\documentclass[11pt]{report}
\usepackage[margin=1in]{geometry}
\usepackage{amsmath,amssymb,amsthm}
\usepackage{enumitem}
\usepackage{xcolor}
\usepackage{listings}
\usepackage{fontspec}
\usepackage{fancyvrb}
\usepackage{longtable}
\usepackage{array}
\newenvironment{mdland}{\clearpage\begingroup
 \pagewidth=11in \pageheight=8.5in
 \setlength\textwidth{9.6in}\setlength\columnwidth{9.6in}\setlength\linewidth{9.6in}\setlength\hsize{9.6in}\setlength\textheight{6.9in}%
 \setlength\oddsidemargin{\dimexpr 0.7in-1in-\hoffset\relax}\setlength\evensidemargin{\oddsidemargin}%
 \setlength\topmargin{\dimexpr 0.7in-1in-\voffset-\headheight-\headsep\relax}}{\clearpage\endgroup}
\usepackage{hyperref}
\hypersetup{colorlinks=true,linkcolor=blue!50!black,urlcolor=blue!50!black,citecolor=blue!50!black}
\setcounter{tocdepth}{1}
\setcounter{secnumdepth}{3}
% ---- fonts used only for the converted Markdown sections (Unicode mathematics in the sources)
\directlua{luaotfload.add_fallback("mdfb",{"Cambria Math:mode=node;script=latn;","Cambria:mode=node;script=latn;"})}
\newfontfamily\mdfont{Cambria}[Ligatures=NoCommon,BoldFont={Cambria Bold},ItalicFont={Cambria Italic},BoldItalicFont={Cambria Bold Italic},RawFeature={fallback=mdfb}]
\newfontfamily\mdtt{Consolas}[Ligatures=NoCommon,RawFeature={fallback=mdfb}]
\setlistdepth{9}
\newlist{mdbul}{itemize}{9}
\setlist[mdbul]{label=\textbullet,leftmargin=1.5em,itemsep=2pt,topsep=2pt,parsep=2pt}
\newlist{mdnum}{itemize}{9}
\setlist[mdnum]{label={},leftmargin=2.3em,labelwidth=1.9em,labelsep=0.4em,align=left,itemsep=2pt,topsep=2pt,parsep=2pt}
\newcommand{\mdlbl}[1]{\textup{#1}}
\newenvironment{mdblock}{\begingroup\mdfont\raggedright\setlength{\parindent}{0pt}\setlength{\parskip}{0.55em}\emergencystretch=3em}{\par\endgroup}
\newenvironment{mdquote}{\begin{list}{}{\setlength{\leftmargin}{1.2em}\setlength{\rightmargin}{0pt}\setlength{\topsep}{2pt}\setlength{\parsep}{0.55em}\setlength{\itemsep}{0pt}}\item\relax}{\end{list}}
\newcommand{\mdhead}[2]{\par\vspace{0.7em}\noindent{\bfseries\ifcase#1\or\large\or\normalsize\or\normalsize\or\small\or\small\or\small\fi #2}\par\nopagebreak\vspace{0.1em}}
\newcommand{\mdhr}{\par\noindent\rule{\linewidth}{0.4pt}\par}
\newcommand{\mdsec}[1]{\phantomsection\section*{#1}\addcontentsline{toc}{section}{#1}}
\newcommand{\mdsubsec}[1]{\phantomsection\subsection*{#1}}
% ---- start of the (verbatim) proof of one chapter: local group, section/equation numbers as in the source proof
\newcommand{\mdproofbegin}{\begingroup
  \renewcommand{\thesection}{\arabic{section}}\renewcommand{\theequation}{\arabic{equation}}%
  \setcounter{section}{0}\setcounter{equation}{0}\setcounter{footnote}{0}%
  \renewcommand{\appendix}{\par\setcounter{section}{0}\renewcommand{\thesection}{\Alph{section}}}%
  \theoremstyle{plain}}
\newcommand{\mdproofend}{\par\endgroup}
\newcommand{\chaptersource}[1]{\par\noindent{\small\itshape\raggedright #1\par}\medskip}

\begin{document}
\begin{center}
{\LARGE\bfseries SMC-BVS v3 — collected proofs of the prove-first list (2026-10-22)\par}
\vspace{0.8em}
{\large Status: proved; per-theorem reviews (Astra x2 + Opus xhigh) and campaign review (Fable 5.1 + Astra xhigh): all ten correct; NOT ADOPTED.\par}
\end{center}

\tableofcontents
\chapter*{Overview}
\addcontentsline{toc}{chapter}{Overview}
\noindent This document collects the ten finished proofs of the prove-first list of the SMC-BVS v3 theory architecture. It is assembled by copying: for each theorem, part (a) is the original conjecture text, part (b) is the \texttt{STATEMENT.md} that was proved, part (c) is the full final proof, and part (d) is the review record (statement-change and review sections of \texttt{FINAL\_REPORT.md}, and every entry of \texttt{CLAIMS\_OF\_FATAL\_FLAWS.md}). Markdown parts are typeset from the unchanged Markdown text; the proofs are the unchanged \LaTeX{} bodies, with the three mechanical changes listed in \texttt{ASSEMBLY\_NOTES.md} (chapter prefixes on labels and on theorem-environment names; no macro had to be renamed, because each proof's own macros are local to its chapter).

\noindent Status of every result in the table: proved, with the stage-A verdict of the lead and Astra \#1 and the stage-B verdicts of Astra \#2 and an Opus xhigh reviewer; the campaign review (stage C) is pending; nothing is adopted. Abbreviations: A, B = review stages; r1, r2, \dots = review rounds; ``0/1/2 = fatal/gap/exposition findings.

\begin{mdland}
{\scriptsize
\begin{longtable}{|>{\raggedright\arraybackslash}p{\dimexpr(\textwidth-14\tabcolsep-8\arrayrulewidth)*9/100\relax}|>{\raggedright\arraybackslash}p{\dimexpr(\textwidth-14\tabcolsep-8\arrayrulewidth)*21/100\relax}|>{\raggedright\arraybackslash}p{\dimexpr(\textwidth-14\tabcolsep-8\arrayrulewidth)*6/100\relax}|>{\raggedright\arraybackslash}p{\dimexpr(\textwidth-14\tabcolsep-8\arrayrulewidth)*16/100\relax}|>{\raggedright\arraybackslash}p{\dimexpr(\textwidth-14\tabcolsep-8\arrayrulewidth)*11/100\relax}|>{\raggedright\arraybackslash}p{\dimexpr(\textwidth-14\tabcolsep-8\arrayrulewidth)*21/100\relax}|>{\raggedright\arraybackslash}p{\dimexpr(\textwidth-14\tabcolsep-8\arrayrulewidth)*16/100\relax}|}
\hline
\textbf{Theorem and parts} & \textbf{Plain-language claim} & \textbf{Outcome} & \textbf{Statement changes} & \textbf{Hypotheses taken from other conjectures} & \textbf{Review history (Stage A, Stage B, judge)} & \textbf{Final proof file and SHA-256} \\ \hline\endhead
\textbf{T09 (T7); all parts (i)--(v)} & The exact start sampler and every move kernel are valid; every round, in any order, mixes at rate at least $\varepsilon_\varphi$; a tiny $\varepsilon_\varphi$ does not imply slow mixing. & Proved (Stage B done, settled); not adopted & None to hypotheses. Two informal sentences, (iii) and (v), written as precise statements; the lead's first reading of (v) was corrected after Astra \#1 round 1. & None (self-contained; uses no P1, P6 or sibling v3 conjecture). & A: Astra \#1 r1 not verified (1 fatal, the lead's own gloss of (v); 1 exposition), r2 verified. B: Astra \#2 r1 and r2 verified; Opus xhigh r1 verified (7 exposition), r2 verified (3 exposition). Judge used: no. & \texttt{T09\_\allowbreak{}T7\_\allowbreak{}kernel\_\allowbreak{}validity/\allowbreak{}proof/\allowbreak{}PROOF\_v3.tex}\newline {\ttfamily\tiny bab55fdf\allowbreak{}1074bdb5\allowbreak{}09f4409f\allowbreak{}c08d9b5e\allowbreak{}d342072d\allowbreak{}a620a2df\allowbreak{}122bdc8d\allowbreak{}03276ba6} \\ \hline
\textbf{T11(iii) only} & For a model $S^*$ fixed by the data and the calibration, $P(S^*$ never among the final run's $Q_f$ iid $\pi_0$ draws $\mid\mathcal C)=(1-\pi_0(S^*))^{Q_f}\le e^{-Q_f\pi_0(S^*)}$. & Proved (Stage B done, settled); not adopted & None to the claim. Setting (data fixed, launch event, EX-S/EX-A/CAL, final-run description) written out; the scoring corollary was restricted to independence proposals after Astra \#1 round 2. & None (EX-A, EX-S, CAL only). & A: Astra \#1 r1 verified (3 gaps, supplementary parts), r2 verified (1 gap, corollary), r3 verified (none). B: Astra \#2 r1, r2, r3 verified; Opus xhigh r1 verified (1 gap F8, 7 exposition), r2 verified (4 exposition), r3 verified (1 note). Judge used: no. & \texttt{T11\_\allowbreak{}LLED\_\allowbreak{}scoring\_\allowbreak{}ledger/\allowbreak{}proof/\allowbreak{}PROOF\_v5.tex}\newline {\ttfamily\tiny 7f4e94c4\allowbreak{}7415f2e4\allowbreak{}bb64c590\allowbreak{}38a4193b\allowbreak{}fbc4fcf7\allowbreak{}1a774fd8\allowbreak{}c429a741\allowbreak{}84886260} \\ \hline
\textbf{T12 (T8); all parts (i)--(iv)} & Given the frozen calibration, each island's $\widehat Z_r$ is exactly unbiased and the recorded RB PIPs satisfy $E[\widehat Z_r\widehat{\mathrm{PIP}}_{rj}]=Z\,\mathrm{PIP}_j$; pooled $0.05\widehat Z\le Z$ w.p.\ $\ge95\%$; $\log\widehat Z$, normalised PIPs and conditioning on diagnostics are not covered (limits proved). & Proved (Stage B done, settled); not adopted & Made precise after an Astra \#1 statement check (approved with changes): arbitrary integer parameters; START imported; negative parts of (iv) as component statements; (iv.d) restricted after the lead's over-extension; target formulas written out after Stage B. No conclusion changed in Stage B. & START (exact $\pi_0$ sampler and $\bar Z_0$), imported from T09(i) / project P1. & A: Astra \#1 statement check approved with changes; review of v2 verified (0 fatal / 0 gap / 3 exposition). B: round 1 on v3: Astra \#2 not verified (0/1/1), Opus xhigh not verified (0/2/6); round 2 on v4: both verified (0/0/2 each). Judge used: no. & \texttt{T12\_\allowbreak{}T8\_\allowbreak{}frozen\_\allowbreak{}run\_\allowbreak{}identities/\allowbreak{}proof/\allowbreak{}PROOF\_v6.tex}\newline {\ttfamily\tiny d6b85fc0\allowbreak{}fea043a8\allowbreak{}df467aeb\allowbreak{}50c969a1\allowbreak{}b47dab1e\allowbreak{}29c42fce\allowbreak{}cff29a15\allowbreak{}496396f6} \\ \hline
\textbf{T13 (T10); all parts (i)--(iv)} & The restricted-ESS test on a candidate set $V$ controls the full-space $1+\chi^2$ once $V$ carries enough mass at the evaluation points; coverage on $[0,1]$ alone is not enough. & Proved (Stage B done, settled); not adopted & Nine precisions, substance unchanged: $\varepsilon_C\in[0,1)$; (iv)(a) avoidance an explicit hypothesis; grid paired with its $V$; witnesses from real method data; coverage at $v$; all real $u,v$; positivity derived; one fixed implementation for (iv)(b); (ii)(d) strengthened. Three Astra \#1 statement checks (one restart). & None from other conjectures (assumption code COVER only). & A: Astra \#1 v1 not verified (0/1/2), v2 verified (0/0/1), v3 verified. B: v3: Astra \#2 not verified (0/1/0), Opus xhigh not verified (0/4/5); v4: both verified (Opus 0/0/3). Judge used: no. & \texttt{T13\_\allowbreak{}T10\_\allowbreak{}restricted\_\allowbreak{}ESS/\allowbreak{}proof/\allowbreak{}PROOF\_v4.tex}\newline {\ttfamily\tiny 4049019c\allowbreak{}683731dd\allowbreak{}56e83c64\allowbreak{}0696bd48\allowbreak{}fc41335d\allowbreak{}f6bbbbf3\allowbreak{}5c0fed01\allowbreak{}9aa60f2d} \\ \hline
\textbf{T14(i) (L-FM, singleton case); part (i) only} & $E|\widehat Z_r/Z-1|\le\frac{2}{1-\varepsilon_*}\sum_t\eta_t$ for every particle count $N$, where $\eta_t$ is the bridge mass outside a fixed model $S^\circ$. & Proved (Stage B done, settled); not adopted & No mathematical change. Notation (numerical subscript = stage index) and the reading of ``endpoint term is essential'' as (i-d). & None (the T12 unbiasedness identity is re-proved inside the proof). & A: Astra \#1 r1 verified (1 exposition), r2 verified. B: Astra \#2 r1 and r2 verified; Opus xhigh r1 verified (3 exposition), r2 verified. Judge used: no. & \texttt{T14\_\allowbreak{}LFM\_\allowbreak{}first\_\allowbreak{}moment\_\allowbreak{}evidence/\allowbreak{}proof/\allowbreak{}PROOF\_v3.tex}\newline {\ttfamily\tiny c12954d1\allowbreak{}7b7f7383\allowbreak{}b026b139\allowbreak{}25c4558e\allowbreak{}917af95c\allowbreak{}14e40b84\allowbreak{}4529f490\allowbreak{}397efdd4} \\ \hline
\textbf{T15 (L-OCC); all parts (i)--(iv)} & The mean fraction of an island's particles outside a fixed set $\mathfrak A$ is at most $\pi(\mathfrak A^c)/c+P(\widehat Z_r/Z<c)$ (and a tail version); the lower tail must be islandwise, not pooled. & Proved (Stage B done, settled); not adopted & None to the inequalities. The T12 dependence made the explicit hypothesis H-T12; the pooled-tail sentence made precise as (iv); after Stage B, (iv) extended to every $a\in(0,cu/4]$ with populations (iv-b); later changes wording only. & H-T12: the conclusion of T12(i), $E_{\mathcal C}[\widehat Z_r\hat q_{r,\tau}(f)]=Z\pi(f)$. & A: Astra \#1 r1 verified (1 exposition), r2 verified. B: Astra \#2 r1 verified, r2 verified (2 exposition), r3 verified; Opus xhigh r1 verified (5 exposition), r2 verified (4), r3 verified (1 optional note). Judge used: no. & \texttt{T15\_\allowbreak{}LOCC\_\allowbreak{}occupancy/\allowbreak{}proof/\allowbreak{}PROOF\_v4.tex}\newline {\ttfamily\tiny 246f0707\allowbreak{}6f0d71fd\allowbreak{}0f08f919\allowbreak{}1ae20e99\allowbreak{}9b4370ef\allowbreak{}e75cdfdf\allowbreak{}7c3e3383\allowbreak{}7ead8b68} \\ \hline
\textbf{T20 (T15); clauses (i)--(ii)} & \texttt{local\_certified} means what it says: no scored neighbour is better; one scan from a neighbour of the true model certifies it on $E_n$, which holds w.p.\ $\to1$ under H03. & Proved (Stage B done, settled); not adopted & No hypothesis changed or added. Method \S6.10 formalised as an algorithm, H03 explicit, EXEC named, seven readings recorded; wording corrections in v2, v3, v3.1. & H03: the conclusion of T03(i), used only in (ii-C). & A: Astra \#1 r1 not verified (1 ``fatal'' in an explanatory note, 1 gap, 1 exposition), r2 verified. B: Astra \#2 verified (1 cosmetic remark); Opus xhigh verified (10 exposition or scope notes); ruling round: both verified. Judge used: no. & \texttt{T20\_\allowbreak{}T15\_\allowbreak{}local\_\allowbreak{}search/\allowbreak{}proof/\allowbreak{}PROOF\_v3.tex}\newline {\ttfamily\tiny 2e34a791\allowbreak{}20c1e92f\allowbreak{}6738acd9\allowbreak{}7dbd54fb\allowbreak{}8051fbc4\allowbreak{}7e38931a\allowbreak{}48d09667\allowbreak{}d2332b24} \\ \hline
\textbf{T22 (T16); parts (i), (ii), (iv)} & Start-side bounds: $\pi_0(S^\circ)\ge(1-q)^{s+1}/(1+e)^s$; an exact product formula for $\pi_0(S_0)/\pi_0(S^\circ)$; $q_\varphi(S_0)\ge\min\{\pi_0(S_0),\pi(S_0)\}$ along the bridge. & Proved (Stage B done, settled); not adopted & Approved by Astra \#1 (APPROVED WITH CORRECTIONS): C1 added hypothesis $S_0\in\Omega$ for (ii) and (iv); C6 generalisation to any real $X,y$ and block partition; C7 scope $u<1$; readings (iv-c) conditional. & (iv-c) only: the conclusions of T03 and of T22(iii) as explicit hypotheses. & A: Astra \#1 statement check approved with corrections; PROOF\_v1 verified (1 exposition). B: Astra \#2 verified (none); Opus xhigh verified (5 exposition); round 2: both verified (Opus 1 new cosmetic note). Judge used: no. & \texttt{T22\_\allowbreak{}T16\_\allowbreak{}start\_\allowbreak{}aggregation/\allowbreak{}proof/\allowbreak{}PROOF\_v3.tex}\newline {\ttfamily\tiny 08c3521c\allowbreak{}5ad28bc7\allowbreak{}d2be26f5\allowbreak{}093b191f\allowbreak{}cb0596d9\allowbreak{}5eae811b\allowbreak{}850ed269\allowbreak{}22f95be1} \\ \hline
\textbf{T23 (T18); part (ii) only} & Pathwise RB PIP error $\le\bar B_r+u_\pi$ on every run, and under BASIN $\le\delta+(1-a)^{\lfloor L/r_0\rfloor}+(L+4)(\eta+\rho)+u_\pi$ in expectation; $P(\widehat{\mathrm{MPM}}\ne S_0)\le2\mathcal B_n$. & Proved (Stage B done, settled); not adopted & None. Five readings recorded in the statement (section 5); EX-S used in the form (F), which both Stage B reviewers ruled faithful. & None (no sibling conjecture; the BASIN conditions are hypotheses of the theorem itself). & A: Astra \#1 on v1 verified (no findings). B: r1 Astra \#2 verified, Opus xhigh verified (6 exposition); r2 both verified (7 further exposition); r3 Astra \#2 verified, Opus verified (2 exposition); r4 both verified (Opus 1 cosmetic note). Judge used: no. & \texttt{T23\_\allowbreak{}T18\_\allowbreak{}basin\_\allowbreak{}prefix/\allowbreak{}proof/\allowbreak{}PROOF\_v5.tex}\newline {\ttfamily\tiny 36d38c99\allowbreak{}b7870cff\allowbreak{}b20d897f\allowbreak{}2f119c6c\allowbreak{}cb656730\allowbreak{}959df9e1\allowbreak{}7b39332d\allowbreak{}63a564ff} \\ \hline
\textbf{T31 (T22); part (i) only} & From unbiasedness alone, $P_c(Z\ge\alpha\widehat Z)\ge1-\alpha$ and $E_c\log\widehat Z\le\log Z$; this is not a relative-error interval, and conditioning on a final E2 flag can destroy coverage. & Proved (Stage B done, settled); not adopted & None to hypotheses or conclusions. ``Under T12'' made the explicit hypothesis H-T12; the two informal cautions written as non-implications (c), (d); explanatory wording revised twice after review. & H-T12: T12(iii), unbiasedness $E_c[\widehat Z]=Z$. & A: Astra \#1 on v1 verified (none). B: r1 Astra \#2 verified (none), Opus xhigh verified (11 exposition); r2 Astra \#2 verified (2), Opus verified (5); r3 Astra \#2 verified (none), Opus verified (3 cosmetic). Judge used: no. & \texttt{T31\_\allowbreak{}T22\_\allowbreak{}certificates/\allowbreak{}proof/\allowbreak{}PROOF\_v3.tex}\newline {\ttfamily\tiny e140d9b8\allowbreak{}1cd01b54\allowbreak{}b1b05201\allowbreak{}1975d5d5\allowbreak{}6a2db13d\allowbreak{}c64095d4\allowbreak{}fe7a6581\allowbreak{}26e2d4cc} \\ \hline
\end{longtable}}
\end{mdland}


\chapter*{Dependencies between the proofs}
\addcontentsline{toc}{chapter}{Dependencies between the proofs}
\noindent The list below is taken from the \texttt{STATEMENT.md} files of the ten proofs (reproduced in each chapter, part (b)). ``Hypothesis'' means that the other result is used only as an explicit assumption of the statement being proved here; nothing about it is proved in the chapter that uses it.
\begin{itemize}
\item \textbf{T09}: no dependency (self-contained).
\item \textbf{T11(iii)}: no dependency (EX-A, EX-S, CAL only).
\item \textbf{T12}: hypothesis START (exact start sampler and exact normaliser $\bar Z_0$), imported from \textbf{T09(i)} (project P1). T12 does not otherwise use a sibling result.
\item \textbf{T13}: no dependency on another result (assumption code COVER is a hypothesis of part (ii), not a result).
\item \textbf{T14(i)}: no dependency (the T12 unbiasedness identity it builds on is re-proved inside its proof).
\item \textbf{T15}: hypothesis H-T12, the conclusion of \textbf{T12(i)}: $E_{\mathcal C}[\widehat Z_r\,\hat q_{r,\tau}(f)]=Z\,\pi(f)$. Parts (i)--(ii) use only the instance $f=\mathbf 1_{\mathfrak A^c}$; part (iv) uses no hypothesis.
\item \textbf{T20}: hypothesis H03, the conclusion of \textbf{T03(i)} (not part of this collection), used only in (ii-C).
\item \textbf{T22}: parts (i), (ii), (iv-a), (iv-b) are deterministic and use no other result; the conditional corollary (iv-c) takes the conclusions of \textbf{T03} (not part of this collection) and of \textbf{T22(iii)} (not part of this collection) as explicit hypotheses.
\item \textbf{T23(ii)}: no dependency on another conjecture; the occupancy step is proved directly from the BASIN conditions and does not cite T15 or any part (i).
\item \textbf{T31(i)}: hypothesis H-T12, taken from \textbf{T12(iii)}: unbiasedness $E_c[\widehat Z]=Z$ for every successful calibration value $c$.
\end{itemize}
\noindent Within this collection the dependency arrows are therefore: T09(i) $\to$ T12; T12(i) $\to$ T15; T12(iii) $\to$ T31(i). The conclusions of T03 and T22(iii) enter T20 and T22 only as hypotheses and are outside the collection.

\chapter*{Shared definitions}
\addcontentsline{toc}{chapter}{Shared definitions}
\chaptersource{Source: \texttt{theorems/00\_COMMON.md}, converted from Markdown; the text is not changed.}
\begin{mdblock}
\mdhead{1}{00\_COMMON.md — shared definitions for every theorem file in this folder}

Status of everything in this folder: \textbf{conjectured, not proven.} Methodology: SMC-BVS v3.0 of 2026-10-19, {\mdtt C:\textbackslash{}\allowbreak{}Users\textbackslash{}\allowbreak{}felip\textbackslash{}\allowbreak{}OneDrive\textbackslash{}\allowbreak{}Documents\textbackslash{}\allowbreak{}Local\_\allowbreak{}to\_\allowbreak{}global\_\allowbreak{}Bayesian\textbackslash{}\allowbreak{}01\_\allowbreak{}Methodology\_\allowbreak{}and\_\allowbreak{}Theory\textbackslash{}\allowbreak{}Methodology\textbackslash{}\allowbreak{}SMC\_\allowbreak{}BVS\_\allowbreak{}v3\_\allowbreak{}20261019\textbackslash{}\allowbreak{}METHODOLOGY\_\allowbreak{}v3.\allowbreak{}md}, SHA-256 {\mdtt 93043758058a6d422cca6ca4b6be8f3a9e8149b6f6de2e1bab1ce31c44b8019c}. The exact algorithm (blocks, tables, composite start, bridge, pilot, candidate grid, final run, local search, statuses) is written out in {\mdtt .\allowbreak{}.\allowbreak{}\textbackslash{}\allowbreak{}THEORY\_\allowbreak{}BRIEF.\allowbreak{}md} §2; the method file is authoritative. Each theorem file {\mdtt T<k>\_\allowbreak{}<name>.\allowbreak{}md} names the definitions and inventory codes it uses; give proof.md the theorem file, this file, THEORY\_BRIEF.md and the method file (all by path).

Felipe's fixed decisions (no theorem changes them): the target (standardised X, y; β\_S \textasciitilde{} N(0, nσ²(G\_S+I)⁻¹); σ² \textasciitilde{} IG(0.01, 0.01); BB(1, p) capped at K = min(p, ⌊n/2⌋)); the HPM over all scored models with exact-equality ties; B = min(⌊√n⌋, ⌊√p⌋, 16); neighbour threshold ρ\_n = √(2 log p/n); no secondary prior arm; the budgets of METHODOLOGY §3.

\mdhead{2}{1. Data, target, decision objects}

\begin{mdbul}
\item Raw data (X\textasciicircum{}r, y\textasciicircum{}r), n ≥ 3, no constant column, non-constant y\textasciicircum{}r. C = I − 11ᵀ/n; D\_x = diag of retained raw column sample SDs. Standardised X = CX\textasciicircum{}rD\_x⁻¹, y = Cy\textasciicircum{}r/s\_y with s\_y² = ‖Cy\textasciicircum{}r‖²/(n−1); so ‖x\_j‖² = ‖y‖² = n − 1. Contrast coordinates: U ∈ ℝ\textasciicircum{}\{n×(n−1)\}, UᵀU = I, UUᵀ = C; y\_c = Uᵀy\textasciicircum{}r, X\_c = UᵀX.
\item State space Ω = \{S ⊆ [p] : |S| ≤ K\}, K = min(p, ⌊n/2⌋). G\_S = X\_SᵀX\_S = (n−1)R\_S; z\_S = X\_Sᵀy; H\_S = (n+1)G\_S + I; D\_S = det H\_S/det(G\_S+I); Q\_S = yᵀy − n z\_SᵀH\_S⁻¹z\_S; ν = 0.01 + n/2; a₀ = b₀ = 0.01.
\item m(S) = (2b₀)\textasciicircum{}\{a₀\}Γ(ν)/(π\textasciicircum{}\{n/2\}Γ(a₀)) · D\_S\textasciicircum{}\{−1/2\}(2b₀ + Q\_S)\textasciicircum{}\{−ν\}; w(k) = ½/C(2p−1, k); W\_K = Σ\_\{k≤K\} C(p,k)w(k); ϖ(S) = w(|S|)/W\_K.
\item F(S) = log[ϖ(S)m(S)]; Z = Σ\_Ω ϖ(S)m(S); π(S) = e\textasciicircum{}\{F(S)\}/Z; BF(S) = m(S)/m(∅).
\item Global HPM S† = argmax\_Ω F (ties: smallest size, then lexicographic). Reported HPM Ŝ = the same rule over 𝒮, the set of all models ever scored (calibration, final run, rejected proposals, local scans); 𝒮\_cal ⊆ 𝒮 is the calibration's scored set. PIP\_j = π(j ∈ S); MPM = \{j : PIP\_j > ½\}.
\item Statistical model: y\textasciicircum{}r = α*1 + X\textasciicircum{}rβ* + ε; μ = CX\textasciicircum{}rβ* = Xθ*, θ* = D\_xβ*; S₀ = supp β*, s₀ = |S₀|. Contrast-scale parameters θ = s\_yβ, v = s\_y²σ². P₀ is over ε; "joint probability" adds the algorithm's randomness. r\_n := √(log n/n). m\_n(ω) := ⌈(s₀ log(2p) + (s₀+1)(log n)/2 + ω)/log 2⌉, m\_n := m\_n(½ log n).
\end{mdbul}

\mdhead{2}{2. Sampler objects}

\begin{mdbul}
\item Blocks C₁, …, C\_J (deterministic in X: complete linkage, top-down cut, next-fit packing), |C\_b| ≤ B = min(⌊√n⌋, ⌊√p⌋, 16). Tables r\_b(A) = BF(A) for A ⊆ C\_b (full response).
\item γ₀(S) = ϖ(S)∏\_b r\_b(S ∩ C\_b); π₀ = γ₀/Z̄₀ (Z̄₀ by exact convolution of block size profiles with the capped size prior; backward sampler gives exact iid π₀ draws). Start mode S̃₀ := argmax π₀.
\item Discrepancy d(S) = log BF(S) − Σ\_b log r\_b(S ∩ C\_b). Bridge γ\_φ = γ₀e\textasciicircum{}\{φd\}, q\_φ = γ\_φ/Z̄\_φ (defined for real φ; production φ ∈ [0,1]); q₀ = π₀, q₁ = π, Z = m(∅)Z̄₁.
\item Round K\_φ: m = p + B + 2 labelled steps in a uniformly random order: p Gibbs flips (o\_j(S, φ) = pre-draw inclusion probability of j at S; a\_j(S,φ) = γ\_φ(S+j)/γ\_φ(S)), one uniform swap, B neighbour swaps (symmetric-union neighbour graph N(·) at threshold ρ\_n), one independence proposal from π₀. Elementary kernels K\_e; reflected kernels K\_e\textasciicircum{}W(x,A) = K\_e(x, A∩W) + K\_e(x, W\textasciicircum{}c)1\_A(x).
\item Calibration 𝓒 := (grid 0 = φ₀ < … < φ\_T = 1, T, L, N\_e, tier e, candidate set V, every V-independent setting), measurable w.r.t. streams 0–9; final streams 10–13 are fresh; escalation is decided before final randomness. Δ\_t = φ\_t − φ\_\{t−1\}; w\_t = 2φ\_t − φ\_\{t−1\}; 𝓔\_grid = \{φ₀,…,φ\_T, w₁,…,w\_T\}. Restricted ESS ρ\_V(u,v) = Z̄\_v(V)²/(Z̄\_u(V)Z̄\_\{2v−u\}(V)).
\item Final run: 4 fresh islands × N\_e iid π₀ draws; frozen grid; multinomial resampling after every reweighting; one round per stage; L burn-in rounds and 4 measured rounds at φ = 1; local search (≤ ℓ moves, scans score every feasible flip and swap of the centre). Outputs: Ẑ\_r (island evidence), Ẑ = ¼Σ\_rẐ\_r, bound 0.05Ẑ; PIP̂\_\{rj\} = Rao–Blackwell average of the pre-draw Gibbs conditionals over the 4 measured rounds; primary PIP̂\_j (equal weight), secondary (Ẑ\_r-weighted); MPM̂; Ŝ; statuses \{exact, completed, schedule\_limit, numerical\_failure, resource\_limit, y\_constant, no\_predictors, invalid\_input\}; flags (E1–E3, low\_empirical\_ESS, local\_certified, global\_certified, …).
\item Budgets: N₀ = max(2500, n, p); H₀ = ⌈512N₀/2500⌉; ℓ = ⌈log₂(n+p+1)⌉; T₀ = max(1000, ℓ²); (N\_e, H\_e, T\_e) = 4\textasciicircum{}e(N₀, H₀, T₀), e ∈ \{0,1\}; L = max(4, ⌈log₂(n+1)⌉). Q\_f := 4N\_e(T + L + 5) iid π₀ draws scored in the final run (initialisation + independence proposals); 4N\_e(L+6) ≤ Q\_f ≤ 4N\_e(T\_e + L + 5).
\item Exact branch: Σ\_\{k≤K\} C(p,k) ≤ (5H₀ + 4N₀)p; it precedes block/Gram construction and scores every feasible model.
\item Stage notation: q\_t = q\_\{φ\_t\}, G\_t = e\textasciicircum{}\{Δ\_t d\}, Z\_t = m(∅)Z̄\_\{φ\_t\}; q̂\_t = empirical population after stage-t resampling and mutation; Ẑ\_\{r,t\} = island r's evidence product through stage t. Reflected normalisation on W: q\_t\textasciicircum{}W = q\_t(·|W), g\_t\textasciicircum{}W = G\_t/[Z̄\_\{φ\_t\}(W)/Z̄\_\{φ\_\{t−1\}\}(W)], 𝒬\textasciicircum{}W\_\{s,t\} = ∏\_\{j=s+1\}\textasciicircum{}t (g\_j\textasciicircum{}W K\_j\textasciicircum{}W); comparator mean with the original start normaliser Z·r\_W, r\_W = q₁(W)/q₀(W).
\item Leakage quantities: u\_\{T,l\} = unit residual of x\_l on X\_T; δ\_\{T,l\} = u\_\{T,l\}ᵀμ; Δ\_T = ‖(I−P\_T)μ‖²/(n−1); v\_T = σ*² + Δ\_T; h\_n = √(2(A+ε) log n); E\_\{T,l\} = min\{δ²\_\{T,l\}/(2Δ\_T), δ²\_\{T,l\}/(2v\_T) + |δ\_\{T,l\}|σ*h\_n/v\_T\} (E\_\{T,l\} = 0 when δ\_\{T,l\} = 0).
\item Truth stability: δ₁ = 1 − π(S₀); u\_π = δ₁/π(S₀); ρ\_* = min\{1, (B+3)u\_π\}.
\item Regimes: R1: n → ∞, s₀ fixed, p ≤ n\textasciicircum{}A. R2: p fixed (SMC statements are component results; the exact branch eventually applies). R3: data and 𝓒 fixed, N or the island count varies (idealised). R4: actual budgets along R1. Fixed-data identities condition on a successful frozen 𝓒, never on passing later diagnostics.
\end{mdbul}

\mdhead{2}{3. Assumption inventory (codes used by the theorem files)}

{\small
\begin{longtable}{|>{\raggedright\arraybackslash}p{\dimexpr(\linewidth-6\tabcolsep-4\arrayrulewidth)*91/1000\relax}|>{\raggedright\arraybackslash}p{\dimexpr(\linewidth-6\tabcolsep-4\arrayrulewidth)*673/1000\relax}|>{\raggedright\arraybackslash}p{\dimexpr(\linewidth-6\tabcolsep-4\arrayrulewidth)*236/1000\relax}|}
\hline
\textbf{Code} & \textbf{Condition} & \textbf{Status} \\ \hline\endhead
G & Fixed design; ε \textasciitilde{} N\_n(0, σ*²I), σ*² > 0 fixed; intercept removed by contrasts. & Standard/project (paper\_draft §2); external precedent UNVERIFIED \\ \hline
R & s₀ fixed; p ≤ n\textasciicircum{}A, A ≥ 1; ‖μ‖²/n ≤ M; λ\_min(R\_\{S₀\}) ≥ φ₀² > 0; eventually s₀ ≤ K. & Project bundle (paper\_draft T6); truth-dependent \\ \hline
SE & λ\_min(R\_T) ≥ φ₀² for every nonempty T with |T| ≤ min(p, m\_n + s₀). & Project (proofs T7\_T8 §1, T9\_C9 §9); strong; not checkable exhaustively \\ \hline
BM-C & min\_\{j∈S₀\} θ*\_j² ≥ M\_*σ*² log n/(nφ₀²), M\_* = 2L\_*+1, L\_* = max\{4C\_*, 32[s₀(A+1)+2]\}, C\_* = 2s₀ + 3[Aj\_* + s₀ + log(j\_*+1) + 1], j\_* = ⌈(2A+1)s₀⌉ + 1. & New candidate fixed constant (J1 audit of the project proof; provisional); the proving run may replace M\_* \\ \hline
BM-∞ & Same inequality with M\_n → ∞ in place of M\_*. & Project sufficient condition (v2's BM) \\ \hline
DX & min retained raw column SD ≥ d₋ > 0. & New; observable \\ \hline
GL & R\_\{S₀\} → Γ₀ ≻ 0 with consistent support ordering. & Project limit form; generic BvM precedent UNVERIFIED \\ \hline
ORTH & G = (n−1)I\_p (so p ≤ n−1); for A4: fixed 1 ≤ s₀ < p, equal magnitudes b\_n on S₀, b\_n²/σ*² bounded. & Project benchmark (P4/Pb1); observable \\ \hline
COH & c\_n = max\_\{i≠j\}|r\_ij|, c\_n log(p∨n) → 0, s₀c\_nθ\_max/θ\_min → 0. & New; optional \\ \hline
DUP & Exact signed duplicate columns; or two Gaussian truths with ‖μ₁−μ₂‖/σ* → 0. & New alternative regimes \\ \hline
NIS & No distinct feasible S ≠ S′ with X\_SX\_Sᵀ = X\_\{S′\}X\_\{S′\}ᵀ (equal sizes then follow from traces). & New; decidable from X \\ \hline
VIS & For U\_b* = S₀∩C\_b ⊄ T ⊆ C\_b: ‖(P\_\{T∪U\_b*\} − P\_T)μ‖² ≥ M\_n\textasciicircum{}V v\_T log n, M\_n\textasciicircum{}V → ∞. & New actual-block visibility; not implied by SE/BM \\ \hline
NL\_avg(κ) & (p−s₀)⁻¹Σ\_\{l∉S₀\} max\_\{U*\_\{b(l)\}⊆T⊆C\_\{b(l)\}∖l\} e\textasciicircum{}\{E\_\{T,l\}\} ≤ n\textasciicircum{}\{½−κ\}, fixed κ > 0. & New; Gaussian-design prevalence UNVERIFIED \\ \hline
NL\_η / TL & Sufficient cases: δ²\_\{T,l\} ≤ ησ*² log n with η < (√(2A+1) − √(2A))², over the block comparisons (NL\_η) or a target-state family (TL). & New explicit threshold \\ \hline
TL\_avg(κ; 𝓡) & For every T ∈ 𝓡, the average of e\textasciicircum{}\{E\_\{T,l\}\} over l ∉ S₀∪T is ≤ n\textasciicircum{}\{½−κ\}. & New; uniform before selection \\ \hline
EXP & Uniform actual-block expansion: log[r\_b(T+l)/r\_b(T)] ≤ −½ log(n+1) + (δ\_\{T,l\} + σ*Z\_\{T,l\})²/(2v\_T) + R\_n, Z\_\{T,l\} = u\_\{T,l\}ᵀε/σ*; R\_n = O\_P(1) (strict rate) or o\_P(log n) (exponent rate); plus a uniform lower gain bound c·gain/v\_T − CB log n for visibility. & New explicit proof obligation (J5) \\ \hline
AGG & For feasible S°, s = |S°|, U\_b* = S°∩C\_b: a\_bj = Σ\_\{U⊇U\_b*, |U∖U\_b*|=j\} r\_b(U)/r\_b(U\_b*); D\_b = max\_\{j≥1\}(j!a\_bj)\textasciicircum{}\{1/j\}; q = (p−s)⁻¹Σ\_bD\_b; e = 2p max\_\{b,U: a=|U\_b*∖U|>0\}[r\_b(U)/r\_b(U∪U\_b*)]\textasciicircum{}\{1/a\}. Hypothesis q < 1. Conventions: empty maxima are 0; q = 0 if s = p; e = 0 if s = 0. & New deterministic table condition \\ \hline
BO & X\_\{C\_b\}ᵀX\_\{C\_b′\} = 0 for b ≠ b′. & New benchmark; observable \\ \hline
ONE & S̃₀ and S₀ differ by one flip, or by a same-block proxy swap with |r\_\{ij′\}| > ρ\_n. 𝔓 = \{S̃₀, S₀\}. & Project case definition \\ \hline
LK\_d & 𝓡\_d = \{U∪P : U ⊆ S₀, P ⊆ S₀\textasciicircum{}c, |P| ≤ d\}, fixed d, s₀+d ≤ K; Λ = max\_\{T∈𝓡\_d\}Σ\_\{l∉S₀∪T\} o\_l(T,1); (L+4)2\textasciicircum{}\{s₀+d\}Λ → 0. & New, truth-dependent \\ \hline
EXIT\_d & ξ = max over starts in 𝓡\_d of the one-round probability that the first exit from 𝓡\_d is caused by a non-Gibbs label; (L+4)ξ → 0. & New; does not bound internal non-Gibbs changes \\ \hline
GAIN\_d & Uniform best missing-true addition gain g⁺ → ∞ and uniform best extra-deletion gain g⁻ → ∞ from proper supersets, over 𝓡\_d. & New; multi-extra deletion not derived \\ \hline
BASIN(𝓡, r₀, a, η, ρ, δ) & Given 𝓒: 𝓡 ∋ S₀; inf\_\{x∈𝓡\}K₁\textasciicircum{}\{r₀\}(x, S₀) ≥ a; one-round basin first exit ≤ η; truth prefix-departure ≤ ρ; expected stage-T outside fraction ≤ δ. & New abstract conditions \\ \hline
BOCC(d) & With probability → 1 all specified pilot/final elementary states have size ≤ d. & New; insufficient alone \\ \hline
EX-A & Exact scores, convolutions, probabilities, comparisons. & Declared idealisation \\ \hline
EX-S & Independent ideal uniform streams with the method's stream separation. & Declared idealisation \\ \hline
CAL & 𝓒 uses streams 0–9 only; streams 10–13 fresh; escalation before final randomness. & Method construction (M §§6.8–6.9) \\ \hline
SR & Each Gibbs update evaluates both feasible configurations (incl. the current one); both count as scored. & Scoring reading of M §§6.5, 6.10 (to be confirmed by Felipe) \\ \hline
COVER(W, ε\_C) & q\_ψ(W) ≥ 1 − ε\_C at every ψ ∈ 𝓔\_grid. & New full-space coverage \\ \hline
HIT & A specified scoring route succeeds with probability → 1 (e.g. Q\_fπ₀(S₀) → ∞, table inclusion, scanned one-neighbour centre). & New conditional route \\ \hline
PX & On no prior discovery, ≥ ω\_n k(p−k) uniform-swap executions at S̃₀ occur with probability → 1, ω\_n → ∞. & New occupancy condition \\ \hline
CPL-E & Common-randomness coupling with the reflected particle system through all populations/ancestors used by the last weight evaluation, failure h\_E; sufficient bound: initial mismatch + Σ\_e E[1\{τ>e⁻\}K\_e(X\_\{e⁻\}, W\textasciicircum{}c)]. & New particle-system obligation \\ \hline
CPL-H / RB & HPM transfer: couple all scored proposals and scans. PIP transfer: couple recorded rewards or bound the reward discrepancy. & New output-specific obligations \\ \hline
MISS-E & Data-measurable 𝓜, ∅ ∉ 𝓜; calibration scores avoid 𝓜; CPL-E for W = 𝓜\textasciicircum{}c with h\_E → 0. & New conditional loss hypothesis \\ \hline
MISS-H & Calibration avoidance plus avoidance of every subsequent scoring opportunity, failure h\_H → 0; RB condition for PIP conclusions. & New \\ \hline
RET / OCC & RET: return visits from neighbours of 𝔓. OCC: expected outside fractions ≤ Cq\_t or specified weighted slacks. & New first-moment conditions; not variance premises \\ \hline
BW2 & ∫q\textasciicircum{}W\_\{t−1\}(dx)[g\_t\textasciicircum{}W(x)]²K\_t\textasciicircum{}W(x,dy)K\_t\textasciicircum{}W(x,dz) ≤ C\_t q\_t\textasciicircum{}W(dy)q\_t\textasciicircum{}W(dz). & New two-lineage domination \\ \hline
BW\_W & sup\_\{s<t\}sup\_x 𝒬\textasciicircum{}W\_\{s,t\}1(x)/inf\_x 𝒬\textasciicircum{}W\_\{s,t\}1(x) ≤ C₀. & New future-weight control \\ \hline
WS & Working set W; terminal initial law λ ≤ wπ(·|W); reflected terminal round has absolute gap g > 0; sup\_\{s<t,x\}𝒬\textasciicircum{}W\_\{s,t\}1(x)/q\_s\textasciicircum{}W(𝒬\textasciicircum{}W\_\{s,t\}1) ≤ a\_n. & New \\ \hline
PW2 & With W\_t = Ẑ\_t/Z\_t: E[W²\_\{t−1\}\{(q̂\_\{t−1\}g\_t)² − 1\}] ≤ (a\_t/N)EW²\_\{t−1\} + ν\_t with controlled slack. & New research obligation \\ \hline
OP-a / OP-s & Preflight passes under the allocation; / schedule failure probability → 0 when final outputs are needed. & New operational conditions \\ \hline
NUM & Simultaneous error bound η for all compared log scores at fixed n. & New validation contract \\ \hline
NDR & No limiting finite bisection comparison equals the ESS threshold; no limiting pilot increment equals the floor. & New nondegeneracy \\ \hline
STIFF & Two-state restricted law, high-discrepancy mass a, gap D: ρ(a, D(b−u)2\textasciicircum{}\{−60\}) < 0.9, ρ(a,x) = [a+(1−a)e\textasciicircum{}\{−x\}]²/[a+(1−a)e\textasciicircum{}\{−2x\}]. & Corrected mass-qualified obstruction \\ \hline
UB & Externally justified upper bound U ≥ Z. & Optional external input \\ \hline
\end{longtable}}

Retired as hypotheses: MASS; AGG-stat (a conclusion); ζ-EXIT; OP's numerics clause; V≤, OT-lb, NA (optional research ideas only). Regimes are not declared universally disjoint; exact duplicates conflict with SE and NIS; ORTH excludes p > n.
\end{mdblock}

\chapter{T09 (T7): exact start, invariant kernels and whole-round minorisation}
\chaptersource{Final proof: \texttt{proofs/\allowbreak{}T09\_\allowbreak{}T7\_\allowbreak{}kernel\_\allowbreak{}validity/\allowbreak{}proof/\allowbreak{}PROOF\_v3.tex}. Sources of parts (a), (b), (d): \texttt{theorems/\allowbreak{}T09\_\allowbreak{}T7\_\allowbreak{}kernel\_\allowbreak{}validity.md}, \texttt{STATEMENT.md}, \texttt{FINAL\_REPORT.md}, \texttt{CLAIMS\_OF\_FATAL\_FLAWS.md}.}
\mdsec{(a) Original conjecture text (theorem file, verbatim)}
\begin{mdblock}
\mdhead{1}{T09 (T7): Exact start, invariant kernels and whole-round minorisation}

\textbf{Status:} conjectured, not proven. Methodology SMC-BVS v3.0 (2026-10-19), SHA-256 93043758058a6d422cca6ca4b6be8f3a9e8149b6f6de2e1bab1ce31c44b8019c. Read with {\mdtt 00\_\allowbreak{}COMMON.\allowbreak{}md}.

\textbf{Category / rank:} algorithmic validity and convergence; \textbf{core} (proposition).

\mdhead{2}{Definitions used}

Blocks, tables r\_b, γ₀, Z̄₀, π₀, backward sampler; d, γ\_φ, q\_φ, Z̄\_φ; round K\_φ with m = p + B + 2 labelled steps (Gibbs, uniform swap, B neighbour swaps on the symmetric-union graph N(·), independence proposal from π₀), elementary kernels K\_e, realised-order kernel K\_σ; ε\_φ := inf\_S π₀(S)/q\_φ(S) (00\_COMMON §2).

\mdhead{2}{Assumptions}

EX-A, EX-S (declared idealisations); the actual tables, cap and move rules.

\mdhead{2}{Precise statement}

For every valid component input and every φ ∈ [0, 1]: (i) Convolving the actual block size profiles with the capped size prior gives Z̄₀ exactly; the backward sampler returns exact iid π₀ draws. (ii) Each elementary kernel is q\_φ-reversible, including the boundaries: addition at |S| = K has Gibbs probability 0; the uniform swap self-loops at |S| ∈ \{0, p\}; the neighbour-swap reverse ratio reads 0 when the reverse insertion is not a neighbour; the stated Hastings ratio supplies the reverse proposal probability. (iii) K\_φ = (1/m!)Σ\_σ K\_\{σ(1)\}⋯K\_\{σ(m)\} is q\_φ-invariant and self-adjoint on L²(q\_φ). An individual realised order need not be reversible. (iv) Every realised order satisfies K\_σ(x, ·) ≥ ε\_φ q\_φ(·) with ε\_φ = (Z̄\_φ/Z̄₀)e\textasciicircum{}\{−φ max\_S d(S)\} > 0; hence ‖K\_φ\textasciicircum{}h(x,·) − q\_φ‖\_TV ≤ (1 − ε\_φ)\textasciicircum{}h and the averaged round's absolute L² gap is ≥ ε\_φ. (v) ε\_φ ≤ π₀(S)/q\_φ(S) for every S. An exponentially small available minorisation constant does not establish exponentially slow mixing of the hybrid kernel; neither a switch nor a large discrepancy range alone supplies such a lower bound.

\mdhead{2}{Dependencies}

Builds on: method §§6.3–6.5; project P1/P6. Used by: T10–T19, T23, T25, T28, T29, T36.

\mdhead{2}{Numerical evidence (provisional)}

provisional evidence (Sol check J2): n = 6, p = 4, K = 3, B = 2: element and round detailed-balance residuals ≤ 2.2·10⁻¹⁶ at φ ∈ \{0, ½, 1\}; minorisation constants 1, 0.786, 0.619 with nonnegative slack (−1.1·10⁻¹⁶ roundoff at φ = 0); addition ratios ≥ 1.09 × floor. Files: judge/sol2/J2\_EXACT\_IDENTITIES.md, J2\_matrices.npz.

\mdhead{2}{Proof-route hint (optional)}

Sum block profiles; verify detailed balance including boundaries; pair each ordering with its reversal for self-adjointness; the independence kernel minorises q\_φ (K\_ind(S,S′) ≥ π₀(S′)min\{1, w(S′)/w(S)\} ≥ q\_φ(S′)/sup w, w = q\_φ/π₀) and its stochastic prefix and invariant suffix preserve the minorisation. Hardest step: boundary and Hastings bookkeeping. \emph{When proof.md is run on this file, its lead considers the statement both with and without these hints; the hints never reach the reviewers.}
\end{mdblock}

\mdsec{(b) STATEMENT.md (verbatim)}
\begin{mdblock}
\mdhead{1}{STATEMENT — T09 (T7): exact start, invariant kernels and whole-round minorisation}

Source: {\mdtt theorems/\allowbreak{}T09\_\allowbreak{}T7\_\allowbreak{}kernel\_\allowbreak{}validity.\allowbreak{}md} (SMC-BVS v3.0, METHODOLOGY\_v3.md SHA-256 93043758058a6d422cca6ca4b6be8f3a9e8149b6f6de2e1bab1ce31c44b8019c). Decision at step 1: \textbf{(a) correct as stated.} No hypothesis is added or changed. Two informal clauses (the last sentence of (iii) and the second sentence of (v)) are given a precise mathematical reading below; the original wording is quoted next to each reading.

Standing idealisations (declared in the source): \textbf{EX-A} exact scores, convolutions, probabilities and comparisons; \textbf{EX-S} independent ideal uniform random streams.

\textbf{Change log.} v2 (after review round 1): the last sentence of the reading of (v) said "(equivalently, no lower bound on its mixing time) that tends to 0", which is wrong as written (a mixing-time lower bound tending to 0 is trivially available). It now speaks of a lower bound tending to infinity and adds the explicit bound t\_mix ≤ 55. The original wording, the hypotheses and the existence claim are unchanged.

\mdhead{2}{1. Setting and definitions}

\textbf{Data.} Integers n ≥ 3, p ≥ 1. X ∈ ℝ\textasciicircum{}\{n×p\} and y ∈ ℝⁿ are standardised data: every column x\_j and y are centred, ‖x\_j‖² = ‖y‖² = n − 1. (These are exactly the outputs of the method's preprocessing; a standardised pair is returned unchanged by it.) K = min(p, ⌊n/2⌋); Ω = \{S ⊆ [p] : |S| ≤ K\}.

\textbf{Scores.} For S ⊆ [p], k = |S|: G\_S = X\_SᵀX\_S, z\_S = X\_Sᵀy, H\_S = (n+1)G\_S + I\_k, D\_S = det H\_S / det(G\_S + I\_k), Q\_S = yᵀy − n z\_SᵀH\_S⁻¹z\_S, ν = 0.01 + n/2, b₀ = 0.01, A\_S := 2b₀ + Q\_S, BF(S) = D\_S\textasciicircum{}\{−1/2\} (A\_∅ / A\_S)\textasciicircum{}\{ν\}  (so BF(∅) = 1; this is m(S)/m(∅)). Size prior: w(k) = ½ / C(2p−1, k), W\_K = Σ\_\{k≤K\} C(p,k) w(k), ϖ(S) = w(|S|)/W\_K for S ∈ Ω.

\textbf{Blocks and tables.} C₁, …, C\_J is the method's block partition of [p] (complete linkage, top-down cut, next-fit packing), |C\_b| ≤ B = min(⌊√n⌋, ⌊√p⌋, 16). Tables r\_b(U) = BF(U) for U ⊆ C\_b. Block size profiles e\_b(j) = Σ\_\{U⊆C\_b, |U|=j\} r\_b(U), j = 0, …, |C\_b|.

\textbf{Start.} γ₀(S) = ϖ(S) ∏\_b r\_b(S ∩ C\_b) for S ∈ Ω; Z̄₀ = Σ\_\{S∈Ω\} γ₀(S); π₀ = γ₀/Z̄₀. Convolution: c\textasciicircum{}\{(0)\} = δ₀ (c\textasciicircum{}\{(0)\}(0) = 1, else 0); c\textasciicircum{}\{(b)\}(t) = Σ\_j e\_b(j) c\textasciicircum{}\{(b−1)\}(t − j) (b = 1, …, J), c := c\textasciicircum{}\{(J)\}. \textbf{Backward sampler} (one draw, fresh uniforms): (1) draw the total size k ∈ \{0, …, K\} with probability w(k)c(k) / Σ\_\{k′≤K\} w(k′)c(k′); (2) set t\_J = k; for b = J, J−1, …, 1 draw j\_b with probability e\_b(j\_b) c\textasciicircum{}\{(b−1)\}(t\_b − j\_b) / c\textasciicircum{}\{(b)\}(t\_b) and set t\_\{b−1\} = t\_b − j\_b; (3) for each b draw U\_b ⊆ C\_b with |U\_b| = j\_b with probability r\_b(U\_b)/e\_b(j\_b); return S = ∪\_b U\_b.

\textbf{Bridge.} d(S) = log BF(S) − Σ\_b log r\_b(S ∩ C\_b); γ\_φ(S) = γ₀(S) e\textasciicircum{}\{φ d(S)\} for S ∈ Ω, γ\_φ(S) := 0 for S ∉ Ω; Z̄\_φ = Σ\_Ω γ\_φ; q\_φ = γ\_φ/Z̄\_φ (a probability on Ω, positive everywhere on Ω).

\textbf{Neighbour graph.} N(j) ⊆ [p] \textbackslash{} \{j\} for each j: the method's symmetric-union graph (l ∈ N(j) ⇔ j ∈ N(l)). Statement (ii) is also asserted, with the stated reading convention, for an arbitrary family N(j) ⊆ [p] \textbackslash{} \{j\}.

\textbf{Elementary kernels on Ω at temperature φ} (all states, all moves in Ω):

\begin{mdbul}
\item \textbf{Gibbs j} (j ∈ [p]): with S⁺ = S ∪ \{j\}, S⁻ = S \textbackslash{} \{j\}: K\_\{G,j\}(S, S\textasciicircum{}±) = γ\_φ(S\textasciicircum{}±) / (γ\_φ(S⁺) + γ\_φ(S⁻)). (When j ∉ S and |S| = K, S⁺ ∉ Ω, so the addition has probability 0.)
\item \textbf{Uniform swap}: if |S| ∈ \{0, p\}, stay. Otherwise pick j uniformly in S and l uniformly in [p] \textbackslash{} S, propose S′ = S \textbackslash{} \{j\} ∪ \{l\}, accept with probability 1 ∧ γ\_φ(S′)/γ\_φ(S); else stay.
\item \textbf{Neighbour swap}: if S = ∅, stay. Otherwise pick j uniformly in S; if N(j) \textbackslash{} S = ∅, stay; else pick l uniformly in N(j) \textbackslash{} S, propose S′ = S \textbackslash{} \{j\} ∪ \{l\}, accept with probability 1 ∧ ([γ\_φ(S′)/γ\_φ(S)] · |N(j) \textbackslash{} S| / |N(l) \textbackslash{} S′|), where this ratio is read as 0 when j ∉ N(l); else stay.
\item \textbf{Independence}: draw S′ \textasciitilde{} π₀ independently of S; accept with probability 1 ∧ e\textasciicircum{}\{φ[d(S′) − d(S)]\}; else stay.
\end{mdbul}

\textbf{Round.} m = p + B + 2 labelled steps: Gibbs 1, …, Gibbs p, one uniform swap, B neighbour swaps, one independence step; their kernels K\_1, …, K\_m (as matrices on Ω). For a permutation σ of [m] (the realised order), K\_σ := K\_\{σ(1)\} K\_\{σ(2)\} ⋯ K\_\{σ(m)\} (step σ(1) is applied first). K\_φ := (1/m!) Σ\_σ K\_σ (order uniform, independent of the state).

\textbf{Analytic objects.} L²(q\_φ): functions on Ω with ⟨f, g⟩ = Σ q\_φ f g. A kernel P is q-reversible if q(x)P(x,y) = q(y)P(y,x) for all x, y; P is self-adjoint on L²(q) if ⟨Pf, g⟩ = ⟨f, Pg⟩ for all f, g. Π\_φ := the kernel Π\_φ(x, ·) = q\_φ(·). Absolute L² gap of P: gap*(P) := 1 − ‖P − Π\_φ‖\_\{L²(q\_φ)→L²(q\_φ)\}. ‖μ − ν‖\_TV = sup\_A |μ(A) − ν(A)|. ε\_φ := min\_\{S∈Ω\} π₀(S)/q\_φ(S). A \textbf{switch} means argmax\_Ω q₀ ≠ argmax\_Ω q₁ (the bridge's mode changes between φ = 0 and φ = 1). The \textbf{discrepancy range} is max\_Ω d − min\_Ω d.

\textbf{Valid component input:} any (n, p, X, y) as above together with the method's blocks, tables and neighbour graph built from it. The kernels are component objects: they are defined for every such input, whether or not the method's exact-enumeration branch would be taken for that (n, p).

\mdhead{2}{2. Claims (for every valid component input and every φ ∈ [0, 1], unless an existence claim says otherwise)}

\textbf{(i)} Σ\_\{k≤K\} (w(k)/W\_K) c(k) = Z̄₀ exactly. Under EX-A and EX-S, the backward sampler is well defined (every conditional probability it uses has a positive denominator and sums to 1) and returns S with probability exactly π₀(S); successive draws with fresh streams are iid π₀.

\textbf{(ii)} Each elementary kernel is a Markov kernel on Ω that is q\_φ-reversible, including the boundary cases: the Gibbs addition at |S| = K has probability 0; the uniform swap self-loops at |S| ∈ \{0, p\}; the neighbour-swap ratio read as 0 when the reverse insertion is not a neighbour (j ∉ N(l)); and the stated Hastings ratios equal γ\_φ(S′)q(S′, S) / (γ\_φ(S)q(S, S′)) with q the actual proposal probabilities (so they supply the reverse proposal probability). For the neighbour swap this holds for any family N(j) ⊆ [p] \textbackslash{} \{j\}; for the method's symmetric graph the "read as 0" case never arises.

\textbf{(iii)} K\_φ is a Markov kernel, q\_φ-invariant, and self-adjoint on L²(q\_φ) (equivalently q\_φ-reversible). \emph{Original wording: "An individual realised order need not be reversible."} Reading: for every σ, K\_σ is q\_φ-invariant, and K\_σ is q\_φ-reversible if and only if K\_σ = K\_\{σ̄\}, σ̄ the reversed order; and there exist a valid component input, φ ∈ [0,1] and an order σ such that K\_σ is not q\_φ-reversible.

\textbf{(iv)} For every σ and every S ∈ Ω: K\_σ(S, ·) ≥ ε\_φ q\_φ(·) (as measures), with ε\_φ = (Z̄\_φ/Z̄₀) e\textasciicircum{}\{−φ max\_Ω d\} ∈ (0, 1]. Hence K\_φ(S, ·) ≥ ε\_φ q\_φ(·), sup\_S ‖K\_φ\textasciicircum{}h(S, ·) − q\_φ‖\_TV ≤ (1 − ε\_φ)\textasciicircum{}h for every integer h ≥ 0, and gap*(K\_φ) ≥ ε\_φ.

\textbf{(v)} ε\_φ ≤ π₀(S)/q\_φ(S) for every S ∈ Ω. \emph{Original wording: "An exponentially small available minorisation constant does not establish exponentially slow mixing of the hybrid kernel; neither a switch nor a large discrepancy range alone supplies such a lower bound."} Reading: there exist constants c > 0, n₀ and, for every n ≥ n₀, a valid component input (with that n) such that at φ = 1 simultaneously: ε₁ ≤ e\textasciicircum{}\{−cn\}; the discrepancy range is ≥ cn; there is a switch; and gap*(K₁) ≥ 1/40 and sup\_S ‖K₁\textasciicircum{}h(S, ·) − q₁‖\_TV ≤ (39/40)\textasciicircum{}h for all h. Consequently, writing t\_mix(K) := min\{h ≥ 0 : sup\_S ‖K\textasciicircum{}h(S, ·) − q\_φ‖\_TV ≤ 1/4\}, the family has t\_mix(K₁) ≤ 55 for every n ≥ n₀; so no upper bound on gap*(K\_φ) that tends to 0, and no lower bound on t\_mix(K\_φ) that tends to infinity (in particular none exponentially large in n), can be derived from ε\_φ being exponentially small, from a switch, or from a large discrepancy range, alone or together.
\end{mdblock}

\mdsec{(c) Final proof, verbatim: body of PROOF\_v3.tex}
\noindent{\small\itshape The proof below is the body of \texttt{PROOF\_v3.tex} from \texttt{\textbackslash begin\{document\}} to \texttt{\textbackslash end\{document\}} without the title block. Its preamble definitions are local to this chapter; labels carry the prefix \texttt{T09:}.}\par\medskip
\mdproofbegin
% ---- preamble of PROOF_v3.tex (theorem environments carry the prefix T09) ----
\newtheorem{T09theorem}{Theorem}
\newtheorem{T09lemma}[T09theorem]{Lemma}
\newtheorem{T09proposition}[T09theorem]{Proposition}
\theoremstyle{remark}
\newtheorem{T09remark}[T09theorem]{Remark}
\newcommand{\BF}{\mathrm{BF}}
\newcommand{\Om}{\Omega}
\newcommand{\eps}{\varepsilon}
\newcommand{\Zb}{\bar Z}
\newcommand{\gs}{\mathrm{gap}^*}
\newcommand{\TV}{\mathrm{TV}}
% ---- body of PROOF_v3.tex ----
All notation is that of STATEMENT.md, which is the statement proved here. Throughout, EX-A (exact arithmetic) and EX-S (ideal independent uniform streams) are in force. Kernels on the finite set $\Om$ are identified with nonnegative $\Om\times\Om$ matrices; a measure is a row vector; $Pf(x)=\sum_y P(x,y)f(y)$ and $\mu P(y)=\sum_x\mu(x)P(x,y)$; $\mathbf 1$ is the constant function~$1$. For $P$ stochastic and $q$ a probability on $\Om$, $P$ is $q$-invariant if $qP=q$.

\section{Preliminaries}

\begin{T09lemma}[Positivity]\label{T09:lem:pos}
For every $S\subseteq[p]$: $Q_S\ge \|y\|^2/(n+1)>0$, $D_S\in[1,\infty)$ is finite, $A_S=2b_0+Q_S>0$ and $\BF(S)\in(0,\infty)$. Consequently $\gamma_\phi(S)\in(0,\infty)$ for all $S\in\Om$ and all real $\phi$, $\Zb_\phi\in(0,\infty)$, and $q_\phi(S)>0$ for every $S\in\Om$.
\end{T09lemma}
\begin{proof}
Let $k=|S|$. If $k=0$, $Q_\emptyset=\|y\|^2$ and $D_\emptyset=1$. Let $k\ge1$ and write the spectral decomposition $G_S=\sum_{i=1}^k s_i b_ib_i^\top$ with $s_i\ge0$ and $(b_i)$ an orthonormal basis of $\mathbb R^k$. Then $H_S=\sum_i((n+1)s_i+1)b_ib_i^\top\succ0$ and $G_S+I=\sum_i(s_i+1)b_ib_i^\top\succ0$, so $D_S=\prod_i\frac{(n+1)s_i+1}{s_i+1}\in[1,\infty)$. For $s_i>0$ put $a_i=X_Sb_i/\sqrt{s_i}$; the $a_i$ are orthonormal, $X_Sb_i=0$ when $s_i=0$, and
$X_SH_S^{-1}X_S^\top=\sum_{i:s_i>0}\frac{s_i}{(n+1)s_i+1}\,a_ia_i^\top$, whose eigenvalues are $<1/(n+1)$. Hence $n z_S^\top H_S^{-1}z_S=n\,y^\top X_SH_S^{-1}X_S^\top y\le \frac{n}{n+1}\|y\|^2$ and $Q_S\ge\|y\|^2/(n+1)=(n-1)/(n+1)>0$. The remaining claims follow from $b_0>0$, the definitions, $\varpi(S)>0$ on $\Om$, and $|\Om|<\infty$.
\end{proof}

\begin{T09lemma}[Metropolis--Hastings form]\label{T09:lem:mh}
Let $q>0$ be a probability on $\Om$, $\mathcal Q$ a proposal kernel on $\Om$, and $P(x,y)=\mathcal Q(x,y)\alpha(x,y)$ for $y\ne x$, $P(x,x)=1-\sum_{y\ne x}P(x,y)$, with $\alpha(x,y)\in[0,1]$. If, for all $x\neq y$ with $\mathcal Q(x,y)>0$,
$\alpha(x,y)=1\wedge r(x,y)$ with $r(x,y)=\dfrac{q(y)\mathcal Q(y,x)}{q(x)\mathcal Q(x,y)}$, then $P$ is a Markov kernel and $q(x)P(x,y)=q(y)P(y,x)$ for all $x,y$.
\end{T09lemma}
\begin{proof}
$P\ge0$ and rows sum to one since $\sum_{y\ne x}\mathcal Q(x,y)\alpha(x,y)\le1$. For $x\ne y$: if $\mathcal Q(x,y)>0$ and $\mathcal Q(y,x)>0$ then $q(x)P(x,y)=\min\{q(x)\mathcal Q(x,y),q(y)\mathcal Q(y,x)\}$, which is symmetric. If exactly one of them is $0$, say $\mathcal Q(y,x)=0<\mathcal Q(x,y)$, then $r(x,y)=0$, so $P(x,y)=0=P(y,x)$. If both are $0$, both sides are $0$. The case $x=y$ is trivial.
\end{proof}

\begin{T09lemma}[Adjoints]\label{T09:lem:adj}
Let $q>0$ be a probability on $\Om$. For a kernel $P$ define $P^*(x,y)=q(y)P(y,x)/q(x)$. Then $\langle Pf,g\rangle_{q}=\langle f,P^*g\rangle_q$ for all $f,g$; $(PR)^*=R^*P^*$; $P$ is self-adjoint on $L^2(q)$ iff $P=P^*$ iff $P$ is $q$-reversible. If $P$ is $q$-reversible and stochastic, then $qP=q$.
\end{T09lemma}
\begin{proof}
$\langle Pf,g\rangle_q=\sum_{x,y}q(x)P(x,y)f(y)g(x)=\sum_{y}q(y)f(y)\sum_x P^*(y,x)g(x)$. The adjoint is unique (the form is nondegenerate because $q>0$), which gives the product rule and the equivalences; $P=P^*$ is the detailed-balance identity. Finally $\sum_xq(x)P(x,y)=\sum_x q(y)P(y,x)=q(y)$.
\end{proof}

\begin{T09lemma}[Doeblin contraction]\label{T09:lem:doeb}
Let $P$ be a stochastic kernel on $\Om$ with $qP=q$, $q>0$ a probability, and suppose $P(x,\cdot)\ge c\,\nu(\cdot)$ for every $x$, for a probability $\nu$ and a constant $c\in(0,1]$. Then:
(a) $\|\mu P\|_1\le(1-c)\|\mu\|_1$ for every signed measure $\mu$ with $\mu(\Om)=0$;
(b) $\sup_x\|P^h(x,\cdot)-q\|_{\TV}\le(1-c)^h$ for every integer $h\ge0$;
(c) if moreover $P$ is self-adjoint on $L^2(q)$, then $\gs(P):=1-\|P-\Pi\|_{L^2(q)\to L^2(q)}\ge c$, where $\Pi(x,\cdot)=q(\cdot)$.
\end{T09lemma}
\begin{proof}
(a) If $c=1$ then $P=\mathbf 1\nu$ and $\mu P=\mu(\Om)\nu=0$. If $c<1$, $R:=(P-c\mathbf 1\nu)/(1-c)$ is nonnegative with unit row sums, and $\mu P=c\mu(\Om)\nu+(1-c)\mu R=(1-c)\mu R$; also $\|\mu R\|_1\le\sum_{x,y}|\mu(x)|R(x,y)=\|\mu\|_1$.
(b) Since $qP^h=q$, $P^h(x,\cdot)-q=(\delta_x-q)P^h$, and $(\delta_x-q)(\Om)=0$. By (a) iterated, $\|P^h(x,\cdot)-q\|_\TV=\tfrac12\|(\delta_x-q)P^h\|_1\le(1-c)^h\tfrac12\|\delta_x-q\|_1=(1-c)^h(1-q(x))\le(1-c)^h$.
(c) Let $L^2_0=\{f:\sum_xq(x)f(x)=0\}$. Since $\sum_x q(x)(Pf)(x)=(qP)f=qf$, $P$ maps $L^2_0$ into itself, and the restriction is self-adjoint; by the spectral theorem (finite dimension) $\|P|_{L^2_0}\|=\max|\lambda|$ over its eigenvalues. Let $Pf=\lambda f$ with $0\ne f\in L^2_0$ and $\mu(x):=q(x)f(x)$. Then $\mu\ne0$, $\mu(\Om)=0$, and by reversibility (Lemma~\ref{T09:lem:adj}) $\mu P(y)=\sum_x q(y)P(y,x)f(x)=q(y)\lambda f(y)=\lambda\mu(y)$. By (a), $|\lambda|\le1-c$. Every $f\in L^2(q)$ splits orthogonally as $f=a\mathbf 1+g$ with $g\in L^2_0$, and $(P-\Pi)f=Pg$ because $P\mathbf 1=\Pi\mathbf 1=\mathbf 1$ and $\Pi g=0$. Hence $\|(P-\Pi)f\|\le(1-c)\|g\|\le(1-c)\|f\|$.
\end{proof}

\section{Part (i): exact normaliser and exact sampler}

Put $e_b(j)=0$ for $j\notin\{0,\dots,|C_b|\}$, $c^{(b)}(t)=0$ for $t<0$, and $P_b=\sum_{i\le b}|C_i|$. These are the conventions under which the recursion $c^{(b)}(t)=\sum_je_b(j)c^{(b-1)}(t-j)$ and the draw of $j_b$ in STATEMENT.md are read: $j$ ranges over all integers, and only $0\le j\le|C_b|$ with $t-j\ge0$ contribute.

\begin{T09lemma}\label{T09:lem:conv}
For $b=0,\dots,J$ and every integer $t$, $c^{(b)}(t)=\sum\prod_{i\le b}r_i(U_i)$, the sum running over tuples $(U_1,\dots,U_b)$ with $U_i\subseteq C_i$ and $\sum_i|U_i|=t$. In particular $c^{(b)}(t)>0$ iff $0\le t\le P_b$.
\end{T09lemma}
\begin{proof}
Induction on $b$. For $b=0$ the only tuple is the empty one, with $t=0$ and product $1$. For the step, group the tuples by $j=|U_b|$: $\sum_{j}\big(\sum_{|U_b|=j}r_b(U_b)\big)c^{(b-1)}(t-j)=\sum_je_b(j)c^{(b-1)}(t-j)$. Positivity: all $r_i>0$ (Lemma~\ref{T09:lem:pos}) and a tuple with $\sum|U_i|=t$ exists iff $0\le t\le P_b$.
\end{proof}

\begin{T09proposition}[Part (i)]
$\sum_{k\le K}\frac{w(k)}{W_K}c(k)=\Zb_0$. The backward sampler is well defined and returns $S$ with probability $\pi_0(S)$ for every $S\in\Om$; draws with fresh streams are iid $\pi_0$.
\end{T09proposition}
\begin{proof}
Because $C_1,\dots,C_J$ partition $[p]$, the map $S\mapsto(S\cap C_1,\dots,S\cap C_J)$ is a bijection from subsets of $[p]$ onto tuples with $U_b\subseteq C_b$, and $|S|=\sum_b|S\cap C_b|$. Since $\varpi(S)=w(|S|)/W_K$ depends on $S$ only through $|S|$,
\[\Zb_0=\sum_{k=0}^{K}\frac{w(k)}{W_K}\sum_{|S|=k}\prod_br_b(S\cap C_b)=\sum_{k=0}^K\frac{w(k)}{W_K}c(k)\]
by Lemma~\ref{T09:lem:conv} with $b=J$. (The values $c^{(b)}(t)$, $t\le K$, depend only on $c^{(b-1)}(t')$ with $t'\le t$, so computing every convolution only up to degree $K$ is exact.)

Well-definedness. Step (1): the denominator $\sum_{k'\le K}w(k')c(k')\ge w(0)c(0)=\tfrac12\prod_br_b(\emptyset)=\tfrac12>0$, so the law of $k$ is a probability, and $k$ is drawn only where $c(k)>0$. Step (2): we show by downward induction that $c^{(b)}(t_b)>0$ whenever $t_b$ is reached with positive probability. This holds for $b=J$. If $c^{(b)}(t_b)>0$, the numbers $e_b(j)c^{(b-1)}(t_b-j)/c^{(b)}(t_b)$ are nonnegative and sum to $1$ over $j$ by the recursion, and $j_b$ has positive probability only if $c^{(b-1)}(t_b-j_b)>0$, i.e.\ $c^{(b-1)}(t_{b-1})>0$; also $e_b(j_b)>0$, so $0\le j_b\le|C_b|$. At the end $c^{(0)}(t_0)>0$ forces $t_0=0$, so $\sum_bj_b=k$. Step (3): $e_b(j_b)>0$ and $\{U\subseteq C_b:|U|=j_b\}$ has total weight $e_b(j_b)$, so the law of $U_b$ is a probability. The output $S=\bigcup_bU_b$ has $|S|=k\le K$, so $S\in\Om$.

Exactness. Fix $S\in\Om$. The only run producing $S$ is $k=|S|$, $j_b=|S\cap C_b|$, $t_b=\sum_{i\le b}j_i$, $U_b=S\cap C_b$; along it all quantities in the denominators are positive (Lemma~\ref{T09:lem:conv}, since $0\le t_b\le P_b$). Its probability is
\[\frac{w(k)c(k)}{\sum_{k'}w(k')c(k')}\prod_{b=1}^J\frac{e_b(j_b)c^{(b-1)}(t_{b-1})}{c^{(b)}(t_b)}\prod_{b=1}^J\frac{r_b(U_b)}{e_b(j_b)}
=\frac{w(k)c(k)}{\sum_{k'}w(k')c(k')}\cdot\frac{c^{(0)}(0)}{c^{(J)}(k)}\prod_br_b(U_b)=\frac{\gamma_0(S)W_K}{W_K\Zb_0}=\pi_0(S),\]
the middle product telescoping. Under EX-A each categorical draw from a finite law with an ideal uniform is exact; under EX-S the uniforms of different draws are independent, and each output is a fixed measurable function of its own uniforms with law $\pi_0$; hence the draws are iid $\pi_0$.
\end{proof}

\section{Part (ii): reversibility of every elementary kernel}

Fix $\phi$ and write $\gamma=\gamma_\phi$, $q=q_\phi$; recall $\gamma(S)=0$ for $S\notin\Om$ and $\gamma>0$ on $\Om$ (Lemma~\ref{T09:lem:pos}). For $S\ne S'$ of equal size with $|S\triangle S'|=2$ write $j=j(S,S')$ for the element of $S\setminus S'$ and $l=l(S,S')$ for the element of $S'\setminus S$; the pair $(j,l)$ determines $S'=S\setminus\{j\}\cup\{l\}$ and conversely.

\paragraph{Gibbs $j$.} For $S\in\Om$, $S^-=S\setminus\{j\}\in\Om$, so the denominator $\gamma(S^+)+\gamma(S^-)\ge\gamma(S^-)>0$ and $K_{G,j}(S,\cdot)$ is a probability on $\{S^+,S^-\}\cap\Om$ (if $S^+\notin\Om$ its probability is $\gamma(S^+)/(\cdot)=0$). This is the boundary case: $S^+\notin\Om$ iff $j\notin S$ and $|S|=K$, and then the addition has probability $0$. Reversibility: let $T\ne S$ with $K_{G,j}(S,T)>0$. Then $T$ is the element of $\{S^+,S^-\}$ other than $S$, so $T\in\Om$ and $\{T^+,T^-\}=\{S^+,S^-\}$, and
$q(S)K_{G,j}(S,T)=\gamma(S)\gamma(T)/(\Zb_\phi(\gamma(S^+)+\gamma(S^-)))$ is symmetric in $(S,T)$; in particular $K_{G,j}(T,S)>0$ too. If $K_{G,j}(S,T)=0$ and $T\neq S$ then either $T\notin\{S^+,S^-\}$ (so $S\notin\{T^+,T^-\}$ and $K_{G,j}(T,S)=0$) or $T=S^+\notin\Om$, which is not a state. Hence detailed balance holds.

\paragraph{Uniform swap.} If $|S|\in\{0,p\}$ the kernel is $\delta_S$ (boundary self-loop). If $1\le|S|\le p-1$, the proposal is $\mathcal Q(S,S')=1/(|S|(p-|S|))$ for each $S'$ with $|S'|=|S|$, $|S\triangle S'|=2$ (distinct pairs $(j,l)$ give distinct $S'$), and $S'\in\Om$. All proposals preserve size, so states of size $0$ or $p$ (each the unique set of its size) are never proposed from other states. For $S'$ proposed from $S$, $S$ is proposed from $S'$ (delete $l$, add $j$) and $\mathcal Q(S',S)=1/(|S'|(p-|S'|))=\mathcal Q(S,S')$, so the Hastings ratio $\gamma(S')\mathcal Q(S',S)/(\gamma(S)\mathcal Q(S,S'))$ equals the stated $\gamma(S')/\gamma(S)$. Lemma~\ref{T09:lem:mh} applies (with $\mathcal Q(S,\cdot)=\delta_S$ on the boundary states).

\paragraph{Neighbour swap (arbitrary $N(j)\subseteq[p]\setminus\{j\}$).} If $S=\emptyset$ the kernel is $\delta_\emptyset$; all moves preserve size, so $\emptyset$ is never proposed from another state. For $S\ne\emptyset$ and $S'\ne S$, the move $S\to S'$ is proposed only by the pair $(j,l)=(j(S,S'),l(S,S'))$, so
\[\mathcal Q(S,S')=\begin{cases}\dfrac{1}{|S|\,|N(j)\setminus S|}&\text{if } l\in N(j),\\[1ex] 0&\text{otherwise,}\end{cases}\qquad
\mathcal Q(S',S)=\begin{cases}\dfrac{1}{|S'|\,|N(l)\setminus S'|}&\text{if } j\in N(l),\\[1ex] 0&\text{otherwise,}\end{cases}\]
(here $j\in S$, $l\notin S$ by definition of $j(S,S')$, $l(S,S')$). The second formula holds because the reverse move deletes $l\in S'$ and must insert $j\notin S'$ from $N(l)\setminus S'$; when $j\in N(l)$ the set $N(l)\setminus S'$ contains $j$ and is nonempty, and when $l\in N(j)$ the set $N(j)\setminus S$ contains $l$, so no denominator vanishes. (Proposals with $N(j)\setminus S=\emptyset$ are self-loops and do not affect Lemma~\ref{T09:lem:mh}.) Let $\mathcal Q(S,S')>0$, so $|S'|=|S|$. If $j\in N(l)$, the Hastings ratio is
\[\frac{\gamma(S')\mathcal Q(S',S)}{\gamma(S)\mathcal Q(S,S')}=\frac{\gamma(S')}{\gamma(S)}\cdot\frac{|S|\,|N(j)\setminus S|}{|S'|\,|N(l)\setminus S'|}=\frac{\gamma(S')}{\gamma(S)}\cdot\frac{|N(j)\setminus S|}{|N(l)\setminus S'|},\]
the stated ratio. If $j\notin N(l)$, the Hastings ratio is $0$, which is the stated ratio under the reading convention. So the stated acceptance is $1\wedge r$ with the true Hastings ratio $r$ in every case, and Lemma~\ref{T09:lem:mh} gives reversibility. For the method's graph, $l\in N(j)\Rightarrow j\in N(l)$, so the convention never applies.

\paragraph{Independence.} $\mathcal Q(S,S')=\pi_0(S')>0$ for all $S,S'\in\Om$. Since $q(S)/\pi_0(S)=(\Zb_0/\Zb_\phi)e^{\phi d(S)}$,
$\dfrac{q(S')\pi_0(S)}{q(S)\pi_0(S')}=e^{\phi[d(S')-d(S)]}$, the stated ratio; Lemma~\ref{T09:lem:mh} applies (a proposal $S'=S$ is a self-loop).

In each case the stated Hastings ratio equals $\gamma(S')\mathcal Q(S',S)/(\gamma(S)\mathcal Q(S,S'))$ with the actual proposal probabilities, i.e.\ it supplies the reverse proposal probability. This proves (ii).

\section{Part (iii): the averaged round}

By (ii) and Lemma~\ref{T09:lem:adj}, each $K_e$ is stochastic, $q$-invariant and self-adjoint, $K_e^*=K_e$. Hence every $K_\sigma$ is stochastic and $qK_\sigma=q$, and so is the average $K_\phi$. By the product rule, $K_\sigma^*=K_{\sigma(m)}^*\cdots K_{\sigma(1)}^*=K_{\sigma(m)}\cdots K_{\sigma(1)}=K_{\bar\sigma}$, $\bar\sigma(i)=\sigma(m+1-i)$. Since $\sigma\mapsto\bar\sigma$ is a bijection of the symmetric group,
$K_\phi^*=\frac1{m!}\sum_\sigma K_{\bar\sigma}=K_\phi$: $K_\phi$ is self-adjoint on $L^2(q_\phi)$, equivalently $q_\phi$-reversible (Lemma~\ref{T09:lem:adj}). Also $K_\sigma$ is $q$-reversible iff $K_\sigma=K_\sigma^*=K_{\bar\sigma}$.

That some realised order is not reversible is shown in Proposition~\ref{T09:prop:nonrev} below, on the family of Section~\ref{T09:sec:family}.

\section{Part (iv): minorisation}

\begin{T09lemma}\label{T09:lem:eps}
$\pi_0(S)/q_\phi(S)=(\Zb_\phi/\Zb_0)e^{-\phi d(S)}$, so for $\phi\ge0$, $\eps_\phi=\min_\Om\pi_0/q_\phi=(\Zb_\phi/\Zb_0)e^{-\phi\max_\Om d}$, and $0<\eps_\phi\le1$.
\end{T09lemma}
\begin{proof}
The identity is $\pi_0=\gamma_0/\Zb_0$ and $q_\phi=\gamma_0e^{\phi d}/\Zb_\phi$; for $\phi\ge0$, $e^{-\phi d}$ is minimal where $d$ is maximal. Positivity by Lemma~\ref{T09:lem:pos}. If $\pi_0(S)>q_\phi(S)$ for all $S$, summing gives $1>1$; so $\eps_\phi\le1$.
\end{proof}

\begin{T09proposition}[Part (iv)]
For every $\sigma$ and $x\in\Om$, $K_\sigma(x,\cdot)\ge\eps_\phi q_\phi(\cdot)$. Hence $K_\phi(x,\cdot)\ge\eps_\phi q_\phi(\cdot)$, $\sup_x\|K_\phi^h(x,\cdot)-q_\phi\|_\TV\le(1-\eps_\phi)^h$ for all $h\ge0$, and $\gs(K_\phi)\ge\eps_\phi$.
\end{T09proposition}
\begin{proof}
Write $\eps=\eps_\phi$, $q=q_\phi$, $K_{\rm ind}$ for the independence kernel. For $u\ne v$,
$K_{\rm ind}(u,v)=\min\{\pi_0(v),\ \pi_0(v)e^{\phi[d(v)-d(u)]}\}$. Here $\pi_0(v)\ge\eps q(v)$, and by Lemma~\ref{T09:lem:eps} $\pi_0(v)e^{\phi[d(v)-d(u)]}=q(v)\,\pi_0(u)/q(u)\ge\eps q(v)$. For $v=u$, $K_{\rm ind}(u,u)\ge\pi_0(u)\ge\eps q(u)$ (the proposal $u$ is a self-loop). So $K_{\rm ind}(u,\cdot)\ge\eps q(\cdot)$ for every $u$.

Let the independence step be at position $i$ of $\sigma$, $A=K_{\sigma(1)}\cdots K_{\sigma(i-1)}$ (identity if $i=1$) and $C=K_{\sigma(i+1)}\cdots K_{\sigma(m)}$ (identity if $i=m$). $A$ is stochastic, $C\ge0$ and $qC=q$ (part (iii)). Then
\[K_\sigma(x,y)=\sum_{u,v}A(x,u)K_{\rm ind}(u,v)C(v,y)\ge\eps\sum_uA(x,u)\sum_vq(v)C(v,y)=\eps q(y).\]
Averaging over $\sigma$ gives $K_\phi(x,\cdot)\ge\eps q$. Since $K_\phi$ is stochastic, $q$-invariant and self-adjoint (part (iii)), Lemma~\ref{T09:lem:doeb} with $\nu=q$, $c=\eps\in(0,1]$ gives the TV bound and $\gs(K_\phi)\ge\eps$.
\end{proof}

\begin{T09remark}
The same argument gives the TV bound for any product of $h$ realised orders (each factor is $q$-invariant and $\ge\eps q$), not only for the averaged kernel.
\end{T09remark}

\section{Part (v), and the non-reversible order of part (iii)}\label{T09:sec:family}

The first sentence of (v) is the definition of $\eps_\phi$ as a minimum. For the second, the asymptotic parameter is the sample size $n$ (with $p=2$ fixed): ``exponentially small'' and ``large range'' refer to $n$; nothing is claimed about growth in $p$.

\paragraph{Notation for this section.} The family has two predictors, and its states are written $\emptyset,\{1\},\{2\},\{1,2\}$; subscripts $1$, $2$, $12$ on $Q$, $D$, $A$, $t$, $R$, $\kappa$ refer to these states. The correlation parameter of the family is $\tau$ and the signal parameter $\rho$; the rate constant of the statement is $\kappa$. Throughout, $q_1=\pi$ is the target at $\phi=1$.

\paragraph{The family.} Fix $\tau\in(0,1)$ and $\rho\in(0,1)$ (for instance $\tau=\tfrac12$, $\rho=\tfrac35$). For $n\ge4$ the space $\mathbf 1^\perp\subset\mathbb R^n$ has dimension $n-1\ge3$; fix orthonormal $h_1,h_2,h_3\in\mathbf 1^\perp$ and put $s=\sqrt{n-1}$,
\[x_1=s h_1,\qquad x_2=s(\tau h_1+\sqrt{1-\tau^2}\,h_2),\qquad y=s(\rho h_1+\sqrt{1-\rho^2}\,h_3).\]
Columns and $y$ are centred with squared norm $n-1$, so this is standardised data in the sense of STATEMENT.md, Section~1. Here $p=2$, $K=\min(2,\lfloor n/2\rfloor)=2$, $\Om=\{\emptyset,\{1\},\{2\},\{1,2\}\}$, and $B=\min(\lfloor\sqrt n\rfloor,1,16)=1$. Since every block has $|C_b|\le B=1$ and the blocks partition $\{1,2\}$, the blocks are $\{1\}$ and $\{2\}$ (in some order; $\gamma_0$ depends only on the partition). $m=p+B+2=5$. The neighbour graph has $N(j)\subseteq\{1,2\}\setminus\{j\}$; no property of it beyond this is used. The prior: $w(0)=\tfrac12$, $w(1)=w(2)=\tfrac16$, so $w(1)/w(0)=\tfrac13$ and $w(2)/w(1)=1$.

\paragraph{Scores.} $x_1^\top x_2=(n-1)\tau$, $z_1=x_1^\top y=(n-1)\rho$, $z_2=(n-1)\tau\rho$. For a singleton, $G=n-1$, $H=(n+1)(n-1)+1=n^2$, $D=n$, $Q=(n-1)-z^2/n$. So
\[Q_\emptyset=n-1,\quad Q_1=(n-1)\Big(1-\rho^2\tfrac{n-1}{n}\Big),\quad Q_2=(n-1)\Big(1-\tau^2\rho^2\tfrac{n-1}{n}\Big),\quad D_1=D_2=n.\]
For $\{1,2\}$, $G=(n-1)R$, $R=\begin{pmatrix}1&\tau\\\tau&1\end{pmatrix}$, with eigenvectors $g_\pm=(1,\pm1)/\sqrt2$ and eigenvalues $1\pm\tau$. Then $H$ has eigenvalues $\alpha_\pm=(n^2-1)(1\pm\tau)+1$ and $G+I$ has $\beta_\pm=(n-1)(1\pm\tau)+1$; since $\alpha_\pm=(n+1)\beta_\pm-n$, $D_{12}=f(1+\tau)f(1-\tau)$ with $f(x)=n+1-\frac{n}{(n-1)x+1}$, increasing in $x>0$ with $f\le n+1$. With $z=(n-1)\rho(1,\tau)^\top$, $z^\top g_\pm=(n-1)\rho(1\pm\tau)/\sqrt2$ and, using $(1\pm\tau)^2/\alpha_\pm=\frac{1\pm\tau}{n^2-1}(1-1/\alpha_\pm)$,
\[t_{12}:=n z^\top H^{-1}z=\frac{n(n-1)^2\rho^2}{2}\sum_\pm\frac{(1\pm\tau)^2}{\alpha_\pm}=\frac{n(n-1)\rho^2}{n+1}\Big(1-\frac\delta2\Big),\qquad \delta=\sum_\pm\frac{1\pm\tau}{\alpha_\pm}\in\Big(0,\frac2{n^2-1}\Big),\]
and $Q_{12}=(n-1)-t_{12}$.

\begin{T09lemma}\label{T09:lem:bounds}
(B1) $|Q_{12}-Q_1|<1/n$. (B2) $Q_{12}>(n-1)/(n+1)\ge\tfrac12$. (B3) If $(n-1)(1-\tau)\ge1$ then $n^2/4\le D_{12}\le(n+1)^2$.
\end{T09lemma}
\begin{proof}
(B1) $t_{12}-\frac{(n-1)^2\rho^2}{n}=(n-1)\rho^2\big[\frac n{n+1}-\frac{n-1}n\big]-\frac{n(n-1)\rho^2\delta}{2(n+1)}$. The first term equals $\frac{(n-1)\rho^2}{n(n+1)}\in[0,1/n)$; the second lies in $[0,\frac{n(n-1)\rho^2}{(n+1)(n^2-1)})=[0,\frac{n\rho^2}{(n+1)^2})\subset[0,1/n)$. A difference of two numbers in $[0,1/n)$ has absolute value $<1/n$. (B2) $t_{12}<\frac{n(n-1)}{n+1}$, so $Q_{12}>\frac{n-1}{n+1}$, which is $\ge\frac12$ for $n\ge3$. (B3) $f\le n+1$ gives the upper bound. If $(n-1)x\ge1$ then $f(x)\ge n+1-n/2\ge n/2$; $f$ increasing and $(n-1)(1-\tau)\ge1$ give $f(1\pm\tau)\ge n/2$.
\end{proof}

\paragraph{Ratios at $\phi=1$.} $\gamma_1=\varpi\BF$ and $\BF(S)=D_S^{-1/2}(A_\emptyset/A_S)^\nu$, $A_S=0.02+Q_S$, $\nu=0.01+n/2$. Define
\[R_1=\frac{\gamma_1(\{1\})}{\gamma_1(\emptyset)}=\frac{\BF(\{1\})}3,\quad R_2=\frac{\gamma_1(\{2\})}{\gamma_1(\emptyset)}=\frac{\BF(\{2\})}3,\]
\[R_{12|1}=\frac{\gamma_1(\{1,2\})}{\gamma_1(\{1\})}=\Big(\frac{D_1}{D_{12}}\Big)^{1/2}\Big(\frac{A_1}{A_{12}}\Big)^\nu,\qquad R_{12|2}=\frac{\gamma_1(\{1,2\})}{\gamma_1(\{2\})}=\Big(\frac{D_2}{D_{12}}\Big)^{1/2}\Big(\frac{A_2}{A_{12}}\Big)^\nu,\]
using $w(2)/w(1)=1$. Let $\kappa_1=\tfrac12\log\frac1{1-\rho^2}$ and $\kappa_2=\tfrac12\log\frac1{1-\tau^2\rho^2}$; then $\kappa_1>\kappa_2>0$ because $0<\tau<1$, $\rho>0$.

\begin{T09lemma}\label{T09:lem:asym}
As $n\to\infty$: $\frac1n\log R_1\to\kappa_1$, $\frac1n\log R_2\to\kappa_2$, $\frac1n\log R_{12|2}\to\kappa_1-\kappa_2>0$, $\frac1n\log\frac{\gamma_1(\{2\})}{\gamma_1(\{1\})}\to-(\kappa_1-\kappa_2)<0$, $\frac1n\log R_{12|1}\to0$, and $R_{12|1}\le 2e^{1.01}n^{-1/2}$ once $(n-1)(1-\tau)\ge1$.
\end{T09lemma}
\begin{proof}
$A_\emptyset/A_1\to1/(1-\rho^2)$, $A_\emptyset/A_2\to1/(1-\tau^2\rho^2)$ and, by (B1), $A_2/A_{12}\to(1-\tau^2\rho^2)/(1-\rho^2)$ and $A_1/A_2\to(1-\rho^2)/(1-\tau^2\rho^2)$, all limits in $(0,\infty)$. Since $\nu/n\to\frac12$, $\frac1n\log(\cdot)^\nu\to\frac12\log(\lim)$. The remaining factors are: $\frac13 D_1^{-1/2}=\frac13n^{-1/2}$ in $R_1$; $\frac13D_2^{-1/2}=\frac13n^{-1/2}$ in $R_2$; and, once $(n-1)(1-\tau)\ge1$, $(D_2/D_{12})^{1/2}\in[n^{1/2}/(n+1),\,2n^{-1/2}]$ in $R_{12|2}$ by (B3). Each lies in $[\frac1{3(n+1)},\,n^{1/2}]$, so contributes $O(\log n)=o(n)$ to the logarithm. This gives the first three limits. Next, $\gamma_1(\{2\})/\gamma_1(\{1\})=(A_1/A_2)^\nu$ since $D_1=D_2$ and $w(1)=w(1)$, which gives the fourth. Since $R_{12|1}=R_{12|2}\cdot\gamma_1(\{2\})/\gamma_1(\{1\})$, the fifth limit is $(\kappa_1-\kappa_2)-(\kappa_1-\kappa_2)=0$. Finally, by (B1)--(B2), $A_1/A_{12}\le1+(1/n)/A_{12}\le1+2/n$, so $(A_1/A_{12})^\nu\le e^{2\nu/n}=e^{1+0.02/n}\le e^{1.01}$; with (B3), $(D_1/D_{12})^{1/2}\le(4/n)^{1/2}$.
\end{proof}

\begin{T09lemma}[Start, target, discrepancy]\label{T09:lem:fam}
(a) $\pi_0(\{1,2\})\to1$, and for all large $n$, $\arg\max_\Om\pi_0=\{1,2\}$. (b) $q_1(\{1\})\to1$, and for all large $n$, $\arg\max_\Om q_1=\{1\}$. (c) $d(\emptyset)=d(\{1\})=d(\{2\})=0$ and $d(\{1,2\})=\log R_{12|1}-\log(3R_2)$, so $\frac1n d(\{1,2\})\to-\kappa_2$; moreover $-d(\{1,2\})\ge\log(3R_2)-\log(2e^{1.01}n^{-1/2})$ once $(n-1)(1-\tau)\ge1$. (d) $\eps_1\le 1/(3R_2\,q_1(\{1\}))$, so $\limsup\frac1n\log\eps_1\le-\kappa_2$.
\end{T09lemma}
\begin{proof}
(a) With blocks $\{1\},\{2\}$, $\gamma_0(S)=\varpi(S)\BF(\{1\})^{\mathbf 1[1\in S]}\BF(\{2\})^{\mathbf 1[2\in S]}$. So $\gamma_0(\{1\})/\gamma_0(\{1,2\})=1/\BF(\{2\})=1/(3R_2)$, $\gamma_0(\{2\})/\gamma_0(\{1,2\})=1/(3R_1)$, $\gamma_0(\emptyset)/\gamma_0(\{1,2\})=3/(\BF(\{1\})\BF(\{2\}))=1/(3R_1R_2)$; all $\to0$ by Lemma~\ref{T09:lem:asym}. (b) $\gamma_1(\emptyset)/\gamma_1(\{1\})=1/R_1$, $\gamma_1(\{1,2\})/\gamma_1(\{1\})=R_{12|1}$, and $\gamma_1(\{2\})/\gamma_1(\{1\})$; all $\to0$. (c) Each singleton or the empty set lies in one block, where the table is the exact $\BF$, so $d=0$ there; and $\BF(\{1,2\})/\BF(\{1\})=R_{12|1}\,w(1)/w(2)=R_{12|1}$, $\BF(\{2\})=3R_2$. The limit and the bound follow from Lemma~\ref{T09:lem:asym}. (d) $\eps_1\le\pi_0(\{1\})/q_1(\{1\})$ and $\pi_0(\{1\})\le\gamma_0(\{1\})/\gamma_0(\{1,2\})=1/(3R_2)$; then use (b) and Lemma~\ref{T09:lem:asym}.
\end{proof}

\begin{T09proposition}[Part (iii), second sentence]\label{T09:prop:nonrev}
In the family, at $\phi=1$ and for all large $n$, the realised order $\sigma=(\text{Gibbs }2,\ \text{Gibbs }1,\ \text{uniform swap},\ \text{neighbour swap},\ \text{independence})$ (Gibbs 2 applied first) gives a kernel $K_\sigma$ that is not $q_1$-reversible.
\end{T09proposition}
\begin{proof}
By part (iii) it suffices to show $K_\sigma\ne K_{\bar\sigma}$. Lower bound along one path:
$K_\sigma(\emptyset,\{1,2\})\ge K_{G,2}(\emptyset,\{2\})\,K_{G,1}(\{2\},\{1,2\})\,K_{\rm us}(\{1,2\},\{1,2\})\,K_{\rm ns}(\{1,2\},\{1,2\})\,K_{\rm ind}(\{1,2\},\{1,2\})$.
Here $K_{G,2}(\emptyset,\{2\})=R_2/(1+R_2)\to1$; $K_{G,1}(\{2\},\{1,2\})=R_{12|2}/(1+R_{12|2})\to1$; $K_{\rm us}(\{1,2\},\{1,2\})=1$ ($|\{1,2\}|=p$); $K_{\rm ns}(\{1,2\},\{1,2\})=1$ (for $j\in\{1,2\}$, $N(j)\setminus\{1,2\}=\emptyset$); $K_{\rm ind}(\{1,2\},\{1,2\})\ge\pi_0(\{1,2\})\to1$. So $K_\sigma(\emptyset,\{1,2\})\to1$.
For $\bar\sigma=(\text{independence},\text{neighbour swap},\text{uniform swap},\text{Gibbs }1,\text{Gibbs }2)$:
$K_{\bar\sigma}(\emptyset,\{1\})\ge K_{\rm ind}(\emptyset,\emptyset)K_{\rm ns}(\emptyset,\emptyset)K_{\rm us}(\emptyset,\emptyset)K_{G,1}(\emptyset,\{1\})K_{G,2}(\{1\},\{1\})$. From $\emptyset$ the independence step moves only if it proposes some $S'\ne\emptyset$ and accepts, with acceptance $1\wedge e^{d(S')}$; so $K_{\rm ind}(\emptyset,\emptyset)\ge1-\pi_0(\{1\})-\pi_0(\{2\})-\pi_0(\{1,2\})e^{d(\{1,2\})}\to1$ by Lemma~\ref{T09:lem:fam}(a),(c). $K_{\rm ns}(\emptyset,\emptyset)=K_{\rm us}(\emptyset,\emptyset)=1$; $K_{G,1}(\emptyset,\{1\})=R_1/(1+R_1)\to1$; $K_{G,2}(\{1\},\{1\})=1/(1+R_{12|1})\to1$. Hence $K_{\bar\sigma}(\emptyset,\{1\})\to1$ and $K_{\bar\sigma}(\emptyset,\{1,2\})\le1-K_{\bar\sigma}(\emptyset,\{1\})\to0$. For large $n$, $K_\sigma(\emptyset,\{1,2\})>\tfrac12>K_{\bar\sigma}(\emptyset,\{1,2\})$.
\end{proof}

\begin{T09proposition}[Part (v), second sentence]\label{T09:prop:v}
Let $\kappa=\kappa_2/2$. There is $n_0$ (depending on $\tau,\rho$) such that for every $n\ge n_0$ the family's input satisfies, at $\phi=1$: $\eps_1\le e^{-\kappa n}$; $\max_\Om d-\min_\Om d\ge\kappa n$; a switch ($\arg\max q_0=\{1,2\}\neq\{1\}=\arg\max q_1$); $K_1(x,\{1\})\ge\frac1{40}$ for every $x\in\Om$; hence $\sup_x\|K_1^h(x,\cdot)-q_1\|_\TV\le(39/40)^h$ for all $h\ge0$, $\gs(K_1)\ge\frac1{40}$, and $t_{\rm mix}(K_1)\le55$, where $t_{\rm mix}(K)=\min\{h\ge0:\sup_x\|K^h(x,\cdot)-q_1\|_\TV\le\frac14\}$.
\end{T09proposition}
\begin{proof}
The first three properties hold for large $n$ by Lemma~\ref{T09:lem:fam}: (d) gives $\frac1n\log\eps_1\le-\kappa_2+o(1)$; (c) gives $\max d-\min d\ge0-d(\{1,2\})\ge\log(3R_2)-\log(2e^{1.01}n^{-1/2})=\kappa_2n+o(n)$; (a),(b) give the switch.

Mixing. Consider the $3!=6$ orders whose last two steps are Gibbs 1 then Gibbs 2 (a fraction $6/120=1/20$ of all orders). For such $\sigma$ and any $x$, let $x'$ be the state before Gibbs 1. Gibbs 1 puts $1$ in with probability $R_1/(1+R_1)$ if $2\notin x'$ and $R_{12|2}/(1+R_{12|2})$ if $2\in x'$ (the conditional depends only on $x'\setminus\{1\}$); then, from a state containing $1$, Gibbs 2 removes or keeps $2$ out with probability $\gamma_1(\{1\})/(\gamma_1(\{1\})+\gamma_1(\{1,2\}))=1/(1+R_{12|1})$, ending at $\{1\}$. So $K_\sigma(x,\{1\})\ge\theta_n:=\min\{\frac{R_1}{1+R_1},\frac{R_{12|2}}{1+R_{12|2}}\}\cdot\frac1{1+R_{12|1}}$, and $K_1(x,\{1\})\ge\theta_n/20$, all other orders contributing $\ge0$. By Lemma~\ref{T09:lem:asym}, $\theta_n\to1$, so $\theta_n\ge\frac12$ for large $n$ and $K_1(x,\cdot)\ge\frac1{40}\delta_{\{1\}}(\cdot)$. $K_1$ is stochastic, $q_1$-invariant and self-adjoint (part (iii)); Lemma~\ref{T09:lem:doeb} with $\nu=\delta_{\{1\}}$ and constant $\frac1{40}$ gives the TV bound and $\gs(K_1)\ge\frac1{40}$. Since $(39/40)^{55}=\exp(-55\log(40/39))<\exp(-1.39)<\frac14$ (as $\log(40/39)>0.0253$ and $55\cdot0.0253>1.39>\log4$), $t_{\rm mix}(K_1)\le55$. Take $n_0$ as the largest of the finitely many thresholds used.
\end{proof}

\begin{T09remark}[Any prescribed exponential rate]\label{T09:rem:rate}
Every argument of this section holds for every fixed $\tau,\rho\in(0,1)$, and $\kappa_2=\frac12\log\frac1{1-\tau^2\rho^2}\to\infty$ as $\tau\rho\uparrow1$. So for any prescribed $K_0>0$, choosing $\tau,\rho$ with $\kappa_2\ge2K_0$ gives a family with $\eps_1\le e^{-K_0n}$ and discrepancy range $\ge K_0n$ for all $n\ge n_0(\tau,\rho)$, a switch, and still $\gs(K_1)\ge\frac1{40}$ and $t_{\rm mix}(K_1)\le55$. Hence even a hypothesis that fixes the exponential rate (``$\eps_\phi\le e^{-K_0n}$'' or ``range $\ge K_0n$'' for a given $K_0$) does not imply slow mixing.
\end{T09remark}

Thus, for the hybrid kernel, an exponentially small minorisation constant $\eps_\phi$, a switch and a discrepancy range linear in $n$ can all hold, at any prescribed exponential rate (Remark~\ref{T09:rem:rate}), while the absolute gap stays $\ge1/40$ and the mixing time stays $\le55$, for every $n\ge n_0$. Hence none of them, alone or together, implies an upper bound on $\gs(K_\phi)$ tending to $0$ or a lower bound on $t_{\rm mix}(K_\phi)$ tending to infinity (in particular none exponentially large in $n$): any such implication would fail on this family. This proves the reading of (v). Together with Sections 2--5 all claims of STATEMENT.md are proved. \qed
\mdproofend
\mdsec{(d) Review record}
\mdsubsec{FINAL\_REPORT.md: statement changes}
\begin{mdblock}
\mdhead{2}{2. Statement changes}

\textbf{No hypothesis was added or changed.} Two informal sentences of the theorem file were written as precise mathematics in STATEMENT.md. The original wording is quoted next to each.

\begin{mdnum}
\item[\mdlbl{1.}] (iii) "An individual realised order need not be reversible" became: a fixed-order round is reversible exactly when it equals its reversed-order round, and an explicit instance exists where it is not.
\item[\mdlbl{2.}] (v) "does not establish exponentially slow mixing …" became an existence claim. It is a family (n → ∞, p = 2) with ε₁ ≤ e\textasciicircum{}\{−κn\}, discrepancy range ≥ κn and a switch, yet gap ≥ 1/40 and mixing time ≤ 55.
\end{mdnum}

One correction to my own reading was made after Astra \#1's round-1 review. My first version said "no lower bound on mixing time that tends to 0", which is wrong; the corrected text says "tends to infinity". The correction is recorded in the STATEMENT.md change log and in CLAIMS\_OF\_FATAL\_FLAWS.md. All three reviewers judged the final readings faithful to the original wording.

\textbf{Dependencies:} none were used. The proof is self-contained: it does not assume the project results P1 or P6 or any sibling v3 conjecture. Neither EXP nor EXP′ is needed.
\end{mdblock}

\mdsubsec{FINAL\_REPORT.md: review history}
\begin{mdblock}
\mdhead{2}{4. Review history}

{\small
\begin{longtable}{|>{\raggedright\arraybackslash}p{\dimexpr(\linewidth-10\tabcolsep-6\arrayrulewidth)*72/1000\relax}|>{\raggedright\arraybackslash}p{\dimexpr(\linewidth-10\tabcolsep-6\arrayrulewidth)*291/1000\relax}|>{\raggedright\arraybackslash}p{\dimexpr(\linewidth-10\tabcolsep-6\arrayrulewidth)*73/1000\relax}|>{\raggedright\arraybackslash}p{\dimexpr(\linewidth-10\tabcolsep-6\arrayrulewidth)*125/1000\relax}|>{\raggedright\arraybackslash}p{\dimexpr(\linewidth-10\tabcolsep-6\arrayrulewidth)*439/1000\relax}|}
\hline
\textbf{Stage} & \textbf{Reviewer} & \textbf{Version} & \textbf{Verdict} & \textbf{Findings} \\ \hline\endhead
A, round 1 & Astra \#1 (gpt-6-astra high, session 01a16420-941b-7872-810a-426ec03c2ee9) & v1 & not verified & 1 fatal (my wrong gloss in the reading of (v)); 1 exposition (0/0 notation). Both accepted and fixed in v2. \\ \hline
A, round 2 & Astra \#1 (same session) & v2 & \textbf{verified} & both findings ruled resolved; no new findings \\ \hline
B, round 1 & Astra \#2 (gpt-6-astra high, new session 01a16428-59e8-7843-847b-ceb1857aa7cb) & v2 & \textbf{verified} & none \\ \hline
B, round 1 & Opus 5.5 xhigh (native agent, Claude Personal, read-only) & v2 & \textbf{verified} & 7 exposition (F1 graded "borderline gap": prescribed rates). All accepted; unified revision v3. \\ \hline
B, round 2 & Astra \#2 & v3 & \textbf{verified} & agrees with all answers; 1 new exposition finding (argmax braces) \\ \hline
B, round 2 & Opus 5.5 xhigh & v3 & \textbf{verified} & agrees with all answers; 3 new exposition findings (notation clashes N1–N3) \\ \hline
\end{longtable}}

\begin{mdbul}
\item CLAIMS\_OF\_FATAL\_FLAWS.md holds the one fatal finding verbatim with my answer, and the borderline F1. Every reviewer who saw the log agreed with the answers and with the fixes.
\item \textbf{Astra judge: not used.} No dispute remained.
\item \textbf{Open exposition notes on PROOF\_v3:} all notation only; no reviewer found any effect on the mathematics. They are not applied, so that the final text is exactly the one all reviewers checked.

\begin{mdbul}
\item The rate symbol K₀ clashes with the cap K and the kernels K\_φ.
\item The vectors h₁, h₂, h₃ clash with the step count h.
\item Subscript 1 means both the state \{1\} and φ = 1.
\item R means both the correlation matrix and the ratios R₁, R₂.
\item "w(1) = w(1)" should say "both singletons have the same prior weight".
\item "argmax = \{1,2\}" should read "the unique maximiser is \{1,2\}".
\end{mdbul}

Fixing them is a five-minute copy-edit for the campaign document.
\end{mdbul}
\end{mdblock}

\mdsubsec{CLAIMS\_OF\_FATAL\_FLAWS.md (every entry, verbatim)}
\begin{mdblock}
\mdhead{1}{CLAIMS\_OF\_FATAL\_FLAWS — T09 (T7)}

Every fatal or gap finding, verbatim, with the lead's answer.

\mdhead{2}{Stage A, Astra \#1 (session 01a16420-941b-7872-810a-426ec03c2ee9), review of PROOF\_v1, finding 1 (fatal)}

\begin{mdquote}
\textbf{Location:} STATEMENT.md, claim (v), final sentence; PROOF\_v1.tex, Section 6, closing paragraph (line 199).

\textbf{Severity:} fatal (the literal concluding claim is false).

\textbf{Claim:} The statement asserts that no upper bound on the absolute gap "(equivalently, no lower bound on its mixing time) that tends to 0" can follow from the listed conditions. The parenthetical does not give an equivalent assertion: a mixing-time lower bound tending to zero is compatible with uniformly fast mixing and is available here.

\textbf{Evidence:} Define the usual worst-case mixing time t\_mix(1/4) = min\{h ≥ 0 : sup\_x ‖K\_φ\textasciicircum{}h(x,·) − q\_φ‖\_TV ≤ 1/4\}. Every permitted state space contains at least ∅ and \{1\}. Therefore sup\_x ‖δ\_x − q\_φ‖\_TV = 1 − min\_x q\_φ(x) ≥ ½, so t\_mix(1/4) ≥ 1. In particular, t\_mix(1/4) ≥ 1/n, a valid lower bound tending to zero, including on the constructed family satisfying all three conditions.

The quoted original wording concerns an \textbf{exponentially large} mixing-time lower bound. The explicit family faithfully demonstrates that such a conclusion need not follow: its geometric TV bound gives a mixing time bounded independently of n. The final parenthetical, read literally, does not faithfully express that conclusion. This finding concerns the concluding statement, not the validity of the constructed family.
\end{mdquote}

\textbf{Lead's answer: accepted and fixed.} The parenthetical was a mis-stated gloss in my reading of the informal clause (a gap upper bound tending to 0 corresponds to a mixing-time lower bound tending to infinity, not to 0). Changes:

\begin{mdbul}
\item STATEMENT.md, reading of (v): the consequence now reads "the family has t\_mix(K₁) ≤ 55 for every n ≥ n₀; so no upper bound on gap*(K\_φ) that tends to 0, and no lower bound on t\_mix(K\_φ) that tends to infinity (in particular none exponentially large in n), can be derived from …", with t\_mix defined. Recorded in the STATEMENT.md change log. The original wording, the hypotheses and the existence claim are unchanged.
\item PROOF\_v2, Section 6: the proposition now also states t\_mix(K₁) ≤ 55, proved from (39/40)\textasciicircum{}55 < 1/4; the closing paragraph is rewritten to match.
\end{mdbul}

(Finding 2 of the same review is exposition: the 0/0 form of the reverse proposal probability for a non-symmetric neighbour family. Fixed in PROOF\_v2 by writing both proposal probabilities by cases.)

\mdhead{2}{Stage B (independent reviews of PROOF\_v2)}

\begin{mdbul}
\item Astra \#2 (session 01a16428-59e8-7843-847b-ceb1857aa7cb): no findings; verdict verified.
\item Opus 5.5 xhigh reviewer: verdict verified; seven findings, all graded exposition (F1 "exposition (borderline gap)"). No fatal or gap finding, so no entry is required here. For the record, F1 is quoted verbatim with its answer, because it was graded borderline:
\end{mdbul}

\begin{mdquote}
\textbf{F1.} Location: Section 6, closing paragraph ("any such implication would fail on this family"); STATEMENT (v), the "Consequently…" sentence. Severity: exposition (borderline gap). Claim: The proof shows the family has ε₁ ≤ e\textasciicircum{}\{−κn\} and discrepancy range ≥ κn with κ = κ\_b/2, where κ\_b = ½log(1/(1−c²ρ²)) for fixed correlation c and ρ. In fact ε₁ ≍ 1/BF(b) ≍ e\textasciicircum{}\{−κ\_b n\}. Take a would-be implication whose hypothesis fixes a rate, such as "ε\_φ ≤ e\textasciicircum{}\{−5n\} ⇒ gap → 0". At the example values c = ½, ρ = 3/5 we get κ\_b ≈ 0.047, so the family does not satisfy that hypothesis and does not refute the implication. Evidence: κ\_b = ½log(1/0.91) ≈ 0.047. The STATEMENT's own wording ("there exist constants c > 0") reads "exponentially small" and "large range" existentially, and under that reading the sentence is proved. The stronger fixed-rate version is easy to repair, but the proof never states the repair: κ\_b → ∞ as cρ ↑ 1, and every argument in Lemmas 10–12 and Propositions 13–14 holds for any fixed c, ρ ∈ (0,1).
\end{mdquote}

\textbf{Lead's answer: accepted; added in PROOF\_v3 as Remark "Any prescribed exponential rate"} (for any K₀ > 0 choose τ, ρ with κ₂ ≥ 2K₀; all arguments hold for fixed τ, ρ; gap ≥ 1/40 and t\_mix ≤ 55 unchanged). The statement is unchanged (its existential reading was already proved). The other exposition findings F2–F7 are all accepted and fixed in PROOF\_v3 (see review/STAGE\_B\_CHANGES\_v2\_to\_v3.md).
\end{mdblock}

\chapter{T11(iii): the final run's iid \texorpdfstring{$\pi_0$}{pi0} draws}
\chaptersource{Final proof: \texttt{proofs/\allowbreak{}T11\_\allowbreak{}LLED\_\allowbreak{}scoring\_\allowbreak{}ledger/\allowbreak{}proof/\allowbreak{}PROOF\_v5.tex}. Sources of parts (a), (b), (d): \texttt{theorems/\allowbreak{}T11\_\allowbreak{}LLED\_\allowbreak{}scoring\_\allowbreak{}ledger.md}, \texttt{STATEMENT.md}, \texttt{FINAL\_REPORT.md}, \texttt{CLAIMS\_OF\_FATAL\_FLAWS.md}.}
\mdsec{(a) Original conjecture text (theorem file, verbatim)}
\begin{mdblock}
\mdhead{1}{T11 (L-LED/T14-D): Unified scoring and discovery ledger}

\textbf{Status:} conjectured, not proven. Methodology SMC-BVS v3.0 (2026-10-19), SHA-256 93043758058a6d422cca6ca4b6be8f3a9e8149b6f6de2e1bab1ce31c44b8019c. Read with {\mdtt 00\_\allowbreak{}COMMON.\allowbreak{}md}.

\textbf{Category / rank:} algorithmic validity and discovery; \textbf{core} (lemma). Merges v2's L-LED and T14-D.

\mdhead{2}{Definitions used}

Scored set 𝒮, 𝒮\_cal, reported HPM Ŝ, global S†; the mechanisms (Gibbs flips, uniform swap, neighbour swaps on N(·), independence proposals and initial draws from π₀, block tables, local scans); execution-time state X\_\{pre-j\} before coordinate j's step; Q\_f = 4N\_e(T+L+5); "opportunity" = a (particle, round, step) at which a mechanism could score the model; first-hit time τ; conditional hazard p\_e (00\_COMMON §2).

\mdhead{2}{Assumptions}

EX-A, EX-S, CAL; SR for Gibbs scoring. Component execution or launch conditions are explicit.

\mdhead{2}{Precise statement}

Fix S* ∈ Ω measurable w.r.t. data and 𝓒 (before final randomness). (i) [per-mechanism rules] Gibbs coordinate j: under SR, S* is evaluated iff X\_\{pre-j\} ∖ \{j\} = S* ∖ \{j\} (this includes the current-state case; before first discovery only the flip-neighbour case gives a new score). Uniform swap from a size-k exchange neighbour: the unique exchange producing S* is proposed with probability 1/[k(p−k)]. Neighbour swap: probability 1/[k|N(j)∖S|] if the insertion is in N(j)∖S, else 0 (degrees may exceed B). Independence proposal / initial draw: proposal probability π₀(S*), regardless of acceptance. Tables: every S* ⊆ C\_b is scored when the tables execute. Local scan: given its centre, every feasible flip and swap is scored; a random centre contributes a probability after averaging, not a deterministic unconditional indicator. (ii) [expected Gibbs evaluations] Σ\_\{r,i,ρ,j\} P(X\textasciicircum{}\{r,i,ρ\}\_\{pre-j\} ∖ \{j\} = S* ∖ \{j\}); if a round starts in stationarity every state-independent ordering prefix preserves q\_φ. Actual particles are not assumed stationary. (iii) [iid draws] P(S* absent from the final iid draws | 𝓒) = (1 − π₀(S*))\textasciicircum{}\{Q\_f\} ≤ e\textasciicircum{}\{−Q\_fπ₀(S*)\}. The adaptive calibration draw count is never a fixed exponent. (iv) [stopped opportunities] For chronological predictable opportunities with success probability ≥ b, with O the count up to first success or the horizon (including the successful one): P(no hit, O ≥ D) ≤ (1−b)\textasciicircum{}D; if P(no hit, O < D) ≤ δ then P(no hit) ≤ δ + (1−b)\textasciicircum{}D. An expected count alone does not supply this premise. Generally P(τ ≤ H) = Σ\_e E[1\{τ > e⁻\}p\_e], and P(hit) ≥ bEO ≥ 1 − E(1−b)\textasciicircum{}O (a weaker bound, not an independent-trials identity). (v) [set exposure] For a fixed set 𝓜: P(𝓜 scored) ≤ 1\{𝒮\_cal ∩ 𝓜 ≠ ∅\} + Q\_fπ₀(𝓜) + EΣ\_e p\_e\textasciicircum{}\{noniid\}, with future Gibbs, swap and local-scan opportunities in the last sum; clip at one. (vi) [launch] On the SMC branch the unconditional final-only miss bound is P(not launched, S* ∉ 𝒮\_cal) + E[1\{launched, S* ∉ 𝒮\_cal\}(1 − π₀(S*))\textasciicircum{}\{Q\_f\}]. Enumeration scores every feasible model. (vii) If S† is scored, then Ŝ = S† (tie rule; NIS not needed).

\mdhead{2}{Dependencies}

Builds on: T09, T10, decision rule. Used by: T20, T21, T24, T27, T28, T29, T30, T31.

\mdhead{2}{Numerical evidence (provisional)}

provisional evidence (Sol check J8): exact stopped Bernoulli tree (b = 0.3): exact hit 0.51 = first-failure sum; the incorrect product expression gives 0.447; P(no hit, O ≥ 2) = 0.49 = bound. Literal Gibbs "iff" without the current-state case: FAIL. provisional evidence (Sol check J7): Q\_f bounds at n = 200, p = 500 tier 0: Q\_min = 140,000, Q\_max = 10,130,000.

\mdhead{2}{Proof-route hint (optional)}

Condition on each execution prefix; fresh proposal randomness; first-hit decomposition and sequential survival bounds; the adaptive union bound with deterministic ceilings for (v). Hardest step: ensuring a claimed opportunity lower bound holds on the no-hit event. \emph{When proof.md is run on this file, its lead considers the statement both with and without these hints; the hints never reach the reviewers.}
\end{mdblock}

\mdsec{(b) STATEMENT.md (verbatim)}
\begin{mdblock}
\mdhead{1}{STATEMENT — T11(iii): iid π₀ draws of the final run}

Source: {\mdtt theorems/\allowbreak{}T11\_\allowbreak{}LLED\_\allowbreak{}scoring\_\allowbreak{}ledger.\allowbreak{}md}, part (iii) only. Method: SMC-BVS v3.0, METHODOLOGY\_v3.md (SHA-256 93043758058a6d422cca6ca4b6be8f3a9e8149b6f6de2e1bab1ce31c44b8019c). Shared definitions: {\mdtt theorems/\allowbreak{}00\_\allowbreak{}COMMON.\allowbreak{}md}.

Original wording (the file's preamble to all parts): "Fix S* ∈ Ω measurable w.r.t. data and 𝓒 (before final randomness)." Original wording of (iii): "P(S* absent from the final iid draws | 𝓒) = (1 − π₀(S*))\textasciicircum{}\{Q\_f\} ≤ e\textasciicircum{}\{−Q\_fπ₀(S*)\}. The adaptive calibration draw count is never a fixed exponent."

Statement changes: \textbf{none.} The items under "Setting" make explicit what the original wording leaves implicit; none adds a hypothesis beyond the file's assumptions (EX-A, EX-S, CAL) and the method's own specification (item 5 transcribes the method's description of the final run; aborts imposed by the computing environment are outside the idealised model, as everywhere in this theory).

\mdhead{2}{Setting}

\begin{mdnum}
\item[\mdlbl{1.}] \textbf{Data.} The data (X, y) are fixed (all probabilities are over the algorithm's randomness given the data). p, n, K, the blocks C\_1, …, C\_J, the tables r\_b, γ₀, Z̄₀ and π₀ = γ₀/Z̄₀ on Ω = \{S ⊆ [p] : |S| ≤ K\} are then fixed (they are functions of the data only; METHODOLOGY §6.3). The definitions (00\_COMMON §§1–2, METHODOLOGY §5, §6.3):

\begin{mdbul}
\item K = min(p, ⌊n/2⌋); the blocks C\_1, …, C\_J partition [p];
\item size prior ϖ(S) = w(|S|)/W\_K for |S| ≤ K, with w(k) = ½/C(2p−1, k) > 0 and W\_K = Σ\_\{k≤K\} C(p,k) w(k);
\item tables r\_b(U) = m(U)/m(∅) for U ⊆ C\_b, where m(·) > 0 is the marginal likelihood, so every r\_b(U) > 0;
\item γ₀(S) = ϖ(S) ∏\_b r\_b(S ∩ C\_b) for S ∈ Ω, Z̄₀ = Σ\_\{S∈Ω\} γ₀(S), π₀ = γ₀/Z̄₀. Hence π₀(S) > 0 for every S ∈ Ω.
\end{mdbul}
\item[\mdlbl{2.}] \textbf{Randomness (EX-S).} The algorithm's randomness is a countable family of independent streams, each an iid sequence of Uniform(0,1) variables, with the method's stream separation: the calibration uses streams 0–9 only and the final run uses streams 10–13 only. Every random choice of the algorithm is made by reading not-yet-read values from the streams; which value is read next is determined by the data and the values read so far. Random operations are executed one at a time (each, e.g. one call of the π₀ sampler, completes before the next begins). All maps are Borel measurable.
\item[\mdlbl{3.}] \textbf{Exactness (EX-A).} Scores, convolutions, probabilities and comparisons are exact. In particular, each call of the π₀ sampler (METHODOLOGY §6.3: total size, then block sizes backwards, then each block's sub-model within its size class) maps the values it reads to its output by a measurable map (which may depend on the run's state when the call starts), and, when the values it reads are iid Uniform(0,1), its output has law exactly π₀.
\item[\mdlbl{4.}] \textbf{Calibration output (CAL).} 𝓒 := (grid 0 = φ₀ < … < φ\_T = 1, T, L, N\_e, tier e, candidate set V, every V-independent setting) is a measurable function of the data and streams 0–9; it includes the indicator of the event Λ = \{the final run is launched\} (SMC branch, preflight passed, a grid was produced; all decided before any final randomness). On Λ, T ≥ 1 is finite.
\item[\mdlbl{5.}] \textbf{Final run} (on Λ; METHODOLOGY §6.8, §6.5). 4 islands, each with N\_e particles started from π₀ draws; the frozen grid; one move round per stage t = 1, …, T, then L burn-in rounds and 4 measured rounds at φ = 1; so each particle of each island goes through T + L + 4 rounds, and each round contains exactly one independence step, whose proposal is a π₀ draw (accepted or not). The \textbf{final iid draws} are these 4N\_e initial draws and 4N\_e(T + L + 4) independence proposals, in total Q\_f := 4N\_e(T + L + 5), a function of 𝓒. A launched final run executes all these rounds (no early stop; aborts imposed by the computing environment, such as wall time, are outside the idealised model).
\item[\mdlbl{6.}] \textbf{Target model.} S* is an Ω-valued random variable that is a measurable function of the data and 𝓒 (the file's "measurable w.r.t. data and 𝓒 (before final randomness)").
\end{mdnum}

\mdhead{2}{Claim (iii)}

On Λ, almost surely,

\hspace*{1.5em}P(S* is not among the Q\_f final iid draws | 𝓒) = (1 − π₀(S*))\textasciicircum{}\{Q\_f\} ≤ exp(−Q\_f π₀(S*)).

\mdhead{2}{Scope note (the clause "the adaptive calibration draw count is never a fixed exponent")}

The identity is asserted only for the final draws, with the 𝓒-measurable exponent Q\_f. Nothing is asserted about the calibration's π₀ draws: their number is adaptive (it depends on the calibration's own randomness, e.g. through the rebuild loop and the escalation decision), so the identity is \textbf{not} claimed with the calibration draw count in place of Q\_f.

\mdhead{2}{Method facts quoted verbatim (METHODOLOGY\_v3.md), used only to justify Setting item 5 and the scoring corollary}

\begin{mdbul}
\item §2, statuses: "{\mdtt numerical\_\allowbreak{}failure} (scoring failed after the fallback; estimates unavailable; the best scored model returned, labelled incomplete) · {\mdtt resource\_\allowbreak{}limit} (preflight predicts the memory or time allocation will be exceeded; estimates of the inputs returned)"; "{\mdtt schedule\_\allowbreak{}limit} (no admissible schedule within the stage ceiling at the highest tier; the best scored model is returned, evidence and PIPs unavailable)".
\item §7 pseudocode: "preflight (§8.6) → maybe return"; "if grid construction failed at e = 1: return schedule\_limit (best scored model)"; then "islands r = 1..4 (streams 10..13): FINAL\_ISLAND(grid, N\_e)    \# bridge + L burn-in + 4 measured rounds".
\item §6.13: "On a Cholesky failure, refactor from scratch, then use an augmented QR on X\_S; if that fails, stop with {\mdtt numerical\_\allowbreak{}failure}".
\item §6.5: "Independence proposal: S′ \textasciitilde{} π₀ (exact draw), accepted with 1 ∧ e\textasciicircum{}\{φ[d(S′) − d(S)]\}. Every proposal is scored, accepted or not."
\end{mdbul}
\end{mdblock}

\mdsec{(c) Final proof, verbatim: body of PROOF\_v5.tex}
\noindent{\small\itshape The proof below is the body of \texttt{PROOF\_v5.tex} from \texttt{\textbackslash begin\{document\}} to \texttt{\textbackslash end\{document\}} without the title block. Its preamble definitions are local to this chapter; labels carry the prefix \texttt{T11:}.}\par\medskip
\mdproofbegin
% ---- preamble of PROOF_v5.tex (theorem environments carry the prefix T11) ----
\newtheorem{T11theorem}{Theorem}
\newtheorem{T11lemma}[T11theorem]{Lemma}
\newtheorem{T11proposition}[T11theorem]{Proposition}
\newtheorem{T11corollary}[T11theorem]{Corollary}
\theoremstyle{remark}
\newtheorem{T11remark}[T11theorem]{Remark}
\newcommand{\cC}{\mathcal{C}}
\newcommand{\cF}{\mathcal{F}}
\newcommand{\Unif}{\mathrm{Unif}(0,1)}
\newcommand{\1}{\mathbf{1}}
% ---- body of PROOF_v5.tex ----
\noindent Notation and the setting are those of \texttt{STATEMENT.md} (items 1--6). Throughout, the data $(X,y)$ are fixed;
$\Omega=\{S\subseteq[p]:|S|\le K\}$; $\pi_0=\gamma_0/\bar Z_0$ is a fixed probability on $\Omega$ with
$\pi_0(S)>0$ for all $S\in\Omega$; $\Lambda$ is the event that the final run is launched; and
$Q_f=4N_e(T+L+5)$.

\section{Two facts about reading fresh uniforms}

Let $(U_\iota)_{\iota\in I}$ be a countable family of independent $\Unif$ variables. A \emph{reading scheme} is a
sequence of measurable maps $f_1,f_2,\dots$ with $f_1\in I$ a constant and $f_k:[0,1]^{k-1}\to I$, such that the
indices
\[
\iota_1=f_1,\qquad \iota_k=f_k(V_1,\dots,V_{k-1}),\qquad V_k:=U_{\iota_k}\quad(k\ge1)
\]
are pairwise distinct for every realisation (each value is read at most once).

\begin{T11lemma}[reading lemma]\label{T11:lem:read}
For every reading scheme, $(V_k)_{k\ge1}$ is an iid $\Unif$ sequence.
\end{T11lemma}

\begin{proof}
It suffices to show, for every $k\ge1$ and Borel sets $A_1,\dots,A_k\subseteq[0,1]$, that
$P(V_1\in A_1,\dots,V_k\in A_k)=\prod_{i\le k}\lambda(A_i)$, with $\lambda$ Lebesgue measure. We induct on $k$; the
case $k=1$ holds because $V_1=U_{f_1}$. Let $k\ge2$ and $E=\{V_1\in A_1,\dots,V_{k-1}\in A_{k-1}\}$. Then
\[
P(E\cap\{V_k\in A_k\})=\sum_{\iota\in I}P\big(E\cap\{\iota_k=\iota\}\cap\{U_\iota\in A_k\}\big).
\]
Fix $\iota$. Let $U^{(\iota)}$ be the family $U$ with $U_\iota$ replaced by the constant $\tfrac12$, and let
$(\iota'_j,V'_j)$ be the indices and values the same scheme produces from $U^{(\iota)}$; these are functions of
$(U_\kappa)_{\kappa\ne\iota}$. On the event that $\iota\notin\{\iota_1,\dots,\iota_{k-1}\}$, an induction on
$j\le k-1$ gives $\iota'_j=\iota_j$ and $V'_j=V_j$ (each index is computed from equal earlier values, and the value read
there is not $U_\iota$), and hence $\iota'_k=\iota_k$. Symmetrically, on the event
$\iota\notin\{\iota'_1,\dots,\iota'_{k-1}\}$ the same equalities hold. Because indices are never repeated,
$\{\iota_k=\iota\}\subseteq\{\iota\notin\{\iota_1,\dots,\iota_{k-1}\}\}$ and likewise for the primed scheme, so
\[
E\cap\{\iota_k=\iota\}=E'\cap\{\iota'_k=\iota\},\qquad E':=\{V'_1\in A_1,\dots,V'_{k-1}\in A_{k-1}\}.
\]
The right side is $\sigma(U_\kappa:\kappa\ne\iota)$-measurable, hence independent of $U_\iota$. Therefore
$P(E\cap\{\iota_k=\iota\}\cap\{U_\iota\in A_k\})=P(E\cap\{\iota_k=\iota\})\,\lambda(A_k)$. Summing over $\iota$
gives $P(E\cap\{V_k\in A_k\})=P(E)\lambda(A_k)$, and the induction hypothesis completes the step.
\end{proof}

Let $\cF_t:=\sigma(V_1,\dots,V_t)$ ($\cF_0$ trivial). A random variable $\tau$ with values in
$\{0,1,2,\dots\}$ is a \emph{stopping time} if $\{\tau=t\}\in\cF_t$ for every $t$.

\begin{T11lemma}[fresh block]\label{T11:lem:fresh}
Let $(V_k)$ be iid $\Unif$ and $\tau$ an a.s.\ finite stopping time (on the null set $\{\tau=\infty\}$ put
$V_{\tau+j}:=\tfrac12$). Then $(V_{\tau+j})_{j\ge1}$ is an iid $\Unif$ sequence independent of
$\cF_\tau:=\{A\in\sigma(V_1,V_2,\dots):A\cap\{\tau=t\}\in\cF_t\ \forall t\}$.
\end{T11lemma}

\begin{proof}
Let $A\in\cF_\tau$ and let $C$ be a measurable subset of $[0,1]^{\mathbb N}$. For each $t$,
$A\cap\{\tau=t\}\in\cF_t$ is independent of $(V_{t+j})_{j\ge1}$, which is iid $\Unif$; hence
$P(A\cap\{\tau=t\}\cap\{(V_{t+j})_j\in C\})=P(A\cap\{\tau=t\})\,\mu(C)$, where $\mu$ is the law of an iid $\Unif$
sequence. Summing over $t$ (using $\tau<\infty$ a.s.) gives
$P(A\cap\{(V_{\tau+j})_j\in C\})=P(A)\mu(C)$. Taking $A$ to be the whole sample space identifies the law as $\mu$, and the product form
gives independence.
\end{proof}

\section{The $\pi_0$ sampler}

Setting item 3 (EX-A) states that one call of the $\pi_0$ sampler maps the values it reads to its output by a
measurable map (possibly depending on the run's state when the call starts) and that, when these values are iid
$\Unif$, its output has law exactly $\pi_0$. For completeness we verify this for the backward sampler of METHODOLOGY \S6.3, written out as follows; only the
conclusion of Lemma~\ref{T11:lem:sampler} is used later. The verification uses the definitions in Setting item 1: the
blocks partition $[p]$, every $r_b(U)>0$, $w(k)>0$, and $\gamma_0(S)=\frac{w(|S|)}{W_K}\prod_br_b(S\cap C_b)$ on
$\Omega$.

Let $e_b(j)=\sum_{U\subseteq C_b,|U|=j}r_b(U)$ for $0\le j\le|C_b|$ (and $e_b(j)=0$ otherwise), $c^{(0)}_k=\1\{k=0\}$,
and $c^{(b)}_k=\sum_j e_b(j)c^{(b-1)}_{k-j}$ for $b=1,\dots,J$ and integers $k$ (with $c^{(b-1)}_m=0$ for $m<0$). Since
the blocks partition $[p]$, $c^{(J)}_k=\sum_{|S|=k}\prod_b r_b(S\cap C_b)$, so
$W_K\bar Z_0=\sum_{k\le K}w(k)c^{(J)}_k=:M$. One call reads values $u_0,u_J,\dots,u_1,u'_J,\dots,u'_1$ (fresh) and:
\begin{enumerate}[label=(\alph*)]
\item draws $k\in\{0,\dots,K\}$ with probability $w(k)c^{(J)}_k/M$ by the inverse distribution function at $u_0$;
\item sets $k_J=k$ and, for $b=J,J-1,\dots,1$, draws $j_b$ with probability $e_b(j)c^{(b-1)}_{k_b-j}/c^{(b)}_{k_b}$ by the
inverse distribution function at $u_b$, and sets $k_{b-1}=k_b-j_b$;
\item for each $b$, draws $U_b\subseteq C_b$ with $|U_b|=j_b$ with probability $r_b(U_b)/e_b(j_b)$ by the inverse
distribution function (in a fixed enumeration of the subsets) at $u'_b$; and returns $S=\bigcup_bU_b$.
\end{enumerate}

\begin{T11lemma}\label{T11:lem:sampler}
Under exact arithmetic, if the values read are iid $\Unif$, the returned $S$ has law $\pi_0$.
\end{T11lemma}

\begin{proof}
All tables are positive, so $e_b(j)>0$ for $0\le j\le|C_b|$ and $c^{(b)}_k>0$ exactly for
$0\le k\le\sum_{b'\le b}|C_{b'}|$. The probabilities in (a)--(c) are nonnegative and sum to one (for (b) by the
recursion defining $c^{(b)}$). A value of $k_b$ with $c^{(b)}_{k_b}=0$ is reached with probability $0$: $k_J=k$ has
$w(k)c^{(J)}_k>0$, and $j_b$ is drawn with positive probability only if $c^{(b-1)}_{k_b-j_b}>0$. At $b=1$ this forces
$k_0=0$, so $\sum_bj_b=k$. An inverse distribution function evaluated at a $\Unif$ value is an exact draw from the
corresponding finite distribution, and the values are independent, so the probability of a path is the product of
its factors. A set $S\in\Omega$ is produced by exactly one path of positive probability: $k=|S|$, $j_b=|S\cap C_b|$, $U_b=S\cap C_b$, with
$k_b=\sum_{b'\le b}j_{b'}$. Its probability is
\[
\frac{w(k)c^{(J)}_k}{M}\prod_{b=1}^J\frac{e_b(j_b)\,c^{(b-1)}_{k_{b-1}}}{c^{(b)}_{k_b}}
\prod_{b=1}^J\frac{r_b(S\cap C_b)}{e_b(j_b)}
=\frac{w(k)}{M}\,c^{(0)}_{0}\prod_b r_b(S\cap C_b)
=\frac{w(|S|)}{W_K}\frac{\prod_b r_b(S\cap C_b)}{\bar Z_0}=\pi_0(S),
\]
the product of $c^{(b-1)}_{k_{b-1}}/c^{(b)}_{k_b}$ telescoping to $c^{(0)}_0/c^{(J)}_k$. Sets outside $\Omega$ are never
returned since $k\le K$.
\end{proof}

\section{The final run is a reading scheme}

\begin{T11remark}[the final run executes in full on $\Lambda$]\label{T11:rem:full}
The main proof uses Setting item 5 as stated. This remark only explains why item 5 is the method's specification,
using the method facts quoted in \texttt{STATEMENT.md}. By the \S7 pseudocode, once launched, the final run consists of
\texttt{FINAL\_ISLAND(grid, $N_e$)} for $r=1,\dots,4$ (``bridge + $L$ burn-in + 4 measured rounds''). Of the statuses
that end a run early, \texttt{resource\_limit} is returned at the preflight and \texttt{schedule\_limit} when grid
construction fails at tier 1, both before the final run; \texttt{numerical\_failure} means ``scoring failed after the
fallback'' (\S2, \S6.13). We read EX-A (``scores are exact'') as: every score computation returns
$\log m(S)$ exactly, and $m(S)$ is a finite positive number for every $S$ (Setting item 1); so scoring does not fail. Hence under EX-A no status interrupts a launched final run.
Aborts imposed by the computing environment (wall time, memory) during the final run are outside the idealised model;
Setting item 5 excludes them.
\end{T11remark}

Write $\Upsilon$ for streams 0--9 and $\Xi$ for streams 10--13. By EX-S, $\Upsilon$ and $\Xi$ are independent. By CAL,
$\cC$ is a measurable function of $\Upsilon$ (the data being fixed), and so are $\Lambda$, $T$, $N_e$, $Q_f$ and $S^*$
(Setting items 4--6). The final run, including the order in which it performs its random operations, is a measurable
function of $(\Upsilon,\Xi)$; it reads values only from $\Xi$.

\begin{T11proposition}\label{T11:prop:iid}
Fix a value $\upsilon$ of $\Upsilon$ with $\upsilon\in\Lambda$, and run the final run with $\Upsilon=\upsilon$ and
$\Xi$ iid as in EX-S. Call a class of the run's $\pi_0$-sampler calls \emph{designated} if whether a call belongs to it
is determined by the run's state when the call starts, and suppose the run makes exactly $Q$ designated calls, $Q$ a
constant (given $\upsilon$). Number their outputs $D_1,\dots,D_Q$ in chronological order. Then $D_1,\dots,D_Q$ are iid
with law $\pi_0$. In particular this holds for
\begin{enumerate}[label=(\alph*)]
\item the final iid draws of Setting item 5 (the $4N_e$ initial draws and the $4N_e(T+L+4)$ independence proposals),
with $Q=Q_f$;
\item the $4N_e(T+L+4)$ independence proposals alone, with $Q=Q_p:=4N_e(T+L+4)$.
\end{enumerate}
\end{T11proposition}

\begin{proof}
With $\upsilon$ fixed, the final run reads not-yet-read values of $\Xi$, each choice of the next index being a
measurable function of the values read so far (Setting item 2). After the run ends, continue reading the unread
values of stream 10 in increasing position; this does not affect the run. This is a reading scheme, so by
Lemma~\ref{T11:lem:read} the read sequence $(V_k)$ is iid $\Unif$, and the whole run is a measurable function of
$(V_k)$.

The $m$-th designated call ($m\le Q$) starts after $\tau_m$ values have been read. The run's state after $t$ reads is
a function of $V_1,\dots,V_t$; whether the run starts a sampler call right after the $t$-th read, whether that call is
designated, and how many designated calls preceded it are all determined by that state, so $\{\tau_m=t\}\in\cF_t$ and
$\tau_m$ is a stopping time. It is finite: the $m$-th designated call does occur (by hypothesis), and since operations
are executed one at a time, each completing before the next begins (Setting item 2), only finitely many values are read
before it starts. The call reads $V_{\tau_m+1},V_{\tau_m+2},\dots$ (the next values read by the run, in order, until the
call returns), so $D_m=G_{Z_m}\big((V_{\tau_m+j})_{j\ge1}\big)$, where $Z_m$ is the run's state when the call starts
and, for each state $z$, $G_z$ maps an input sequence to the sampler's output. Measurability: take the state to be
the history $z=(t,v_1,\dots,v_t)\in\bigsqcup_{t\ge0}[0,1]^t$ (any coarser state is a measurable function of it). Then
$Z_m=(\tau_m,V_1,\dots,V_{\tau_m})$ is $\cF_{\tau_m}$-measurable, and, since the whole run is a Borel function of the
values it reads (Setting item 2), $(z,w)\mapsto G_z(w)$ is jointly measurable on each piece $[0,1]^t\times[0,1]^{\mathbb N}$,
hence on the disjoint union.
By Setting item 3 (for the sampler of METHODOLOGY \S6.3, Lemma~\ref{T11:lem:sampler}, where $G_z$ does not depend on $z$),
$G_z$ applied to an iid $\Unif$ sequence has law $\pi_0$ for every $z$. The calls do not overlap, so
$D_1,\dots,D_{m-1}$ are functions of the values read before the $m$-th designated call starts, that is, they are
$\cF_{\tau_m}$-measurable. By Lemma~\ref{T11:lem:fresh}, $W:=(V_{\tau_m+j})_{j\ge1}$ is iid $\Unif$ and independent of
$\cF_{\tau_m}$; hence, for $A\in\cF_{\tau_m}$ and $s\in\Omega$, Fubini's theorem gives
$P(A\cap\{G_{Z_m}(W)=s\})=E\big[\1_A\,P(G_z(W)=s)|_{z=Z_m}\big]=P(A)\,\pi_0(s)$. Taking
$A=\{D_1=s_1,\dots,D_{m-1}=s_{m-1}\}$, for all $s_1,\dots,s_Q\in\Omega$,
\[
P(D_1=s_1,\dots,D_m=s_m)=P(D_1=s_1,\dots,D_{m-1}=s_{m-1})\,\pi_0(s_m),\qquad m=1,\dots,Q,
\]
and induction gives $P(D_1=s_1,\dots,D_Q=s_Q)=\prod_m\pi_0(s_m)$.

For (a) and (b): whether a sampler call is an initial draw or the independence step of some particle's round is part
of the run's state when the call starts, so both classes are designated; by Setting item 5 they contain exactly $Q_f$
and $Q_p$ calls respectively, constants given $\upsilon$.
\end{proof}

\section{The theorem}

\begin{T11theorem}[T11(iii)]\label{T11:thm:main}
Under Setting items 1--6, on $\Lambda$, almost surely,
\[
P\big(S^*\notin\{D_1,\dots,D_{Q_f}\}\,\big|\,\cC\big)=(1-\pi_0(S^*))^{Q_f}\le\exp\{-Q_f\pi_0(S^*)\}.
\]
\end{T11theorem}

\begin{proof}
Let $A=\Lambda\cap\{S^*\notin\{D_1,\dots,D_{Q_f}\}\}$. Its indicator is $\1_A=h(\Upsilon,\Xi)$ for a measurable $h$, and
for fixed $\upsilon$ the map $\xi\mapsto h(\upsilon,\xi)$ is the indicator that the final run started from $\upsilon$
avoids $s^*(\upsilon)$ (the value of $S^*$ at $\upsilon$), or $0$ if $\upsilon\notin\Lambda$. Because $\Upsilon$ and
$\Xi$ are independent, Fubini's theorem gives $E[h(\Upsilon,\Xi)\mid\Upsilon]=g(\Upsilon)$ a.s., with
$g(\upsilon):=E\,h(\upsilon,\Xi)$. For $\upsilon\in\Lambda$, Proposition~\ref{T11:prop:iid}(a) gives
\[
g(\upsilon)=P(D_m\ne s^*(\upsilon)\ \text{for } m=1,\dots,Q_f(\upsilon))=(1-\pi_0(s^*(\upsilon)))^{Q_f(\upsilon)},
\]
and $g(\upsilon)=0$ otherwise. So $g(\Upsilon)=\1_\Lambda(1-\pi_0(S^*))^{Q_f}$, which is a function of $\cC$ (Setting
items 4--6). Since $\sigma(\cC)\subseteq\sigma(\Upsilon)$, the tower property gives
$P(A\mid\cC)=E[g(\Upsilon)\mid\cC]=\1_\Lambda(1-\pi_0(S^*))^{Q_f}$ a.s. On $\Lambda\in\sigma(\cC)$,
$P(S^*\notin\{D_m\}\mid\cC)=P(A\mid\cC)$ a.s., which is the identity. The inequality follows from
$0\le1-x\le e^{-x}$ for $x=\pi_0(S^*)\in(0,1]$, raised to the power $Q_f\ge0$.
\end{proof}

\section{Remarks}

\begin{T11remark}[scope: calibration draws]\label{T11:rem:scope}
The argument uses two facts about the final draws that fail for the calibration's $\pi_0$ draws. (1) The exponent
$Q_f$ is fixed by $\cC$ before any final value is read; the calibration's number of draws depends on its own
randomness (rebuild loop, escalation). (2) The conditioning $\sigma$-field $\sigma(\cC)$ is independent of the final
streams, but in general not of the calibration streams ($\cC$ is computed from them), so given $\cC$ the calibration
draws need not be iid $\pi_0$. With a random count, $(1-\pi_0(S^*))^N$ is in general not the absence probability, even in expectation: with
$D_1,D_2$ iid $\pi_0$, $x=\pi_0(s)\in(0,1)$, and $N=1$ if $D_1\ne s$, $N=2$ otherwise,
$P(s\notin\{D_1,\dots,D_N\})=1-x$ while $E(1-x)^N=(1-x)^2(1+x)<1-x$. Accordingly, Theorem~\ref{T11:thm:main} is asserted
only with the exponent $Q_f$.
\end{T11remark}

\begin{T11remark}[unconditional form]
Taking expectations, $P(\Lambda,\ S^*\notin\{D_1,\dots,D_{Q_f}\})=E[\1_\Lambda(1-\pi_0(S^*))^{Q_f}]$.
\end{T11remark}

\begin{T11corollary}[scoring; beyond the stated part]
Read ``scored'' as ``its score $\log m(S)$ is computed''. By the method fact quoted in \texttt{STATEMENT.md}
(\S6.5: ``Every proposal is scored, accepted or not''), each of the $Q_p=4N_e(T+L+4)$ final independence proposals is
scored. Applying the argument of Theorem~\ref{T11:thm:main} to the proposals alone (by Proposition~\ref{T11:prop:iid}(b) they
are iid $\pi_0$ given $\Upsilon$, and $Q_p$ is a function of $\cC$), on $\Lambda$,
$P(S^*\text{ is not scored in the final run}\mid\cC)\le(1-\pi_0(S^*))^{Q_p}$ a.s. If the initial draws are also scored
(the quoted text does not say so explicitly), the exponent $Q_p$ may be replaced by $Q_f$.
\end{T11corollary}
\mdproofend
\mdsec{(d) Review record}
\mdsubsec{FINAL\_REPORT.md: statement changes}
\begin{mdblock}
\mdhead{2}{Statement changes}

None to the claim. STATEMENT.md writes out what the original one-line wording left implicit:

\begin{mdbul}
\item The data are fixed. The probability is over the algorithm's randomness.
\item The formula holds on the event that the final run is launched.
\item The stream model (EX-S), exact sampling (EX-A) and calibration-measurability (CAL) are spelled out.
\item The final-run description of METHODOLOGY §6.5, §6.8 and §7 is written out. A launched run executes all its rounds; aborts by the computing environment are outside the idealised model.
\item S* is "measurable w.r.t. data and 𝓒 (before final randomness)", quoting the theorem file's own preamble.
\end{mdbul}

During review, definitions (w, W\_K, ϖ, the tables, γ₀) and short verbatim method quotes were added so that reviewers could check them. No new hypothesis was added. EXP / EXP′ is not needed.
\end{mdblock}

\mdsubsec{FINAL\_REPORT.md: review history}
\begin{mdblock}
\mdhead{2}{Review history}

\textbf{Stage A: lead with Astra \#1} (GPT-6 Astra, high; session {\mdtt 01a1641a-\allowbreak{}4997-\allowbreak{}7370-\allowbreak{}a99a-\allowbreak{}43d4f2a4f245}).

{\small
\begin{longtable}{|>{\raggedright\arraybackslash}p{\dimexpr(\linewidth-12\tabcolsep-7\arrayrulewidth)*53/1000\relax}|>{\raggedright\arraybackslash}p{\dimexpr(\linewidth-12\tabcolsep-7\arrayrulewidth)*53/1000\relax}|>{\raggedright\arraybackslash}p{\dimexpr(\linewidth-12\tabcolsep-7\arrayrulewidth)*85/1000\relax}|>{\raggedright\arraybackslash}p{\dimexpr(\linewidth-12\tabcolsep-7\arrayrulewidth)*53/1000\relax}|>{\raggedright\arraybackslash}p{\dimexpr(\linewidth-12\tabcolsep-7\arrayrulewidth)*47/1000\relax}|>{\raggedright\arraybackslash}p{\dimexpr(\linewidth-12\tabcolsep-7\arrayrulewidth)*708/1000\relax}|}
\hline
\textbf{Round} & \textbf{Proof} & \textbf{Verdict} & \textbf{Fatal} & \textbf{Gap} & \textbf{Outcome} \\ \hline\endhead
1 & v1 & verified & 0 & 3 & All gaps were in supplementary parts, because the packet lacked definitions and method text. All three were accepted and fixed. \\ \hline
2 & v2 & verified & 0 & 1 & The gap concerned the scoring corollary, which is outside (iii): the quoted text does not show that initial draws are scored. Accepted; the corollary was restricted to proposals. Two exposition points were also fixed. \\ \hline
3 & v3 & verified & 0 & 0 & No findings. Astra \#1 agreed with every answer. Stage A closed. \\ \hline
\end{longtable}}

\textbf{Stage B: two fresh independent reviews in parallel} (amendment 1).

{\small
\begin{longtable}{|>{\raggedright\arraybackslash}p{\dimexpr(\linewidth-8\tabcolsep-5\arrayrulewidth)*68/1000\relax}|>{\raggedright\arraybackslash}p{\dimexpr(\linewidth-8\tabcolsep-5\arrayrulewidth)*68/1000\relax}|>{\raggedright\arraybackslash}p{\dimexpr(\linewidth-8\tabcolsep-5\arrayrulewidth)*384/1000\relax}|>{\raggedright\arraybackslash}p{\dimexpr(\linewidth-8\tabcolsep-5\arrayrulewidth)*479/1000\relax}|}
\hline
\textbf{Round} & \textbf{Proof} & \textbf{Astra \#2 (session {\mdtt 01a16421-\allowbreak{}b0ef-\allowbreak{}7281-\allowbreak{}984e-\allowbreak{}2af293f78cf1})} & \textbf{Opus 5.5 xhigh (native agent, Claude Personal)} \\ \hline\endhead
1 & v3 & verified, no findings & verified; 0 fatal; 1 gap (F8); 7 exposition \\ \hline
2 & v4 & verified; agreed with all answers & verified; agreed with all answers; 4 exposition notes \\ \hline
3 & v5 & verified, no findings & verified; 1 exposition note, no change required \\ \hline
\end{longtable}}

\begin{mdbul}
\item \textbf{Opus F8 (round 1)} was in Corollary 9, outside (iii). The proposals form a subfamily of the draws whose position in the chronological order may depend on the run, so the iid property of all draws did not directly transfer to them. It was accepted and fixed: Proposition 5 now covers any class of "designated" calls decided when each call starts.
\item \textbf{Opus exposition notes (round 2)} were applied in v5.
\end{mdbul}

\textbf{Disputes:} none. No Astra judge was needed.

The full log, including every fatal or gap claim verbatim with the lead's answer and the reviewers' rulings, is in {\mdtt CLAIMS\_\allowbreak{}OF\_\allowbreak{}FATAL\_\allowbreak{}FLAWS.\allowbreak{}md}.
\end{mdblock}

\mdsubsec{CLAIMS\_OF\_FATAL\_FLAWS.md (every entry, verbatim)}
\begin{mdblock}
\mdhead{1}{CLAIMS\_OF\_FATAL\_FLAWS — T11(iii)}

Every fatal or gap finding, copied verbatim, with the lead's answer.

\mdhead{2}{Stage A, round 1 — Astra \#1 (session 01a1641a-4997-7370-a99a-43d4f2a4f245), review of PROOF\_v1 (verdict: verified)}

\mdhead{3}{A1-F1 (gap)}

\begin{mdquote}
\textbf{Finding 1}

\begin{mdnum}
\item[\mdlbl{1.}] \textbf{Location:} Section 2, sampler definitions and final equality in Lemma \(\text{\mdtt\textbackslash ref\{lem:sampler\}}\).
\item[\mdlbl{2.}] \textbf{Severity:} \textbf{gap} — confined to the supplementary sampler verification.
\item[\mdlbl{3.}] \textbf{Claim:} The identification of the displayed sampler distribution with \(\pi_0\) is not established from the supplied setting.
\item[\mdlbl{4.}] \textbf{Evidence:} The setting does not define \(w(k)\), \(W_K\), or the product representation 
\[
\gamma_0(S)=\frac{w(|S|)}{W_K}\prod_b r_b(S\cap C_b).
\]
 Nor does it explicitly supply positivity of every table entry. The algebra is correct given these additional facts, but their applicability is unverified. This does not undermine Claim (iii), because EX-A already supplies sampler exactness.
\end{mdnum}
\end{mdquote}

\textbf{Lead's answer: accepted and fixed.} The packet's STATEMENT.md did not carry the definitions of w, W\_K, ϖ, γ₀ and the positivity of the tables, which the sampler verification uses. Change: Setting item 1 of STATEMENT.md now lists these definitions (taken from 00\_COMMON §§1–2 and METHODOLOGY §5, §6.3: w(k) = ½/C(2p−1,k) > 0, W\_K = Σ\_\{k≤K\}C(p,k)w(k), ϖ(S) = w(|S|)/W\_K, r\_b(U) = m(U)/m(∅) > 0, γ₀(S) = ϖ(S)∏\_b r\_b(S∩C\_b)); these are definitions, not new hypotheses. PROOF\_v2 §2 states that the verification uses exactly these. As the finding notes, the main claim uses only the exactness in Setting item 3.

\mdhead{3}{A1-F2 (gap)}

\begin{mdquote}
\textbf{Finding 2}

\begin{mdnum}
\item[\mdlbl{1.}] \textbf{Location:} Section 3, Remark \(\text{\mdtt\textbackslash ref\{rem:full\}}\).
\item[\mdlbl{2.}] \textbf{Severity:} \textbf{gap} — confined to the supplementary implementation justification.
\item[\mdlbl{3.}] \textbf{Claim:} The assertions about all possible interruption statuses and all matrices factorised cannot be verified from the supplied files.
\item[\mdlbl{4.}] \textbf{Evidence:} The setting does not specify the cited failure mechanisms or matrix inventory. Positive definiteness follows for the displayed matrices, but their exhaustiveness is unsupported here. Full execution is already stipulated by Setting item 5, so the main proof does not require these assertions.
\end{mdnum}
\end{mdquote}

\textbf{Lead's answer: accepted and fixed.} STATEMENT.md now has a section "Method facts quoted verbatim" (statuses of §2, the §7 pseudocode, §6.13's failure path). Remark 4 of PROOF\_v2 ("the final run executes in full") is rewritten: it says the main proof uses Setting item 5 as stated and only explains why item 5 is the method's specification. The argument no longer relies on an inventory of factorised matrices: {\mdtt numerical\_\allowbreak{}failure} is "scoring failed after the fallback", and under EX-A scores are computed exactly, so scoring does not fail; the positive-definiteness of H\_S and G\_S + I is kept only as an illustration.

\mdhead{3}{A1-F3 (gap)}

\begin{mdquote}
\textbf{Finding 3}

\begin{mdnum}
\item[\mdlbl{1.}] \textbf{Location:} Final corollary, “scoring; beyond the stated part.”
\item[\mdlbl{2.}] \textbf{Severity:} \textbf{gap}.
\item[\mdlbl{3.}] \textbf{Claim:} The scoring conclusion requires implementation facts absent from Setting items 1–6.
\item[\mdlbl{4.}] \textbf{Evidence:} Those items specify initial draws and independence proposals, but do not require every proposal’s score to be computed or specify the stage-1 weighting calculation. The corollary invokes unavailable methodology sections for those premises. Its set-inclusion argument is valid if the premises hold; the premises remain unverified. This corollary is explicitly outside Claim (iii).
\end{mdnum}
\end{mdquote}

\textbf{Lead's answer: accepted and fixed.} The premises are now quoted verbatim in STATEMENT.md (§6.5: "Every proposal is scored, accepted or not."; §6.12: the evidence product, whose t = 1 factor contains d(S\_i\textasciicircum{}\{(0)\}), which involves log m(S\_i\textasciicircum{}\{(0)\})). The corollary in PROOF\_v2 cites these quotes. It remains labelled as beyond the stated part (iii).

\mdhead{2}{Stage A, round 2 — Astra \#1 rulings on round-1 answers and review of PROOF\_v2 (verdict: verified)}

Rulings (verbatim):

\begin{mdquote}
\textbf{Rulings on the three previous findings}

\begin{mdnum}
\item[\mdlbl{1.}] \textbf{A1-F1 — sampler definitions: agree; the fix works.} Updated Setting item 1 supplies positivity of the tables, the size-prior normalisation, and the product formula for \(\gamma_0\). These justify \(M=W_K\bar Z_0\) and the sampler proof’s final identification with \(\pi_0\).
\end{mdnum}

\begin{mdnum}
\item[\mdlbl{2.}] \textbf{A1-F2 — full execution: agree with the substantive answer; the fix works for the main argument.} The proof explicitly takes full execution from Setting item 5. The supplied quotations support the timing of preflight and schedule failures and identify the scoring-failure status. The argument no longer requires an exhaustive matrix inventory. One illustrative implementation assertion remains unsupported; see Finding 1 below.
\end{mdnum}

\begin{mdnum}
\item[\mdlbl{3.}] \textbf{A1-F3 — scoring: partially agree; the fix is incomplete.} The new quotation establishes that every independence proposal is scored. However, the evidence-product formula does not, by itself, establish that each initial draw’s \(\log m(S)\) is computed during the final run. See Finding 2 below.
\end{mdnum}
\end{mdquote}

\mdhead{3}{A1-R2-F2 (gap; outside Claim (iii))}

\begin{mdquote}
\textbf{Finding 2}

\begin{mdnum}
\item[\mdlbl{1.}] \textbf{Location:} Final scoring corollary, assertion that each initial draw is scored “by the evidence product.”
\item[\mdlbl{2.}] \textbf{Severity:} \textbf{gap} — outside Claim (iii).
\item[\mdlbl{3.}] \textbf{Claim:} An expression containing \(d(S)\) does not necessarily require computing \(\log m(S)\) during the final run.
\item[\mdlbl{4.}] \textbf{Evidence:} Consider the permitted case of one block, \(C_1=[p]\). Then 
\[
r_1(S)=m(S)/m(\emptyset),\qquad
d(S)=\log\!\frac{m(S)}{m(\emptyset)}-\log r_1(S)=0.
\]
 Thus every evidence-product factor can be evaluated as \(1\) without computing an initial model’s score during that run. The quoted formula is compatible with this evaluation, while the explicit scoring instruction covers proposals only.

Consequently, the quotations do not establish that the event “\(S^*\) is not scored” is contained in the event “\(S^*\) is absent from all initial draws and proposals.” The main absence theorem remains valid.
\end{mdnum}
\end{mdquote}

\textbf{Lead's answer: accepted and fixed.} The quoted method text guarantees scoring only for the independence proposals. PROOF\_v3's corollary now uses the proposals alone: each of the Q\_p = 4N\_e(T+L+4) final proposals is scored (§6.5 quote), they are iid π₀ given the calibration streams (Proposition 5) and Q\_p is 𝓒-measurable, so P(S* not scored in the final run | 𝓒) ≤ (1 − π₀(S*))\textasciicircum{}\{Q\_p\} on Λ; the exponent Q\_f is stated only conditionally on the initial draws also being scored. The §6.12 quote was removed from STATEMENT.md (no longer used). Claim (iii) itself is unaffected (Astra agrees).

Exposition findings of round 2 (not fatal/gap; handled in PROOF\_v3): Finding 1 (matrix example in Remark 4 unsupported) — parenthetical removed. Finding 3 (σ(𝓒) "not independent" of calibration streams too categorical) — reworded to "in general not ... (𝓒 is computed from them)".

\mdhead{2}{Stage A, round 3 — Astra \#1 rulings and review of PROOF\_v3 (verdict: verified; no new findings). Stage A closed: lead and Astra \#1 both hold PROOF\_v3 correct.}

\begin{mdquote}
\textbf{Rulings on the author’s responses}

\begin{mdnum}
\item[\mdlbl{1.}] \textbf{Round-2 Finding 2 — scoring gap: agree; the fix works.} The revised corollary uses only the \(Q_p=4N_e(T+L+4)\) independence proposals, which the quoted method text explicitly requires to be scored. Their absence probability gives the stated upper bound on the probability of not being scored. The stronger exponent \(Q_f\) is now explicitly conditional on the initial draws also being scored. Removing the unused §6.12 quotation is consistent with this change.
\end{mdnum}

\begin{mdnum}
\item[\mdlbl{2.}] \textbf{Round-2 Finding 1 — unsupported matrix illustration: agree; the fix works.} The parenthetical implementation assertion has been removed. The main proof continues to use full execution as stipulated by Setting item 5.
\end{mdnum}

\begin{mdnum}
\item[\mdlbl{3.}] \textbf{Round-2 Finding 3 — categorical non-independence: agree; the fix works.} “In general not” correctly allows degenerate cases, including constant calibration output, while explaining why independence from calibration randomness is not guaranteed.
\end{mdnum}
\end{mdquote}

\mdhead{2}{Stage B, round 1 — Astra \#2 (session 01a16421-b0ef-7281-984e-2af293f78cf1): verified, no findings. Opus 5.5 xhigh reviewer: verified, no fatal findings, one gap (F8, outside (iii)) and seven exposition findings. All addressed in PROOF\_v4.}

\mdhead{3}{OB-F8 (gap; Corollary 9, outside Claim (iii)) — Opus reviewer}

\begin{mdquote}
\textbf{F8}

\begin{mdbul}
\item \textbf{Location:} Corollary 9, lines 211–218. This is outside claim (iii).
\item \textbf{Severity:} gap (in the corollary only; it does not affect the verdict on (iii)).
\item \textbf{Claim:} The step "by Proposition 5 they [the proposals alone] are iid π₀ given Υ" does not follow from Proposition 5 as stated.
\item \textbf{Evidence:}

\begin{mdbul}
\item Proposition 5 proves that the chronologically numbered D\_1..D\_\{Q\_f\} are iid. The proposals are a sub-family whose positions in that order may depend on Ξ: item 2 lets the order of operations depend on the values read so far, for example through how the islands are interleaved.
\item A sub-family chosen with data-dependent positions is not iid in general. For example, with D\_1, D\_2 iid, choosing D\_1 if D\_1 = s and D\_2 otherwise does not give a π₀-distributed result.
\item The corollary can be repaired by rerunning Proposition 5's stopping-time argument over proposal calls only. Whether a call is a proposal is decided at its start, so that event is 𝓕\_\{τ\_m\}-measurable. As written, though, the citation does not supply this.
\end{mdbul}
\end{mdbul}
\end{mdquote}

\textbf{Lead's answer: accepted and fixed.} Proposition 5 is restated in PROOF\_v4 for any class of \emph{designated} sampler calls (membership decided by the run's state when the call starts) with a constant number Q of calls; its stopping-time argument is run over the designated calls only, exactly as the reviewer suggests (\{m-th designated call starts after read t\} ∈ 𝓕\_t). Parts (a) (all final draws, Q = Q\_f) and (b) (independence proposals alone, Q = Q\_p) are stated explicitly; Theorem 6 cites (a) and Corollary 9 cites (b).

Exposition findings (not fatal/gap), handled in PROOF\_v4 / STATEMENT.md:

\begin{mdbul}
\item F1 (constant 0 outside (0,1)): replacement constant is now 1/2.
\item F2 (𝓕\_τ, τ = ∞): 𝓕\_τ now requires A ∈ σ(V); V\_\{τ+j\} := 1/2 on the null set \{τ = ∞\}.
\item F3 (finiteness of τ\_m): now justified by item 5 (the m-th designated call occurs) and item 2 (one operation at a time).
\item F4 (other sampler calls): covered by the designated-call formulation.
\item F5 (single map G): the proof now allows a state-dependent map G\_\{Z\_m\} with Z\_m 𝓕\_\{τ\_m\}-measurable, and Setting item 3 is worded accordingly (measurable map, possibly state-dependent; exact law π₀ on iid uniform input); the step uses Lemma 2 plus Fubini.
\item F6: "exactly one path of positive probability".
\item F7 (Remark 4): the remark now notes m(S) is finite and positive (item 1) and that environment-imposed aborts are outside the idealised model; Setting item 5 now says explicitly that a launched final run executes all its rounds, and the "Statement changes" paragraph says item 5 transcribes the method's specification.
\item Setting item 6 faithfulness: STATEMENT.md now quotes the theorem file's own definition of S* ("Fix S* ∈ Ω measurable w.r.t. data and 𝓒 (before final randomness)").
\end{mdbul}

\mdhead{2}{Stage B, round 2 — rulings on PROOF\_v4}

\begin{mdbul}
\item Astra \#2: agrees with every logged answer (A1-F1, A1-F2, A1-F3 as resolved by v4, A1-R2-F2, OB-F8) and with the handling of all exposition items; no new findings; \textbf{VERDICT: verified} (review/astra2\_stageB\_round2\_review.md).
\item Opus reviewer: agrees with every logged answer; OB-F8 closed by Proposition 5(b); F1–F7 and the item-6 concern adequately handled; four new exposition findings N1–N4, no fatal or gap; \textbf{VERDICT: verified} (review/opus\_stageB\_round2\_review.md).
\item No open fatal or gap claim. Exposition N1–N4 handled in PROOF\_v5 (diff against v4 is exposition only): N1 "A = Ω" (clash with the model space) → "A the whole sample space"; N2 Section 2 paraphrase of item 3 updated to the reworded item 3; N3 joint measurability of (z, w) ↦ G\_z(w) written out (state = history in ⊔\_t [0,1]\textasciicircum{}t; the run is a Borel function of the values read, Setting item 2); N4 Remark 4 now states the reading of EX-A it uses ("every score computation returns log m(S) exactly").
\end{mdbul}

\mdhead{2}{Stage B, round 3 — confirmation of PROOF\_v5 (final)}

\begin{mdbul}
\item Astra \#2: v5 differs from v4 only at N1–N4 (plus title); N3 measurability argument correct; no findings; \textbf{VERDICT: verified} (review/astra2\_stageB\_round3\_review.md).
\item Opus reviewer: same diff check; N3 correct step by step; one exposition note E1 ("for every z" is slightly broader than needed: G\_z is only defined at histories where a call starts; harmless, "no change is required"); \textbf{VERDICT: verified} (review/opus\_stageB\_round3\_review.md). Lead: E1 noted, no change (the Fubini step evaluates G\_z only at Z\_m, where it is defined).
\item Final state: no open fatal or gap claim; all three reviewers (Astra \#1, Astra \#2, Opus xhigh) hold the proof correct; no Astra judge needed.
\end{mdbul}
\end{mdblock}

\chapter{T12 (T8): frozen-run evidence and recorded Rao--Blackwell identities}
\chaptersource{Final proof: \texttt{proofs/\allowbreak{}T12\_\allowbreak{}T8\_\allowbreak{}frozen\_\allowbreak{}run\_\allowbreak{}identities/\allowbreak{}proof/\allowbreak{}PROOF\_v6.tex}. Sources of parts (a), (b), (d): \texttt{theorems/\allowbreak{}T12\_\allowbreak{}T8\_\allowbreak{}frozen\_\allowbreak{}run\_\allowbreak{}identities.md}, \texttt{STATEMENT.md}, \texttt{FINAL\_REPORT.md}, \texttt{CLAIMS\_OF\_FATAL\_FLAWS.md}.}
\mdsec{(a) Original conjecture text (theorem file, verbatim)}
\begin{mdblock}
\mdhead{1}{T12 (T8): Frozen-run evidence and recorded-RB identities}

\textbf{Status:} conjectured, not proven. Methodology SMC-BVS v3.0 (2026-10-19), SHA-256 93043758058a6d422cca6ca4b6be8f3a9e8149b6f6de2e1bab1ce31c44b8019c. Read with {\mdtt 00\_\allowbreak{}COMMON.\allowbreak{}md}.

\textbf{Category / rank:} algorithmic validity; \textbf{core} (theorem).

\mdhead{2}{Definitions used}

Stage notation q\_t, Z\_t, q̂\_\{r,t\}, Ẑ\_\{r,t\}; the final run (4 islands × N\_e, multinomial resampling after every reweighting, one round per stage, L burn-in + 4 measured rounds); PIP̂\_\{rj\} = pre-draw Gibbs conditionals averaged over particles and the 4 measured rounds; pooled Ẑ = ¼Σ\_rẐ\_r; primary/secondary PIPs (00\_COMMON §2).

\mdhead{2}{Assumptions}

EX-A, EX-S, CAL; prescribed multinomial resampling and independent particle mutations.

\mdhead{2}{Precise statement}

Conditional on a successful 𝓒, for every island r, N ≥ 1, stage t and bounded fixed f: (i) E[Ẑ\_\{r,t\}q̂\_\{r,t\}(f)] = Z\_tq\_t(f); in particular EẐ\_r = Z. Appending terminal invariant rounds with unit potentials preserves the identity; it holds at fixed elementary prefixes and immediately before a specified Gibbs label in a state-independent random permutation; not at arbitrary state-dependent stopping times. (ii) For the actual pre-draw RB recordings: E[Ẑ\_r PIP̂\_\{rj\}] = Z·PIP\_j. (iii) For the pooled evidence: EẐ = Z; E log Ẑ ≤ log Z; P(Z ≥ αẐ) ≥ 1 − α for 0 < α < 1 (Jensen also islandwise). (iv) Primary and secondary normalised PIPs and log Ẑ are not unbiased in general. Conditional island independence holds given 𝓒, not necessarily after averaging over shared calibration. Averaging over a 𝓒-measurable launch event preserves evidence unbiasedness; conditioning on favourable final diagnostics does not guarantee any of (i)–(iii). NUM alone supplies no unbiasedness or Markov guarantee for implemented floating-point evidence.

\mdhead{2}{Dependencies}

Builds on: T09; project P7/P11. Used by: T14, T15, T16, T17, T19, T25, T28, T31.

\mdhead{2}{Numerical evidence (provisional)}

provisional evidence (Sol check J2): exact weighted-law propagation (N = 1, T = 1, 2, 3, through terminal resampling, burn-in and measured rounds with pre-draw rewards): max relative residual of EẐ = Z and E[Ẑ PIP̂\_j] = Z·PIP\_j is 1.3·10⁻¹⁵. Files: judge/sol2/J2\_EXACT\_IDENTITIES.md.

\mdhead{2}{Proof-route hint (optional)}

Induct on the unnormalised empirical measure through weighting, resampling and mutation (Feynman–Kac unbiasedness); append unit-potential invariant kernels; treat the pre-coordinate reward as a marked transition and replace the indicator by o\_j (project P11); tower property. Hardest step: the recording position within a round; terminal indicators cannot replace recorded conditionals. \emph{When proof.md is run on this file, its lead considers the statement both with and without these hints; the hints never reach the reviewers.}
\end{mdblock}

\mdsec{(b) STATEMENT.md (verbatim)}
\begin{mdblock}
\mdhead{1}{STATEMENT — T12 (T8): frozen-run evidence and recorded Rao–Blackwell identities}

Methodology: SMC-BVS v3.0 (METHODOLOGY\_v3.md, SHA-256 93043758058a6d422cca6ca4b6be8f3a9e8149b6f6de2e1bab1ce31c44b8019c), §§6.3–6.5, 6.8, 6.11, 6.12. Source conjecture: {\mdtt theorems/\allowbreak{}T12\_\allowbreak{}T8\_\allowbreak{}frozen\_\allowbreak{}run\_\allowbreak{}identities.\allowbreak{}md}. This file is the precise statement to be proved; nothing else is assumed.

\mdhead{2}{0. Data and target (fixed, deterministic objects)}

A \textbf{valid input} is (X\textasciicircum{}r, y\textasciicircum{}r) with finite entries, n ≥ 3, at least one retained (non-constant) column, and y\textasciicircum{}r not constant. X and y are the standardised data (columns and y centred and scaled so that every squared norm is n − 1). p ≥ 1 is the number of retained columns, K = min(p, ⌊n/2⌋) ≥ 1, and the state space is Ω = \{S ⊆ [p] : |S| ≤ K\} (finite, |Ω| ≥ 2).

For S ∈ Ω, with k = |S|, G\_S = X\_SᵀX\_S, z\_S = X\_Sᵀy, a₀ = b₀ = 0.01 and ν = a₀ + n/2 (the v3 target, METHODOLOGY §5 and 00\_COMMON §1):

\begin{mdbul}
\item H\_S = (n + 1)G\_S + I\_k, D\_S = det H\_S / det(G\_S + I\_k), Q\_S = yᵀy − n z\_SᵀH\_S⁻¹z\_S (for S = ∅: D\_∅ = 1, Q\_∅ = yᵀy);
\item m(S) = C\_n D\_S\textasciicircum{}\{−1/2\}(2b₀ + Q\_S)\textasciicircum{}\{−ν\} with C\_n = (2b₀)\textasciicircum{}\{a₀\}Γ(ν)/(π\textasciicircum{}\{n/2\}Γ(a₀)) > 0, so BF(S) := m(S)/m(∅) = D\_S\textasciicircum{}\{−1/2\}[(2b₀ + Q\_S)/(2b₀ + yᵀy)]\textasciicircum{}\{−ν\};
\item ϖ(S) = w(k)/W\_K with w(k) = B(1 + k, 2p − k)/B(1, p) = ½/C(2p − 1, k) and W\_K = Σ\_\{j≤K\} C(p, j)w(j) (the BB(1, p) prior capped at K);
\item Z = Σ\_\{S∈Ω\} ϖ(S)m(S), π(S) = ϖ(S)m(S)/Z, PIP\_j = π(\{S : j ∈ S\}).
\end{mdbul}

m > 0 (Q\_S > 0, as Q\_S ≥ yᵀy/(n + 1)) and ϖ > 0. Everything except the explicit instance of (iv.b)/(iv.e) uses only m > 0 and ϖ > 0. The instance uses the formulas above.

The following are deterministic functions of the data (method §§6.3–6.4):

\begin{mdbul}
\item blocks C₁, …, C\_J partitioning [p]; tables r\_b(U) = BF(U) > 0 for U ⊆ C\_b;
\item γ₀(S) = ϖ(S)∏\_b r\_b(S ∩ C\_b), Z̄₀ = Σ\_Ω γ₀, π₀ = γ₀/Z̄₀ (so π₀(S) > 0 for every S ∈ Ω);
\item the discrepancy d(S) = log BF(S) − Σ\_b log r\_b(S ∩ C\_b) (so d(∅) = 0); for real φ, γ\_φ = γ₀e\textasciicircum{}\{φd\}, Z̄\_φ = Σ\_Ω γ\_φ, q\_φ = γ\_φ/Z̄\_φ; note q₀ = π₀ and q₁ = π, and m(∅)Z̄₁ = Z;
\item the neighbour graph N(·), symmetric: l ∈ N(j) ⇔ j ∈ N(l); B = min(⌊√n⌋, ⌊√p⌋, 16).
\end{mdbul}

\mdhead{2}{1. Elementary steps and rounds (method §6.5)}

At temperature φ there are m = p + B + 2 labelled elementary steps, each a Markov kernel on Ω:

\begin{mdbul}
\item \textbf{Gibbs coordinate j} (j ∈ [p]): at S, with S⁺ = S ∪ \{j\} and S⁻ = S \textbackslash{} \{j\}: if |S⁺| ≤ K, move to S⁺ with probability o\_j(S, φ) = γ\_φ(S⁺)/(γ\_φ(S⁺) + γ\_φ(S⁻)) and to S⁻ otherwise; if |S⁺| > K, set o\_j(S, φ) = 0 and move to S⁻ (= S). The value o\_j(S, φ) is the \textbf{pre-draw recording} of the step.
\item \textbf{Uniform swap:} at S with |S| ∈ \{0, p\}, stay. Otherwise choose j ∈ S and l ∈ [p] \textbackslash{} S uniformly and independently, propose S′ = S \textbackslash{} \{j\} ∪ \{l\}, accept with probability 1 ∧ γ\_φ(S′)/γ\_φ(S), else stay.
\item \textbf{Neighbour swap} (B copies, each a separate label): at S = ∅, stay. Otherwise choose j ∈ S uniformly; if N(j) \textbackslash{} S = ∅, stay; else choose l uniformly in N(j) \textbackslash{} S, propose S′ = S \textbackslash{} \{j\} ∪ \{l\}, accept with probability 1 ∧ [γ\_φ(S′)/γ\_φ(S)]·|N(j) \textbackslash{} S|/|N(l) \textbackslash{} S′|, else stay.
\item \textbf{Independence proposal:} draw S′ \textasciitilde{} π₀, accept with probability 1 ∧ e\textasciicircum{}\{φ[d(S′) − d(S)]\}, else stay.
\end{mdbul}

A \textbf{round at φ} applied to a particle: draw an order σ (a bijection from positions \{1, …, m\} to the m labels) uniformly at random, independently of the particle's state, and apply the elementary steps σ(1), …, σ(m) in turn. The states X\_\{[0]\} (round start), X\_\{[1]\}, …, X\_\{[m]\} (round end) form the particle's \textbf{within-round path}. K\_φ denotes the resulting round kernel.

\mdhead{2}{2. The frozen run with fixed parameters}

\textbf{Frozen parameters} c consist of: a grid 0 = φ₀ < φ₁ < … < φ\_T = 1 with integer T ≥ 1 (Δ\_t = φ\_t − φ\_\{t−1\}, G\_t = e\textasciicircum{}\{Δ\_t d\}); an integer particle number N ≥ 1; an integer number L′ ≥ 0 of rounds at φ = 1 run after the stage-T round; an integer number M ≥ 1 of measured rounds; and a flag saying whether a round at φ₀ = 0 is applied at stage 0. \textbf{The method} (METHODOLOGY §6.8: "one move round per stage, then L burn-in rounds and 4 measured rounds") is the case N = N\_e, M = 4, L′ = L, flag unset. Other values are generalisations, covered by the same statements.

\textbf{One island I(c):} 1. Stage 0: X₀\textasciicircum{}i iid \textasciitilde{} π₀ (i = 1, …, N); Ẑ₀ = Z₀ := m(∅)Z̄₀; if the flag is set, each particle performs one round at φ₀. Call the population X\textasciicircum{}\{(0)\}. 2. Stage t = 1, …, T: weights w\textasciicircum{}i = G\_t(X\textasciicircum{}\{(t−1),i\}); Ẑ\_t = Ẑ\_\{t−1\}·(1/N)Σ\_i w\textasciicircum{}i; multinomial resampling (ancestors A¹, …, A\textasciicircum{}N iid with P(A = k) = w\textasciicircum{}k/Σ\_l w\textasciicircum{}l given the past); each resampled particle X\textasciicircum{}\{(t−1),A\textasciicircum{}i\} performs one round at φ\_t; call the population X\textasciicircum{}\{(t)\}. q̂\_t := (1/N)Σ\_i δ\_\{X\textasciicircum{}\{(t),i\}\}; q̂₀ := (1/N)Σ\_i δ\_\{X\textasciicircum{}\{(0),i\}\}. 3. After stage T: L′ rounds at φ = 1, then M measured rounds at φ = 1. In measured round u, Y\textasciicircum{}i\_\{u,j\} is particle i's state immediately before its Gibbs-j step. Island evidence Ẑ := Ẑ\_T; island RB estimate PIP̂\_j := (MN)⁻¹Σ\_\{u=1\}\textasciicircum{}\{M\}Σ\_\{i=1\}\textasciicircum{}\{N\} o\_j(Y\textasciicircum{}i\_\{u,j\}, 1).

Stage notation: q\_t = q\_\{φ\_t\}, Z\_t = m(∅)Z̄\_\{φ\_t\}; so Z\_T = Z.

\textbf{Randomness (from EX-S and the prescribed operations).} (R1) the stage-0 draws are iid π₀; (R2) given everything before a resampling, the ancestors are iid with the stated law; (R3) given the σ-field H of everything before a round (including the states entering it), the particles' (order, within-round path) pairs are conditionally independent across particles, each order is uniform and independent of H, and given H and its order σ, a particle's within-round path is a Markov chain started at its entering state with kernels K\_\{σ(1)\}, …, K\_\{σ(m)\} (each step uses fresh randomness; the independence proposal is a fresh exact π₀ draw).

\textbf{Hypotheses.}

\begin{mdbul}
\item \textbf{EX-A:} exact arithmetic: every score, table, Z̄₀, probability and comparison is exact (no numerical-failure status occurs).
\item \textbf{START} (imported dependency: T09(i) / project P1, not re-proved here): the start sampler returns exact iid π₀ draws, the reported Z̄₀ equals Σ\_Ω γ₀, and the independence proposals use the same exact sampler. This is not an extra restriction on inputs: the definition of I(c) above already uses these properties, and START is what makes the method's final island an instance of I(c). EX-A alone does not give it.
\item \textbf{EX-S:} independent ideal uniform streams with the method's stream separation; in particular (R1)–(R3). The stream model includes measurability: the map (c, u) ↦ Φ(c, u) below, from (calibration value, stream) to the island record, is jointly measurable, and the record takes values in a Polish space. This holds for any algorithmic realisation from uniform streams.
\item \textbf{CAL:} the calibration output 𝓒 (grid, T, L, N\_e, tier, candidate set and every other setting the final run uses) is measurable with respect to the data and streams 0–9 only; the final islands r = 1, …, 4 use streams 10–13, which are independent of each other and of streams 0–9; escalation is decided before any final randomness.
\item \textbf{Success:} 𝒜 ∈ σ(𝓒) is the event that calibration succeeds (a grid is constructed at the final tier, so the final run executes). On 𝒜 the final islands are O\_r = Φ(𝓒, U\_r), r = 1, …, 4, where U\_r is stream 9 + r and Φ(c, ·) realises I(c) with the method's parameters (N = N\_e, M = 4, L′ = L, flag unset). "Given 𝓒" means conditional on σ(𝓒) on the event 𝒜. Island r's outputs are Ẑ\_r, PIP̂\_\{rj\}, and its stage quantities Ẑ\_\{r,t\}, q̂\_\{r,t\}.
\item \textbf{Pooled estimators:} Ẑ = ¼Σ\_rẐ\_r; primary P̄\_j = ¼Σ\_r PIP̂\_\{rj\}; secondary P̃\_j = Σ\_rẐ\_rPIP̂\_\{rj\}/Σ\_rẐ\_r.
\end{mdbul}

\mdhead{2}{3. Conclusions}

Let E\_c denote expectation for the frozen run with fixed parameters c. For every valid input:

\textbf{(i) Evidence identities.}

\begin{mdbul}
\item (i.a) For every c, t ∈ \{0, …, T\} and f : Ω → ℝ: E\_c[Ẑ\_t q̂\_t(f)] = Z\_t q\_t(f). In particular E\_c[Ẑ] = Z.
\item (i.b) \emph{Terminal invariant rounds with unit potentials.} After stage T, let any fixed finite sequence of operations be applied, each either a round at φ = 1 or a multinomial resampling with unit weights (which leaves Ẑ unchanged). For the population q̂ after any of these operations, E\_c[Ẑ q̂(f)] = Z π(f).
\item (i.c) \emph{State-independent prefixes.} Consider any round of the run (a stage-t round, with Ẑ\_t and target q\_t, or a post-terminal round, with Ẑ\_T and target π) and, for each particle i, a prefix length κ\_i = κ(σ\_i) ∈ \{0, …, m\} that is a fixed function of that particle's order σ\_i alone (examples: a fixed k; κ(σ) = σ⁻¹(ℓ) − 1, the position just before a specified label ℓ). Then E\_c[Ẑ\_t (1/N)Σ\_i f(X\textasciicircum{}i\_\{[κ\_i]\})] = Z\_t q\_t(f).
\item (i.d) \emph{Not at state-dependent stopping times.} For every c and every round of the run (with its Ẑ\_t, q\_t as in (i.c)), there are f and a stopping time τ\_i of particle i's within-round path for which E\_c[Ẑ\_t (1/N)Σ\_i f(X\textasciicircum{}i\_\{[τ\_i]\})] > Z\_t q\_t(f). Namely τ\_i = min\{k ∈ \{0, …, m\} : X\textasciicircum{}i\_\{[k]\} = ∅\} ∧ m, with min ∅ = ∞, and f = 1\{∅\}.
\end{mdbul}

\textbf{(ii) Recorded RB identity.} For every c and j ∈ [p]: E\_c[Ẑ·PIP̂\_j] = Z·PIP\_j.

\textbf{(iii) Pooled evidence.} Under EX-A, START, EX-S, CAL, given 𝓒 on 𝒜: for every r and every stage t ≤ T of 𝓒's grid, E[Ẑ\_\{r,t\}q̂\_\{r,t\}(f) | 𝓒] = Z\_tq\_t(f) (Z\_t and q\_t are those of 𝓒's grid, so this right side is 𝓒-measurable) and E[Ẑ\_rPIP̂\_\{rj\} | 𝓒] = Z·PIP\_j; E[Ẑ | 𝓒] = Z; E[log Ẑ | 𝓒] ≤ log Z and E[log Ẑ\_r | 𝓒] ≤ log Z; and for every 0 < α < 1, P(Z ≥ αẐ | 𝓒) ≥ 1 − α and P(Z ≥ αẐ\_r | 𝓒) ≥ 1 − α.

\textbf{(iv) Limits.}

\begin{mdbul}
\item (iv.a) \emph{log Ẑ.} For every c: if d ≡ 0 on Ω then Ẑ = Z almost surely; otherwise E\_c[log Ẑ] < log Z, and, for the four islands of the method given 𝓒, E[log Ẑ | 𝓒] < log Z and E[log Ẑ\_r | 𝓒] < log Z.
\item (iv.b) \emph{Normalised PIPs.} Given 𝓒, E[P̄\_j | 𝓒] − PIP\_j = −Cov(PIP̂\_\{1j\}, Ẑ₁ | 𝓒)/Z. \emph{Component example.} There is a valid input and frozen parameters c for which both normalised estimators are biased: n = 20, p = 2 with the data of the appendix, T = 1, N = 1, M = 4, L′ = L = max(4, ⌈log₂(n + 1)⌉) = 5, flag unset, four independent islands I(c), and j = 1. Then E[P̄₁] − PIP₁ < −0.17 and E[P̃₁] − PIP₁ < −0.15. The same bounds hold with the flag set and/or L′ = L − 1 (generalisations). Scope: at p = 2 the full method would take the exact branch, and N = 1 is not the production budget. So this shows bias of the frozen SMC component, not bias of the full method's output on this input.
\item (iv.c) \emph{Island independence.} Given 𝓒, the islands O₁, …, O₄ are independent, each with the law of I(𝓒). Assume P(𝒜) > 0. Under P(· | 𝒜), for every bounded measurable h of an island's outputs, Cov(h(O₁), h(O₂)) = Var(E[h(O₁) | 𝓒]). Hence the islands are independent under P(· | 𝒜) if and only if the conditional law of an island given 𝓒 is a.s. constant, which none of the hypotheses ensure. In particular the island law does depend on c: if d has a unique maximiser on Ω, the law of Ẑ under (grid, N) differs from that under (grid, 4N). (Not claimed: that the actual calibration procedure produces, on some specific input, two values of 𝓒 with different island laws, each with positive probability.)
\item (iv.d) \emph{Launch events.} Let A ∈ σ(𝓒) with A ⊆ 𝒜 and P(A) > 0. The conclusions of (iii) whose right side does not depend on 𝓒 hold under P(· | A): E[Ẑ | A] = E[Ẑ\_r | A] = Z, E[Ẑ\_rPIP̂\_\{rj\} | A] = Z·PIP\_j, E[log Ẑ | A] ≤ log Z, E[log Ẑ\_r | A] ≤ log Z, P(Z ≥ αẐ | A) ≥ 1 − α and P(Z ≥ αẐ\_r | A) ≥ 1 − α. For an intermediate stage t and A ⊆ \{T ≥ t\}, the averaged form holds: E[Ẑ\_\{r,t\}q̂\_\{r,t\}(f) | A] = E[Z\_tq\_t(f) | A], where Z\_tq\_t(f) is computed on 𝓒's grid. The left side is a number, whereas Z\_tq\_t(f) is a random variable, and it is non-constant when, for example, φ\_t takes two values with positive probability under P(· | A) and f = 1\{S\} with d(S) ≠ 0.
\item (iv.e) \emph{Favourable final diagnostics.} Let F be the event that the final E2 check passes: the standard deviation of \{log Ẑ\_r\}\_\{r=1\}\textasciicircum{}\{4\} is ≤ 1 nat (either the 1/4 or the 1/3 normalisation). In the component example of (iv.b) (same scope), P(F) > 0, E[Ẑ\_r | F] < Z, E[Ẑ | F] < Z and E[Ẑ\_r PIP̂\_\{r1\} | F] ≠ Z·PIP₁. So conditioning on a favourable final diagnostic breaks an identity in each of (i) (E Ẑ\_r = Z), (ii) and (iii) (E Ẑ = Z). This is not a claim that the Jensen or Markov inequalities of (iii) fail under F; in this example they still hold, because E[Ẑ | F] < Z.
\item (iv.f) \emph{NUM alone.} A NUM-compatible perturbed-score model with bound η > 0 is the same algorithm in exact arithmetic with every evaluated log m(S) replaced by a value within η of it. (This does not assert that any particular floating-point implementation makes this error; it shows that NUM alone does not exclude it.) For every η > 0 there is such an implementation (log m(S) + η for every S) with E\_c[Ẑ\textasciicircum{}impl] = e\textasciicircum{}ηZ ≠ Z for every c and valid input; and on every valid input with d ≡ 0 (for example p = 1), P(Z ≥ αẐ\textasciicircum{}impl) = 0 for every α ∈ (e\textasciicircum{}\{−η\}, 1).
\end{mdbul}

\mdhead{2}{Appendix: data for (iv.b) and (iv.e)}

n = 20, p = 2, i = 0, …, 19: x1\_i = (7i mod 20) − 9.5; x2\_i = x1\_i + ((5i² mod 7) − 3); y\_i = x1\_i − x2\_i + ((3i mod 5) − 2). (Raw data; standardised as in §0.) Then B = 1, the blocks are \{1\}, \{2\}, N(1) = \{2\}, N(2) = \{1\}, K = 2, m = 5, L = 5.

\mdhead{2}{Statement history}

\begin{mdbul}
\item v1 (2026-10-22): written by the lead from the conjecture. The formalisation choices are listed in review/packets/STATEMENT\_REPORT\_A1\_call1.md.
\item v2 (2026-10-22): statement check by Astra \#1 (review/A1\_call1\_statement\_check\_REPLY.md), ruling APPROVED WITH CHANGES. All four required changes adopted: 1. (iv.d) repaired. v1 claimed every conclusion of (iii) under P(· | A). That is false for intermediate stages, whose targets depend on the random grid; only the 𝓒-free conclusions survive, and intermediate stages hold in averaged form. 2. The method is identified as L′ = L, flag unset (METHODOLOGY §6.8); other values are labelled generalisations; count parameters are declared integers. 3. START is labelled an imported dependency (T09(i)), not a new restriction. 4. Component scope stated beside (iv.b)/(iv.e). (iv.e) is stated as breaking the identities, not the Jensen/Markov inequalities. Also, following Astra's item 5: min ∅ = ∞ and time 0 are made explicit in (i.d).
\item v2.1 (2026-10-22): after Astra \#1's proof review (verdict verified; three exposition findings), the assumption P(𝒜) > 0 was made explicit in (iv.c). This is exposition only.
\item v2.2 (2026-10-22): after the Stage B reviews (Astra \#2 and Opus xhigh; both raised the same gap). §0 now writes out the target formulas (m(S), BF, ϖ, w, W\_K, a₀ = b₀ = 0.01, ν = a₀ + n/2), copied from METHODOLOGY §5 / 00\_COMMON §1, Felipe's fixed target. EX-S now states that joint measurability of Φ is part of the stream model. In (iv.d), 'only … in general random' is replaced by an exact remark. No conclusion changed.
\end{mdbul}
\end{mdblock}

\mdsec{(c) Final proof, verbatim: body of PROOF\_v6.tex}
\noindent{\small\itshape The proof below is the body of \texttt{PROOF\_v6.tex} from \texttt{\textbackslash begin\{document\}} to \texttt{\textbackslash end\{document\}} without the title block. Its preamble definitions are local to this chapter; labels carry the prefix \texttt{T12:}.}\par\medskip
\mdproofbegin
% ---- preamble of PROOF_v6.tex (theorem environments carry the prefix T12) ----
\lstset{basicstyle=\ttfamily\scriptsize,breaklines=true,columns=fullflexible,frame=single}

\newtheorem{T12theorem}{Theorem}
\newtheorem{T12lemma}{Lemma}[section]
\newtheorem{T12corollary}[T12lemma]{Corollary}
\theoremstyle{remark}
\newtheorem{T12remark}[T12lemma]{Remark}

\newcommand{\E}{\mathbb{E}}
\renewcommand{\P}{\mathbb{P}}
\newcommand{\R}{\mathbb{R}}
\newcommand{\cC}{\mathcal{C}}
\newcommand{\cA}{\mathcal{A}}
\newcommand{\cF}{\mathcal{F}}
\newcommand{\cH}{\mathcal{H}}
\newcommand{\Zh}{\widehat Z}
\newcommand{\qh}{\widehat q}
\newcommand{\PIPh}{\widehat{\mathrm{PIP}}}
\newcommand{\PIP}{\mathrm{PIP}}
\newcommand{\Zb}{\bar Z}
\newcommand{\ind}{\mathbf 1}
\newcommand{\Om}{\Omega}
% ---- body of PROOF_v6.tex ----
\noindent We prove every conclusion of \texttt{STATEMENT.md}, with its definitions, notation and hypotheses (EX-A, START, EX-S, CAL, the randomness model (R1)--(R3), the success event $\cA$). Section numbers refer to this document. All functions on $\Om$ are bounded because $\Om$ is finite, and every random quantity below is bounded, so all expectations exist. For a function $h$ on $\Om$ and a Markov kernel $P$ on $\Om$, $Ph(x)=\sum_y P(x,y)h(y)$; for a measure $\mu$ on $\Om$, $\mu(h)=\sum_x\mu(x)h(x)$ and $\mu P(y)=\sum_x\mu(x)P(x,y)$. We write $J_j=\{S\in\Om: j\in S\}$, so $\PIP_j=\pi(J_j)=\pi(\ind_{J_j})$. Parts (i), (ii) and (iv.a)--(iv.b) are proved for the frozen run $I(c)$ with arbitrary integer parameters. The method's final island is the case $N=N_e$, $M=4$, $L'=L$ with the stage-0 flag unset; other values are generalisations. START is an imported dependency (T09(i)): it is what makes the stage-0 draws and the independence proposals exact $\pi_0$ draws and $\Zb_0$ exact, as the definition of $I(c)$ requires.

\section{Deterministic facts}

\begin{T12lemma}[bridge identities]\label{T12:lem:bridge}
For every real $\varphi$, $\Zb_\varphi\in(0,\infty)$ and $q_\varphi(S)>0$ for every $S\in\Om$. For $t=1,\dots,T$ and every $h:\Om\to\R$,
\[ Z_{t-1}\,q_{t-1}(G_t h)=Z_t\,q_t(h). \]
Moreover $m(\emptyset)\Zb_1=Z$, so $Z_T=Z$; and $Z_0=m(\emptyset)\Zb_0$.
\end{T12lemma}
\begin{proof}
$\gamma_\varphi(S)=\gamma_0(S)e^{\varphi d(S)}>0$ because $\varpi>0$ and $r_b>0$; $\Om$ is finite, so $\Zb_\varphi\in(0,\infty)$. Since $\varphi_{t-1}+\Delta_t=\varphi_t$,
$q_{t-1}(x)G_t(x)=\gamma_0(x)e^{\varphi_{t-1}d(x)}e^{\Delta_t d(x)}/\Zb_{\varphi_{t-1}}=\gamma_{\varphi_t}(x)/\Zb_{\varphi_{t-1}}$. Hence
$Z_{t-1}q_{t-1}(G_th)=m(\emptyset)\Zb_{\varphi_{t-1}}\sum_x \gamma_{\varphi_t}(x)h(x)/\Zb_{\varphi_{t-1}}=m(\emptyset)\Zb_{\varphi_t}q_t(h)=Z_tq_t(h)$.
Finally $\gamma_0(S)e^{d(S)}=\varpi(S)\prod_b r_b(S\cap C_b)\cdot \mathrm{BF}(S)/\prod_b r_b(S\cap C_b)=\varpi(S)m(S)/m(\emptyset)$, so $m(\emptyset)\Zb_1=\sum_S\varpi(S)m(S)=Z$.
\end{proof}

\begin{T12lemma}[elementary steps are $q_\varphi$-reversible]\label{T12:lem:rev}
Fix a real $\varphi$. Each elementary step of STATEMENT \S1 (Gibbs coordinate $j$, uniform swap, neighbour swap, independence proposal) defines a Markov kernel $K_e$ on $\Om$ that satisfies detailed balance with respect to $q_\varphi$: $q_\varphi(S)K_e(S,S')=q_\varphi(S')K_e(S',S)$ for all $S,S'\in\Om$. Consequently $q_\varphi K_e=q_\varphi$.
\end{T12lemma}
\begin{proof}
Write $q=q_\varphi$, $\gamma=\gamma_\varphi$, $\Zb=\Zb_\varphi$. Detailed balance for $S=S'$ is trivial. Summing it over $S$ gives $q K_e(S')=q(S')\sum_SK_e(S',S)=q(S')$. So it suffices to check $S\ne S'$, after checking that each step stays in $\Om$ with total probability one.

\emph{Gibbs coordinate $j$.} At $S$ the step moves to $S^+=S\cup\{j\}$ (only if $|S^+|\le K$) or to $S^-=S\setminus\{j\}$, with probabilities $o_j(S,\varphi)$ and $1-o_j(S,\varphi)$. Both targets are in $\Om$ whenever they are used, since $|S^-|\le |S|\le K$. The only transitions between distinct states are between some pair $\{U,U\cup\{j\}\}$ with $j\notin U$ and $|U|+1\le K$. From either member of such a pair, $S^+=U\cup\{j\}$ and $S^-=U$, so $K_e(U,U\cup\{j\})=\gamma(U\cup\{j\})/(\gamma(U\cup\{j\})+\gamma(U))$ and $K_e(U\cup\{j\},U)=\gamma(U)/(\gamma(U\cup\{j\})+\gamma(U))$. Then
$q(U)K_e(U,U\cup\{j\})=\gamma(U)\gamma(U\cup\{j\})/(\Zb(\gamma(U\cup\{j\})+\gamma(U)))=q(U\cup\{j\})K_e(U\cup\{j\},U)$. No other pair of distinct states communicates, and every other $K_e(S,S')$ with $S\ne S'$ is $0$ in both directions.

\emph{Uniform swap.} If $|S|\in\{0,p\}$ the step stays at $S$. Otherwise, with $k=|S|\in(0,p)$, the pair $(j,l)\in S\times([p]\setminus S)$ has probability $1/(k(p-k))$. It proposes $S'=S\setminus\{j\}\cup\{l\}$, with $|S'|=k\le K$, so $S'\in\Om$ and $S'\neq S$. The map $(j,l)\mapsto S'$ is injective (its inverse is $S'\mapsto(S\setminus S',\,S'\setminus S)$). Hence for $S'\ne S$, $K_e(S,S')=\frac{1}{k(p-k)}\min\{1,\gamma(S')/\gamma(S)\}$ if $|S\triangle S'|=2$ and $|S'|=|S|$, and $K_e(S,S')=0$ otherwise. This condition is symmetric in $(S,S')$ and $|S'|=k$, so
$q(S)K_e(S,S')=\frac{1}{\Zb\,k(p-k)}\min\{\gamma(S),\gamma(S')\}=q(S')K_e(S',S)$.

\emph{Neighbour swap.} At $S=\emptyset$ the step stays. Otherwise, with $k=|S|\ge1$, choose $j\in S$ with probability $1/k$. Let $A=N(j)\setminus S$. If $A=\emptyset$, stay. Else choose $l\in A$ with probability $1/|A|$ and propose $S'=S\setminus\{j\}\cup\{l\}$; then $|S'|=k$, $S'\in\Om$, $S'\ne S$. Put $A'=N(l)\setminus S'$. Since $l\in N(j)$, the symmetry of the graph gives $j\in N(l)$, and $j\notin S'$, so $j\in A'$ and $|A'|\ge1$; the acceptance probability $\min\{1,\gamma(S')|A|/(\gamma(S)|A'|)\}$ is well defined. The map $(j,l)\mapsto S'$ is injective as before. Hence for $S'\neq S$, $K_e(S,S')=\frac1{k|A|}\min\{1,\frac{\gamma(S')|A|}{\gamma(S)|A'|}\}$ when $S'=S\setminus\{j\}\cup\{l\}$ with $j\in S$ and $l\in N(j)\setminus S$, and $0$ otherwise. If $S'$ arises from $S$ through $(j,l)$, then $S$ arises from $S'$ through $(l,j)$: $l\in S'$, $j\in N(l)\setminus S'=A'$, and the sets used in the reverse move are $N(l)\setminus S'=A'$ and $N(j)\setminus S=A$. Hence $K_e(S',S)=\frac1{k|A'|}\min\{1,\frac{\gamma(S)|A'|}{\gamma(S')|A|}\}$, and
\[ q(S)K_e(S,S')=\frac{1}{\Zb k}\min\Big\{\frac{\gamma(S)}{|A|},\frac{\gamma(S')}{|A'|}\Big\}=q(S')K_e(S',S). \]
If $S'$ does not arise from $S$, then $S$ does not arise from $S'$ (by the same reversal), and both sides are $0$.

\emph{Independence proposal.} $\pi_0$ is supported on $\Om$. For $S'\ne S$, $K_e(S,S')=\pi_0(S')\min\{1,e^{\varphi(d(S')-d(S))}\}$. Since $q(S)=\pi_0(S)e^{\varphi d(S)}\Zb_0/\Zb$,
$q(S)K_e(S,S')=\frac{\Zb_0}{\Zb}\pi_0(S)\pi_0(S')\min\{e^{\varphi d(S)},e^{\varphi d(S')}\}$, which is symmetric.
\end{proof}

\begin{T12lemma}[state-independent prefixes; recordings]\label{T12:lem:prefix}
Fix a real $\varphi$, let $\Sigma$ be the set of orders (bijections from positions $\{1,\dots,m\}$ to labels), and let $\kappa:\Sigma\to\{0,\dots,m\}$. Define
\[ P_\kappa=\frac1{m!}\sum_{\sigma\in\Sigma}K_{\sigma(1)}K_{\sigma(2)}\cdots K_{\sigma(\kappa(\sigma))}\qquad(\text{empty product}=\text{identity}). \]
Then $q_\varphi P_\kappa=q_\varphi$. In particular the round kernel $K_\varphi=P_m$ is $q_\varphi$-invariant. Moreover, if $K_{g_j}$ is the Gibbs-$j$ kernel, then $K_{g_j}\ind_{J_j}=o_j(\cdot,\varphi)$ and $q_\varphi(o_j(\cdot,\varphi))=q_\varphi(J_j)$.
\end{T12lemma}
\begin{proof}
By Lemma~\ref{T12:lem:rev}, every $K_e$ preserves $q_\varphi$, so every finite product does, and so does every average of products. For the second claim, at $S$ the Gibbs-$j$ step ends in a state containing $j$ exactly when it moves to $S^+$, which happens with probability $o_j(S,\varphi)$ (this is $0$ when $|S^+|>K$, where the step moves to $S^-$). Hence $K_{g_j}\ind_{J_j}=o_j(\cdot,\varphi)$, and $q_\varphi(o_j)=(q_\varphi K_{g_j})(\ind_{J_j})=q_\varphi(J_j)$.
\end{proof}

\begin{T12lemma}[positivity]\label{T12:lem:pos}
Fix a real $\varphi$ and let $R_d=\max_\Om d-\min_\Om d$.
\begin{enumerate}[label=(\alph*)]
\item For every order $\sigma$ in which the independence label is last, and all $x,S\in\Om$, $K_{\sigma(1)}\cdots K_{\sigma(m)}(x,\{S\})\ge \pi_0(S)e^{-|\varphi|R_d}$. Hence $K_\varphi(x,\{S\})\ge \pi_0(S)e^{-|\varphi|R_d}/m$. In particular, for every $\varphi\in[0,1]$ (all temperatures used by the run), $K_\varphi(x,\{S\})\ge\varepsilon_S/m$ with the temperature-free constant $\varepsilon_S:=\pi_0(S)e^{-R_d}>0$.
\item For every order $\sigma$ with $\sigma(1)=$ Gibbs coordinate $1$, $K_{\sigma(1)}\cdots K_{\sigma(m)}(\emptyset,\Om\setminus\{\emptyset\})\ge \delta>0$ for a constant $\delta$ that depends only on $\varphi$ and the input. Hence $K_\varphi(\emptyset,\Om\setminus\{\emptyset\})\ge \delta/m>0$.
\end{enumerate}
\end{T12lemma}
\begin{proof}
(a) Whatever the state $y$ before the last step, the independence step proposes $S$ with probability $\pi_0(S)$. If $S\ne y$ it accepts with probability $\min\{1,e^{\varphi(d(S)-d(y))}\}\ge e^{-|\varphi|R_d}$; if $S=y$ the state stays $S$ whether or not it accepts. Hence $K_{\rm ind}(y,\{S\})\ge\pi_0(S)e^{-|\varphi|R_d}$ for every $y$, and the earlier steps are stochastic. For $\varphi\in[0,1]$, $e^{-|\varphi|R_d}\ge e^{-R_d}$. Exactly $(m-1)!$ of the $m!$ orders have the independence label last.

(b) At $\emptyset$, Gibbs coordinate $1$ moves to $\{1\}$, which is in $\Om$ since $K\ge1$, with probability $o_1(\emptyset,\varphi)=\gamma_\varphi(\{1\})/(\gamma_\varphi(\{1\})+\gamma_\varphi(\emptyset))>0$. For every elementary step $e$ and every non-empty $y$, $c_e(y):=K_e(y,\Om\setminus\{\emptyset\})>0$. Gibbs $j$ with $j\notin y$ ends in $y$ or $y\cup\{j\}$, so $c_e(y)=1$. Gibbs $j$ with $j\in y$ stays at $y$ with probability $o_j(y,\varphi)>0$, because $y^+=y\in\Om$. Swaps preserve $|y|$ or stay, so $c_e(y)=1$. The independence step stays at $y$ with probability at least $\pi_0(y)>0$. Let $c=\min_e\min_{y\ne\emptyset}c_e(y)>0$ (a minimum over finitely many positive numbers). By the Markov property along the fixed order, the probability of being non-empty after step 1 and of remaining non-empty after each later step is at least $\delta:=o_1(\emptyset,\varphi)\,c^{m-1}$. Exactly $(m-1)!$ orders have Gibbs coordinate $1$ first.
\end{proof}

\section{Conditional means of the island operations}

All random variables of one island live on one probability space. For a population $(X^1,\dots,X^N)$ entering an operation, $\cF$ denotes the $\sigma$-field of everything generated before the operation.

\begin{T12lemma}[resampling]\label{T12:lem:res}
Let the population $(X^1,\dots,X^N)$ and weights $w^k=G(X^k)>0$ be $\cF$-measurable, and let the ancestors $A^1,\dots,A^N$ be conditionally iid given $\cF$ with $\P(A^i=k\mid\cF)=w^k/\sum_lw^l$ (R2). Then for every $h:\Om\to\R$,
\[ \E\Big[\Big(\tfrac1N\textstyle\sum_lw^l\Big)\tfrac1N\textstyle\sum_i h(X^{A^i})\,\Big|\,\cF\Big]=\tfrac1N\textstyle\sum_k G(X^k)h(X^k). \]
\end{T12lemma}
\begin{proof}
$\E[h(X^{A^i})\mid\cF]=\sum_kh(X^k)w^k/\sum_lw^l$ for each $i$. Multiply by the $\cF$-measurable factor $\frac1N\sum_lw^l$ and average over $i$.
\end{proof}

\begin{T12lemma}[one round, state-independent prefix]\label{T12:lem:round}
Let a population $(\tilde X^1,\dots,\tilde X^N)$ enter a round at temperature $\varphi$, let $\cH$ be the $\sigma$-field of everything before the round, and let $X^i_{[k]}$ be particle $i$'s within-round path. For $\kappa:\Sigma\to\{0,\dots,m\}$ and every $f:\Om\to\R$, under (R3), $\E[f(X^i_{[\kappa(\sigma_i)]})\mid\cH]=(P_\kappa f)(\tilde X^i)$.
\end{T12lemma}
\begin{proof}
By (R3), given $\cH$ and $\sigma_i=\sigma$, the path is a Markov chain started at $X^i_{[0]}=\tilde X^i$ with kernels $K_{\sigma(1)},\dots,K_{\sigma(m)}$. Hence $\E[f(X^i_{[k]})\mid\cH,\sigma_i=\sigma]=(K_{\sigma(1)}\cdots K_{\sigma(k)}f)(\tilde X^i)$ for each fixed $k$, and in particular for $k=\kappa(\sigma)$. Since $\sigma_i$ is uniform on the finite set $\Sigma$ and independent of $\cH$, averaging over $\sigma$ gives the claim.
\end{proof}

\section{Proof of (i.a)--(i.c)}

Write $X^{(t-1),A^i}=:\tilde X^{(t),i}$ for the resampled particles of stage $t\ge1$, and $\tilde X^{(0),i}:=X_0^i$.

\paragraph{Claim $\mathsf C_t$ ($t=0,\dots,T$).} For every $h:\Om\to\R$: (a) $\E[\Zh_t\frac1N\sum_ih(\tilde X^{(t),i})]=Z_tq_t(h)$, and (b) $\E[\Zh_t\frac1N\sum_ih(X^{(t),i})]=Z_tq_t(h)$. Moreover (c), for every round of stage $t$ (the stage-$t$ round, or the stage-0 round if the flag is set) and every $\kappa$, $\E[\Zh_t\frac1N\sum_i f(X^i_{[\kappa(\sigma_i)]})]=Z_tq_t(f)$.

\emph{$t=0$.} $\Zh_0=Z_0$ is deterministic and the $X_0^i$ are iid $\pi_0=q_0$ by (R1) and START, so (a) holds. If no stage-0 round is applied, $X^{(0)}=X_0$ and (b) is (a). If a round at $\varphi_0=0$ is applied, then for its prefix with $\kappa$, Lemma~\ref{T12:lem:round} and (a) give $\E[Z_0\frac1N\sum_if(X^i_{[\kappa]})]=\E[Z_0\frac1N\sum_i(P_\kappa f)(X_0^i)]=Z_0q_0(P_\kappa f)=Z_0q_0(f)$ by Lemma~\ref{T12:lem:prefix}. This is (c), and $\kappa\equiv m$ gives (b).

\emph{$t-1\to t$.} Let $\cF$ be everything before the stage-$t$ resampling. $\Zh_t=\Zh_{t-1}\frac1N\sum_lw^l$ with $w^l=G_t(X^{(t-1),l})$, and $\Zh_{t-1}$ is $\cF$-measurable. By Lemma~\ref{T12:lem:res},
$\E[\Zh_t\frac1N\sum_ih(\tilde X^{(t),i})\mid\cF]=\Zh_{t-1}\frac1N\sum_k(G_th)(X^{(t-1),k})$.
Taking expectations and applying $\mathsf C_{t-1}$(b) to the function $G_th$ gives $Z_{t-1}q_{t-1}(G_th)=Z_tq_t(h)$ by Lemma~\ref{T12:lem:bridge}. This is (a). For the stage-$t$ round with prefix $\kappa$, $\Zh_t$ is measurable with respect to $\cH$ (everything before the round), so Lemma~\ref{T12:lem:round} gives
$\E[\Zh_t\frac1N\sum_if(X^i_{[\kappa(\sigma_i)]})\mid\cH]=\Zh_t\frac1N\sum_i(P_\kappa f)(\tilde X^{(t),i})$.
By (a) with $h=P_\kappa f$ and Lemma~\ref{T12:lem:prefix} (at $\varphi_t$), its expectation is $Z_tq_t(P_\kappa f)=Z_tq_t(f)$. This is (c), and $\kappa\equiv m$ gives (b). \qed

Claim $\mathsf C_t$(b) is (i.a). With $t=T$ and $h\equiv1$, $\E\Zh=Z_Tq_T(1)=Z$ by Lemma~\ref{T12:lem:bridge}. Claim $\mathsf C_t$(c) is (i.c) for the rounds of stages $0,\dots,T$.

\paragraph{Post-terminal operations: (i.b) and (i.c).} Let $\mathcal O_1,\dots,\mathcal O_k$ be the fixed sequence of operations after stage $T$. Each is a round at $\varphi=1$ or a multinomial resampling with unit weights. Write $\qh^{(s)}$ for the population after $\mathcal O_s$ and $\qh^{(0)}=\qh_T$. We show by induction on $s$ that $\E[\Zh\,\qh^{(s)}(h)]=Z\pi(h)$ for all $h$; for $s=0$ this is (i.a) at $t=T$, since $q_T=\pi$. A unit-weight resampling leaves $\Zh$ unchanged (its factor is $\frac1N\sum_l1=1$), and by Lemma~\ref{T12:lem:res} with $G\equiv1$, $\E[\qh^{(s)}(h)\mid\cF]=\qh^{(s-1)}(h)$, with $\Zh$ $\cF$-measurable. A round at $\varphi=1$ with prefix $\kappa$ gives, by Lemma~\ref{T12:lem:round}, $\E[\Zh\frac1N\sum_if(X^i_{[\kappa(\sigma_i)]})\mid\cH]=\Zh\,\qh^{(s-1)}(P_\kappa f)$. By the induction hypothesis and Lemma~\ref{T12:lem:prefix} at $\varphi=1$, its expectation is $Z\pi(P_\kappa f)=Z\pi(f)$. This is (i.c) for that round, and $\kappa\equiv m$ gives the step of the induction. This proves (i.b), and (i.c) for every post-terminal round. In particular it covers the $L'$ rounds and the $M$ measured rounds of $I(c)$, which are post-terminal rounds at $\varphi=1$. \qed

\section{Proof of (i.d): state-dependent stopping times}

Fix $c$ and a round $\rho$ of the run, with prefactor $\Zh_\rho$ (that is, $\Zh_t$ for a stage-$t$ round and $\Zh_T$ for a post-terminal round) and target $q_\rho$ (that is, $q_t$ or $\pi$). Let $\cH$ be everything before $\rho$, $\tilde X^i$ the entering states, and $\tau_i=\min\{k\in\{0,\dots,m\}: X^i_{[k]}=\emptyset\}\wedge m$ (with $\min\emptyset=\infty$, so $\tau_i=m$ when the path never visits $\emptyset$). Then $\tau_i$ is a stopping time for $\cF^i_k=\sigma(\cH,\sigma_i,X^i_{[0]},\dots,X^i_{[k]})$. Let $f=\ind\{\emptyset\}$.

\emph{Step 1: $D_i:=f(X^i_{[\tau_i]})-f(X^i_{[m]})\ge \ind_{E_i}\ge 0$, where $E_i=\{X^i_{[0]}=\emptyset,\,X^i_{[m]}\ne\emptyset\}$.} If the path visits $\emptyset$, then $f(X^i_{[\tau_i]})=1\ge f(X^i_{[m]})$. If it does not, $\tau_i=m$ and $D_i=0$. On $E_i$, $\tau_i=0$, $f(X^i_{[0]})=1$ and $f(X^i_{[m]})=0$.

\emph{Step 2: $\P(\tilde X^1=\emptyset)>0$.} Let $V$ be the event that every particle is at $\emptyset$ after every operation of the run before $\rho$, and at stage 0. $\P(X_0^i=\emptyset\ \forall i)=\pi_0(\emptyset)^N>0$ by (R1). If all particles are at $\emptyset$, every resampled particle is at $\emptyset$. Given $\cH'$ (everything before a round) on the event that all entering particles are at $\emptyset$, the particles end the round at $\emptyset$ with conditional probability $K_\varphi(\emptyset,\{\emptyset\})^N\ge(\varepsilon_\emptyset/m)^N>0$, by the conditional independence in (R3) and Lemma~\ref{T12:lem:pos}(a). Finitely many operations precede $\rho$, so $\P(V)>0$, and $V\subseteq\{\tilde X^1=\emptyset\}$.

\emph{Step 3.} By (R3) and Lemma~\ref{T12:lem:pos}(b) at the temperature of $\rho$, $\P(E_1\mid\cH)=\ind\{\tilde X^1=\emptyset\}K_\varphi(\emptyset,\Om\setminus\{\emptyset\})\ge\ind\{\tilde X^1=\emptyset\}\delta/m$. Since $\Zh_\rho>0$ surely (it is a product of positive factors) and is $\cH$-measurable,
\[ \E\Big[\Zh_\rho\tfrac1N\textstyle\sum_iD_i\Big]\ \ge\ \tfrac1N\E\big[\Zh_\rho\,\ind\{\tilde X^1=\emptyset\}\big]\,\delta/m\ >\ 0 . \]
By (i.c) with $\kappa\equiv m$, $\E[\Zh_\rho\frac1N\sum_if(X^i_{[m]})]=Z_\rho q_\rho(f)$, where $Z_\rho$ is $Z_t$ or $Z$. Hence $\E[\Zh_\rho\frac1N\sum_if(X^i_{[\tau_i]})]>Z_\rho q_\rho(f)$. \qed

\section{Proof of (ii)}

Measured round $u\in\{1,\dots,M\}$ is a post-terminal round at $\varphi=1$. Let $\kappa_j(\sigma)=\sigma^{-1}(g_j)-1$, where $g_j$ is the Gibbs-$j$ label. Then $Y^i_{u,j}=X^i_{[\kappa_j(\sigma_i)]}$ in that round. By (i.c) with $f=o_j(\cdot,1)$, and then Lemma~\ref{T12:lem:prefix} at $\varphi=1$ (where $q_1=\pi$),
\[ \E\Big[\Zh\,\tfrac1N\textstyle\sum_io_j(Y^i_{u,j},1)\Big]=Z\,\pi(o_j(\cdot,1))=Z\,\pi(J_j)=Z\,\PIP_j . \]
Averaging over $u=1,\dots,M$ gives $\E_c[\Zh\,\PIPh_j]=Z\,\PIP_j$. \qed

\begin{T12remark}
Lemma~\ref{T12:lem:prefix} is used only through $\pi(o_j)=\pi(J_j)$. The recorded value is a conditional probability at a state-independent prefix. So the identity holds whatever the position of the Gibbs-$j$ label in the order, because that position is a function of $\sigma_i$ alone.
\end{T12remark}

\section{Embedding in the method: conditioning on $\cC$}

\begin{T12lemma}[freezing]\label{T12:lem:freeze}
Let $\cC$ be a random element of a measurable space $(E,\mathcal E)$, and let $U=(U_1,\dots,U_4)$ be a random element of $(\Upsilon,\mathcal U)$ independent of $\cC$. For every bounded measurable $g:E\times\Upsilon\to\R$, $\E[g(\cC,U)\mid\cC]=\bar g(\cC)$ a.s., where $\bar g(c)=\E[g(c,U)]$ is measurable.
\end{T12lemma}
\begin{proof}
For $g=\ind_{A\times B}$ and $C\in\mathcal E$: $\E[g(\cC,U)\ind_C(\cC)]=\P(\cC\in A\cap C)\P(U\in B)=\E[\bar g(\cC)\ind_C(\cC)]$, by independence, with $\bar g(c)=\ind_A(c)\P(U\in B)$. The class of bounded measurable $g$ for which $\bar g$ is measurable and the identity holds for all $C$ contains the indicators of the $\pi$-system of measurable rectangles. It is closed under linear combinations and bounded monotone limits (monotone convergence). By the monotone class theorem it contains every bounded $\mathcal E\otimes\mathcal U$-measurable $g$.
\end{proof}

By CAL and EX-S, $\cC$ (with the event $\cA\in\sigma(\cC)$ that the final run executes) is a function of the data and streams 0--9. The final streams $U_r$ (stream $9+r$) are independent of each other and of streams 0--9, and are identically distributed. On $\cA$, island $r$'s record is $O_r=\Phi(\cC,U_r)$. Here $\Phi(c,\cdot)$ realises $I(c)$ with the method's parameters ($N=N_e$, $M=4$, $L'=L$, flag unset), and $\Phi$ is jointly measurable with values in a Polish space (part of EX-S in STATEMENT.md). Applying Lemma~\ref{T12:lem:freeze} to $g\,\ind_\cA(\cC)$, for every bounded measurable function $g$ of $(O_1,\dots,O_4)$,
\begin{equation}\label{T12:eq:freeze}
\E[g(O_1,\dots,O_4)\mid\cC]=\E_c\big[g(O_1,\dots,O_4)\big]\big|_{c=\cC}\quad\text{a.s. on }\cA,
\end{equation}
where on the right the four islands are independent copies of $I(c)$. Every conclusion stated for fixed $c$ and for the four independent islands of $I(c)$ therefore holds given $\cC$ on $\cA$. This includes the identities of (i) and (ii), which hold for every $c$, and so hold almost surely given $\cC$.

\section{Proof of (iii)}

Fix $c$ and four independent islands of $I(c)$. By (i.a) and (ii) for each island, $\E_c[\Zh_{r,t}\qh_{r,t}(f)]=Z_tq_t(f)$ and $\E_c[\Zh_r\PIPh_{rj}]=Z\PIP_j$. Hence $\E_c\Zh=\frac14\sum_r\E_c\Zh_r=Z$. Each $\Zh_r$ and $\Zh$ takes values in the compact interval $[Z_0e^{\min d},Z_0e^{\max d}]\subset(0,\infty)$, because each factor $\frac1N\sum_ie^{\Delta_td(\cdot)}$ lies between $e^{\Delta_t\min d}$ and $e^{\Delta_t\max d}$ and $\sum_t\Delta_t=1$. So $\log\Zh$ and $\log\Zh_r$ are bounded, and Jensen's inequality for the concave $\log$ gives $\E_c\log\Zh\le\log\E_c\Zh=\log Z$, and likewise for $\Zh_r$. For $0<\alpha<1$, Markov's inequality for the non-negative $\Zh$ gives $\P_c(\alpha\Zh>Z)\le\alpha\E_c\Zh/Z=\alpha$, that is $\P_c(Z\ge\alpha\Zh)\ge1-\alpha$; likewise for $\Zh_r$. By \eqref{T12:eq:freeze} applied to the bounded functions involved (indicators, $\Zh_{r,t}\qh_{r,t}(f)$, $\Zh_r\PIPh_{rj}$, $\Zh$, $\log\Zh$, $\log\Zh_r$), all of these hold given $\cC$ on $\cA$. For the stage identities, fix $t$ and apply \eqref{T12:eq:freeze} on $\cA\cap\{T(\cC)\ge t\}$ (an event in $\sigma(\cC)$). This gives $\E[\Zh_{r,t}\qh_{r,t}(f)\mid\cC]=Z_tq_t(f)$, with $Z_t,q_t$ computed on $\cC$'s grid; the right side is a $\sigma(\cC)$-measurable random variable. \qed

\section{Proof of (iv)}

\subsection{(iv.a) $\log\Zh$}
If $d\equiv0$ on $\Om$, then $G_t\equiv1$, every factor of $\Zh$ is $1$, and $\Zh=Z_0=m(\emptyset)\Zb_0=m(\emptyset)\Zb_1=Z$ (Lemma~\ref{T12:lem:bridge}, since $\gamma_1=\gamma_0$).

Suppose $d\not\equiv0$; since $d(\emptyset)=0$, pick $S$ with $d(S)\ne0$. For $U\in\{\emptyset,S\}$, let $E_U$ be the event that all stage-0 draws equal $U$, and that after every round preceding the stage-$T$ resampling (the stage-0 round if the flag is set, and the stage-$t$ rounds for $t\le T-1$) every particle is at $U$. As in Step 2 of \S4 (using Lemma~\ref{T12:lem:pos}(a) for $U$ and the fact that a population entirely at $U$ resamples to $U$), $\P_c(E_U)\ge \pi_0(U)^N(\varepsilon_U/m)^{NT}>0$, with at most $T$ rounds involved. On $E_U$, every weight at stage $t$ equals $e^{\Delta_td(U)}$, so $\Zh=Z_0e^{d(U)\sum_t\Delta_t}=Z_0e^{d(U)}$. Thus $\Zh$ takes the two distinct values $Z_0$ and $Z_0e^{d(S)}$ with positive probability, so $\P_c(\Zh\ne Z)>0$. For the bounded positive $W=\Zh$ with $\E_cW=Z$, the inequality $\log w\le\log Z+(w-Z)/Z$ holds for all $w>0$, with equality only at $w=Z$. Taking expectations, with strict inequality on an event of positive probability, gives $\E_c\log\Zh<\log Z$. For the method's four islands given $\cC$: by \eqref{T12:eq:freeze} the islands are independent copies of $I(\cC)$. On the event that all four islands are in their copy of $E_U$ (conditional probability $\P_\cC(E_U)^4>0$), $\Zh=\frac14\sum_r\Zh_r=Z_0e^{d(U)}$. So the same argument gives $\E[\log\Zh\mid\cC]<\log Z$, and also $\E[\log\Zh_r\mid\cC]<\log Z$. \qed

\subsection{(iv.b) normalised PIPs}
Given $\cC$, the islands are identically distributed (by \eqref{T12:eq:freeze}), so $\E[\bar P_j\mid\cC]=\E[\PIPh_{1j}\mid\cC]$. By (ii) and (i.a), $\E[\Zh_1\PIPh_{1j}\mid\cC]=Z\PIP_j$ and $\E[\Zh_1\mid\cC]=Z$. Hence
$\mathrm{Cov}(\PIPh_{1j},\Zh_1\mid\cC)=Z\PIP_j-Z\,\E[\PIPh_{1j}\mid\cC]$, which is the stated formula.

\emph{The instance.} Take the appendix data of STATEMENT.md: $n=20$, $p=2$. As recorded there, $K=\min(2,10)=2$, so $\Om$ is all four subsets. Further $B=\min(4,1,16)=1$, so the blocks are $\{1\},\{2\}$ and $r_b=\mathrm{BF}$ on singletons. The neighbour graph has $N(1)=\{2\}$ and $N(2)=\{1\}$, and $m=5$, $L=\max(4,\lceil\log_2 21\rceil)=5$. The prior of STATEMENT \S0 gives $w(k)=\tfrac12/\binom{3}{k}$, that is $w(0)=\frac12$ and $w(1)=w(2)=\frac16$, and $W_K=\binom20\frac12+\binom21\frac16+\binom22\frac16=1$.

\emph{The scores.} By STATEMENT \S0, $C_n$ cancels in $\mathrm{BF}(S)=D_S^{-1/2}\big[(2b_0+Q_S)/(2b_0+y^\top y)\big]^{-\nu}$, with $\nu=a_0+n/2=10.01$ and $a_0=b_0=0.01$.
\begin{itemize}
\item For $S=\{j\}$: $G=x_j^\top x_j$, $H=(n+1)G+1$, $D_S=H/(G+1)$, $Q_S=y^\top y-n(x_j^\top y)^2/H$.
\item For $S=\{1,2\}$: write $g_{ab}$ for the entries of $G_S$, $z=(z_1,z_2)=X_S^\top y$ and $h_{ab}=(n+1)g_{ab}+\ind\{a=b\}$. Then $\det H_S=h_{11}h_{22}-h_{12}^2$, $\det(G_S+I)=(g_{11}+1)(g_{22}+1)-g_{12}^2$, and $z^\top H_S^{-1}z=(h_{22}z_1^2-2h_{12}z_1z_2+h_{11}z_2^2)/\det H_S$ (the $2\times2$ inverse).
\end{itemize}
These are exactly the expressions in \texttt{log\_bf} of Appendix~\ref{T12:app:code}. The function \texttt{standardise} centres each raw column and $y$ and divides by the sample standard deviation with denominator $n-1$, as in STATEMENT \S0. So $\gamma_0=(\frac12,\frac{b_1}6,\frac{b_2}6,\frac{b_1b_2}6)$ and $\gamma_1=(\frac12,\frac{b_1}6,\frac{b_2}6,\frac{b_{12}}6)$ on $(\emptyset,\{1\},\{2\},\{1,2\})$, with $b_S=\mathrm{BF}(S)$. Moreover $d$ vanishes except at $\{1,2\}$, where $d(\{1,2\})=D:=\log(b_{12}/(b_1b_2))$. With $T=1$ and $N=1$, $\Zh=Z_0e^{D\,B}$ with $B=\ind\{X^{(0)}=\{1,2\}\}$. Resampling of a single particle is trivial. The particle then performs $1+L'$ rounds at $\varphi=1$ before the $M=4$ measured rounds. The recording $o_1(Y,1)$ equals $c_0=\gamma_1(\{1\})/(\gamma_1(\{1\})+\gamma_1(\emptyset))$ if $2\notin Y$, and $c_1=\gamma_1(\{1,2\})/(\gamma_1(\{1,2\})+\gamma_1(\{2\}))$ if $2\in Y$. Hence $\PIPh_1=c_0+(c_1-c_0)k/4$, where $k$ is the number of measured rounds whose pre-Gibbs-1 state contains $2$.

Under the stage-0 flag, $X^{(0)}=X_0K_0$ has law $\pi_0K_0=\pi_0$ (Lemma~\ref{T12:lem:prefix}). Everything after stage 0 depends on the past only through $X^{(0)}$. Hence the joint law of $(B,k)$, and so of $(\Zh,\PIPh_1)$, is the same under both stage-0 readings. It is computed by propagating a law on triples (bit $b$, state $x$, count $\kappa$):
\begin{enumerate}[label=(\arabic*)]
\item start from $\mu_0(b,x,0)=\pi_0(x)\ind\{b=\ind\{x=\{1,2\}\}\}$;
\item apply the round kernel $K_1$ to the state coordinate $1+L'$ times;
\item apply four measured rounds, each mapping $(b,x,\kappa)$ to $(b,z,\kappa+\ind\{2\in y\})$ with weight
$\mathsf J(x,y,z)=\frac1{120}\sum_\sigma[K_{\sigma(1)}\cdots K_{\sigma(\kappa_1(\sigma))}](x,y)\,[K_{\sigma(\kappa_1(\sigma)+1)}\cdots K_{\sigma(5)}](y,z)$,
where $y$ is the pre-Gibbs-1 state;
\item marginalise over the final state.
\end{enumerate}
This is the law of $(B,k)$ by (R3) and Lemma~\ref{T12:lem:round}. Within a round, the path splits at the pre-Gibbs-1 position into the prefix product and the suffix product of the same order. The $4\times4$ elementary kernels at $\varphi=1$ are given by STATEMENT \S1. At $p=2$, $B=1$ the uniform and neighbour swaps coincide: both are the Metropolis swap $\{1\}\leftrightarrow\{2\}$ with acceptance $\min\{1,b_2/b_1\}$ (resp.\ $\min\{1,b_1/b_2\}$), and they stay at $\emptyset$ and $\{1,2\}$. The expectations then are finite sums: $\E[\bar P_1]=\E[\PIPh_1]$, and
$\E[\tilde P_1]=\sum\prod_{r=1}^4\P(B_r,k_r)\cdot\frac{\sum_re^{DB_r}\PIPh(k_r)}{\sum_re^{DB_r}}$ over the $10^4$ joint values of four independent islands.

We evaluated every quantity in 60-digit interval arithmetic (mpmath \texttt{iv}). Every input ($n$, the data, $a_0=b_0=0.01$) enters as an enclosing interval, every operation is rigorous interval arithmetic, and $\min\{1,x\}$ is enclosed by $[\min(1,\underline x),\min(1,\overline x)]$. The listing is in Appendix~\ref{T12:app:code}. Every enclosure has width below $10^{-52}$: the last \texttt{CHECK} line of each block of the log (Appendix~\ref{T12:app:log}) subtracts the two exact endpoints of each enclosure in ordinary \texttt{mpf} arithmetic and compares the largest difference with the floating-point constant \texttt{1e-52}; that subtraction rounds by at most one part in $2^{53}$, against a margin of a factor five, and no conclusion of the proof uses this width. The five decisive \texttt{CHECK} lines compare an exact interval endpoint with a floating-point constant whose rounding cannot cross the threshold it stands for. Rounded to 12 digits, they are:
\begin{center}
\begin{tabular}{lcc}
 & $L'=L$ (method; $1+L'=6$) & $L'=L-1$ (generalisation; $1+L'=5$)\\\hline
$D$ & \multicolumn{2}{c}{$8.14676240634$}\\
$a=\pi_0(\{1,2\})$ & \multicolumn{2}{c}{$0.0279141173179$}\\
$c_0,\ c_1$ & \multicolumn{2}{c}{$0.0694711126908,\ 0.998708337452$}\\
$\PIP_1$ & \multicolumn{2}{c}{$0.990618012243$}\\
$(\E\Zh-Z)/Z_0$ & interval $\ni0$ & interval $\ni0$\\
$(\E[\Zh\PIPh_1]-Z\PIP_1)/Z_0$ & interval $\ni0$ & interval $\ni0$\\
$\E[\bar P_1]-\PIP_1$ & $-0.173974177850$ & $-0.212065526919$\\
$\E[\tilde P_1]-\PIP_1$ & $-0.159677356166$ & $-0.194638440548$\\
\end{tabular}
\end{center}
The first two difference rows check (i.a) and (ii) on the instance; they are not used in the argument. In both columns and under both stage-0 readings, the biases are certified to be below $-0.17$ (primary) and $-0.15$ (secondary). \emph{Scope.} This is a component example. At $p=2$ the full method would take the exact branch, and $N=1$ is not the production budget. It shows that the frozen SMC component's normalised PIP estimators are biased, not that the full method's output on this input is. \qed

\subsection{(iv.c) island independence}
By \eqref{T12:eq:freeze} with $g=\prod_rh_r(O_r)$, $\E[\prod_rh_r(O_r)\mid\cC]=\prod_r\psi_{h_r}(\cC)$ on $\cA$, where $\psi_h(c)=\E_c[h(O)]$ is the same function for every island. This is conditional independence with conditional law that of $I(\cC)$. Assume $\P(\cA)>0$, so that $\P(\cdot\mid\cA)$ is defined. Under $\P(\cdot\mid\cA)$ and for bounded $h$: $\E[h(O_1)h(O_2)]=\E[\psi_h(\cC)^2]$ and $\E h(O_r)=\E\psi_h(\cC)$, so $\mathrm{Cov}(h(O_1),h(O_2))=\mathrm{Var}(\psi_h(\cC))$.

\emph{Characterisation.} The record $O_r$ is a finite collection of real and finite-valued variables, so it takes values in a Polish space $\mathsf O$. Choose a countable $\pi$-system $\{B_k\}$ generating its Borel $\sigma$-field. If $O_1,O_2$ are independent under $\P(\cdot\mid\cA)$, then $\mathrm{Var}(\psi_{\ind_{B_k}}(\cC))=\mathrm{Cov}(\ind_{B_k}(O_1),\ind_{B_k}(O_2))=0$ for every $k$. So, outside one null set, $\P(O_1\in B_k\mid\cC)=\mu(B_k)$ for all $k$, where $\mu$ is the law of $O_1$ under $\P(\cdot\mid\cA)$. Two probability measures that agree on a generating $\pi$-system are equal, so the conditional law is a.s.\ $\mu$. Conversely, if the conditional law is a.s.\ a fixed $\mu$, then $\E[h_1(O_1)h_2(O_2)]=\E[\psi_{h_1}\psi_{h_2}(\cC)]=\mu(h_1)\mu(h_2)=\E h_1(O_1)\E h_2(O_2)$; so the islands are independent, and the same holds for all four by the same computation. None of the hypotheses restricts the law of $\cC$ beyond its measurability in streams 0--9. So they do not exclude, for example, a calibration law that puts positive mass on two values $c\ne c'$ with different island laws, and then the islands are dependent under $\P(\cdot\mid\cA)$. The next paragraph shows such pairs exist: $(\text{grid},N)$ and $(\text{grid},4N)$ on any input where $d$ has a unique maximiser. The instance of (iv.b) is one, since there $d=(0,0,0,D)$ with $D>0$. Whether the method's actual calibration procedure produces such a law on some input is not claimed.

\emph{The island law depends on $c$.} Suppose $d$ has a unique maximiser $S^*$ on $\Om$, and let $d^*=d(S^*)$. Then $\Zh/Z_0=\prod_{t=1}^T\frac1N\sum_ie^{\Delta_td(X^{(t-1),i})}\le e^{d^*}$, with equality iff $X^{(t-1),i}=S^*$ for all $i$ and all $t=1,\dots,T$. This event has probability $a^N$, where $a=\pi_0(S^*)\prod_{t=1}^{T-1}K_{\varphi_t}(S^*,\{S^*\})$ does not depend on $N$. Here we used (R1)--(R3): $X^{(0),i}$ has law $\pi_0$ under either stage-0 reading (by Lemma~\ref{T12:lem:prefix}), independently across $i$; a population at $S^*$ resamples to $S^*$; and the rounds act independently across particles. Now $0<a\le\pi_0(S^*)<1$ by Lemma~\ref{T12:lem:pos}(a) and $|\Om|\ge2$. So $\P_{(\text{grid},N)}(\Zh=Z_0e^{d^*})=a^N\ne a^{4N}=\P_{(\text{grid},4N)}(\Zh=Z_0e^{d^*})$. \qed

\subsection{(iv.d) launch events}
Let $A\in\sigma(\cC)$, $A\subseteq\cA$, $\P(A)>0$. By the tower property, $\E[Y\ind_A]=\E[\E[Y\mid\cC]\ind_A]$ for every bounded $Y$. For the conclusions of (iii) whose right side $v$ is a constant (fixed by the data), this gives $\E[Y\mid A]=v$ for $Y=\Zh$, $Y=\Zh_r$ and $Y=\Zh_r\PIPh_{rj}$ (with $v=Z$, $Z$, $Z\PIP_j$). Likewise $\E[\log\Zh\mid A]\le\log Z$ and $\E[\log\Zh_r\mid A]\le\log Z$. With $Y=\ind\{Z\ge\alpha\Zh\}$ (resp.\ $\Zh_r$), $\P(Z\ge\alpha\Zh\mid A)=\E[\P(Z\ge\alpha\Zh\mid\cC)\ind_A]/\P(A)\ge1-\alpha$. For an intermediate stage $t$ and $A\subseteq\{T(\cC)\ge t\}$, the right side $Z_tq_t(f)$ of (iii) is a $\sigma(\cC)$-measurable random variable. The same computation then gives only $\E[\Zh_{r,t}\qh_{r,t}(f)\ind_A]=\E[Z_tq_t(f)\ind_A]$, which is the averaged form. The left side divided by $\P(A)$ is a number. The right side's integrand $Z_tq_t(f)$ is non-constant when, for example, $\varphi_t$ takes two values $u\ne u'$ with positive probability under $\P(\cdot\mid A)$ and $f=\ind\{S\}$ with $d(S)\ne0$. Indeed $Z_tq_t(\ind\{S\})=m(\emptyset)\gamma_0(S)e^{\varphi_td(S)}$ by the definitions, and this takes different values at $u$ and $u'$. \qed

\subsection{(iv.e) favourable final diagnostics}
In the instance of (iv.b), $\log\Zh_r=\log Z_0+DB_r$ with $B_r$ iid Bernoulli($a$) across the four independent islands. If exactly $k\in\{1,2,3\}$ of the $B_r$ equal $1$, the four values $DB_r$ have standard deviation $D\sqrt{k(4-k)}/4\ge D\sqrt3/4$ (normalisation $1/4$) or $D\sqrt{k(4-k)/12}\ge D/2$ (normalisation $1/3$). Both exceed $1$ because $D>8>4/\sqrt3$. If $k\in\{0,4\}$ the standard deviation is $0$. Hence $F=\{B_1=B_2=B_3=B_4\}$ under either convention, and $\P(F)=(1-a)^4+a^4>0$. The four islands are iid and $F$ is invariant under permuting them, so the pairs $(\Zh_r,\PIPh_{r1})$ have the same law under $\P(\cdot\mid F)$ for every $r$. Hence $\E[\Zh_r\mid F]=\E[\Zh\mid F]=\E[\Zh_1\mid F]$ and $\E[\Zh_r\PIPh_{r1}\mid F]=\E[\Zh_1\PIPh_{11}\mid F]$ for every $r$. Moreover
\[ \E[\Zh_1\mid F]=Z_0\,\frac{(1-a)^4+a^4e^D}{(1-a)^4+a^4}=Z_0\big(1+(e^D-1)\theta\big),\qquad \theta=\frac{a^4}{(1-a)^4+a^4}. \]
Since $a<\frac12$, $\theta<a$ (because $a^4<a(1-a)^4+a^5\iff a^3<(1-a)^3$). Also $e^D>1$ and $Z=\E\Zh_1=Z_0(1+(e^D-1)a)$. Hence $\E[\Zh_1\mid F]<Z$ and $\E[\Zh\mid F]<Z$. So the identity $\E\Zh_r=Z$ of (i) and the identity $\E\Zh=Z$ of (iii) fail under $\P(\cdot\mid F)$. For (ii),
\[ \E[\Zh_1\PIPh_1\mid F]=\frac{\sum_{b,k}\P(B=b,k)\,Z_0e^{Db}\,\PIPh(k)\,\P(B=b)^3}{(1-a)^4+a^4}, \]
which the same interval computation evaluates as $(\E[\Zh_1\PIPh_1\mid F]-Z\PIP_1)/Z_0=-95.609622724\ldots$ ($L'=L$) and $-95.648818292\ldots$ ($L'=L-1$), under either stage-0 reading; both are certified non-zero. This is a component calculation (scope as in (iv.b)). It breaks an identity in each of (i), (ii) and (iii). It does not break the Jensen or Markov inequalities of (iii): since $\E[\Zh\mid F]<Z$, they still hold under $\P(\cdot\mid F)$. \qed

\subsection{(iv.f) NUM alone}
Fix $\eta>0$ and the NUM-compatible perturbed-score model in which every evaluated $\log m(S)$ is replaced by $\log m(S)+\eta$. (This is a model NUM does not exclude, not a claim about a particular floating-point implementation.) It satisfies the NUM bound with equality. Every quantity the run uses other than the prefactor $m(\emptyset)$ depends on scores only through ratios $m(S)/m(\emptyset)$ or $m(S)/m(S')$: the tables $r_b$, $\mathrm{BF}$, $d$, $\gamma_\varphi$ up to the common factor $\varpi$, $\pi_0$, $\Zb_0$, all acceptance and Gibbs probabilities, and the weights. These are unchanged, so the run has the same law, and $\Zh^{\rm impl}=e^\eta m(\emptyset)\Zb_0\prod_t(\cdots)=e^\eta\Zh$. By (i.a), $\E_c\Zh^{\rm impl}=e^\eta Z\neq Z$. If $d\equiv0$ (for example $p=1$: a single block $\{1\}$, so $d(S)=\log\mathrm{BF}(S)-\log r_1(S)=0$), then $\Zh=Z$ surely by (iv.a), so $\Zh^{\rm impl}=e^\eta Z$, and $\P(Z\ge\alpha\Zh^{\rm impl})=\ind\{\alpha e^\eta\le1\}=0$ for $\alpha\in(e^{-\eta},1)$. So NUM alone gives neither unbiasedness nor the Markov bound. \qed

\appendix
\section{Listing of the interval computation}\label{T12:app:code}
The certified values of \S8.2 and \S8.5 are produced by \texttt{certify\_instance.py iv}, which calls \texttt{p2\_exact.py} with \texttt{IvArith(60)} (mpmath 1.3.0 interval arithmetic at 60 digits). The listings follow.

\subsection*{p2\_exact.py}
\lstinputlisting[language=Python]{build/helpers/p2_exact.py}
\subsection*{certify\_instance.py}
\lstinputlisting[language=Python]{build/helpers/certify_instance.py}

\section{Output of the certified run}\label{T12:app:log}
The output of \texttt{python -I certify\_instance.py iv} (file \texttt{helpers/certify\_iv.log}) follows. Each enclosure is printed as its two exact endpoints, rounded to 25 digits for display. The \texttt{CHECK} lines compare exact endpoints. The thresholds $-0.1701$ and $-0.1501$ lie below $-0.17$ and $-0.15$ whatever their binary rounding, and $2.31>4/\sqrt3$. Rows \texttt{identity\_Z} and \texttt{identity\_ZP} are $(\E\Zh-Z)/Z_0$ and $(\E[\Zh\PIPh_1]-Z\PIP_1)/Z_0$. Rows \texttt{condF\_Z\_minus\_Z} and \texttt{condF\_ZP\_minus\_ZPIP} are $(\E[\Zh_1\mid F]-Z)/Z_0$ and $(\E[\Zh_1\PIPh_{11}\mid F]-Z\PIP_1)/Z_0$. Row \texttt{condF\_Elog\_minus\_logZ} is the unscaled difference $\E[\log\Zh_1\mid F]-\log Z$ (not used in the proof). The last \texttt{CHECK} line of each block computes, for every enclosure in that block, the difference of its two exact endpoints in ordinary \texttt{mpf} arithmetic and compares the largest with the floating-point constant \texttt{1e-52}; this is a display check and is not used by the proof. \texttt{rounds\_before\_measured}$=1+L'$.
\lstinputlisting{build/helpers/certify_iv.log}
\mdproofend
\mdsec{(d) Review record}
\mdsubsec{FINAL\_REPORT.md: statement changes}
\begin{mdblock}
\mdhead{2}{Statement changes, and why}

\begin{mdnum}
\item[\mdlbl{1.}] \textbf{Made precise (Astra \#1 statement check: APPROVED WITH CHANGES).}

\begin{mdbul}
\item The frozen run is stated for arbitrary integer parameters. The method is the case N = N\_e, M = 4, L′ = L, with no stage-0 round.
\item START, exactness of the π₀ sampler and of Z̄₀, is an imported dependency (T09(i)), not an extra restriction.
\item The negative parts of (iv) are component statements, with the scope stated: the p = 2 example would take the exact branch in the full method, and N = 1 is not the production budget.
\item "Not necessarily independent after averaging over 𝓒" is formalised as an exact characterisation: unconditional independence holds iff the island law given 𝓒 is a.s. constant. No claim is made about the actual calibration law.
\item "Favourable diagnostics" is formalised as the final E2 check passing. "NUM alone" is formalised as a perturbed-score model that NUM does not exclude.
\end{mdbul}
\item[\mdlbl{2.}] \textbf{One error of mine, caught by Astra \#1.} My first draft extended (iv.d) to \emph{every} conclusion of (iii) under a launch event. That is false for intermediate stages, whose targets depend on the random grid. It was repaired: only the 𝓒-free conclusions are kept, and intermediate stages hold in averaged form. This was an over-extension beyond the conjecture, so it is not counted as a step-1 restart.
\item[\mdlbl{3.}] \textbf{Stage B round 1.} Both reviewers independently found that the statement pointed to 00\_COMMON for the marginal-likelihood formula instead of writing it out. STATEMENT §0 now carries Felipe's fixed target formulas. Joint measurability of the stream map is now stated explicitly under EX-S. No conclusion changed.
\end{mdnum}
\end{mdblock}

\mdsubsec{FINAL\_REPORT.md: review history}
\begin{mdblock}
\mdhead{2}{Review history}

{\small
\begin{longtable}{|>{\raggedright\arraybackslash}p{\dimexpr(\linewidth-10\tabcolsep-6\arrayrulewidth)*137/1000\relax}|>{\raggedright\arraybackslash}p{\dimexpr(\linewidth-10\tabcolsep-6\arrayrulewidth)*212/1000\relax}|>{\raggedright\arraybackslash}p{\dimexpr(\linewidth-10\tabcolsep-6\arrayrulewidth)*137/1000\relax}|>{\raggedright\arraybackslash}p{\dimexpr(\linewidth-10\tabcolsep-6\arrayrulewidth)*228/1000\relax}|>{\raggedright\arraybackslash}p{\dimexpr(\linewidth-10\tabcolsep-6\arrayrulewidth)*286/1000\relax}|}
\hline
\textbf{Stage / round} & \textbf{Reviewer} & \textbf{Object} & \textbf{Verdict} & \textbf{Fatal / gap / exposition} \\ \hline\endhead
A, statement check & Astra \#1 (high) & STATEMENT v1 & APPROVED WITH CHANGES & 1 false clause, (iv.d), plus 3 required clarifications; all adopted \\ \hline
A, review & Astra \#1 (high), resumed & PROOF\_v2 & \textbf{verified} & 0 / 0 / 3 (fixed in v3) \\ \hline
B, round 1 & Astra \#2 (high), new session & PROOF\_v3 & not verified & 0 / 1 / 1 \\ \hline
B, round 1 & Opus 5.5 (xhigh), Claude Personal & PROOF\_v3 & not verified & 0 / 2 / 6 \\ \hline
B, round 2 & Astra \#2 & PROOF\_v4 + claims log & \textbf{verified}; agrees with all 3 answers & 0 / 0 / 2 (fixed in v5) \\ \hline
B, round 2 & Opus 5.5 (xhigh) & PROOF\_v4 + claims log & \textbf{verified}; agrees with all 3 answers & 0 / 0 / 2 (fixed in v5) \\ \hline
\end{longtable}}

\begin{mdbul}
\item CLAIMS\_OF\_FATAL\_FLAWS.md has every gap finding verbatim, with my answer and both reviewers' rulings. No fatal finding was raised at any stage.
\item The Astra judge was not needed.
\item The independent checks of the numerical example agree to 12+ digits:

\begin{mdbul}
\item a Sonnet assistant wrote its own implementation (helpers/sonnet\_check/);
\item the Opus reviewer re-ran the certificate, wrote its own float implementation, and ran a single-seed Monte Carlo of the literal algorithm with 200,000 islands, getting −0.1752 ± 0.0008 against the certified −0.17397.
\end{mdbul}
\end{mdbul}
\end{mdblock}

\mdsubsec{CLAIMS\_OF\_FATAL\_FLAWS.md (every entry, verbatim)}
\begin{mdblock}
\mdhead{1}{CLAIMS\_OF\_FATAL\_FLAWS — T12 (T8)}

Every fatal or gap finding, copied verbatim, with the lead's answer. Stage A (Astra \#1 on PROOF\_v2) had no fatal or gap findings, only 3 exposition findings, all fixed in PROOF\_v3. The entries below are from Stage B round 1 on PROOF\_v3. The fixes are in PROOF\_v4 and STATEMENT v2.2.

\mdhr

\mdhead{2}{Claim 1 (Astra \#2, finding 1, gap)}

\begin{mdquote}
\textbf{Location:} §8.2, {\mdtt PROOF\_\allowbreak{}v3.\allowbreak{}tex} lines 185–214; §8.5, lines 227–231; Appendix A, {\mdtt helpers/\allowbreak{}p2\_\allowbreak{}exact.\allowbreak{}py} lines 58–84.

\textbf{Severity:} gap.

\textbf{Claim:} The computational examples are not sufficiently linked to the marginal likelihood specified in the theorem.

\textbf{Evidence:} The helper's {\mdtt log\_\allowbreak{}bf} function asserts that it computes the v3 target's Bayes factors. It uses \(a_0=b_0=0.01\), \(H=(n+1)X_S^\top X_S+I\), a determinant ratio involving \(X_S^\top X_S+I\), and exponent \(a_0+n/2\). None of the permitted documents supplies the target formula or derives these expressions from it. {\mdtt STATEMENT.\allowbreak{}md} refers to {\mdtt 00\_\allowbreak{}COMMON §1}, whose contents are not supplied in the permitted materials, and says that only positivity of \(m\) and \(\varpi\) is used.

Positivity alone does not determine these Bayes factors or the numerical conclusions. In particular, the finite-state enumeration and saved intervals establish results for the scores implemented by {\mdtt log\_\allowbreak{}bf}; they do not establish that those scores equal the theorem's \(m(S)/m(\emptyset)\). Consequently, the bias bounds in \textbf{(iv.b)} and the instance-specific diagnostic conclusions in \textbf{(iv.e)} remain unverified against the stated target. This is not evidence that the implemented formula is false; its required identification is unsupported here.
\end{mdquote}

\mdhead{2}{Claim 2 (Opus xhigh reviewer, G1, gap). Same issue as Claim 1, found independently.}

\begin{mdquote}
\textbf{G1}

\begin{mdbul}
\item \textbf{Location:} §8.2, "The instance" and the interval-computation paragraph; §8.5; {\mdtt helpers/\allowbreak{}p2\_\allowbreak{}exact.\allowbreak{}py}, {\mdtt log\_\allowbreak{}bf} (lines 58–84) and {\mdtt setup} (line 98).
\item \textbf{Severity:} gap.
\item \textbf{Claim:} The (iv.b) bias bounds (−0.17 and −0.15) and the (iv.e) result E[Ẑ\_r·PIP̂\_\{r1\} | F] ≠ Z·PIP₁ depend on the exact formula for m(S), on a₀ = b₀ = 0.01, and on the prior weights w(k). None of these is stated or justified in the proof or in STATEMENT.md, so I cannot verify that the computation is about the v3 target.
\item \textbf{Evidence:}

\begin{mdbul}
\item The formula exists only in code, as log BF(S) = −½ log(det H / det(G+I)) − ν[log(2b₀ + yᵀy − n zᵀH⁻¹z) − log(2b₀ + yᵀy)], with H = (n+1)G + I and ν = a₀ + n/2. Its only justification is the comment "for the v3 target".
\item STATEMENT §0 defers the formula to 00\_COMMON §1, which is not in the packet. It also says "only m > 0 and ϖ > 0 are used", which is not true for this instance.
\item The prior values w(0) = ½, w(1) = w(2) = 1/6, W\_K = 1 match the standard BB(1,2) beta-binomial (I recomputed them), but that definition is also not in the packet.
\item The margins are small: at the method column, 0.0040 for the primary bound and 0.0097 for the secondary.
\item A diagnostic only: with ν = a₀ + (n−1)/2 the bounds still hold (primary −0.179106, secondary −0.164769, D = 7.7534). So the conclusion may well survive, but I could not verify the formula itself.
\item The qualitative (iv.e) facts (P(F) > 0, E[Ẑ\_r | F] < Z, E[Ẑ | F] < Z) need only D > 4/√3 and a < ½. Both of those numbers come from the same unverified formula.
\end{mdbul}
\end{mdbul}
\end{mdquote}

\textbf{Lead's answer to Claims 1–2: accepted and fixed.} The finding is right. STATEMENT.md delegated the target formulas to a file outside the packet, and its remark "only m > 0 and ϖ > 0 are used" was false for the instance. Fixes: 1. STATEMENT §0 (v2.2) now writes out the fixed v3 target, copied from METHODOLOGY\_v3 §5 and 00\_COMMON §1: H\_S, D\_S, Q\_S, m(S) with C\_n, BF(S) = D\_S\textasciicircum{}\{−1/2\}[(2b₀+Q\_S)/(2b₀+yᵀy)]\textasciicircum{}\{−ν\}, ν = a₀ + n/2, a₀ = b₀ = 0.01, w(k) = B(1+k, 2p−k)/B(1,p) = ½/C(2p−1,k), W\_K. The positivity remark is restricted to everything except the instance. This adds no assumption and changes no conclusion: these are Felipe's fixed target, which the statement always referred to. 2. PROOF\_v4 §8.2 has a new paragraph, "The scores". It derives the 1×1 and 2×2 expressions that {\mdtt log\_\allowbreak{}bf} computes from the STATEMENT §0 formula (C\_n cancels; explicit 2×2 determinant and inverse; ν = 10.01), states that {\mdtt standardise} implements the §0 standardisation, and computes the prior values from w(k) = ½/C(3,k). 3. The certificate was re-run after cosmetic helper changes (E6), with no change to any value. It now prints exact endpoint comparisons (CHECK lines), and its output is reproduced in a new Appendix B.

\mdhead{2}{Claim 3 (Opus xhigh reviewer, G2, gap, minor)}

\begin{mdquote}
\textbf{G2}

\begin{mdbul}
\item \textbf{Location:} §6, the paragraph after Lemma 6.1 ("the joint measurability of Φ is part of the stream model").
\item \textbf{Severity:} gap (minor).
\item \textbf{Claim:} The proof assumes Φ is jointly measurable in (c, u). STATEMENT.md does not state this; EX-S, CAL and Success only say that Φ(c, ·) realises I(c).
\item \textbf{Evidence:} Two steps need it. Lemma 6.1 is applied to g(Φ(c,u₁),…,Φ(c,u₄))·1\_𝒜(c), which must be measurable on E×Υ. The (iv.c) characterisation needs c ↦ law of I(c) to be a probability kernel. The assumption is standard and true for any algorithmic implementation, but it is imported rather than given.
\end{mdbul}
\end{mdquote}

\textbf{Lead's answer: accepted and fixed.} STATEMENT (v2.2) now says explicitly, under EX-S, that the stream model includes joint measurability of (c, u) ↦ Φ(c, u) with values in a Polish space; PROOF\_v4 §6 cites it. This states what EX-S ("independent ideal uniform streams") already presupposes, since the island record is produced from uniform streams by finitely many measurable operations. It restricts no input and changes no conclusion.

\mdhr

\mdhead{2}{Exposition findings (not fatal/gap; recorded for completeness)}

\begin{mdbul}
\item Astra \#2 finding 2 = Opus E1 (broken \textbackslash{}ref, log not in appendix): fixed. The reference is repaired and the log is reproduced in Appendix B.
\item Opus E2 (ε\_U temperature): fixed. Lemma 1.4(a) now defines the temperature-free ε\_S = π₀(S)e\textasciicircum{}\{−R\_d\}, valid for every φ ∈ [0, 1].
\item Opus E3 ((iv.c) implicit step): fixed. Added the explicit sentence: the hypotheses admit a calibration law with mass on (grid, N) and (grid, 4N); the (iv.b) instance has a unique maximiser of d; no claim is made about the actual calibration procedure.
\item Opus E4 ("only the averaged form"): fixed. STATEMENT (iv.d) now gives an exact remark instead of "only … in general random". PROOF\_v4 §8.4 proves that Z\_tq\_t(1\{S\}) = m(∅)γ₀(S)e\textasciicircum{}\{φ\_t d(S)\} is non-constant when φ\_t takes two values.
\item Opus E5 (block/graph facts rely on method rules): the proof now cites the STATEMENT appendix, where these facts are recorded.
\item Opus E6 (helper cosmetics): fixed. Endpoint display, the prior comment and the {\mdtt \_\allowbreak{}\_\allowbreak{}main\_\allowbreak{}\_\allowbreak{}} block are corrected; the certificate was re-run with identical values (certify\_iv.log; previous log kept as helpers/certify\_iv\_v1.log).
\item Opus, (iv.e) "extension from r = 1 to every r by exchangeability is implicit": made explicit in §8.5.
\end{mdbul}

\mdhead{2}{Rulings by the Stage B reviewers (round 2, on PROOF\_v4 / STATEMENT v2.2)}

\begin{mdbul}
\item \textbf{Astra \#2} (review/A2\_call2\_stageB\_r2\_ruling\_REPLY.md): agrees with the answers to Claims 1, 2 and 3; each fix works. No remaining fatal or gap findings. Two exposition items: (a) the log displayed rounded endpoints, so the width claim was not visible; (b) the Appendix B description wrongly called the log row condF\_Elog "divided by Z₀". \textbf{Verdict: verified.}
\item \textbf{Opus xhigh} (review/OPUS\_stageB\_r2\_ruling.md): agrees with all three answers; each fix works, checked term by term and numerically. Two exposition items: X1, certify\_iv\_v1.log was not in its packet (packaging; the file is in the project's helpers/); X2, an unneeded remark about the method's max-abs scaling. \textbf{Verdict: verified.}
\item Lead: all four exposition items were fixed in PROOF\_v5 (exposition only, no mathematical change from the verified v4). The certificate now prints an exact width CHECK (max width 1.92·10⁻⁵³ < 10⁻⁵²); the Appendix B row descriptions are corrected; the X2 remark is removed. No dispute remains, so the Astra judge was not needed.
\end{mdbul}
\end{mdblock}

\chapter{T13 (T10): restricted ESS and the omitted-mass identity}
\chaptersource{Final proof: \texttt{proofs/\allowbreak{}T13\_\allowbreak{}T10\_\allowbreak{}restricted\_\allowbreak{}ESS/\allowbreak{}proof/\allowbreak{}PROOF\_v4.tex}. Sources of parts (a), (b), (d): \texttt{theorems/\allowbreak{}T13\_\allowbreak{}T10\_\allowbreak{}restricted\_\allowbreak{}ESS.md}, \texttt{STATEMENT.md}, \texttt{FINAL\_REPORT.md}, \texttt{CLAIMS\_OF\_FATAL\_FLAWS.md}.}
\mdsec{(a) Original conjecture text (theorem file, verbatim)}
\begin{mdblock}
\mdhead{1}{T13 (T10): Restricted ESS and the omitted-mass identity}

\textbf{Status:} conjectured, not proven. Methodology SMC-BVS v3.0 (2026-10-19), SHA-256 93043758058a6d422cca6ca4b6be8f3a9e8149b6f6de2e1bab1ce31c44b8019c. Read with {\mdtt 00\_\allowbreak{}COMMON.\allowbreak{}md}.

\textbf{Category / rank:} algorithmic validity; \textbf{core} (lemma).

\mdhead{2}{Definitions used}

Bridge q\_ψ for real ψ; candidate set V (frozen, nonempty); restricted ESS ρ\_V(u, v) = Z̄\_v(V)²/(Z̄\_u(V)Z̄\_\{2v−u\}(V)); a(ψ) := q\_ψ(V); for W ⊆ V, a\_W(ψ) := q\_ψ(W), q\textasciicircum{}W\_ψ = q\_ψ(·|W); χ²(q\_v‖q\_u) = ∫(dq\_v/dq\_u)²dq\_u − 1; 𝓔\_grid (00\_COMMON §2).

\mdhead{2}{Assumptions}

EX-A; fixed V; COVER(V, ε\_C) or COVER(W, ε\_C) for the corollaries.

\mdhead{2}{Precise statement}

For nonempty frozen V, u < v, w = 2v − u: (i) [identity] 1 + χ²(q\_v‖q\_u) = ρ\_V(u,v)⁻¹·a(v)²/(a(u)a(w)). (ii) [COVER corollary] An accepted interval (ρ\_V ≥ 0.9) with a(u), a(w) ≥ 1 − ε\_C satisfies 1 + χ²(q\_v‖q\_u) ≤ 10/(9(1−ε\_C)²). COVER supplies this at every grid evaluation point (φ\_\{t−1\}, φ\_t, w\_t); coverage on [0,1] alone does not supply it when w\_t > 1; the largest w\_t is not automatically the worst coverage point. (iii) [restricted working set] For nonempty W ⊆ V: 1 + χ²(q\_v\textasciicircum{}W‖q\_u\textasciicircum{}W) = ρ\_V(u,v)⁻¹·(a\_V(v)²/a\_W(v)²)·(a\_W(u)a\_W(w))/(a\_V(u)a\_V(w)); COVER(W, ε\_C) gives the same bound as (ii). (iv) [restricted grid] If ∅ ∉ 𝓜 and calibration never scores 𝓜, then V ∩ 𝓜 = ∅; with the same pilot temperatures and the same V, restricting γ\_φ to 𝓜\textasciicircum{}c leaves the grid construction unchanged. No assertion derives particle mixing from ESS.

\mdhead{2}{Dependencies}

Builds on: bridge definitions; project L4/L7. Used by: T18, T25, T28, T31, T36.

\mdhead{2}{Numerical evidence (provisional)}

None specific (the identity is algebraic); provisional evidence (Sol check J6) uses it in the reflected-pair product check.

\mdhead{2}{Proof-route hint (optional)}

Substitute Z̄\_ψ(V) = a(ψ)Z̄\_ψ into the two second-moment identities (project L4, L7; valid for every real ψ, so over-tempered points are covered); (iv): ρ\_V depends only on γ\_φ restricted to V. Hardest step: none for the identity; establishing coverage for the actual scored set is the substantive obstacle (T25). \emph{When proof.md is run on this file, its lead considers the statement both with and without these hints; the hints never reach the reviewers.}
\end{mdblock}

\mdsec{(b) STATEMENT.md (verbatim)}
\begin{mdblock}
\mdhead{1}{STATEMENT — T13 (T10): restricted ESS and the omitted-mass identity}

Source: {\mdtt Theory\_\allowbreak{}architecture\_\allowbreak{}v3/\allowbreak{}theorems/\allowbreak{}T13\_\allowbreak{}T10\_\allowbreak{}restricted\_\allowbreak{}ESS.\allowbreak{}md} (all parts), read with {\mdtt 00\_\allowbreak{}COMMON.\allowbreak{}md}, {\mdtt THEORY\_\allowbreak{}BRIEF.\allowbreak{}md} §2.10 and the method file {\mdtt METHODOLOGY\_\allowbreak{}v3.\allowbreak{}md} §6.7 (SHA-256 93043758058a6d422cca6ca4b6be8f3a9e8149b6f6de2e1bab1ce31c44b8019c). Version 3 (after statement checks 1 and 2). Changes against the source are listed in §2 with reasons.

\mdhead{2}{0. Setting and definitions}

\begin{mdbul}
\item \textbf{Data and model objects} are those of 00\_COMMON §1: standardised X, y; Ω = \{S ⊆ [p] : |S| ≤ K\}, K = min(p, ⌊n/2⌋); m(S), BF(S) = m(S)/m(∅); ϖ(S) = w(|S|)/W\_K with w(k) = ½/C(2p−1, k).
\item \textbf{Bridge} (00\_COMMON §2, METHODOLOGY §6.3): blocks C₁, …, C\_J; tables r\_b(U) = BF(U), U ⊆ C\_b; γ₀(S) = ϖ(S)∏\_b r\_b(S ∩ C\_b); d(S) = log BF(S) − Σ\_b log r\_b(S ∩ C\_b). For every real ψ: γ\_ψ = γ₀e\textasciicircum{}\{ψd\}, Z̄\_ψ = Σ\_\{S∈Ω\}γ\_ψ(S), q\_ψ = γ\_ψ/Z̄\_ψ. (The proof shows from these definitions that γ₀ > 0 and d is finite on Ω, so every q\_ψ has full support.) For nonempty A ⊆ Ω: Z̄\_ψ(A) = Σ\_\{S∈A\}γ\_ψ(S), q\_ψ(A) = Z̄\_ψ(A)/Z̄\_ψ, q\textasciicircum{}A\_ψ = q\_ψ(· | A).
\item \textbf{Candidate set.} V is a nonempty subset of Ω, held fixed. a(ψ) = a\_V(ψ) := q\_ψ(V); for nonempty W ⊆ V, a\_W(ψ) := q\_ψ(W).
\item \textbf{χ² divergence.} For laws μ, λ on a finite set with λ > 0 wherever μ > 0: χ²(μ‖λ) = Σ\_\{S:λ(S)>0\} λ(S)(μ(S)/λ(S))² − 1.
\item \textbf{Restricted ESS.} For real u, v and w := 2v − u: ρ\_V(u, v) = Z̄\_v(V)²/(Z̄\_u(V)Z̄\_w(V)) = exp\{2F\_V(v) − F\_V(u) − F\_V(w)\}, F\_V(t) = log Σ\_\{S∈V\}γ₀(S)e\textasciicircum{}\{td(S)\}. An interval (u, v), u < v, is \textbf{accepted} if ρ\_V(u, v) ≥ 0.9.
\item \textbf{Grid.} 0 = φ₀ < φ₁ < … < φ\_T = 1, T ≥ 1; w\_t = 2φ\_t − φ\_\{t−1\}; 𝓔\_grid = \{φ₀, …, φ\_T, w₁, …, w\_T\}.
\item \textbf{Candidate-grid construction} (METHODOLOGY §6.7). Inputs: an increasing list of pilot temperatures in (0, 1] ending at 1, the candidate set V, T\_e. Start at u = 0; for the next pilot temperature b: if ρ\_V(u, b) ≥ 0.90 append b; else bisect on (u, b] with u fixed, at most 60 halvings, and append the largest tested endpoint that passes (ρ\_V(u, ·) ≥ 0.90); continue from the appended point; validate strict progress and the inequality at every appended point; non-progress, failed validation or more than T\_e intervals ⇒ failure. A \textbf{successful construction} outputs a grid. An \textbf{implementation} of the construction is a fixed deterministic procedure realising these rules (it fixes the bisection points, stopping and endpoint choices); it accesses the bridge only through evaluations of ρ\_V.
\item \textbf{Candidate set of the method's form.} For a set 𝒫 ⊆ Ω of scored models (the \emph{pool}), V(𝒫) := top-20p of 𝒫 by log γ₀ ∪ top-20p of 𝒫 by log γ₁ ∪ \{∅\} (ties lexicographic; if |𝒫| ≤ 20p each top list is all of 𝒫).
\item \textbf{COVER(W, ε\_C)}, for nonempty W ⊆ Ω, a grid, and \textbf{0 ≤ ε\_C < 1}: q\_ψ(W) ≥ 1 − ε\_C for every ψ ∈ 𝓔\_grid.
\item \textbf{EX-A}: exact scores, sums, probabilities and comparisons.
\item \textbf{Method-class witness} (used in (ii)(c)–(d)): a raw data set (X\textasciicircum{}r, y\textasciicircum{}r) satisfying the method's input conditions (n ≥ 3, no constant column, y\textasciicircum{}r non-constant) on which the exact-branch condition Σ\_\{k≤K\}C(p, k) ≤ (5H₀ + 4N₀)p fails (so the SMC branch runs); the method's blocks, tables, γ₀ and d for these data; a candidate set V = V(𝒫) of the method's form for some pool 𝒫 ⊆ Ω; a tier e ∈ \{0, 1\} with T\_e = 4\textasciicircum{}eT₀ computed from the data's (n, p) (00\_COMMON §2 budgets); a fixed deterministic implementation of the construction as in this section; a list of pilot temperatures; and the grid that this implementation outputs, successfully and under EX-A, from V, that list and T\_e. (No claim is made about the probability that the random pilot outputs that list, or that a calibration reaches the pool 𝒫.)
\end{mdbul}

\mdhead{2}{1. Statement}

Throughout, V ⊆ Ω is nonempty and fixed, u, v are real, w = 2v − u.

\textbf{(i) Identity.} For all real u, v (in particular u < v):\\
\hspace*{1.5em}1 + χ²(q\_v ‖ q\_u) = ρ\_V(u, v)⁻¹ · a(v)²/(a(u)a(w)).

\textbf{(ii) Coverage corollary.} (a) Let 0 ≤ ε\_C < 1. If (u, v) is accepted and a(u) ≥ 1 − ε\_C and a(w) ≥ 1 − ε\_C, then 1 + χ²(q\_v ‖ q\_u) ≤ 10/(9(1 − ε\_C)²). (b) Assume EX-A. Let (φ\_t)\_\{t=0\}\textasciicircum{}T be the output of a successful candidate-grid construction run with candidate set V (in particular the frozen grid of a calibration, paired with the V of the construction that produced it). Then every interval (φ\_\{t−1\}, φ\_t) is accepted; its evaluation points φ\_\{t−1\}, φ\_t, w\_t lie in 𝓔\_grid; and if COVER(V, ε\_C) holds, then 1 + χ²(q\_\{φ\_t\} ‖ q\_\{φ\_\{t−1\}\}) ≤ 10/(9(1 − ε\_C)²) for every t = 1, …, T. (c) Coverage on [0, 1] alone does not supply the hypothesis of (a) when w\_t > 1: there is a method-class witness with grid T = 1 (φ₀ = 0, φ₁ = 1, w₁ = 2) such that, with ε\_C = 0.1, the interval (0, 1) is accepted and q\_ψ(V) ≥ 1 − ε\_C for every ψ ∈ [0, 1], but q\_\{w₁\}(V) < 1 − ε\_C and 1 + χ²(q₁ ‖ q₀) > 10/(9(1 − ε\_C)²). (d) The largest w\_t is not automatically the worst coverage point: there is a method-class witness with grid T = 2 (both intervals accepted) such that the minimum of ψ ↦ q\_ψ(V) over 𝓔\_grid is attained only at φ₀, not at max\_t w\_t, and the minimum over \{w₁, …, w\_T\} is not attained at max\_t w\_t.

\textbf{(iii) Restricted working set.} For every nonempty W ⊆ V and all real u, v:\\
\hspace*{1.5em}1 + χ²(q\textasciicircum{}W\_v ‖ q\textasciicircum{}W\_u) = ρ\_V(u, v)⁻¹ · (a\_V(v)²/a\_W(v)²) · (a\_W(u)a\_W(w))/(a\_V(u)a\_V(w)). In particular (W = V) ρ\_V(u, v) = 1/(1 + χ²(q\textasciicircum{}V\_v ‖ q\textasciicircum{}V\_u)). Moreover, with 0 ≤ ε\_C < 1, if (u, v) is accepted and a\_W(v) ≥ 1 − ε\_C, then 1 + χ²(q\textasciicircum{}W\_v ‖ q\textasciicircum{}W\_u) ≤ 10/(9(1 − ε\_C)²). Hence, for a grid as in (ii)(b), COVER(W, ε\_C) gives at every interval t both 1 + χ²(q\textasciicircum{}W\_\{φ\_t\} ‖ q\textasciicircum{}W\_\{φ\_\{t−1\}\}) ≤ 10/(9(1 − ε\_C)²) and (because W ⊆ V, COVER(W, ε\_C) implies COVER(V, ε\_C)) the bound of (ii)(b).

\textbf{(iv) Restricted grid.} Let 𝓜 ⊆ Ω with ∅ ∉ 𝓜. (a) Let V = V(𝒫) be of the method's form. If 𝒫 ∩ 𝓜 = ∅ (pool avoidance), then V ∩ 𝓜 = ∅. Pool avoidance holds, in particular, when calibration never scores 𝓜 and every member of 𝒫 was scored by the calibration or belongs to an inherited set of scored models (tables, leaders carried across tiers) that is assumed to avoid 𝓜. (b) Suppose V ∩ 𝓜 = ∅. Let γ′\_φ := γ\_φ1\_\{𝓜\textasciicircum{}c\} for real φ (the bridge on Ω′ = Ω ∖ 𝓜 ∋ ∅ with γ₀, d restricted to Ω′; q′\_φ = q\_φ(· | 𝓜\textasciicircum{}c)). Then F′\_V = F\_V and ρ′\_V = ρ\_V on all real arguments, and, under EX-A, one and the same implementation of the candidate-grid construction, run on the restricted and on the original bridge with the same pilot temperatures, the same V and the same T\_e, has the same outcome (the same grid, or failure in both).

No assertion derives particle mixing from ESS.

\mdhead{2}{2. Changes against the source statement, with reasons}

\begin{mdnum}
\item[\mdlbl{1.}] \textbf{ε\_C ∈ [0, 1)} added (approved in check 1): for ε\_C ≥ 1 the bound is undefined or false.
\item[\mdlbl{2.}] \textbf{(iv)(a) pool avoidance} made the explicit hypothesis (modification of check 1): 𝒫 ∩ 𝓜 = ∅ is sufficient for V ∩ 𝓜 = ∅. The source's wording "calibration never scores 𝓜" is kept as a sufficient condition, with the provenance condition that every pool member was scored by the calibration or comes from an inherited set assumed to avoid 𝓜 (the method does not say whether tables or carried leaders feed V's top lists).
\item[\mdlbl{3.}] \textbf{(ii)(b) the V paired with a grid} is the candidate set of the construction that produced it (approved in check 1).
\item[\mdlbl{4.}] \textbf{(ii)(c)–(d) witnesses are method-class} (replaces the rejected abstract-bridge version): real data run through the method's blocks, tables, γ₀ and d, on the SMC branch; V of the method's form; a grid output by the construction. The pilot's random output is not claimed; the pilot temperatures enter as construction inputs, as they do in the source's (iv).
\item[\mdlbl{5.}] \textbf{(iii) the bound uses coverage at v} (a\_W(v) ≥ 1 − ε\_C), supplied by COVER(W, ε\_C) because φ\_t ∈ 𝓔\_grid; both the restricted and the full-space bound are stated (approved in check 1; the condition is sufficient, not claimed necessary).
\item[\mdlbl{6.}] \textbf{(i), (iii) for all real u, v} (approved in check 1).
\item[\mdlbl{7.}] \textbf{Positivity}: γ₀ > 0 and d finite on Ω are derived in the proof from the definitions of 00\_COMMON §1 (ϖ > 0, m(S) ∈ (0, ∞)), not assumed.
\item[\mdlbl{8.}] \textbf{(iv)(b) implementation}: the "unchanged construction" conclusion is for one fixed deterministic implementation applied to both bridges (check 1: EX-A alone does not fix bisection and endpoint choices).
\item[\mdlbl{9.}] \textbf{(ii)(d) strengthened} (recorded at check 2): the witness must have the grid-evaluation minimum attained \textbf{only} at φ₀ (the source asks only that the largest w\_t not be the worst point). This is a strengthening, not a weakening.
\item[\mdlbl{10.}] \textbf{Witness parameters} (modification of check 2): a witness names its tier e, uses T\_e = 4\textasciicircum{}eT₀ from its own data, and fixes a deterministic implementation.
\end{mdnum}
\end{mdblock}

\mdsec{(c) Final proof, verbatim: body of PROOF\_v4.tex}
\noindent{\small\itshape The proof below is the body of \texttt{PROOF\_v4.tex} from \texttt{\textbackslash begin\{document\}} to \texttt{\textbackslash end\{document\}} without the title block. Its preamble definitions are local to this chapter; labels carry the prefix \texttt{T13:}.}\par\medskip
\mdproofbegin
% ---- preamble of PROOF_v4.tex (theorem environments carry the prefix T13) ----
\newfontfamily\qfont{Cambria Math}
\lstset{basicstyle=\ttfamily\footnotesize,breaklines=true,columns=fullflexible}
\newtheorem{T13lemma}{Lemma}
\newtheorem{T13proposition}[T13lemma]{Proposition}
\theoremstyle{remark}
\newtheorem{T13remark}[T13lemma]{Remark}
\newcommand{\Zb}{\bar Z}
\newcommand{\E}{\mathcal{E}_{\mathrm{grid}}}
\newcommand{\M}{\mathcal{M}}
\newcommand{\Pool}{\mathcal{P}}
\newcommand{\chisq}{\chi^2}
% ---- body of PROOF_v4.tex ----
All notation is that of STATEMENT.md (version 3). In particular $\Omega=\{S\subseteq[p]:|S|\le K\}$, $\gamma_\psi=\gamma_0e^{\psi d}$, $\Zb_\psi(A)=\sum_{S\in A}\gamma_\psi(S)$, $\Zb_\psi=\Zb_\psi(\Omega)$, $q_\psi=\gamma_\psi/\Zb_\psi$, $q_\psi^A=q_\psi(\cdot\mid A)$, $a_W(\psi)=q_\psi(W)$, $a=a_V$, and $w=2v-u$ whenever $u,v$ are given. Throughout, $V\subseteq\Omega$ is nonempty and fixed.

STATEMENT.md defines the model and sampler objects by reference to \texttt{00\_COMMON.md} and \texttt{METHODOLOGY\_v3.md}. Every definition from those files that this proof uses is quoted verbatim in Appendix~B (with the files' SHA-256 checksums); citations of the form [B: file, \S, lines] point there.

\section{Preliminaries}

\begin{T13lemma}[Positivity]\label{T13:lem:pos}
For every $S\in\Omega$, $\varpi(S)>0$ and $0<m(S)<\infty$. Consequently $\gamma_0(S)\in(0,\infty)$ and $d(S)\in\mathbb{R}$ for all $S\in\Omega$; for every real $\psi$ and nonempty $A\subseteq\Omega$, $\gamma_\psi>0$ on $\Omega$, $\Zb_\psi(A)\in(0,\infty)$, and $q_\psi$ and $q_\psi^A$ give positive mass to every point of $\Omega$ and of $A$ respectively.
\end{T13lemma}
\begin{proof}
The definitions used are [B: 00\_COMMON.md, \S1, lines 10--11] (also [B: METHODOLOGY\_v3.md, \S5]): $G_S=X_S^\top X_S$, $z_S=X_S^\top y$, $H_S=(n+1)G_S+I$, $D_S=\det H_S/\det(G_S+I)$, $Q_S=y^\top y-nz_S^\top H_S^{-1}z_S$, $\nu=0.01+n/2$, $a_0=b_0=0.01$,
$m(S)=(2b_0)^{a_0}\Gamma(\nu)/(\pi^{n/2}\Gamma(a_0))\cdot D_S^{-1/2}(2b_0+Q_S)^{-\nu}$, $w(k)=\tfrac12/\binom{2p-1}{k}$, $W_K=\sum_{k\le K}\binom pk w(k)$, $\varpi(S)=w(|S|)/W_K$; and $\mathrm{BF}(S)=m(S)/m(\emptyset)$, $r_b(U)=\mathrm{BF}(U)$ for $U\subseteq C_b$, $\gamma_0$, $d$ as in [B: 00\_COMMON.md, \S2, lines 18--20].

$\varpi(S)=w(|S|)/W_K$ with $w(k)=\tfrac12/\binom{2p-1}{k}>0$ for $0\le k\le K\le p\le 2p-1$, and $W_K=\sum_{k\le K}\binom pk w(k)\ge w(0)>0$ is finite.
For $m(S)$, and in fact for $m(U)$ with any $U\subseteq[p]$ (the argument does not use $|U|\le K$; we write it for $S$): $G_S=X_S^\top X_S$ is positive semidefinite, so $H_S=(n+1)G_S+I$ and $G_S+I$ are positive definite (for $S=\emptyset$ all matrices are $0\times 0$, determinants are $1$ and $Q_\emptyset=y^\top y$). Hence $D_S=\det H_S/\det(G_S+I)\in(0,\infty)$. With $\hat\beta=nH_S^{-1}z_S$,
\[
\|y-X_S\hat\beta\|^2+\tfrac1n\bigl(\|X_S\hat\beta\|^2+\|\hat\beta\|^2\bigr)
= y^\top y-2\hat\beta^\top z_S+\tfrac1n\hat\beta^\top\bigl((n+1)G_S+I\bigr)\hat\beta
\]\[
= y^\top y-2n z_S^\top H_S^{-1}z_S+n z_S^\top H_S^{-1}z_S=Q_S,
\]
so $Q_S\ge0$ and $2b_0+Q_S\ge 0.02>0$. The constant $(2b_0)^{a_0}\Gamma(\nu)/(\pi^{n/2}\Gamma(a_0))$ is positive and finite. Hence $m(S)=\text{const}\cdot D_S^{-1/2}(2b_0+Q_S)^{-\nu}\in(0,\infty)$, so $\mathrm{BF}(S)=m(S)/m(\emptyset)\in(0,\infty)$, every table entry $r_b(U)=\mathrm{BF}(U)$ ($U\subseteq C_b$) is in $(0,\infty)$ by the same argument applied to $U$, $\gamma_0(S)=\varpi(S)\prod_b r_b(S\cap C_b)\in(0,\infty)$, and $d(S)=\log\mathrm{BF}(S)-\sum_b\log r_b(S\cap C_b)$ is finite. The remaining claims follow because $\Omega$ is finite.
\end{proof}

\begin{T13lemma}[Restricted second-moment identity]\label{T13:lem:sm}
For every nonempty $A\subseteq\Omega$ and all real $u,v$,
\[
1+\chisq\bigl(q_v^A\,\big\|\,q_u^A\bigr)=\frac{\Zb_u(A)\,\Zb_w(A)}{\Zb_v(A)^2}.
\]
For $A=\Omega$ this reads $1+\chisq(q_v\|q_u)=\Zb_u\Zb_w/\Zb_v^2$.
\end{T13lemma}
\begin{proof}
By Lemma~\ref{T13:lem:pos}, $q_u^A>0$ exactly on $A$ and $q_v^A$ vanishes off $A$, so the $\chi^2$ divergence is defined and its sum runs over $A$. For $S\in A$, $q_v^A(S)/q_u^A(S)=e^{(v-u)d(S)}\Zb_u(A)/\Zb_v(A)$. Hence
\[
\sum_{S\in A}q_u^A(S)\Bigl(\frac{q_v^A(S)}{q_u^A(S)}\Bigr)^2
=\frac{\Zb_u(A)^2}{\Zb_v(A)^2}\sum_{S\in A}\frac{\gamma_0(S)e^{ud(S)}e^{2(v-u)d(S)}}{\Zb_u(A)}
=\frac{\Zb_u(A)\Zb_w(A)}{\Zb_v(A)^2},
\]
since $u+2(v-u)=2v-u=w$. Subtracting $1$ gives the claim. For $A=\Omega$, $q^\Omega_\psi=q_\psi$.
\end{proof}

\begin{T13lemma}[Omitted mass]\label{T13:lem:om}
For nonempty $A\subseteq\Omega$ and real $\psi$: $\Zb_\psi(A)=a_A(\psi)\Zb_\psi$ with $a_A(\psi)=q_\psi(A)\in(0,1]$. If $\emptyset\neq W\subseteq V$ then $\Zb_\psi(W)=\Zb_\psi(V)\,a_W(\psi)/a_V(\psi)$ and $a_W(\psi)\le a_V(\psi)$.
\end{T13lemma}
\begin{proof}
The first identity is the definition $q_\psi(A)=\Zb_\psi(A)/\Zb_\psi$; positivity is Lemma~\ref{T13:lem:pos}. The second follows by dividing the first for $A=W$ by the first for $A=V$. Monotonicity: $W\subseteq V$ and $\gamma_\psi\ge0$.
\end{proof}

\section{Part (i)}
By Lemma~\ref{T13:lem:sm} with $A=\Omega$ and Lemma~\ref{T13:lem:om} with $A=V$ (so $\Zb_\psi=\Zb_\psi(V)/a(\psi)$ for $\psi\in\{u,v,w\}$),
\[
1+\chisq(q_v\|q_u)=\frac{\Zb_u\Zb_w}{\Zb_v^2}
=\frac{\Zb_u(V)\Zb_w(V)}{\Zb_v(V)^2}\cdot\frac{a(v)^2}{a(u)a(w)}
=\rho_V(u,v)^{-1}\,\frac{a(v)^2}{a(u)a(w)},
\]
for all real $u,v$. All quantities are finite and positive by Lemma~\ref{T13:lem:pos}. \qed

\section{Part (iii)}
Let $\emptyset\ne W\subseteq V$. By Lemma~\ref{T13:lem:sm} with $A=W$ and Lemma~\ref{T13:lem:om},
\[
1+\chisq(q_v^W\|q_u^W)=\frac{\Zb_u(W)\Zb_w(W)}{\Zb_v(W)^2}
=\frac{\Zb_u(V)\Zb_w(V)}{\Zb_v(V)^2}\cdot\frac{a_W(u)}{a_V(u)}\cdot\frac{a_W(w)}{a_V(w)}\cdot\frac{a_V(v)^2}{a_W(v)^2},
\]
which is the identity of (iii), since the first factor is $\rho_V(u,v)^{-1}$. For $W=V$ the last three factors equal $1$, so $\rho_V(u,v)=1/(1+\chisq(q_v^V\|q_u^V))$.

\emph{Bound.} Let $0\le\varepsilon_C<1$, $\rho_V(u,v)\ge 0.9$ and $a_W(v)\ge1-\varepsilon_C>0$. By Lemma~\ref{T13:lem:om}, $a_W(u)/a_V(u)\le1$ and $a_W(w)/a_V(w)\le1$; also $a_V(v)\le1$. Hence
$1+\chisq(q_v^W\|q_u^W)\le (10/9)\cdot 1/(1-\varepsilon_C)^2$.

\emph{Grid consequence.} Let $(\varphi_t)$ be as in (ii)(b) and assume COVER$(W,\varepsilon_C)$. By (ii)(b) (proved below, without using (iii)) every interval $(\varphi_{t-1},\varphi_t)$ is accepted, and $\varphi_t\in\E$, so $a_W(\varphi_t)\ge1-\varepsilon_C$; the bound gives $1+\chisq(q^W_{\varphi_t}\|q^W_{\varphi_{t-1}})\le 10/(9(1-\varepsilon_C)^2)$. Since $q_\psi(V)\ge q_\psi(W)$ for every $\psi$, COVER$(W,\varepsilon_C)$ implies COVER$(V,\varepsilon_C)$, and (ii)(b) gives the full-space bound. \qed

\section{Part (ii)(a)--(b)}
\emph{(a).} By (i), $a(v)\le1$, $a(u)a(w)\ge(1-\varepsilon_C)^2>0$ and $\rho_V(u,v)^{-1}\le 10/9$:
$1+\chisq(q_v\|q_u)\le (10/9)(1-\varepsilon_C)^{-2}$.

\emph{(b).} Assume EX-A and let $0=\varphi_0<\dots<\varphi_T=1$ be the output of a successful construction with candidate set $V$. We read the construction of STATEMENT \S0 (quoted from [B: METHODOLOGY\_v3.md, \S6.7]) as follows, which is the only reading consistent with ``start at $u=0$ \dots\ continue from the appended point'' and with an output $0=\varphi_0<\dots<\varphi_T=1$: the current left point $u$ is always the last appended point (initially $u=0$), and the output grid is $0$ followed by the appended points in the order appended. (The parenthetical case of the frozen grid of a calibration is the special case in which the frozen grid is such an output, paired with its $V$.) So $\varphi_t$ was appended while $u=\varphi_{t-1}$. A point $b'$ is appended only in two cases: $b'=b$ is the next pilot temperature and $\rho_V(u,b)\ge0.90$; or $b'$ is a tested bisection endpoint that passes, i.e.\ $\rho_V(u,b')\ge0.90$. In both cases $\rho_V(\varphi_{t-1},\varphi_t)\ge 0.9$, and under EX-A this comparison (and the validation of the inequality at every appended point) is exact. Strict progress gives $\varphi_{t-1}<\varphi_t$. Hence every interval is accepted. By definition of $\E$, $\varphi_{t-1},\varphi_t,w_t\in\E$. Under COVER$(V,\varepsilon_C)$, $a(\varphi_{t-1})\ge1-\varepsilon_C$ and $a(w_t)\ge1-\varepsilon_C$, and (a) with $(u,v,w)=(\varphi_{t-1},\varphi_t,w_t)$ gives the bound for every $t$. \qed

\section{Part (ii)(c)--(d): method-class witnesses}

\subsection{The data set}
Let $n=42$, $p=40$. Let $e_1,\dots,e_{41}$ be an orthonormal basis of $\mathbf 1^\perp=\{x\in\mathbb R^{42}:\sum_i x_i=0\}$ (dimension 41). Put $x_j=\sqrt{41}\,e_j$ ($j\in[40]$), $X^r=X=(x_1,\dots,x_{40})$.

\emph{Blocks.} We use three properties of the method's blocks, all from [B: METHODOLOGY\_v3.md, \S6.3] and [B: 00\_COMMON.md, \S2, line 18].
(a) \emph{They depend on $X$ only}: ``Blocks $C_1,\dots,C_J$ (deterministic in X: complete linkage, top-down cut, next-fit packing)''; the clustering distance $1-|r_{jl}|$ is ``computed from X'', and the cut and the packing use only the dendrogram, $B$ and index order. So the blocks can be computed from $X$ before $y$ is chosen.
(b) \emph{They partition $[p]$}: the dendrogram's root contains all $p$ retained columns; the cut repeatedly replaces a node by its two children (which partition it) until every retained cluster has at most $B$ members, so the retained clusters partition $[p]$; packing places whole retained clusters (``a retained cluster is never split'') into blocks, each cluster into exactly one block. Hence every column lies in exactly one block.
(c) \emph{Their size is capped}: $|C_b|\le B=\min(\lfloor\sqrt n\rfloor,\lfloor\sqrt p\rfloor,16)=\min(6,6,16)=6<40$.
By (a) the blocks are fixed by $X$. Let $j_1=1$; by (c) the block of $j_1$ does not contain all 40 columns, so we can pick $j_2\in[40]$ outside it; by (b) $j_1$ and $j_2$ lie in different blocks. Put
\[
y^r=y=\sqrt{41}\,(c_1e_{j_1}+c_2e_{j_2}+c_3e_{41}),\qquad c_1=0.97,\ c_2=0.10,\ c_3=\sqrt{1-c_1^2-c_2^2}=\sqrt{0.0491}.
\]
\emph{Input conditions and standardisation.} All entries are finite, $n=42\ge3$, $p=40\ge1$; each $x_j$ and $y$ is centred and nonzero, hence non-constant, so no column is dropped and $y$ is not constant [B: METHODOLOGY\_v3.md, \S4, items 1--3]. Both conventions in the sources use the sample SD with denominator $n-1$: [B: 00\_COMMON.md, \S1, line 9] (``$s_y^2=\|Cy^r\|^2/(n-1)$; so $\|x_j\|^2=\|y\|^2=n-1$'') and [B: METHODOLOGY\_v3.md, \S4, item 4] (``divide by its maximum absolute value, compute the mean, centre, then divide by the sample SD (denominator n $-$ 1)''). Here $\|x_j\|^2=\|y\|^2=41=n-1$ and the vectors are centred, so every sample SD equals $1$. Under 00\_COMMON's form, $D_x=I$, $s_y=1$, $CX^rD_x^{-1}=X$ and $Cy^r/s_y=y$. Under the method's form, a column $x$ is divided by $c=\max_i|x_i|>0$, stays centred, and has sample SD $1/c$, so the final division returns $x$; the same holds for $y$. Either way the standardised data are $X$ and $y$ themselves.
\emph{SMC branch.} By [B: 00\_COMMON.md, \S2, lines 24--25] and [B: METHODOLOGY\_v3.md, \S3, \S6.2]: $N_0=\max(2500,n,p)=2500$ and $H_0=\lceil512N_0/2500\rceil=\lceil512\cdot2500/2500\rceil=512$, so $(5H_0+4N_0)p=12560\cdot40=502400$; $K=\min(p,\lfloor n/2\rfloor)=\min(40,21)=21$ and $\sum_{k\le21}\binom{40}k\ge\binom{40}{20}=137846528820>502400$. The exact branch is not taken.
\emph{Budget.} Same sources: $\ell=\lceil\log_2(n+p+1)\rceil=\lceil\log_2 83\rceil=7$, $T_0=\max(1000,\ell^2)=1000$. Both witnesses use tier $e=0$, $T_e=T_0=1000$.

\begin{T13lemma}[Closed forms]\label{T13:lem:cf}
Let $\nu=0.01+n/2=21.01$, $t_i=(n-1)^2c_i^2/\bigl(n(n-1+2b_0)\bigr)=1681c_i^2/(42\cdot41.02)$ ($i=1,2$), $g_i=(1-t_i)^{-\nu}$, and for $S\in\Omega$ let $I(S)=\{i\in\{1,2\}:j_i\in S\}$. Then $t_1+t_2<1$ and
\[
\mathrm{BF}(S)=n^{-|S|/2}\Bigl(1-\sum_{i\in I(S)}t_i\Bigr)^{-\nu},\qquad
\gamma_0(S)=\varpi(S)\,n^{-|S|/2}\prod_{i\in I(S)}g_i,\qquad
d(S)=D\,\mathbf 1\{I(S)=\{1,2\}\},
\]
with $D=\nu\log\dfrac{(1-t_1)(1-t_2)}{1-t_1-t_2}>0$. These forms do not depend on the block partition beyond $j_1,j_2$ lying in different blocks.
\end{T13lemma}
\begin{proof}
For $|S|=k$: $G_S=41I_k$ (orthogonality), so $H_S=(43\cdot41+1)I_k=n^2I_k$, $G_S+I=nI_k$, $D_S=n^{2k}/n^k=n^k$. $z_j=x_j^\top y=41c_i$ if $j=j_i$ and $0$ otherwise (as $e_{41}\perp x_j$). So $Q_S=y^\top y-n\sum_{j\in S}z_j^2/n^2=41-\sum_{i\in I(S)}1681c_i^2/42$, and
$\mathrm{BF}(S)=D_S^{-1/2}\bigl((2b_0+Q_S)/(2b_0+Q_\emptyset)\bigr)^{-\nu}=n^{-k/2}(1-\sum_{i\in I(S)}t_i)^{-\nu}$ since $2b_0+Q_\emptyset=41.02$. Next $t_1+t_2=1681\cdot0.9509/1722.84<0.928<1$. A block $C_b$ contains at most one of $j_1,j_2$, so $r_b(U)=n^{-|U|/2}\prod_{i\in I(U)}g_i$ for $U\subseteq C_b$; multiplying over blocks ($\sum_b|S\cap C_b|=|S|$, and each $j_i\in S$ lies in exactly one block, both by the partition property (b)) gives $\gamma_0$. Finally $d(S)=-\nu\log(1-\sum_{I(S)}t_i)+\nu\sum_{I(S)}\log(1-t_i)$, which is $0$ when $|I(S)|\le1$ and $D$ when $I(S)=\{1,2\}$; $D>0$ because $(1-t_1)(1-t_2)=1-t_1-t_2+t_1t_2>1-t_1-t_2>0$.
\end{proof}

\emph{Grouped sums.} Let $N=[40]\setminus\{j_1,j_2\}$ ($|N|=38$) and $\sigma(a)=\sum_{m=0}^{\min(38,K-a)}\binom{38}{m}\varpi(a+m)n^{-(a+m)/2}$. By Lemma~\ref{T13:lem:cf}, grouping $S$ by $I(S)$ and $|S\cap N|$,
\[
\Zb_\psi=P+R\,e^{\psi D},\qquad P=\sigma(0)+(g_1+g_2)\sigma(1),\qquad R=g_1g_2\,\sigma(2),
\]
with $P$ the $\gamma_0$-mass of models not containing both $j_1,j_2$ and $R$ the $\gamma_0$-mass of the others. Write $r=R/P$.

\emph{Numerical values and their certification.} Every number used below is a finite expression in exact rationals ($\varpi(k)$, $t_i$, $\nu$, binomial coefficients) and the operations $+,-,\times,/,\sqrt{\cdot},\log,\exp$. The script of Appendix~A computes the rationals exactly (integer and fraction arithmetic) and evaluates each expression in interval arithmetic with outward rounding (\texttt{mpmath.iv}, 50 digits), so each internal output interval is guaranteed to contain the exact real value (for display, its endpoints are rounded outward, so the printed intervals contain it too). Every inequality between transcendental quantities used in \S\S5.2--5.3 is decided by comparing interval endpoints (lower endpoint of the larger side against upper endpoint of the smaller side); the script prints these seven checks ($D>0$, $a_c(1)\ge0.9$, $a_c(2)<0.9$, $1+\chi^2>1000/729$, $\beta>\alpha$, and the two acceptance inequalities), and all pass (output in Appendix~A). The remaining facts used are exact: $t_1+t_2<1$ (exact rational arithmetic, asserted in the script and shown by hand in Lemma~\ref{T13:lem:cf}); the set sizes $|V_c|=117$, $|V_d|=782\le800$, the counts $1,2\le T_e$ and the exact-branch inequality (integer arithmetic); $P>0$ and $R>0$ (Lemma~\ref{T13:lem:pos}). All enclosures have width below $10^{-20}$. Below, ``$x\approx c$'' means that the certified enclosure of $x$ lies within one unit of the last displayed digit of $c$ (these displayed values are for orientation only; the proof uses only the seven certified checks):
\begin{center}
\begin{tabular}{ll}
$t_1\approx0.91804979$, $t_2\approx0.0097571452$ & $D\approx2.457388918$, $e^D\approx11.6742892$\\
$r=R/P\approx5.27195546\times10^{-3}$ & $re^D\approx0.0615463326$, $re^{2D}\approx0.718509685$
\end{tabular}
\end{center}

\subsection{Witness for (ii)(c)}
$V_c=\{\emptyset\}\cup\{\{j\}:j\in[40]\}\cup\{\{j_i,k\}:i=1,2,\ k\in N\}$, $|V_c|=117$. All its members have size $\le2\le K$, so $V_c\subseteq\Omega$. With pool $\Pool=V_c$: $|\Pool|\le 20p=800$, so both top-$20p$ lists equal $\Pool$, and $V(\Pool)=\Pool\cup\{\emptyset\}=V_c$; $V_c$ is of the method's form.
No member of $V_c$ contains both $j_1,j_2$, so $d\equiv0$ on $V_c$ and $\Zb_\psi(V_c)=\Zb_0(V_c)=:Z_c$ for all $\psi$. Hence $F_{V_c}$ is constant and $\rho_{V_c}(u,v)=1$ for all $u,v$.
\emph{Grid.} An implementation exists: e.g.\ test the pilot temperature $b$ first; if it fails, test the midpoints of $(u,b]$ successively (halving towards $u$ after a failure and towards $b$ after a success, at most 60 halvings), append the largest tested point that passed, and fail if none passed. In both witnesses (\S5.2 and \S5.3) every pilot temperature passes on its first test, so no bisection, endpoint choice or stopping rule is ever used, and every implementation satisfying STATEMENT \S0 gives the same output. Pilot list $(1)$, tier $0$: at $u=0$, $b=1$, $\rho_{V_c}(0,1)=1\ge0.90$, so $b$ is appended (strict progress $1>0$; inequality validated; $1\le T_e$ intervals); the list is exhausted and the output is the grid $(0,1)$, $T=1$, $w_1=2$. The interval $(0,1)$ is accepted.
\emph{Coverage.} $a_c(\psi)=Z_c/(P+Re^{\psi D})$ is strictly decreasing in $\psi$ ($R>0$, $D>0$), where
$Z_c=\varpi(0)+\varpi(1)n^{-1/2}(38+g_1+g_2)+38\,\varpi(2)n^{-1}(g_1+g_2)$, with certified $Z_c/P\approx0.984162996$. Hence, by the certified checks,
\[
\min_{\psi\in[0,1]}a_c(\psi)=a_c(1)=\frac{Z_c/P}{1+re^{D}}\approx0.92710319\ \ge0.9,\qquad
a_c(2)=\frac{Z_c/P}{1+re^{2D}}\approx0.57268400\ <0.9 .
\]
\emph{Divergence.} By Lemma~\ref{T13:lem:sm} with $A=\Omega$,
\[
1+\chisq(q_1\|q_0)=\frac{\Zb_0\Zb_2}{\Zb_1^2}=\frac{(1+r)(1+re^{2D})}{(1+re^D)^2}\approx1.53305467\ >\ \frac{10}{9\cdot0.81}=\frac{1000}{729}\approx1.37174211
\]
(certified: the lower endpoint of the enclosure of the left side exceeds the upper endpoint of the enclosure of $1000/729$).
So with $\varepsilon_C=0.1$ the interval is accepted and coverage holds on $[0,1]$, while $q_{w_1}(V_c)<1-\varepsilon_C$ and the bound of (ii)(a) fails. This proves (ii)(c).

\subsection{Witness for (ii)(d)}
$V_d=\{\emptyset\}\cup\{\{j_1\}\}\cup\{\{j_1,k\}:k\in N\}\cup\{S:\{j_1,j_2\}\subseteq S,\ |S|\le4\}$, $|V_d|=1+1+38+(1+38+703)=782\le800$; sizes $\le4\le K$; as above $V(V_d)=V_d$ is of the method's form. By Lemma~\ref{T13:lem:cf}, $\Zb_\psi(V_d)=A+Be^{\psi D}$ with
\[
A=\varpi(0)+g_1\bigl(\varpi(1)n^{-1/2}+38\varpi(2)n^{-1}\bigr),\qquad
B=g_1g_2\bigl(\varpi(2)n^{-1}+38\varpi(3)n^{-3/2}+703\varpi(4)n^{-2}\bigr).
\]
Let $\alpha=A/P\approx0.984162996$ (this agrees with $Z_c/P$ to the digits shown but is a different number: $Z_c-A=\varpi(1)n^{-1/2}(38+g_2)+38\varpi(2)n^{-1}g_2>0$ is tiny relative to $P$; the certified enclosures in Appendix~A are disjoint), $\beta=B/R\approx0.996250785$ (certified: $\beta-\alpha>0$, enclosure $\approx0.0120877894$) and $\lambda(\psi)=Re^{\psi D}/(P+Re^{\psi D})$, strictly increasing in $\psi$. Then
$a_d(\psi)=(A+Be^{\psi D})/(P+Re^{\psi D})=(1-\lambda(\psi))\alpha+\lambda(\psi)\beta$, strictly increasing because $\beta>\alpha$.
\emph{Grid.} Pilot list $(\tfrac12,1)$, tier $0$, any implementation. With $s=B/A\approx5.33670722\times10^{-3}$,
\[
\rho_{V_d}(u,v)=\frac{(1+se^{vD})^2}{(1+se^{uD})(1+se^{(2v-u)D})},\qquad
\rho_{V_d}(0,\tfrac12)\approx0.97081344,\quad \rho_{V_d}(\tfrac12,1)\approx0.91376273,
\]
and the certified checks give $\rho_{V_d}(0,\tfrac12)\ge0.9$ and $\rho_{V_d}(\tfrac12,1)\ge0.9$. So the construction appends $\tfrac12$ and then $1$ without bisection (strict progress, validated inequalities, $2\le T_e$ intervals) and outputs $(0,\tfrac12,1)$, $T=2$, $w_1=1$, $w_2=\tfrac32$, $\E=\{0,\tfrac12,1,\tfrac32\}$; both intervals are accepted.
\emph{Coverage.} Since $a_d$ is strictly increasing, the minimum over $\E$ is attained only at $\varphi_0=0$, and the minimum over $\{w_1,w_2\}$ is attained at $w_1=1<\max_tw_t=w_2$. (For illustration: $a_d(0)\approx0.98422639$, $a_d(\tfrac12)\approx0.98437688$, $a_d(1)\approx0.98486382$, $a_d(\tfrac32)\approx0.98626327$; COVER$(V_d,\varepsilon_C)$ holds with $\varepsilon_C=0.016$, and the worst point is $\varphi_0$, not the largest $w_t$. Only $\beta>\alpha$ and the two acceptance inequalities are used.) This proves (ii)(d). \qed

\section{Part (iv)}
\emph{(a).} Both top lists of $V(\Pool)$ are subsets of $\Pool$, so $V\subseteq\Pool\cup\{\emptyset\}$ and $V\cap\M\subseteq(\Pool\cap\M)\cup(\{\emptyset\}\cap\M)=\emptyset$ under pool avoidance and $\emptyset\notin\M$. If calibration never scores a model of $\M$ and every member of $\Pool$ was scored by the calibration or belongs to an inherited set of scored models that avoids $\M$, then every member of $\Pool$ lies outside $\M$, i.e.\ $\Pool\cap\M=\emptyset$.

\emph{(b).} Suppose $V\cap\M=\emptyset$ and let $\gamma'_\psi=\gamma_\psi\mathbf 1_{\M^c}$. On $\Omega'=\Omega\setminus\M$ (nonempty, as $\emptyset\in\Omega'$), $\gamma'_\psi=\gamma_0|_{\Omega'}\,e^{\psi d|_{\Omega'}}$ with $\gamma_0|_{\Omega'}>0$ (Lemma~\ref{T13:lem:pos}), so $\gamma'$ is the bridge on $\Omega'$ with restricted $\gamma_0,d$, and $q'_\psi=\gamma_\psi\mathbf 1_{\M^c}/\Zb_\psi(\M^c)=q_\psi(\cdot\mid\M^c)$. Since $V\subseteq\M^c$, $\gamma'_0(S)e^{td'(S)}=\gamma_0(S)e^{td(S)}$ for every $S\in V$, hence $F'_V(t)=F_V(t)$ for all real $t$ and $\rho'_V(u,v)=\exp\{2F'_V(v)-F'_V(u)-F'_V(2v-u)\}=\rho_V(u,v)$ for all real $u,v$.
Fix one implementation $\mathcal I$ and the inputs (pilot list, $V$, $T_e$). By definition, $\mathcal I$ is deterministic and accesses the bridge only through evaluations of $\rho_V$; each of its steps (the next query point, a comparison with $0.90$, a validation, appending, failing, stopping) is a function of the inputs and of the answers to the previous queries. Run $\mathcal I$ on both bridges. By induction on $k$: if the first $k-1$ queries and their answers coincide, then the $k$-th query point coincides (it is determined by the same inputs and answers), and its answer coincides because $\rho'_V=\rho_V$ at every real pair; under EX-A every comparison of the answer with $0.90$ is exact, so the same branch is taken. Hence the whole transcripts coincide, and so do the outcomes: the same grid, or failure in both. \qed

\begin{T13remark}
Nothing above concerns particles: the statements are identities and inequalities between exact bridge quantities, and no assertion derives particle mixing from ESS.
\end{T13remark}

\appendix
\section{Certified computation of the witness inequalities}
Script \texttt{witness\_certified.py} (run with \texttt{python -I witness\_certified.py}; Python 3 with mpmath):
\begin{lstlisting}
"""Certified enclosures for the T13 witnesses (proof Section 5), by interval arithmetic.

All rational inputs (varpi(k), t_i, nu, binomial counts) are computed EXACTLY with Python integers / Fractions and
converted to intervals that contain them; every transcendental operation (sqrt, log, exp) uses mpmath.iv, whose
operations return intervals guaranteed to contain the exact result (outward rounding). Each internal interval
therefore contains the exact real value, and each inequality is decided by comparing internal interval endpoints.
For display, the internal binary endpoints are converted exactly to rationals and rounded OUTWARD to 22 decimals
(lower endpoint down, upper endpoint up), so every printed interval also contains the exact value.
Run: python -I witness_certified.py
"""
from fractions import Fraction as Fr
from math import comb
from mpmath import iv

iv.dps = 50
n, p = 42, 40
K = min(p, n // 2)                       # 21
nu = Fr(1, 100) + Fr(n, 2)               # 21.01 = 2101/100
b0 = Fr(1, 100)
c1, c2 = Fr(97, 100), Fr(10, 100)
w = lambda k: Fr(1, 2) / comb(2 * p - 1, k)
WK = sum(comb(p, k) * w(k) for k in range(K + 1))
vp = lambda k: w(k) / WK                 # exact rationals
t = [Fr((n - 1) ** 2) * c * c / (n * (n - 1 + 2 * b0)) for c in (c1, c2)]   # exact rationals
assert t[0] + t[1] < 1

def I(q):  # interval containing the rational q
    return iv.mpf(q.numerator) / iv.mpf(q.denominator)

NU, N = I(nu), iv.mpf(n)
rs = 1 / iv.sqrt(N)                      # n^{-1/2}
g = [iv.exp(-NU * iv.log(I(1 - ti))) for ti in t]                 # (1 - t_i)^{-nu}
D = NU * iv.log(I((1 - t[0]) * (1 - t[1]) / (1 - t[0] - t[1])))   # D
eD = iv.exp(D)
Nn = p - 2
def sigma(a):
    tot = iv.mpf(0)
    for m in range(0, min(Nn, K - a) + 1):
        tot += iv.mpf(comb(Nn, m)) * I(vp(a + m)) * rs ** (a + m)
    return tot
s0, s1, s2 = sigma(0), sigma(1), sigma(2)
P = s0 + (g[0] + g[1]) * s1
R = g[0] * g[1] * s2
r = R / P
# witness (c)
Zc = I(vp(0)) + I(vp(1)) * rs * (Nn + g[0] + g[1]) + I(vp(2)) / N * Nn * (g[0] + g[1])
zc = Zc / P
ac1 = zc / (1 + r * eD)
ac2 = zc / (1 + r * eD * eD)
chi = (1 + r) * (1 + r * eD * eD) / (1 + r * eD) ** 2
bound = iv.mpf(1000) / iv.mpf(729)       # 10/(9*0.81) = 1000/729 exactly
# witness (d)
A = I(vp(0)) + g[0] * (I(vp(1)) * rs + Nn * I(vp(2)) / N)
B = g[0] * g[1] * (I(vp(2)) / N + Nn * I(vp(3)) * rs ** 3 + comb(Nn, 2) * I(vp(4)) / N ** 2)
alpha, beta, s = A / P, B / R, B / A
def rho(u2, v2):  # u = u2/2, v = v2/2 (half-integers); rho = (1+s e^{vD})^2 / ((1+s e^{uD})(1+s e^{(2v-u)D}))
    E = lambda k2: iv.exp(D * iv.mpf(k2) / 2)
    return (1 + s * E(v2)) ** 2 / ((1 + s * E(u2)) * (1 + s * E(2 * v2 - u2)))
ad = lambda k2: (A + B * iv.exp(D * iv.mpf(k2) / 2)) / (P + R * iv.exp(D * iv.mpf(k2) / 2))

def exact_ends(x):  # the interval's two binary endpoints as exact rationals
    def fr(t):
        sign, man, exp, _ = t
        q = Fr(man) * (Fr(2) ** exp)
        return -q if sign else q
    return fr(x._mpi_[0]), fr(x._mpi_[1])
def show(name, x, k=22):  # outward-rounded decimal display: printed [lo, hi] contains the internal interval
    lo, hi = exact_ends(x)
    s = 10 ** k
    L = Fr((lo.numerator * s) // lo.denominator)                    # floor(lo * 10^k)
    H = Fr(-((-hi.numerator * s) // hi.denominator))                # ceil(hi * 10^k)
    fmt = lambda q: (str(q.numerator // s) if q >= 0 else "-" + str((-q.numerator) // s)) + "." + str(abs(q.numerator) % s).zfill(k)
    print(f"{name:18s} in [{fmt(L)}, {fmt(H)}]  (decimals outward-rounded; width < 1e-20: {hi - lo < Fr(1, 10**20)})")
for nm, x in [("t1", I(t[0])), ("t2", I(t[1])), ("D", D), ("e^D", eD), ("r = R/P", r), ("r e^D", r * eD), ("r e^{2D}", r * eD * eD), ("Z_c/P", zc),
              ("a_c(1)", ac1), ("a_c(2)", ac2), ("1+chi2(q1||q0)", chi), ("1000/729", bound),
              ("alpha = A/P", alpha), ("beta = B/R", beta), ("beta - alpha", beta - alpha), ("s = B/A", s),
              ("rho_Vd(0,1/2)", rho(0, 1)), ("rho_Vd(1/2,1)", rho(1, 2)),
              ("a_d(0)", ad(0)), ("a_d(1/2)", ad(1)), ("a_d(1)", ad(2)), ("a_d(3/2)", ad(3))]:
    show(nm, x)
checks = {
    "D > 0": D.a > 0,
    "a_c(1) >= 0.9": ac1.a >= (iv.mpf(9) / 10).b,
    "a_c(2) < 0.9": ac2.b < (iv.mpf(9) / 10).a,
    "1+chi2 > 1000/729": chi.a > bound.b,
    "beta > alpha": (beta - alpha).a > 0,
    "rho(0,1/2) >= 0.9": rho(0, 1).a >= (iv.mpf(9) / 10).b,
    "rho(1/2,1) >= 0.9": rho(1, 2).a >= (iv.mpf(9) / 10).b,
}
for k, v in checks.items(): print(f"CHECK {k:22s}: {v}")
print("ALL CHECKS PASS" if all(checks.values()) else "SOME CHECK FAILS")
\end{lstlisting}
Output (each printed interval is the internal enclosure rounded outward to 22 decimals, so it also contains the exact value):
\begin{lstlisting}
t1                 in [0.9180497898818230363817, 0.9180497898818230363818]  (decimals outward-rounded; width < 1e-20: True)
t2                 in [0.0097571451788906688955, 0.0097571451788906688956]  (decimals outward-rounded; width < 1e-20: True)
D                  in [2.4573889178992225124231, 2.4573889178992225124232]  (decimals outward-rounded; width < 1e-20: True)
e^D                in [11.6742891816087945733827, 11.6742891816087945733828]  (decimals outward-rounded; width < 1e-20: True)
r = R/P            in [0.0052719554605095421403, 0.0052719554605095421404]  (decimals outward-rounded; width < 1e-20: True)
r e^D              in [0.0615463325985499584319, 0.0615463325985499584320]  (decimals outward-rounded; width < 1e-20: True)
r e^{2D}           in [0.7185096848229484693078, 0.7185096848229484693079]  (decimals outward-rounded; width < 1e-20: True)
Z_c/P              in [0.9841629956846136023762, 0.9841629956846136023763]  (decimals outward-rounded; width < 1e-20: True)
a_c(1)             in [0.9271031941446113197618, 0.9271031941446113197619]  (decimals outward-rounded; width < 1e-20: True)
a_c(2)             in [0.5726839972892024642265, 0.5726839972892024642266]  (decimals outward-rounded; width < 1e-20: True)
1+chi2(q1||q0)     in [1.5330546700532342666175, 1.5330546700532342666176]  (decimals outward-rounded; width < 1e-20: True)
1000/729           in [1.3717421124828532235939, 1.3717421124828532235940]  (decimals outward-rounded; width < 1e-20: True)
alpha = A/P        in [0.9841629956846136023757, 0.9841629956846136023758]  (decimals outward-rounded; width < 1e-20: True)
beta = B/R         in [0.9962507850715599508558, 0.9962507850715599508559]  (decimals outward-rounded; width < 1e-20: True)
beta - alpha       in [0.0120877893869463484800, 0.0120877893869463484801]  (decimals outward-rounded; width < 1e-20: True)
s = B/A            in [0.0053367072216948640454, 0.0053367072216948640455]  (decimals outward-rounded; width < 1e-20: True)
rho_Vd(0,1/2)      in [0.9708134408716465566174, 0.9708134408716465566175]  (decimals outward-rounded; width < 1e-20: True)
rho_Vd(1/2,1)      in [0.9137627308499763841343, 0.9137627308499763841344]  (decimals outward-rounded; width < 1e-20: True)
a_d(0)             in [0.9842263877716183656109, 0.9842263877716183656110]  (decimals outward-rounded; width < 1e-20: True)
a_d(1/2)           in [0.9843768807694762001798, 0.9843768807694762001799]  (decimals outward-rounded; width < 1e-20: True)
a_d(1)             in [0.9848638215300282733429, 0.9848638215300282733430]  (decimals outward-rounded; width < 1e-20: True)
a_d(3/2)           in [0.9862632653219948237377, 0.9862632653219948237378]  (decimals outward-rounded; width < 1e-20: True)
CHECK D > 0                 : True
CHECK a_c(1) >= 0.9         : True
CHECK a_c(2) < 0.9          : True
CHECK 1+chi2 > 1000/729     : True
CHECK beta > alpha          : True
CHECK rho(0,1/2) >= 0.9     : True
CHECK rho(1/2,1) >= 0.9     : True
ALL CHECKS PASS
\end{lstlisting}
\section{Verbatim quotations of the definitions used}
The passages below are copied mechanically (by line number) from \texttt{00\_COMMON.md} (SHA-256 \texttt{a7d05134350823efd62add9d66a8ec1e8ff78185f2efc78415a504cc819e023d}) and \texttt{METHODOLOGY\_v3.md} (SHA-256 \texttt{93043758058a6d422cca6ca4b6be8f3a9e8149b6f6de2e1bab1ce31c44b8019c}), the files STATEMENT.md refers to. For display only, lines longer than 100 characters are wrapped at spaces, and continuation lines begin with ``\texttt{\ \ |\ }''; removing these markers and rejoining recovers the source text (checked by the extraction script). Markdown markup is kept as in the source.
\begin{Verbatim}[fontsize=\footnotesize,formatcom=\qfont]
--- 00_COMMON.md, §1 Data, target (first three bullets), lines 9-11 ---
- Raw data (X^r, y^r), n ≥ 3, no constant column, non-constant y^r. C = I − 11ᵀ/n; D_x = diag of
  | retained raw column sample SDs. Standardised X = CX^rD_x⁻¹, y = Cy^r/s_y with s_y² = ‖Cy^r‖²/(n−1);
  | so ‖x_j‖² = ‖y‖² = n − 1. Contrast coordinates: U ∈ ℝ^{n×(n−1)}, UᵀU = I, UUᵀ = C; y_c = Uᵀy^r, X_c
  | = UᵀX.
- State space Ω = {S ⊆ [p] : |S| ≤ K}, K = min(p, ⌊n/2⌋). G_S = X_SᵀX_S = (n−1)R_S; z_S = X_Sᵀy; H_S
  | = (n+1)G_S + I; D_S = det H_S/det(G_S+I); Q_S = yᵀy − n z_SᵀH_S⁻¹z_S; ν = 0.01 + n/2; a₀ = b₀ =
  | 0.01.
- m(S) = (2b₀)^{a₀}Γ(ν)/(π^{n/2}Γ(a₀)) · D_S^{−1/2}(2b₀ + Q_S)^{−ν}; w(k) = ½/C(2p−1, k); W_K =
  | Σ_{k≤K} C(p,k)w(k); ϖ(S) = w(|S|)/W_K.

--- 00_COMMON.md, §2 Sampler objects (blocks, start, bridge), lines 18-20 ---
- Blocks C₁, …, C_J (deterministic in X: complete linkage, top-down cut, next-fit packing), |C_b| ≤
  | B = min(⌊√n⌋, ⌊√p⌋, 16). Tables r_b(A) = BF(A) for A ⊆ C_b (full response).
- γ₀(S) = ϖ(S)∏_b r_b(S ∩ C_b); π₀ = γ₀/Z̄₀ (Z̄₀ by exact convolution of block size profiles with
  | the capped size prior; backward sampler gives exact iid π₀ draws). Start mode S̃₀ := argmax π₀.
- Discrepancy d(S) = log BF(S) − Σ_b log r_b(S ∩ C_b). Bridge γ_φ = γ₀e^{φd}, q_φ = γ_φ/Z̄_φ
  | (defined for real φ; production φ ∈ [0,1]); q₀ = π₀, q₁ = π, Z = m(∅)Z̄₁.

--- 00_COMMON.md, §2 Budgets and exact branch, lines 24-25 ---
- Budgets: N₀ = max(2500, n, p); H₀ = ⌈512N₀/2500⌉; ℓ = ⌈log₂(n+p+1)⌉; T₀ = max(1000, ℓ²); (N_e,
  | H_e, T_e) = 4^e(N₀, H₀, T₀), e ∈ {0,1}; L = max(4, ⌈log₂(n+1)⌉). Q_f := 4N_e(T + L + 5) iid π₀ draws
  | scored in the final run (initialisation + independence proposals); 4N_e(L+6) ≤ Q_f ≤ 4N_e(T_e + L +
  | 5).
- Exact branch: Σ_{k≤K} C(p,k) ≤ (5H₀ + 4N₀)p; it precedes block/Gram construction and scores every
  | feasible model.

--- METHODOLOGY_v3.md, §3 Notation and budgets (table rows), lines 62-73 ---
| Symbol | Definition | n = 200, p = 500 | n = 200, p = 1,000 |
|---|---|---|---|
| p | number of retained (non-constant) predictors | 500 | 1,000 |
| K | size cap, min(p, ⌊n/2⌋) | 100 | 100 |
| B | block size, min(⌊√n⌋, ⌊√p⌋, B*), B* = 16 (implementation limit) | 14 | 14 |
| ℓ | ⌈log₂(n + p + 1)⌉ (local-search and rebuild budget) | 10 | 11 |
| L | max(4, ⌈log₂(n + 1)⌉) (burn-in rounds at φ = 1) | 8 | 8 |
| N₀ | max(2,500, n, p) final particles per island | 2,500 | 2,500 |
| H₀ | ⌈512·N₀/2,500⌉ calibration particles (pilot and each diagnostic island) | 512 | 512 |
| T₀ | max(1,000, ℓ²) stage ceiling | 1,000 | 1,000 |
| R | 4 islands (final and diagnostic) | 4 | 4 |
| tier e | e ∈ {0, 1}; (N_e, H_e, T_e) = 4^e (N₀, H₀, T₀) | | |

--- METHODOLOGY_v3.md, §4 Preprocessing (items 1–4), lines 80-87 ---
1. Stop with `invalid_input` if any entry of X or y is not finite, or n < 3, or p < 1.
2. Drop column j only if all its entries are **exactly** equal; record it; it is never selected. p
  | := number kept. If p = 0,
   return Ŝ = ∅, PIPs and evidence unavailable, status `no_predictors`.
3. If all entries of y are exactly equal, return Ŝ = ∅ (the exact HPM: every Q_S = 0, D_S ≥ 1 and w
  | is decreasing), PIPs and
   evidence unavailable, status `y_constant`. (The exact PIPs are not 0; they are not computed.)
4. Standardise stably: for each column, divide by its maximum absolute value, compute the mean,
  | centre, then divide by the
   sample SD (denominator n − 1), so ‖x_j‖² = n − 1. Same for y, recording s_y² = ‖Cy‖²/(n − 1) with
  | C = I − 11ᵀ/n.
   Original indices, means and scales are kept for reporting and back-transformation.

--- METHODOLOGY_v3.md, §5 Target posterior, lines 97-108 ---
On the standardised data (X, y), for S ⊆ {1, …, p} with k = |S| ≤ K, G_S = X_SᵀX_S, z_S = X_Sᵀy:

```text
y | β_S, σ², S ~ N_n(X_S β_S, σ² I_n)                   (no intercept; y is centred)
β_S | σ², S    ~ N_k(0, n σ² (G_S + I_k)⁻¹)              g = n (unit information), τ = 1 (one
  | pseudo-observation)
σ²             ~ IG(a₀, b₀),  a₀ = b₀ = 0.01
ϖ(S)           = w(k) 1{k ≤ K} / W_K,   w(k) = B(1+k, 2p−k)/B(1, p) = ½ / C(2p−1, k),   W_K =
  | Σ_{j≤K} C(p, j) w(j)

H_S = (n + 1) G_S + I_k,   Q_S = yᵀy − n z_Sᵀ H_S⁻¹ z_S,
log m(S) = log C_n − ½ [log det H_S − log det(G_S + I_k)] − (a₀ + n/2) log(2b₀ + Q_S),
C_n = (2b₀)^{a₀} Γ(a₀ + n/2) / (π^{n/2} Γ(a₀)),
π(S | y) = ϖ(S) m(S) / Z,   Z = Σ_{|S|≤K} ϖ(S) m(S).

--- METHODOLOGY_v3.md, §6.2 Exact branch, lines 154-161 ---
### 6.2 Exact branch

If E = Σ_{k≤K} C(p, k) ≤ (5H₀ + 4N₀)·p (12,560·p at baseline; counted with saturating integer
  | arithmetic), score every model
with |S| ≤ K in a streaming pass: running best (strict order of §6.10), log-sum-exp accumulation of
  | Z and of the PIP numerators.
No blocks are built. Status `exact`; `global_certified` when all comparisons are numerically
  | validated. This applies for
p ≤ 17 for n ≥ 22 (p ≤ 18 at n = 20, 21), and for larger p when the cap binds (for example p ≤ 36 at
  | n = 10; all p ≤ 3,000 at n ≤ 5). Because the
threshold grows with N₀ ≥ n, every fixed-p problem is eventually enumerated as n grows.


--- METHODOLOGY_v3.md, §6.3 Blocks, tables and the composite start, lines 162-177 ---
### 6.3 Blocks, tables and the composite start

- **Clustering:** complete linkage on the distance 1 − |r_jl| computed from X; ties between
  | candidate merges broken
  lexicographically on the sorted original-index lists.
- **Cut:** top-down from the root; any node with more than B members is split into its two children
  | (visited smallest index
  first) until every retained cluster has at most B members.
- **Packing:** next-fit in dendrogram leaf order, whole clusters only (a retained cluster is never
  | split), into blocks
  C₁, …, C_J of at most B members. J ≥ ⌈p/B⌉ (packing leaves gaps).
- **Tables:** for each block, every sub-model U ⊆ C_b is scored against the full y; r_b(U) =
  | m(U)/m(∅) (table exponent 1).
  Since B ≤ K for n ≥ 3, no sub-model exceeds the cap. Every table score counts for the HPM. Table
  | scores are not inserted in
  the score cache.
- **Start:** γ₀(S) = ϖ(S) ∏_b r_b(S ∩ C_b), π₀ = γ₀/Z̄₀. Z̄₀ is computed exactly by convolving the
  | block size profiles
  e_b(j) = Σ_{|U|=j} r_b(U) up to K with the size prior (valid because ϖ depends on S only through
  | |S|). Exact iid draws:
  draw the total size, then block sizes backwards, then each block's sub-model within its size
  | class.
- **Discrepancy:** d(S) = log BF(S) − Σ_b log r_b(S ∩ C_b), BF(S) = m(S)/m(∅); bridge γ_φ(S) = γ₀(S)
  | e^{φ d(S)}, φ ∈ [0, 1];
  γ₁ ∝ the posterior; Z = m(∅) Σ_S γ₁(S).

--- METHODOLOGY_v3.md, §6.7 Candidate set and candidate grid (first four bullets), lines 211-222 ---
### 6.7 Candidate set and candidate grid

- **V:** the union of the top 20p scored models by log γ₀, the top 20p scored models by log γ₁ (the
  | target), and ∅;
  deduplicated; list ties lexicographic; frozen during each construction.
- **Restricted log-normaliser:** F_V(t) = log Σ_{S∈V} γ₀(S) e^{t d(S)}.
- **Restricted ESS** between temperatures u < v: ρ_V(u, v) = exp{2F_V(v) − F_V(u) − F_V(2v − u)},
  | the population ESS fraction
  of the weights e^{(v−u)d} under the restricted law q_u(· | V) (evaluating F_V beyond 1 is a
  | finite-sum moment, not a
  production temperature).
- **Construction:** start at u = 0 with the pilot temperatures as targets. For the next pilot
  | temperature b: if
  ρ_V(u, b) ≥ 0.90 append b; otherwise bisect on (u, b] with u fixed, at most 60 halvings, and
  | append the largest tested
  endpoint that passes; continue from the appended point. Validate strict numerical progress and the
  | inequality at every
  appended point. Non-progress, failed validation, or more than T_e intervals → the construction
  | fails.
\end{Verbatim}
\mdproofend
\mdsec{(d) Review record}
\mdsubsec{FINAL\_REPORT.md: statement changes}
\begin{mdblock}
\mdhead{2}{2. Changes to the statement, and why}

All changes were checked by Astra \#1 in three statement checks ({\mdtt review\textbackslash{}\allowbreak{}astra\_\allowbreak{}01.\allowbreak{}.\allowbreak{}03}). Check 1 rejected one reading, which counts as one restart of step 1 (the limit is 3). The final statement is {\mdtt STATEMENT.\allowbreak{}md} §2.

\begin{mdnum}
\item[\mdlbl{1.}] \textbf{ε\_C must be in [0, 1).} For ε\_C ≥ 1 the bound is undefined or false.
\item[\mdlbl{2.}] \textbf{(iv)(a): "calibration never scores 𝓜" became an explicit hypothesis.} The hypothesis is that the scored pool V is drawn from avoids 𝓜. The method does not say whether table scores or leaders carried across tiers feed V's top lists. The source wording is kept as a sufficient condition, with that provenance proviso.
\item[\mdlbl{3.}] \textbf{(ii)(b): each grid is paired with the V of the construction that produced it.}
\item[\mdlbl{4.}] \textbf{(ii)(c)–(d): the two cautionary examples must come from real data run through the method.} That means the method's blocks, tables, γ₀ and d, on the SMC branch, with a V of the method's form and a grid the construction actually outputs. My first version used abstract weights; Astra rejected that as a weakening (restart 1). The witness also names its tier, its T\_e and a deterministic implementation.
\item[\mdlbl{5.}] \textbf{(iii): the bound uses coverage at v.} COVER supplies it.
\item[\mdlbl{6.}] \textbf{(i), (iii) hold for all real u, v,} a free strengthening.
\item[\mdlbl{7.}] \textbf{Positivity is derived, not assumed.}
\item[\mdlbl{8.}] \textbf{(iv)(b) is for one fixed deterministic implementation.}
\item[\mdlbl{9.}] \textbf{(ii)(d) is strengthened:} the minimum is attained \emph{only} at φ₀.
\end{mdnum}

None of these changes the substance of the source theorem. EXP / EXP′ is not needed anywhere.
\end{mdblock}

\mdsubsec{FINAL\_REPORT.md: review history}
\begin{mdblock}
\mdhead{2}{4. Review history and my judgments}

{\small
\begin{longtable}{|>{\raggedright\arraybackslash}p{\dimexpr(\linewidth-12\tabcolsep-7\arrayrulewidth)*47/1000\relax}|>{\raggedright\arraybackslash}p{\dimexpr(\linewidth-12\tabcolsep-7\arrayrulewidth)*75/1000\relax}|>{\raggedright\arraybackslash}p{\dimexpr(\linewidth-12\tabcolsep-7\arrayrulewidth)*66/1000\relax}|>{\raggedright\arraybackslash}p{\dimexpr(\linewidth-12\tabcolsep-7\arrayrulewidth)*113/1000\relax}|>{\raggedright\arraybackslash}p{\dimexpr(\linewidth-12\tabcolsep-7\arrayrulewidth)*103/1000\relax}|>{\raggedright\arraybackslash}p{\dimexpr(\linewidth-12\tabcolsep-7\arrayrulewidth)*597/1000\relax}|}
\hline
\textbf{Stage} & \textbf{Reviewer} & \textbf{Version} & \textbf{Verdict} & \textbf{Findings (fatal / gap / exposition)} & \textbf{My answer} \\ \hline\endhead
A & Astra \#1 & v1 & not verified & 0 / 1 / 2 & Gap accepted: the witness inequalities had no certified error bound. Fixed by interval arithmetic (v2). Exposition fixed. \\ \hline
A & Astra \#1 & v2 & verified & 0 / 0 / 1 & Printed endpoints were rounded to nearest; fixed with outward rounding (v3). \\ \hline
A & Astra \#1 & v3 & verified & none & Stage A proof candidate. \\ \hline
B & Astra \#2 (fresh) & v3 & not verified & 0 / 1 / 0 & Gap accepted: the method definitions were not in the packet. Fixed by Appendix B (v4). \\ \hline
B & Opus xhigh (fresh) & v3 & not verified & 0 / 4 / 5 & The four gaps are the same issue in detail (m(S) formula, n − 1 standardisation, block properties, budgets). All accepted and fixed in v4, with the exposition points. \\ \hline
B & Astra \#2 & v4 & \textbf{verified} & none & Agrees with every logged answer. \\ \hline
B & Opus xhigh & v4 & \textbf{verified} & 0 / 0 / 3 & The provenance note is closed by my mechanical check (9/9 quoted passages match the source lines; checksums match). The other two notes need no change (the reviewer agrees). \\ \hline
\end{longtable}}

\begin{mdbul}
\item No reviewer found a mathematical error at any stage. Every finding was about rigour or completeness of the record.
\item Every gap finding, verbatim with my answer and the reviewers' rulings, is in {\mdtt CLAIMS\_\allowbreak{}OF\_\allowbreak{}FATAL\_\allowbreak{}FLAWS.\allowbreak{}md}.
\item Raw replies are in {\mdtt review\textbackslash{}\allowbreak{}}: the Astra logs {\mdtt astra\_\allowbreak{}0*.\allowbreak{}log} and {\mdtt stageB\_\allowbreak{}astra2\_\allowbreak{}0*.\allowbreak{}log}, and the Opus replies {\mdtt stageB\_\allowbreak{}opus\_\allowbreak{}0*.\allowbreak{}md}.
\item \textbf{Astra judge: not used.} There was no dispute: both Stage B reviewers verified v4 in their second round.
\end{mdbul}
\end{mdblock}

\mdsubsec{CLAIMS\_OF\_FATAL\_FLAWS.md (every entry, verbatim)}
\begin{mdblock}
\mdhead{1}{CLAIMS OF FATAL FLAWS — T13 (T10) restricted ESS}

Every fatal or gap finding, copied verbatim, with the lead's answer. Exposition findings are listed briefly at the end of each round.

\mdhead{2}{Stage A, round 1 — Astra \#1 review of PROOF\_v1 (review\textbackslash{}astra\_05\_proof\_review\_v1.log)}

\mdhead{3}{A1-1 (gap)}

\begin{mdquote}
\textbf{Location:} §5.1, “Numerical values”; §§5.2–5.3, numerical comparisons; Appendix A, particularly {\mdtt mp.\allowbreak{}dps = 40}.

\textbf{Severity:} gap.

\textbf{Claim:} The exact inequalities establishing the witnesses lack a justified numerical error bound.

\textbf{Evidence:} Appendix A evaluates noninteger powers, logarithms, exponentials, sums, and ratios using ordinary finite-precision {\mdtt mpmath} arithmetic. It supplies neither certified enclosures nor an error analysis. Consequently, the displayed approximations alone do not rigorously establish \(a_c(1)\ge0.9\), \(a_c(2)<0.9\), the divergence-bound violation, \(\beta>\alpha\), or the two acceptance inequalities. These comparisons are essential to parts (ii)(c)–(d). EX-A specifies the idealized construction’s arithmetic; it does not certify the appendix’s floating-point calculations. Independent agreement and the substantial numerical margins make this plausibly fixable, but do not supply the missing certification.
\end{mdquote}

\textbf{Lead's answer: accepted and fixed in PROOF\_v2.} All rational inputs (ϖ(k), t\_i, ν, binomial counts) are now computed exactly with integer and fraction arithmetic and converted to enclosing intervals. Every transcendental step (√, log, exp) is evaluated in interval arithmetic with outward rounding ({\mdtt mpmath.\allowbreak{}iv}, 50 digits), so each output interval contains the exact real value. The seven inequalities the witnesses use (D > 0, a\_c(1) ≥ 0.9, a\_c(2) < 0.9, 1+χ² > 1000/729, β > α, ρ\_\{V\_d\}(0,½) ≥ 0.9, ρ\_\{V\_d\}(½,1) ≥ 0.9) are decided by comparing interval endpoints, and all pass. All enclosures are narrower than 10⁻²⁰. The script {\mdtt witness\_\allowbreak{}certified.\allowbreak{}py} and its full output are in Appendix A of PROOF\_v2. §5.1 now states the certification rule, and the displayed decimals are marked "≈" and described as orientation only.

\mdhead{3}{Exposition findings, round 1 (not fatal or gap; handled in PROOF\_v2)}

\begin{mdbul}
\item A1-2: the claimed independent validation (Cholesky check, reduced enumeration) could not be audited from the packet. \textbf{Fixed:} the claim is removed from the proof. The proof does not rely on it; the check itself is in {\mdtt helpers\textbackslash{}\allowbreak{}check\_\allowbreak{}sonnet\textbackslash{}\allowbreak{}}.
\item A1-3: rounded values were written as exact equalities, and 10/(9·0.81) is exactly 1000/729. \textbf{Fixed:} "≈" with a stated meaning, and 1000/729 written exactly.
\end{mdbul}

\mdhead{2}{Stage A, round 2 — Astra \#1 ruling on the log and review of PROOF\_v2 (review\textbackslash{}astra\_06\_ruling\_v2.log)}

Rulings: A1-1 agree, fix works; A1-2 agree, fix works; A1-3 agree for the main text, incomplete in Appendix A. Verdict on PROOF\_v2: \textbf{verified}. No fatal or gap finding. One exposition finding: printed interval endpoints were rounded to nearest, so a printed interval could fail to contain the exact value (e.g. the printed 1000/729 singleton). \textbf{Fixed in PROOF\_v3:} the endpoints are converted exactly to rationals and rounded outward to 22 decimals, and the docstring and §5.1 now say so. The decisive checks were already on internal endpoints and are unchanged.

\mdhead{2}{Stage A, round 3 — Astra \#1 confirmation of PROOF\_v3 (review/astra\_07\_confirm\_v3.log)}

Fix confirmed; no remaining findings; verdict \textbf{verified}. PROOF\_v3 = Stage A proof candidate.

\mdhead{2}{Stage B, round 1 — independent reviews of PROOF\_v3}

\mdhead{3}{B-A2-1 (Astra \#2, gap) — review/stageB\_astra2\_01\_review\_v3.log}

\begin{mdquote}
\begin{mdbul}
\item \textbf{Location:} §1, Lemma 1; §5.1, “The data set” and Lemma 4.
\item \textbf{Severity:} gap.
\item \textbf{Claim:} The identification of the proof’s likelihood, block construction and budgets with the method incorporated into STATEMENT.md cannot be verified from the permitted files.
\item \textbf{Evidence:} STATEMENT.md imports definitions from {\mdtt 00\_\allowbreak{}COMMON.\allowbreak{}md} and {\mdtt METHODOLOGY\_\allowbreak{}v3.\allowbreak{}md} without reproducing the marginal-likelihood formula, hyperparameters, block-size rule, dependence of blocks on \(X\) alone, or budget formulas. The proof supplies these as facts. Their conformity with the imported method is essential: the likelihood determines positivity and the witness scores; the block rule permits choosing \(j_2\) before defining \(y\); and the budgets establish the SMC branch and allowable grids. The subsequent calculations are consistent, but the correspondence to the referenced method remains unchecked. This is a verification limitation, not evidence that those assertions are false.
\end{mdbul}
\end{mdquote}

\textbf{Lead's answer: accepted and fixed in PROOF\_v4.} New Appendix B quotes verbatim, copied mechanically by line number with both files' SHA-256 checksums, every definition the proof uses: 00\_COMMON §1 (data, standardisation, Ω, G\_S, z\_S, H\_S, D\_S, Q\_S, ν, a₀ = b₀ = 0.01, m(S), w, W\_K, ϖ), 00\_COMMON §2 (blocks deterministic in X with |C\_b| ≤ B, tables, γ₀, d, bridge; budgets; exact branch), and METHODOLOGY §3 (budget table), §4 items 1–4 (input checks, standardisation with the n − 1 SD), §5 (target), §6.2 (exact branch), §6.3 (clustering, cut, packing, tables, start, discrepancy) and §6.7 (V, F\_V, ρ\_V, construction). Lemma 1 now restates the definitions with citations. §5.1 now derives the three block properties (X-only, partition, size cap) and the standardisation and budget facts from the quotes.

\mdhead{3}{B-O-1 (Opus, gap) — review/stageB\_opus\_01\_review\_v3.md, finding 1}

\begin{mdquote}
\textbf{Claim:} The proof uses a specific formula for m(S) that I cannot check against its source. STATEMENT.md defines m(S) only by reference to 00\_COMMON §1, which is not in the packet. [...] The proof never states the definition of Q\_S or the values of a₀ and b₀. [...] The only open question is whether this is the m(S) of 00\_COMMON.
\end{mdquote}

(Full text verbatim in review/stageB\_opus\_01\_review\_v3.md, finding 1.)

\textbf{Lead's answer: accepted and fixed in PROOF\_v4.} The m(S) formula with a₀ = b₀ = 0.01, ν = 0.01 + n/2, the definition of Q\_S, H\_S and D\_S are now stated in the proof of Lemma 1, cited to Appendix B (00\_COMMON §1 lines 10–11; METHODOLOGY §5). They match what v3 used.

\mdhead{3}{B-O-2 (Opus, gap) — finding 2}

\begin{mdquote}
\textbf{Claim:} The step "‖x\_j‖² = 41 = n−1, so every sample SD equals 1 … standardisation returns X and y unchanged" assumes 00\_COMMON standardises with the n−1 denominator. That convention is not in the packet.
\end{mdquote}

\textbf{Lead's answer: accepted and fixed in PROOF\_v4.} Both sources use the n − 1 denominator: 00\_COMMON §1 line 9 ("s\_y² = ‖Cy\textasciicircum{}r‖²/(n−1); so ‖x\_j‖² = ‖y‖² = n − 1") and METHODOLOGY §4 item 4 ("divide by the sample SD (denominator n − 1)"). Both are quoted in Appendix B. §5.1 now shows that both forms of the standardisation (00\_COMMON's CX\textasciicircum{}rD\_x⁻¹, and the method's divide by max |x|, centre, divide by SD) return X and y unchanged. Lemma 4 is therefore correct as written.

\mdhead{3}{B-O-3 (Opus, gap) — finding 3}

\begin{mdquote}
\textbf{Claim:} The proof uses three properties of the method's block construction that the packet does not contain. [...] (a) The blocks depend on X only, not on y. [...] (b) The blocks partition [p]. [...] (c) Block size is capped at B = min(⌊√n⌋, ⌊√p⌋, 16).
\end{mdquote}

\textbf{Lead's answer: accepted and fixed in PROOF\_v4.} §5.1 "Blocks" derives the three properties from the quotes:

\begin{mdbul}
\item (a) 00\_COMMON §2 line 18 ("deterministic in X") and METHODOLOGY §6.3 (distance "computed from X"; the cut and packing use only the dendrogram, B and index order).
\item (b) The root holds all p columns; the cut replaces a node by its two children until every retained cluster has at most B members; packing places whole clusters ("a retained cluster is never split") into blocks.
\item (c) |C\_b| ≤ B = min(⌊√n⌋, ⌊√p⌋, 16).
\end{mdbul}

Lemma 4's proof cites (b).

\mdhead{3}{B-O-4 (Opus, gap, low risk) — finding 4}

\begin{mdquote}
\textbf{Claim:} The formulas N₀ = max(2500, n, p), H₀ = ⌈512·2500/2500⌉, ℓ = ⌈log₂(n+p+1)⌉ and T₀ = max(1000, ℓ²), and the constants in the exact-branch condition, come from 00\_COMMON §2 and cannot be checked in the packet. The H₀ expression as printed (2500/2500) does not show which of the two 2500s is N₀.
\end{mdquote}

\textbf{Lead's answer: accepted and fixed in PROOF\_v4.} The formulas are cited to 00\_COMMON §2 lines 24–25 and METHODOLOGY §3 and §6.2, all quoted in Appendix B. H₀ is now written ⌈512N₀/2500⌉ = ⌈512·2500/2500⌉.

\mdhead{3}{Exposition findings, Stage B round 1 (handled in PROOF\_v4)}

\begin{mdbul}
\item B-O-5 (no implementation named): §5.2 "Grid" now names one (test b first, then midpoint bisection), and notes that no bisection is ever used in either witness, so every implementation gives the same output.
\item B-O-6 (reading of the output grid; the frozen-grid parenthetical): §4 (ii)(b) now states the reading explicitly and why it is the only consistent one; the parenthetical is the special case.
\item B-O-7 (a) (the seven checks do not cover every inequality): §5.1 now lists the exact integer and rational facts separately. (b) (re\textasciicircum{}D and re\textasciicircum{}\{2D\} had no printed enclosure): the script now prints certified enclosures for both.
\item B-O-8 (Z\_c/P and α display identically): §5.3 notes that they differ (Z\_c − A > 0, about 5·10⁻²² relative) and that their enclosures are disjoint.
\item B-O-9 (B, C\_b used before definition; |U| ≤ K step unneeded): Lemma 1 now argues m(U) ∈ (0, ∞) for every U ⊆ [p], without using the size cap.
\end{mdbul}

\mdhead{2}{Stage B, round 2 — rulings on this log and review of PROOF\_v4}

\begin{mdbul}
\item \textbf{Astra \#2} (review/stageB\_astra2\_02\_ruling\_v4.log): agrees with the lead's answer on every logged item (A1-1..3, Stage A round 2, B-A2-1, B-O-1..9), and every fix works. No new or remaining findings. \textbf{Verdict: verified.}
\item \textbf{Opus reviewer} (review/stageB\_opus\_02\_ruling\_v4.md): agrees with the lead's answer on every logged item, and every fix works. No fatal or gap finding. \textbf{Verdict: verified.} Three exposition notes: 1. Provenance of the Appendix B quotes could not be checked from the packet. \textbf{Closed by the lead:} helpers/proof\_build/check\_appendix\_B.py compares the quotes inside PROOF\_v4.tex with the cited source lines. Result: 9/9 passages match, and both SHA-256 checksums in the proof equal those of the current files (output in helpers/proof\_build/check\_appendix\_B\_output.txt). 2. The example implementation in §5.2 is a sketch. \textbf{No change:} the reviewer agrees that existence is not in doubt, and that no unspecified branch is reached in either witness. 3. Lemma 1 extends the m formula to every U ⊆ [p], beyond the source's |S| ≤ K. \textbf{No change:} the reviewer rates it harmless; positivity is needed only where m is defined.
\end{mdbul}

No fatal or gap finding is open. Astra judge not needed.
\end{mdblock}

\chapter{T14(i) (L-FM, singleton case): first-moment evidence consistency}
\chaptersource{Final proof: \texttt{proofs/\allowbreak{}T14\_\allowbreak{}LFM\_\allowbreak{}first\_\allowbreak{}moment\_\allowbreak{}evidence/\allowbreak{}proof/\allowbreak{}PROOF\_v3.tex}. Sources of parts (a), (b), (d): \texttt{theorems/\allowbreak{}T14\_\allowbreak{}LFM\_\allowbreak{}first\_\allowbreak{}moment\_\allowbreak{}evidence.md}, \texttt{STATEMENT.md}, \texttt{FINAL\_REPORT.md}, \texttt{CLAIMS\_OF\_FATAL\_FLAWS.md}.}
\mdsec{(a) Original conjecture text (theorem file, verbatim)}
\begin{mdblock}
\mdhead{1}{T14 (L-FM): First-moment evidence consistency}

\textbf{Status:} conjectured, not proven. Methodology SMC-BVS v3.0 (2026-10-19), SHA-256 93043758058a6d422cca6ca4b6be8f3a9e8149b6f6de2e1bab1ce31c44b8019c. Read with {\mdtt 00\_\allowbreak{}COMMON.\allowbreak{}md}.

\textbf{Category / rank:} algorithmic convergence (consistency of the evidence); \textbf{core} (lemma).

\mdhead{2}{Definitions used}

Stage laws q\_t, G\_t = e\textasciicircum{}\{Δ\_td\}, empirical population q̂\_\{t−1\} reweighted at stage t, island evidence Ẑ\_r (00\_COMMON §2). For a model S° fixed given 𝓒: η\_t = 1 − q\_t(S°), ε\_* = max\_\{0≤t≤T\}η\_t; R\_t(S) = G\_t(S)/G\_t(S°); A\_t = ∏\_\{v=1\}\textasciicircum{}t q̂\_\{v−1\}(R\_v); b\_t = q₀(S°)/q\_t(S°). For a set 𝔄: ε\_t = 1 − q\_t(𝔄); c\_t = q\_\{t−1\}(G\_t1\_𝔄)/q\_\{t−1\}(𝔄); R\_t = G\_t/c\_t; A\_t, b\_t = q₀(𝔄)/q\_t(𝔄) analogously; Φ\_t = E[A\_\{t−1\}|q̂\_\{t−1\}((R\_t − 1)1\_𝔄)|]/b\_\{t−1\}.

\mdhead{2}{Assumptions}

EX-A, EX-S, CAL; exact start; prescribed multinomial resampling and q\_φ-invariant rounds, one round per stage.

\mdhead{2}{Precise statement}

Conditional on a successful 𝓒: (i) [singleton] Ẑ\_r/Z = A\_T/b\_T, and for every N ≥ 1 (ε\_* < 1):\\
\hspace*{1.5em}E|Ẑ\_r/Z − 1| ≤ [Σ\_\{t=1\}\textasciicircum{}T(b\_\{t−1\}η\_\{t−1\} + b\_tη\_t) + |1 − b\_T|]/b\_T ≤ (2/(1−ε\_*))Σ\_\{t=0\}\textasciicircum{}T η\_t. The endpoint term |1 − b\_T| is essential. Convexity gives the same bound for the pooled Ẑ. (ii) [consequence] Along launched families, Σ\_\{t=0\}\textasciicircum{}T η\_t → 0 in data/calibration probability implies Ẑ/Z → 1 in joint probability. No over-tempered condition and no N\textasciicircum{}\{−1/2\} rate is asserted. (iii) [general set] Ẑ\_r/Z = A\_T/b\_T with the set definitions, and E|Ẑ\_r/Z − 1| ≤ [2Σ\_\{t=0\}\textasciicircum{}T ε\_t + Σ\_\{t=1\}\textasciicircum{}T Φ\_t]/(1 − max\_tε\_t). The signed expression inside Φ\_t's absolute value has weighted mean zero; Φ\_t ≥ 0 need not vanish (it vanishes for a singleton); Φ\_t ≤ q\_\{t−1\}(|R\_t − 1|1\_𝔄), which need not be summably small. Set mass alone does not control within-set weight dispersion; reflection is one possible further route, not an automatic transfer.

\mdhead{2}{Dependencies}

Builds on: T12. Used by: T15, T25 (first-moment alternative), T26.

\mdhead{2}{Numerical evidence (provisional)}

provisional evidence (Sol check J6): exact finite-state laws (N ≤ 2, T = 3): a rare state with singleton deficits (10⁻¹², 3·10⁻⁸, 10⁻³) has exact L¹ error 0.0020 ≤ telescoping-with-endpoint 0.00200 ≤ bound 0.00200, while E(Ẑ/Z)² ≈ 10⁶; a two-state set with full set mass has L¹ error 1.996 = Φ. Files: judge/sol2/J6\_FIRST\_MOMENTS.md.

\mdhead{2}{Proof-route hint (optional)}

Use E[A\_\{t−1\}q̂\_\{t−1\}(f)] = b\_\{t−1\}q\_\{t−1\}(f) (T12), the exact identity q\_\{t−1\}(R\_t1\_\{S°ᶜ\}) = (b\_t/b\_\{t−1\})q\_t(S°ᶜ) and R\_t(S°) = 1 to get E|A\_t − A\_\{t−1\}| ≤ b\_\{t−1\}η\_\{t−1\} + b\_tη\_t; telescope A\_t; retain |1 − b\_T|; note b\_t/b\_T = q\_T(S°)/q\_t(S°) and |1 − b\_T|/b\_T = |q\_T(S°) − q₀(S°)|/q₀(S°). For sets, choose c\_t so that R\_t has q\_\{t−1\}(·|𝔄)-mean 1 and separate inside/outside contributions. Hardest step: Φ\_t for larger sets. \emph{When proof.md is run on this file, its lead considers the statement both with and without these hints; the hints never reach the reviewers.}
\end{mdblock}

\mdsec{(b) STATEMENT.md (verbatim)}
\begin{mdblock}
\mdhead{1}{STATEMENT — T14(i) (L-FM, singleton case): first-moment evidence consistency}

Source: {\mdtt Theory\_\allowbreak{}architecture\_\allowbreak{}v3/\allowbreak{}theorems/\allowbreak{}T14\_\allowbreak{}LFM\_\allowbreak{}first\_\allowbreak{}moment\_\allowbreak{}evidence.\allowbreak{}md}, part (i) only. Method: SMC-BVS v3.0, METHODOLOGY\_v3.md (SHA-256 93043758058a6d422cca6ca4b6be8f3a9e8149b6f6de2e1bab1ce31c44b8019c). Parts (ii) and (iii) of the file are not part of this run.

Statement analysis outcome: \textbf{(a) correct as stated.} No assumption was changed or added. EXP / EXP′ is not needed. No other v3 conjecture is used as a hypothesis (the unbiasedness identity of T12 that the theorem builds on is re-derived inside the proof).

\mdhead{2}{1. Setting (from 00\_COMMON.md §§1–2 and METHODOLOGY §§6.3, 6.5, 6.6, 6.8, 6.12)}

\begin{mdbul}
\item Ω = \{S ⊆ [p] : |S| ≤ K\} is a finite set. For S ∈ Ω: γ₀(S) = ϖ(S)∏\_b r\_b(S ∩ C\_b) > 0 (the size prior ϖ and the block tables r\_b = BF are strictly positive), d(S) ∈ ℝ, and for real φ, γ\_φ(S) = γ₀(S)e\textasciicircum{}\{φ d(S)\}, Z̄\_φ = Σ\_\{S∈Ω\} γ\_φ(S), q\_φ = γ\_φ/Z̄\_φ. Hence q\_φ(S) > 0 for every S ∈ Ω. At temperature φ = 0 the bridge is π₀, at φ = 1 it is π, and Z = m(∅)·Z̄\_φ evaluated at φ = 1, with m(∅) > 0. (Below, a numerical subscript on q or Z̄ is a stage index; since φ\_T = 1, Z = m(∅)Z̄\_T and q\_T = π.)
\item The calibration 𝓒 fixes the grid 0 = φ₀ < φ₁ < … < φ\_T = 1 (so T ≥ 1) and the per-island particle count N (= N\_e). Write Δ\_t = φ\_t − φ\_\{t−1\}, q\_t = q\_\{φ\_t\}, Z̄\_t = Z̄\_\{φ\_t\}, G\_t(S) = e\textasciicircum{}\{Δ\_t d(S)\} (t = 1, …, T).
\item \textbf{Final island r} (one island; N ≥ 1 particles), given the data and 𝓒: 1. ξ₀¹, …, ξ₀ᴺ are iid with law q₀ (exact start). 2. For t = 1, …, T: (reweighting) particle i gets weight G\_t(ξ\_\{t−1\}ⁱ); (multinomial resampling) ζ\_t¹, …, ζ\_tᴺ are drawn, conditionally on the past, iid from Σ\_i G\_t(ξ\_\{t−1\}ⁱ)δ\_\{ξ\_\{t−1\}ⁱ\} / Σ\_i G\_t(ξ\_\{t−1\}ⁱ); (mutation, one round per stage) conditionally on the past and on ζ\_t, the ξ\_tʲ are independent with ξ\_tʲ \textasciitilde{} K\_t(ζ\_tʲ, ·), where K\_t = K\_\{φ\_t\} is a Markov kernel on Ω leaving q\_t invariant (q\_t K\_t = q\_t). 3. q̂\_t = N⁻¹Σ\_i δ\_\{ξ\_tⁱ\} (empirical population after stage-t resampling and mutation; q̂₀ = the initial draws). 4. Island evidence: Ẑ\_r = m(∅) Z̄₀ ∏\_\{t=1\}\textasciicircum{}T q̂\_\{t−1\}(G\_t), computed exactly. Pooled evidence: Ẑ = ¼Σ\_\{r=1\}\textasciicircum{}4 Ẑ\_r over the four final islands. The burn-in rounds, measured rounds and local search that follow stage T do not enter Ẑ\_r.
\item \textbf{Reference model.} S° ∈ Ω is fixed given 𝓒: a function of the data and 𝓒 only, not of the final-run randomness. Define

\begin{mdbul}
\item η\_t = 1 − q\_t(S°) (t = 0, …, T), ε\_* = max\_\{0≤t≤T\} η\_t;
\item R\_t(S) = G\_t(S)/G\_t(S°) (t = 1, …, T), so R\_t(S°) = 1;
\item A\_t = ∏\_\{v=1\}\textasciicircum{}t q̂\_\{v−1\}(R\_v) (A₀ = 1);
\item b\_t = q₀(S°)/q\_t(S°) (b₀ = 1).
\end{mdbul}
\end{mdbul}

\mdhead{2}{2. Hypotheses}

\begin{mdbul}
\item \textbf{EX-A:} exact scores, normalisers (Z̄₀ in particular), probabilities and arithmetic.
\item \textbf{EX-S + CAL:} ideal independent uniform streams; 𝓒 is measurable with respect to calibration streams 0–9 and the data; the final islands use streams 10–13, independent of 𝓒. Consequently, conditionally on (data, 𝓒), each final island evolves exactly as in §1 item 2.
\item \textbf{Exact start:} the initial particles are exact iid π₀ = q₀ draws.
\item \textbf{Prescribed resampling and moves:} multinomial resampling after every reweighting; one q\_\{φ\_t\}-invariant move round per stage, applied to each particle as in §1 item 2.
\item \textbf{Successful 𝓒:} the frozen grid exists (no {\mdtt schedule\_\allowbreak{}limit}); everything below is conditional on (data, 𝓒). E denotes expectation over the final-run randomness given (data, 𝓒).
\end{mdbul}

\mdhead{2}{3. Conclusion (to prove)}

For every island r and every N ≥ 1 (note ε\_* < 1 always, since q\_t(S°) > 0):

(i-a) Ẑ\_r/Z = A\_T/b\_T, surely.

(i-b) E|Ẑ\_r/Z − 1| ≤ [Σ\_\{t=1\}\textasciicircum{}T (b\_\{t−1\}η\_\{t−1\} + b\_tη\_t) + |1 − b\_T|]/b\_T ≤ (2/(1 − ε\_*)) Σ\_\{t=0\}\textasciicircum{}T η\_t.

(i-c) The same two bounds hold for the pooled Ẑ in place of Ẑ\_r.

(i-d) Reading of the file's sentence "the endpoint term |1 − b\_T| is essential": the middle expression cannot in general be replaced by Σ\_\{t=1\}\textasciicircum{}T (b\_\{t−1\}η\_\{t−1\} + b\_tη\_t)/b\_T. Precisely: there is an instance of the abstract setting of §1 (finite Ω, positive γ₀, real d, a grid, q\_t-invariant kernels K\_t, N ≥ 1, S° ∈ Ω) for which E|Ẑ\_r/Z − 1| > Σ\_\{t=1\}\textasciicircum{}T (b\_\{t−1\}η\_\{t−1\} + b\_tη\_t)/b\_T. (This sentence is an informal remark in the theorem file; this is the reading adopted here. It is not claimed that such an instance arises from an actual SMC-BVS data set.)

\mdhead{2}{4. Change log}

\begin{mdbul}
\item 2026-10-22, after Astra review round 1 (exposition finding 1): notation only. "q₁ = π, Z = m(∅)Z̄₁" in §1 was rewritten so that a numerical subscript is always a stage index (q\_t = q\_\{φ\_t\}); the endpoint is written q\_T, Z̄\_T. No hypothesis or conclusion changed.
\end{mdbul}
\end{mdblock}

\mdsec{(c) Final proof, verbatim: body of PROOF\_v3.tex}
\noindent{\small\itshape The proof below is the body of \texttt{PROOF\_v3.tex} from \texttt{\textbackslash begin\{document\}} to \texttt{\textbackslash end\{document\}} without the title block. Its preamble definitions are local to this chapter; labels carry the prefix \texttt{T14:}.}\par\medskip
\mdproofbegin
% ---- preamble of PROOF_v3.tex (theorem environments carry the prefix T14) ----
\newtheorem{T14theorem}{Theorem}
\newtheorem{T14lemma}[T14theorem]{Lemma}
\newtheorem{T14proposition}[T14theorem]{Proposition}
\theoremstyle{remark}
\newtheorem{T14remark}[T14theorem]{Remark}

\newcommand{\E}{\mathbb{E}}
\newcommand{\F}{\mathcal{F}}
\newcommand{\Om}{\Omega}
\newcommand{\qh}{\widehat q}
\newcommand{\Zh}{\widehat Z}
\newcommand{\Zb}{\bar Z}
\newcommand{\So}{S^{\circ}}
% ---- body of PROOF_v3.tex ----
\section{Setting and statement}

Everything is conditional on the data and on a successful calibration $\mathcal C$; $\E$ denotes expectation over the
final-run randomness given (data, $\mathcal C$). By EX-S and CAL the final islands use streams independent of $\mathcal C$,
so given (data, $\mathcal C$) each island evolves as described below; by EX-A every quantity is computed exactly.

\paragraph{Bridge.} $\Om$ is a finite set. For $S\in\Om$, $\gamma_0(S)>0$ and $d(S)\in\mathbb R$; for real $\varphi$,
$\gamma_\varphi(S)=\gamma_0(S)e^{\varphi d(S)}$, $\Zb_\varphi=\sum_{S\in\Om}\gamma_\varphi(S)$ and
$q_\varphi=\gamma_\varphi/\Zb_\varphi$. Thus $q_\varphi(S)>0$ for all $S\in\Om$. The target evidence is
$Z=m(\emptyset)\sum_{S\in\Om}\gamma_0(S)e^{d(S)}$ (the bridge normaliser at temperature $\varphi=1$) with $m(\emptyset)>0$.

\paragraph{Grid.} $0=\varphi_0<\varphi_1<\dots<\varphi_T=1$ (so $T\ge1$), $\Delta_t=\varphi_t-\varphi_{t-1}$,
$q_t=q_{\varphi_t}$, $\Zb_t=\Zb_{\varphi_t}$, and $G_t(S)=e^{\Delta_t d(S)}>0$ for $t=1,\dots,T$. Since $\varphi_T=1$, $Z=m(\emptyset)\Zb_T$. From here on a numerical subscript on $q$ or $\Zb$ is always a stage index ($q_t=q_{\varphi_t}$), never a temperature.

\paragraph{One final island ($N\ge1$ particles).}
\begin{enumerate}[label=(\alph*)]
\item $\xi_0^1,\dots,\xi_0^N$ are iid with law $q_0$.
\item For $t=1,\dots,T$: conditionally on $\F_{t-1}$ (defined below), $\zeta_t^1,\dots,\zeta_t^N$ are iid with law
  \[
  \mu_t:=\frac{\sum_{i=1}^N G_t(\xi_{t-1}^i)\,\delta_{\xi_{t-1}^i}}{\sum_{i=1}^N G_t(\xi_{t-1}^i)}
  \]
  (multinomial resampling); conditionally on $\F_{t-1}$ and $\zeta_t=(\zeta_t^1,\dots,\zeta_t^N)$, the $\xi_t^j$ are independent
  with $\xi_t^j\sim K_t(\zeta_t^j,\cdot)$, where $K_t$ is a Markov kernel on $\Om$ with $q_tK_t=q_t$ (one move round at
  $\varphi_t$). If these laws are given conditionally on a larger past (for instance one that also contains other
  islands), they hold in particular given $\F_{t-1}$, respectively given $(\F_{t-1},\zeta_t)$, by the tower property,
  because $\mu_t$ is $\F_{t-1}$-measurable and the law of $\xi_t^j$ depends only on $\zeta_t^j$.
\item $\F_{t}=\sigma(\xi_0,\zeta_1,\xi_1,\dots,\zeta_t,\xi_t)$, $\F_0=\sigma(\xi_0)$, and
  $\qh_t=N^{-1}\sum_{i=1}^N\delta_{\xi_t^i}$, so $\qh_t(f)=N^{-1}\sum_i f(\xi_t^i)$.
\item $\Zh_r=m(\emptyset)\Zb_0\prod_{t=1}^T\qh_{t-1}(G_t)$. The pooled evidence is $\Zh=\tfrac14\sum_{r=1}^4\Zh_r$ over four
  islands, each of the above form with the same grid, kernels and $N$.
\end{enumerate}
Only $\qh_0,\dots,\qh_{T-1}$ enter $\Zh_r$; what the island does after stage $T$ (resampling, burn-in and measured rounds,
local search) is irrelevant here. Since $\Om$ is finite, every function $f:\Om\to\mathbb R$ is bounded and all
expectations below are of bounded random variables.

\paragraph{Reference model.} $\So\in\Om$ is fixed given (data, $\mathcal C$). Put
$\eta_t=1-q_t(\So)$ ($0\le t\le T$), $\varepsilon_*=\max_{0\le t\le T}\eta_t$,
$R_t=G_t/G_t(\So)$ ($1\le t\le T$), $A_t=\prod_{v=1}^t\qh_{v-1}(R_v)$ with $A_0=1$, and $b_t=q_0(\So)/q_t(\So)$ with
$b_0=1$. Note $R_t(\So)=1$, $R_t>0$, $A_t>0$, $b_t>0$, and $\varepsilon_*<1$ because $q_t(\So)>0$ for every $t$.

\begin{T14theorem}[T14(i)]\label{T14:thm:main}
For every $N\ge1$:
\begin{enumerate}[label=(i-\alph*)]
\item $\Zh_r/Z=A_T/b_T$ surely;
\item
  \[
  \E\Bigl|\frac{\Zh_r}{Z}-1\Bigr|\;\le\;\frac{\sum_{t=1}^T\bigl(b_{t-1}\eta_{t-1}+b_t\eta_t\bigr)+|1-b_T|}{b_T}
  \;\le\;\frac{2}{1-\varepsilon_*}\sum_{t=0}^T\eta_t ;
  \]
\item the same two bounds hold with the pooled $\Zh$ in place of $\Zh_r$;
\item the endpoint term cannot be dropped in general: there is an instance of the setting above with
  $\E|\Zh_r/Z-1|>\sum_{t=1}^T(b_{t-1}\eta_{t-1}+b_t\eta_t)/b_T$.
\end{enumerate}
\end{T14theorem}

\section{Three exact identities}

\begin{T14lemma}[change of temperature]\label{T14:lem:change}
For $1\le t\le T$ and every $h:\Om\to\mathbb R$,
\[
q_{t-1}(G_t h)=\frac{\Zb_t}{\Zb_{t-1}}\,q_t(h).
\]
\end{T14lemma}
\begin{proof}
For $S\in\Om$,
$q_{t-1}(S)G_t(S)=\gamma_0(S)e^{\varphi_{t-1}d(S)}e^{\Delta_t d(S)}/\Zb_{t-1}=\gamma_{\varphi_t}(S)/\Zb_{t-1}
=(\Zb_t/\Zb_{t-1})\,q_t(S)$. Multiply by $h(S)$ and sum over $\Om$.
\end{proof}

\begin{T14lemma}[ratio of the $b$'s]\label{T14:lem:b}
For $1\le t\le T$, \ $b_t/b_{t-1}=\dfrac{\Zb_t/\Zb_{t-1}}{G_t(\So)}$.
\end{T14lemma}
\begin{proof}
$b_t/b_{t-1}=q_{t-1}(\So)/q_t(\So)$, and by Lemma~\ref{T14:lem:change} with $h=\mathbf 1_{\{\So\}}$,
$q_{t-1}(\So)G_t(\So)=(\Zb_t/\Zb_{t-1})q_t(\So)$. Divide by $G_t(\So)q_t(\So)>0$.
\end{proof}

\begin{T14lemma}[one stage in conditional mean]\label{T14:lem:step}
For $1\le t\le T$ and every $f:\Om\to\mathbb R$,
\[
\E\bigl[\qh_t(f)\bigm|\F_{t-1}\bigr]=\frac{\qh_{t-1}(G_t\,K_tf)}{\qh_{t-1}(G_t)},
\qquad K_tf(x):=\sum_{y\in\Om}K_t(x,y)f(y).
\]
\end{T14lemma}
\begin{proof}
Fix $j$. Given $\F_{t-1}$ and $\zeta_t$, $\xi_t^j\sim K_t(\zeta_t^j,\cdot)$, so
$\E[f(\xi_t^j)\mid\F_{t-1},\zeta_t]=K_tf(\zeta_t^j)$. Given $\F_{t-1}$, $\zeta_t^j\sim\mu_t$, so by the tower property
\[
\E[f(\xi_t^j)\mid\F_{t-1}]=\mu_t(K_tf)=\frac{\sum_iG_t(\xi_{t-1}^i)K_tf(\xi_{t-1}^i)}{\sum_iG_t(\xi_{t-1}^i)}
=\frac{\qh_{t-1}(G_tK_tf)}{\qh_{t-1}(G_t)}.
\]
Average over $j=1,\dots,N$.
\end{proof}

(Only the conditional marginal laws of the $\zeta_t^j$ and $\xi_t^j$ were used; the independence across particles is not
needed for first moments.)

\section{The weighted mean identity}

\begin{T14proposition}\label{T14:prop:unbiased}
For $0\le t\le T-1$ and every $f:\Om\to\mathbb R$,
\[
\E\bigl[A_t\,\qh_t(f)\bigr]=b_t\,q_t(f).
\]
\end{T14proposition}
\begin{proof}
Induction on $t$, for all $f$ at once. For $t=0$: $A_0=b_0=1$ and $\E\qh_0(f)=q_0(f)$ because the $\xi_0^i$ have law $q_0$.

Let $1\le t\le T-1$ and assume the claim at $t-1$ for all functions. $A_t=A_{t-1}\qh_{t-1}(R_t)$ is
$\F_{t-1}$-measurable, so by Lemma~\ref{T14:lem:step},
\[
\E[A_t\qh_t(f)]=\E\Bigl[A_{t-1}\,\qh_{t-1}(R_t)\,\frac{\qh_{t-1}(G_tK_tf)}{\qh_{t-1}(G_t)}\Bigr]
=\E\bigl[A_{t-1}\,\qh_{t-1}(R_tK_tf)\bigr],
\]
since $R_t=G_t/G_t(\So)$ gives $\qh_{t-1}(R_t)/\qh_{t-1}(G_t)=1/G_t(\So)$ and $\qh_{t-1}(G_tK_tf)/G_t(\So)
=\qh_{t-1}(R_tK_tf)$. The induction hypothesis with the function $R_tK_tf$, then Lemma~\ref{T14:lem:change} with
$h=K_tf$, invariance $q_t(K_tf)=q_t(f)$, and Lemma~\ref{T14:lem:b} give
\[
\E[A_t\qh_t(f)]=b_{t-1}\,q_{t-1}(R_tK_tf)=\frac{b_{t-1}}{G_t(\So)}\cdot\frac{\Zb_t}{\Zb_{t-1}}\,q_t(K_tf)=b_t\,q_t(f).
\qedhere
\]
\end{proof}

\section{Proof of Theorem~\ref{T14:thm:main}}

\paragraph{(i-a).} Since $G_t=G_t(\So)R_t$, $\qh_{t-1}(G_t)=G_t(\So)\qh_{t-1}(R_t)$, and
$\prod_{t=1}^TG_t(\So)=e^{d(\So)\sum_t\Delta_t}=e^{d(\So)}$ because $\varphi_T-\varphi_0=1$. Hence
\[
\frac{\Zh_r}{Z}=\frac{\Zb_0}{\Zb_T}\,e^{d(\So)}\,A_T .
\]
Now, since $\varphi_T=1$, $q_T(\So)/q_0(\So)=\bigl(\gamma_0(\So)e^{d(\So)}/\Zb_T\bigr)\big/\bigl(\gamma_0(\So)/\Zb_0\bigr)
=e^{d(\So)}\Zb_0/\Zb_T$, so $\Zh_r/Z=A_T\,q_T(\So)/q_0(\So)=A_T/b_T$.

\paragraph{One telescoping increment.} Fix $1\le t\le T$. Since $R_t(\So)=1$,
\[
\qh_{t-1}(R_t)-1=\qh_{t-1}(R_t-1)=\qh_{t-1}\bigl((R_t-1)\mathbf 1_{\Om\setminus\{\So\}}\bigr),
\]
and, because $R_t>0$,
$|\qh_{t-1}(R_t)-1|\le\qh_{t-1}(R_t\mathbf 1_{\Om\setminus\{\So\}})+\qh_{t-1}(\mathbf 1_{\Om\setminus\{\So\}})$.
Multiplying by $A_{t-1}>0$,
\[
|A_t-A_{t-1}|=A_{t-1}|\qh_{t-1}(R_t)-1|
\le A_{t-1}\qh_{t-1}\bigl(R_t\mathbf 1_{\Om\setminus\{\So\}}\bigr)+A_{t-1}\qh_{t-1}\bigl(\mathbf 1_{\Om\setminus\{\So\}}\bigr).
\]
Take expectations and apply Proposition~\ref{T14:prop:unbiased} at $t-1\in\{0,\dots,T-1\}$:
\[
\E|A_t-A_{t-1}|\le b_{t-1}\,q_{t-1}\bigl(R_t\mathbf 1_{\Om\setminus\{\So\}}\bigr)+b_{t-1}\eta_{t-1}.
\]
By Lemma~\ref{T14:lem:change} with $h=\mathbf 1_{\Om\setminus\{\So\}}$ and then Lemma~\ref{T14:lem:b},
\[
b_{t-1}\,q_{t-1}\bigl(R_t\mathbf 1_{\Om\setminus\{\So\}}\bigr)
=\frac{b_{t-1}}{G_t(\So)}\cdot\frac{\Zb_t}{\Zb_{t-1}}\,q_t(\Om\setminus\{\So\})=b_t\eta_t .
\]
Therefore
\begin{equation}\label{T14:eq:incr}
\E|A_t-A_{t-1}|\le b_{t-1}\eta_{t-1}+b_t\eta_t\qquad(1\le t\le T).
\end{equation}

\paragraph{(i-b), first inequality.} By (i-a),
\[
\frac{\Zh_r}{Z}-1=\frac{A_T-1}{b_T}+\frac{1-b_T}{b_T},
\qquad |A_T-1|=|A_T-A_0|\le\sum_{t=1}^T|A_t-A_{t-1}| .
\]
Taking expectations and using \eqref{T14:eq:incr},
\[
\E\Bigl|\frac{\Zh_r}{Z}-1\Bigr|\le\frac{\sum_{t=1}^T(b_{t-1}\eta_{t-1}+b_t\eta_t)+|1-b_T|}{b_T}.
\]

\paragraph{(i-b), second inequality.} For $0\le s\le T$,
\[
\frac{b_s}{b_T}=\frac{q_T(\So)}{q_s(\So)}\le\frac1{q_s(\So)}=\frac1{1-\eta_s}\le\frac1{1-\varepsilon_*}.
\]
Hence
\[
\frac{\sum_{t=1}^T(b_{t-1}\eta_{t-1}+b_t\eta_t)}{b_T}\le\frac{\sum_{t=1}^T(\eta_{t-1}+\eta_t)}{1-\varepsilon_*}
=\frac{2\sum_{t=0}^T\eta_t-\eta_0-\eta_T}{1-\varepsilon_*}.
\]
For the endpoint term, $b_T=q_0(\So)/q_T(\So)$ gives
\[
\frac{|1-b_T|}{b_T}=\frac{|q_T(\So)-q_0(\So)|}{q_0(\So)}=\frac{|\eta_0-\eta_T|}{1-\eta_0}
\le\frac{\eta_0+\eta_T}{1-\varepsilon_*},
\]
using $|\eta_0-\eta_T|\le\max(\eta_0,\eta_T)\le\eta_0+\eta_T$ (both are in $[0,1)$). Adding the two displays gives the
bound $2(1-\varepsilon_*)^{-1}\sum_{t=0}^T\eta_t$.

\paragraph{(i-c).} Both right-hand sides of (i-b) depend only on $(q_t)_{0\le t\le T}$ and $\So$, which are fixed by the
data and $\mathcal C$; they are therefore the same for every island. Call the middle expression $M$. By the triangle inequality,
$|\Zh/Z-1|=\bigl|\tfrac14\sum_r(\Zh_r/Z-1)\bigr|\le\tfrac14\sum_r|\Zh_r/Z-1|$, so
$\E|\Zh/Z-1|\le\tfrac14\sum_r\E|\Zh_r/Z-1|\le M\le 2(1-\varepsilon_*)^{-1}\sum_t\eta_t$. No independence between islands
is used.

\paragraph{(i-d).} Take $\Om=\{\So,x\}$, $T=1$, $N=1$, $m(\emptyset)=1$, $K_1$ the identity kernel (it leaves every law
invariant), and fix $a=0.01$, $c=0.1$. Put $\gamma_0(\So)=1-a$, $\gamma_0(x)=a$, $d(\So)=0$,
$d(x)=\log\frac{c(1-a)}{a(1-c)}$. Then $q_0=(1-a,a)$, and $\gamma_1=(1-a,\;c(1-a)/(1-c))=\frac{1-a}{1-c}(1-c,c)$, so
$q_1=(1-c,c)$. With $T=N=1$ (so stage $1$ is temperature $1$), $\Zh_r/Z=\Zb_0G_1(\xi_0^1)/\Zb_1=q_1(\xi_0^1)/q_0(\xi_0^1)$ (Lemma~\ref{T14:lem:change} with
$h=\mathbf 1_{\{y\}}$ gives $\Zb_0 G_1(y)/\Zb_1=q_1(y)/q_0(y)$), and $\xi_0^1\sim q_0$, so
\[
\E\Bigl|\frac{\Zh_r}{Z}-1\Bigr|=\sum_{y\in\Om}|q_1(y)-q_0(y)|=2|c-a|=0.18 .
\]
Here $\eta_0=a$, $\eta_1=c$, $b_0=1$, $b_1=(1-a)/(1-c)=1.1$, and
\[
\frac{b_0\eta_0+b_1\eta_1}{b_1}=\frac{0.01+0.11}{1.1}=0.1090\ldots<0.18 .
\]
So the bound without the endpoint term fails, while the full middle expression of (i-b) equals
$0.1090\ldots+|1-1.1|/1.1=0.2$ and the right-hand side equals $(2/0.9)(0.11)=0.2444\ldots$, both $\ge0.18$ as the theorem
requires. \qed

\begin{T14remark}
Proposition~\ref{T14:prop:unbiased} is equivalent to the frozen-run identity $\E[\Zh_{r,t}\qh_t(f)]=Z_tq_t(f)$ (the form in
which the project's conjecture T12(i) states it; that file is not needed here) rewritten through
$\Zh_{r,t}=m(\emptyset)\Zb_0e^{\varphi_td(\So)}A_t$ and $Z_t=m(\emptyset)\Zb_t$; it is proved above directly, so the
argument uses nothing from T12. Nothing about the internal structure of the move round is used beyond $q_tK_t=q_t$, and
nothing about EXP/EXP$'$ or the statistical model enters.
\end{T14remark}
\mdproofend
\mdsec{(d) Review record}
\mdsubsec{FINAL\_REPORT.md: statement changes}
\begin{mdblock}
\mdhead{2}{3. Statement changes}

There were no mathematical changes. The statement analysis concluded it was correct as stated (proof.md step 1(a)).

\begin{mdbul}
\item \textbf{Notation (after Astra \#1 finding 1).} "q₁ = π, Z = m(∅)Z̄₁" clashed with stage indexing. It was rewritten so that a numerical subscript is always a stage index (Z = m(∅)Z̄\_T). This is logged in STATEMENT.md §4.
\item \textbf{Reading of an informal sentence.} "The endpoint term |1 − b\_T| is essential" is read as: in the abstract setting, dropping the term breaks the bound for some instance. This reading is stated as (i-d), and it is not claimed that such an instance comes from a real SMC-BVS data set. All three reviewers judged this reading fair.
\end{mdbul}
\end{mdblock}

\mdsubsec{FINAL\_REPORT.md: review history}
\begin{mdblock}
\mdhead{2}{4. Review history}

{\small
\begin{longtable}{|>{\raggedright\arraybackslash}p{\dimexpr(\linewidth-10\tabcolsep-6\arrayrulewidth)*61/1000\relax}|>{\raggedright\arraybackslash}p{\dimexpr(\linewidth-10\tabcolsep-6\arrayrulewidth)*374/1000\relax}|>{\raggedright\arraybackslash}p{\dimexpr(\linewidth-10\tabcolsep-6\arrayrulewidth)*98/1000\relax}|>{\raggedright\arraybackslash}p{\dimexpr(\linewidth-10\tabcolsep-6\arrayrulewidth)*98/1000\relax}|>{\raggedright\arraybackslash}p{\dimexpr(\linewidth-10\tabcolsep-6\arrayrulewidth)*369/1000\relax}|}
\hline
\textbf{Stage} & \textbf{Reviewer} & \textbf{Input} & \textbf{Verdict} & \textbf{Findings} \\ \hline\endhead
A & Astra \#1 (gpt-6-astra, high), session 01a16418-6bd3-78f1-ad35-0348b23c7968 & PROOF\_v1 & verified & 1 exposition (stage/temperature subscript clash): accepted, fixed in v2 \\ \hline
A & Astra \#1, same session (resume) & PROOF\_v2 + log & verified & agrees with the fix; no new findings \\ \hline
B & Astra \#2 (gpt-6-astra, high), new session 01a1641c-418a-7463-a810-25b59b68de23 & PROOF\_v2 & verified & none \\ \hline
B & Opus 5.5 xhigh (Claude Personal, T3 delegate\_task) & PROOF\_v2 & verified & 3 exposition: conditioning on F\_\{t−1\}; (i-c) wording; remark citing T12. All accepted, unified into v3 \\ \hline
B & Astra \#2 (resume) & PROOF\_v3 + log & verified & agrees with all answers; none \\ \hline
B & Opus 5.5 xhigh (new round task) & PROOF\_v3 + log & verified & agrees with all three answers; none \\ \hline
\end{longtable}}

\begin{mdbul}
\item Fatal or gap findings: none, in any round. CLAIMS\_OF\_FATAL\_FLAWS.md records this and logs the exposition findings with the answers.
\item Astra judge: not used.
\item All replies are saved in {\mdtt review/\allowbreak{}}. Stage B Opus scratch scripts are in {\mdtt review/\allowbreak{}stageB\_\allowbreak{}opus\_\allowbreak{}scratch/\allowbreak{}}.
\item Temp packet folders were deleted at the end; a copy of the Stage A packet is kept in {\mdtt review/\allowbreak{}packet\_\allowbreak{}stageA\_\allowbreak{}copy/\allowbreak{}}.
\end{mdbul}
\end{mdblock}

\mdsubsec{CLAIMS\_OF\_FATAL\_FLAWS.md (every entry, verbatim)}
\begin{mdblock}
\mdhead{1}{CLAIMS\_OF\_FATAL\_FLAWS — T14(i)}

Reviewer: Astra (gpt-6-astra, reasoning high), session 01a16418-6bd3-78f1-ad35-0348b23c7968, round 1 on PROOF\_v1 (2026-10-22). Verdict: verified.

\mdhead{2}{Fatal or gap findings}

None. Astra raised no fatal and no gap finding.

\mdhead{2}{Exposition findings (recorded for completeness)}

\textbf{Finding 1 (exposition), verbatim claim:} "The notation \(q_1,\bar Z_1\) is used both for temperature \(1\) and for stage \(1\). These meanings generally differ when \(T>1\). In particular, the assertion “\(q_T=q_1\)” requires the temperature interpretation of the right-hand side."

\textbf{Lead's answer: accepted and fixed.} PROOF\_v2.tex defines Z = m(∅)Σγ₀e\textasciicircum{}\{d\} (the bridge normaliser at φ = 1), states that a numerical subscript on q or Z̄ is always a stage index, and writes Z = m(∅)Z̄\_T and q\_T(S°)/q₀(S°) = e\textasciicircum{}\{d(S°)\}Z̄₀/Z̄\_T in the proof of (i-a). The (i-d) example keeps q₁, Z̄₁ because there T = 1, as now noted. STATEMENT.md §1 was reworded the same way (notation only; change log §4).

\mdhr

\mdhead{1}{Stage B (Protocol Amendment 1), round 1 on PROOF\_v2, 2026-10-22}

Reviewers: Astra \#2 (gpt-6-astra, high, new session 01a1641c-418a-7463-a810-25b59b68de23): verdict verified, no findings. Opus 5.5 (xhigh, Claude Personal, clientRequestId T14i-stageB-opus-personal-r1): verdict verified, three exposition findings.

\mdhead{2}{Fatal or gap findings}

None. Neither Stage B reviewer raised a fatal or gap finding.

\mdhead{2}{Exposition findings (Opus), all accepted and fixed in the unified PROOF\_v3.tex}

\textbf{Opus finding 1, verbatim claim:} "the statement gives the resampling and move laws "conditionally on the past". The proof conditions on the smaller σ-field F\_\{t−1\} = σ(ξ₀, ζ₁, …, ξ\_\{t−1\}) without saying why that is allowed." \textbf{Lead's answer: accepted.} §1 item (b) now says that laws given on a larger past hold in particular given F\_\{t−1\} (resp. (F\_\{t−1\}, ζ\_t)) by the tower property, since μ\_t is F\_\{t−1\}-measurable and the law of ξ\_t\textasciicircum{}j depends only on ζ\_t\textasciicircum{}j.

\textbf{Opus finding 2, verbatim claim:} "the proof says the islands "share the grid, the kernels, N and S°". Only the grid, γ₀, d and S° are actually needed." \textbf{Lead's answer: accepted.} (i-c) now says both right-hand sides depend only on (q\_t)\_\{0≤t≤T\} and S°, which the data and 𝓒 fix.

\textbf{Opus finding 3, verbatim claim:} "it cites T12(i) (E[Ẑ\_\{r,t\} q̂\_t(f)] = Z\_t q\_t(f)), which is not in the packet, so I cannot check it against its source." \textbf{Lead's answer: accepted.} The remark now says Proposition 5 is equivalent to the identity in the form in which the project's conjecture T12(i) states it, that the file is not needed, and that the argument uses nothing from T12. (The remark is not used by the proof.)
\end{mdblock}

\chapter{T15 (L-OCC): occupancy from an evidence lower tail}
\chaptersource{Final proof: \texttt{proofs/\allowbreak{}T15\_\allowbreak{}LOCC\_\allowbreak{}occupancy/\allowbreak{}proof/\allowbreak{}PROOF\_v4.tex}. Sources of parts (a), (b), (d): \texttt{theorems/\allowbreak{}T15\_\allowbreak{}LOCC\_\allowbreak{}occupancy.md}, \texttt{STATEMENT.md}, \texttt{FINAL\_REPORT.md}, \texttt{CLAIMS\_OF\_FATAL\_FLAWS.md}.}
\mdsec{(a) Original conjecture text (theorem file, verbatim)}
\begin{mdblock}
\mdhead{1}{T15 (L-OCC): Occupancy from an evidence lower tail}

\textbf{Status:} conjectured, not proven. Methodology SMC-BVS v3.0 (2026-10-19), SHA-256 93043758058a6d422cca6ca4b6be8f3a9e8149b6f6de2e1bab1ce31c44b8019c. Read with {\mdtt 00\_\allowbreak{}COMMON.\allowbreak{}md}.

\textbf{Category / rank:} algorithmic convergence (source of BASIN's δ); \textbf{core} (lemma).

\mdhead{2}{Definitions used}

Island evidence Ẑ\_r; f\_r(𝔄ᶜ) = the fraction of island r's particles outside a set 𝔄 (fixed given 𝓒) at stage T or at a fixed terminal recording prefix (as in T12(i)); W := Ẑ\_r/Z.

\mdhead{2}{Assumptions}

EX-A, EX-S, CAL; prefixes as specified in T12.

\mdhead{2}{Precise statement}

Conditional on a successful 𝓒, for every c, u ∈ (0, 1):\\
\hspace*{1.5em}E f\_r(𝔄ᶜ) ≤ min\{1, π(𝔄ᶜ)/c + P(Ẑ\_r/Z < c)\},  P(f\_r(𝔄ᶜ) > u) ≤ π(𝔄ᶜ)/(cu) + P(Ẑ\_r/Z < c). If E|Ẑ\_r/Z − 1| ≤ B then the lower-tail term is ≤ B/(1−c). These are calibration-conditional bounds; conditional on an already realised population its actual outside fraction must be used. The lower tail must be islandwise; a pooled-evidence lower tail does not substitute.

\mdhead{2}{Dependencies}

Builds on: T12 (and optionally T14 or T25 for the lower tail). Used by: T23, T24, T26.

\mdhead{2}{Numerical evidence (provisional)}

provisional evidence (Sol check J6): exact finite-particle check of E[Wf] = π(𝔄ᶜ) (weighted fraction 0.01980 = target outside 0.01980) and of the linear bound. Files: judge/sol2/J6\_FIRST\_MOMENTS.md.

\mdhead{2}{Proof-route hint (optional)}

T12(i) with f = 1\_\{𝔄ᶜ\} gives E[Wf\_r] = π(𝔄ᶜ); split at W = c using 0 ≤ f\_r ≤ 1; Markov for the probability bound. No substantive hardest step; the content is the lower-tail input. \emph{When proof.md is run on this file, its lead considers the statement both with and without these hints; the hints never reach the reviewers.}
\end{mdblock}

\mdsec{(b) STATEMENT.md (verbatim)}
\begin{mdblock}
\mdhead{1}{STATEMENT — T15 (L-OCC): occupancy from an evidence lower tail}

Source: {\mdtt theorems/\allowbreak{}T15\_\allowbreak{}LOCC\_\allowbreak{}occupancy.\allowbreak{}md} (SMC-BVS v3.0, METHODOLOGY\_v3.md SHA-256 93043758058a6d422cca6ca4b6be8f3a9e8149b6f6de2e1bab1ce31c44b8019c), read with {\mdtt theorems/\allowbreak{}00\_\allowbreak{}COMMON.\allowbreak{}md}. Parts to prove: all.

\mdhead{2}{1. Setting (compact; the method file is authoritative)}

\begin{mdbul}
\item \textbf{Target.} Ω = \{S ⊆ [p] : |S| ≤ K\} is a finite set of models. Each S ∈ Ω has a strictly positive unnormalised posterior weight γ₁(S) = ϖ(S)m(S) (prior × marginal likelihood). Z := Σ\_\{S∈Ω\} γ₁(S) ∈ (0, ∞) is the evidence and π := γ₁/Z the posterior on Ω. For A ⊆ Ω, π(A) = Σ\_\{S∈A\} π(S).
\item \textbf{Calibration.} 𝓒 is the frozen output of the calibration phase (temperature grid 0 = φ₀ < φ₁ < … < φ\_T = 1, budgets, N = N\_e, and all other settings). It is measurable with respect to calibration random streams 0–9 and the data. "Successful 𝓒" means the calibration ended with a constructed grid and the final run is launched. P\_𝓒 and E\_𝓒 denote probability and expectation over the final-run randomness (streams 10–13) given 𝓒 and the data.
\item \textbf{Final run, island r ∈ \{1, 2, 3, 4\}.} N particles drawn i.i.d. from the start law π₀; for t = 1, …, T: reweight by G\_t(S) = exp\{(φ\_t − φ\_\{t−1\})d(S)\}, multinomial resampling, one move round at temperature φ\_t; then L burn-in rounds and 4 measured rounds at φ = 1 (no reweighting in these terminal rounds). A move round is a list of elementary steps applied, for each particle, in a uniformly random order drawn independently of the particle's state.
\item \textbf{Island evidence.} Ẑ\_r := m(∅)Z̄₀ ∏\_\{t=1\}\textasciicircum{}\{T\} (1/N) Σ\_\{i=1\}\textasciicircum{}\{N\} G\_t(S\_i\textasciicircum{}\{(t−1)\}), the product over all stages (the terminal rounds contribute no factor). Here m(∅) > 0 is the marginal likelihood of the empty model, π₀ = γ₀/Z̄₀ is the composite start law with normaliser Z̄₀ = Σ\_Ω γ₀ > 0 (γ₀ > 0), and d(S) is the real-valued discrepancy (all defined in {\mdtt theorems/\allowbreak{}00\_\allowbreak{}COMMON.\allowbreak{}md} §§1–2). Since every factor is a finite positive number (each G\_t(S) = e\textasciicircum{}\{(φ\_t−φ\_\{t−1\})d(S)\} ∈ (0, ∞)), Ẑ\_r is a finite, strictly positive random variable under EX-A. The proof uses only W\_r ≥ 0. Write \textbf{W\_r := Ẑ\_r / Z}. Pooled evidence: Ẑ := ¼ Σ\_r Ẑ\_r, and W̄ := Ẑ/Z = ¼ Σ\_r W\_r.
\item \textbf{Recording positions.} 𝒯 denotes the set of positions τ in island r's terminal part at which T12(i) is stated: (a) τ = "end of stage T" (after the stage-T resampling and move round); (b) any fixed elementary prefix within the L + 4 terminal rounds; (c) "immediately before a specified Gibbs label" in a terminal round, under the state-independent random order. q̂\_\{r,τ\} is the empirical measure (1/N)Σ\_i δ\_\{S\_i(τ)\} of island r's particle states at position τ.
\item \textbf{Outside set and outside fraction.} 𝔄 ⊆ Ω is a set that is fixed given 𝓒 (𝓒- and data-measurable, independent of streams 10–13). f\_r := q̂\_\{r,τ\}(1\_\{𝔄ᶜ\}) ∈ [0, 1], the fraction of island r's particles outside 𝔄 at a fixed τ ∈ 𝒯.
\end{mdbul}

\mdhead{2}{2. Hypotheses}

\begin{mdbul}
\item \textbf{EX-A} (exact arithmetic), \textbf{EX-S} (independent ideal uniform streams, with the method's stream separation), \textbf{CAL} (𝓒 uses streams 0–9 only; streams 10–13 are fresh; escalation is decided before any final randomness). These are standing assumptions of the setting.
\item \textbf{(H-T12) Frozen-run identity} (the conclusion of the sibling conjecture T12(i), used here only as an explicit hypothesis): conditional on a successful 𝓒, for island r, each τ ∈ 𝒯, and every bounded function f on Ω that is fixed given 𝓒,

E\_𝓒[ Ẑ\_r · q̂\_\{r,τ\}(f) ] = Z · π(f).

(At these positions the target is q\_T = q₁ = π and Z\_T = Z.) Parts (i)–(ii) use only the instance f = 1\_\{𝔄ᶜ\}, and only the inequality E\_𝓒[W\_r f\_r] ≤ π(𝔄ᶜ). The first claim of (iii), P\_𝓒(W\_r < c) ≤ B/(1 − c), uses only its own L¹ hypothesis; its "consequently" clause rests on (i)–(ii) and hence on the same instance of (H-T12). Part (iv) is a statement about abstract laws and uses no hypothesis; it shows that even the identity of (H-T12) for every test function does not yield the pooled form.
\end{mdbul}

\mdhead{2}{3. Conclusions}

Fix a successful 𝓒, an island r, a position τ ∈ 𝒯, and a set 𝔄 fixed given 𝓒. Let f\_r = q̂\_\{r,τ\}(1\_\{𝔄ᶜ\}) and W\_r = Ẑ\_r/Z.

\textbf{(i) Mean bound.} For every c ∈ (0, 1):\\
\hspace*{1.5em}E\_𝓒 f\_r ≤ min\{ 1, π(𝔄ᶜ)/c + P\_𝓒(W\_r < c) \}.

\textbf{(ii) Tail bound.} For every c, u ∈ (0, 1):\\
\hspace*{1.5em}P\_𝓒(f\_r > u) ≤ π(𝔄ᶜ)/(cu) + P\_𝓒(W\_r < c).

\textbf{(iii) Lower tail from an L¹ bound.} If E\_𝓒|W\_r − 1| ≤ B for a constant B ≥ 0, then for every c ∈ (0, 1), P\_𝓒(W\_r < c) ≤ B/(1 − c). Consequently (i) and (ii) hold with P\_𝓒(W\_r < c) replaced by B/(1 − c).

\textbf{(iv) The lower tail must be islandwise.} The bounds (i)–(ii) with the pooled lower tail P\_𝓒(W̄ < c) in place of the islandwise P\_𝓒(W\_r < c) do not follow from the structure used, even with the identity of (H-T12) for every test function. (This is a non-implication about direct substitution; it does not assert that every conceivable occupancy bound must use islandwise evidence.) Precisely, fix c, u ∈ (0, 1) and a ∈ (0, cu/4].

(iv-a) There are four i.i.d. random pairs (W\_r, f\_r), r = 1, …, 4, with W\_r > 0 (strictly), f\_r ∈ \{0, 1\}, E W\_r = 1 and E[W\_r f\_r] = a, such that, with W̄ = ¼ Σ\_r W\_r,\\
\hspace*{1.5em}E f\_1 > a/c + P(W̄ < c)  and  P(f\_1 > u) > a/(cu) + P(W̄ < c); in particular, since min\{1, x\} ≤ x, also E f\_1 > min\{1, a/c + P(W̄ < c)\}.

(iv-b) Moreover, let Ω′ be any finite set, π′ any probability on Ω′, 𝔄′ ⊆ Ω′ with π′(𝔄′ᶜ) = a, and N ≥ 1. (For example, the method's Ω and π whenever some 𝔄 ⊆ Ω has π(𝔄ᶜ) = a; such an 𝔄 exists for a := the smallest posterior mass of a single model whenever |Ω| ≥ 4/(cu), since that mass lies in (0, 1/|Ω|] ⊆ (0, cu/4].) Then the pairs of (iv-a) can be realised jointly with random populations (S\_\{r,1\}, …, S\_\{r,N\}) ∈ Ω′\textasciicircum{}N, r = 1, …, 4, such that the four pairs (W\_r, (S\_\{r,i\})\_\{i≤N\}) are i.i.d., f\_r = q̂\_r(1\_\{𝔄′ᶜ\}) for the empirical measure q̂\_r = (1/N) Σ\_i δ\_\{S\_\{r,i\}\}, and

\hspace*{1.5em}E[ W\_r · q̂\_r(g) ] = π′(g)  for every function g : Ω′ → ℝ,

which is the identity of (H-T12) with Ẑ\_r := Z′ W\_r for any constant Z′ > 0, for every test function g and not only for g ∈ \{1, 1\_\{𝔄′ᶜ\}\}. Each island has one population; using it at every recording position (q̂\_\{r,τ\} := q̂\_r for all τ ∈ 𝒯) gives the identity at every position as well. The four islands of this example are i.i.d.; that property only strengthens the example and is not asserted here as a property of the method.

\mdhead{2}{4. Scope remarks (from the theorem file; interpretations, not separate claims)}

\begin{mdbul}
\item (i)–(iii) are statements about the law of the final run given 𝓒 (calibration-conditional). They give no bound on the outside fraction of an already realised population; for a realised population its actual f\_r must be used.
\item Part (iv) formalises the theorem file's sentence "a pooled-evidence lower tail does not substitute" as a non-implication from the properties listed (W\_r > 0, f\_r a fraction in [0,1], i.i.d. islands, and the identity of (H-T12) for every test function and every recording position). It does not claim that the SMC final run, for a particular data set and calibration, realises the counterexample law.
\end{mdbul}

\mdhead{2}{5. Changes relative to the theorem file}

\begin{mdbul}
\item None to the inequalities, their constants or their quantifiers.
\item The theorem file's dependency on T12 ("Builds on: T12"; "prefixes as specified in T12") is stated as the explicit hypothesis (H-T12), as the run rules require for sibling conjectures. The positions 𝒯 are exactly those named in T12(i).
\item The sentence on pooled lower tails is made precise as part (iv) (a non-implication), and the sentence on realised populations is kept as a scope remark (§4).
\item Assumption EXP (and its replacement EXP′) is not used.
\item v2 (after Astra round 1, exposition finding 1): in (iv) the example pairs are required to have W\_r > 0 strictly (v1 allowed W\_r ≥ 0), matching the strict positivity of actual island evidence, and a sentence fixes the scope of the heading. The inequalities of (i)–(iii) are unchanged.
\item v3 (after Stage B, Opus exposition findings 1–5; no fatal or gap findings in any review): (iv) now holds for every a ∈ (0, cu/4] (v2: a = cu/4 only) and adds (iv-b), realising the example with particle populations that satisfy the full identity (H-T12) for every test function; the min\{1,·\} form is refuted explicitly; the sentence presenting island i.i.d. as a method property is removed; Ẑ\_r's symbols and positivity are spelled out; the hypothesis-use note in §2 is made exact. (i)–(iii) unchanged.
\item v4 (after Stage B round 2; both reviewers verified v3 with exposition findings only): wording only — the hypothesis-use note for (iii) separates its two claims; "full identity" replaced by "identity for every test function", with the single-population/every-position remark; the method's-Ω example states when a set of mass a exists. No mathematical content changed.
\item Post-review alignment (after Stage B round 3, Opus finding 1, marked optional): the last "full identity" in the §4 scope remark replaced by the v4 wording. Wording only, in a scope remark; not re-sent for review.
\end{mdbul}
\end{mdblock}

\mdsec{(c) Final proof, verbatim: body of PROOF\_v4.tex}
\noindent{\small\itshape The proof below is the body of \texttt{PROOF\_v4.tex} from \texttt{\textbackslash begin\{document\}} to \texttt{\textbackslash end\{document\}} without the title block. Its preamble definitions are local to this chapter; labels carry the prefix \texttt{T15:}.}\par\medskip
\mdproofbegin
% ---- preamble of PROOF_v4.tex (theorem environments carry the prefix T15) ----
\newtheorem{T15theorem}{Theorem}
\newtheorem{T15lemma}{Lemma}
\theoremstyle{remark}
\newtheorem{T15remark}{Remark}
\newcommand{\Pc}{\mathbb{P}_{\mathcal C}}
\newcommand{\Ec}{\mathbb{E}_{\mathcal C}}
\newcommand{\A}{\mathfrak{A}}
% ---- body of PROOF_v4.tex ----
\section{Setting and hypotheses}

We use the notation of STATEMENT.md. $\Omega$ is the finite model space, $\pi=\gamma_1/Z$ the posterior with
evidence $Z\in(0,\infty)$, and $\mathcal C$ a successful frozen calibration; $\Pc$ and $\Ec$ are the probability and
expectation over the final-run randomness given $\mathcal C$ and the data. For an island $r$, $\hat Z_r$ is its
evidence estimate, which under EX-A is a finite, strictly positive random variable, and
\[
  W_r:=\hat Z_r/Z \in (0,\infty).
\]
Fix a recording position $\tau\in\mathcal T$ and a set $\A\subseteq\Omega$ fixed given $\mathcal C$, and let
\[
  f_r:=\hat q_{r,\tau}(1_{\A^c})=\frac1N\sum_{i=1}^N 1\{S_i(\tau)\notin\A\}\in[0,1].
\]
Since $\A$ is fixed given $\mathcal C$, the function $f=1_{\A^c}$ is bounded and fixed given $\mathcal C$, so
hypothesis (H-T12) applies to it and gives
\begin{equation}\label{T15:eq:id}
  \Ec[W_r f_r]=\frac{\Ec[\hat Z_r\,\hat q_{r,\tau}(1_{\A^c})]}{Z}=\frac{Z\,\pi(1_{\A^c})}{Z}=\pi(\A^c).
\end{equation}
For parts (i)--(iii) only the inequality $\Ec[W_rf_r]\le\pi(\A^c)$ is used. The random variable $W_rf_r$ is
nonnegative, and by \eqref{T15:eq:id} it is integrable.

Throughout, $r$, $\tau$ and $\A$ are fixed, and we drop the subscript $r$: $W=W_r$, $f=f_r$.

\section{Two elementary facts}

\begin{T15lemma}[Splitting at $c$]\label{T15:lem:split}
Let $W\ge0$ and $f\in[0,1]$ be random variables on a common probability space, with $\mathbb E[Wf]\le a<\infty$.
Then for every $c>0$,
\[
  \mathbb E\big[f\,1\{W\ge c\}\big]\le a/c .
\]
\end{T15lemma}
\begin{proof}
On the event $\{W\ge c\}$ we have $W/c\ge1$ and $f\ge0$, hence $f\le fW/c$. Off this event,
$f\,1\{W\ge c\}=0\le fW/c$ because $f,W\ge0$. So pointwise
$f\,1\{W\ge c\}\le fW/c$, and taking expectations gives $\mathbb E[f1\{W\ge c\}]\le \mathbb E[fW]/c\le a/c$.
\end{proof}

\begin{T15lemma}[Markov]\label{T15:lem:markov}
If $V\ge0$ is a random variable and $s>0$, then $\mathbb P(V\ge s)\le \mathbb E V/s$.
\end{T15lemma}
\begin{proof}
$s\,1\{V\ge s\}\le V$ pointwise; take expectations.
\end{proof}

\section{Proof of (i): the mean bound}

Fix $c\in(0,1)$. Since $0\le f\le1$, we have $\Ec f\le1$. Write $f=f\,1\{W\ge c\}+f\,1\{W<c\}$. By
Lemma~\ref{T15:lem:split} with $a=\pi(\A^c)$ (justified by \eqref{T15:eq:id}), $\Ec[f1\{W\ge c\}]\le\pi(\A^c)/c$. Since
$f\le1$, $\Ec[f1\{W<c\}]\le\Pc(W<c)$. Adding,
\[
  \Ec f\le \pi(\A^c)/c+\Pc(W<c).
\]
Together with $\Ec f\le1$ this is $\Ec f\le\min\{1,\ \pi(\A^c)/c+\Pc(W<c)\}$. \qed

\section{Proof of (ii): the tail bound}

Fix $c,u\in(0,1)$. Decompose
\[
  \{f>u\}\subseteq\big(\{f>u\}\cap\{W\ge c\}\big)\cup\{W<c\}.
\]
On $\{f>u\}\cap\{W\ge c\}$ we have $fW\ge fc>uc$, using $W\ge c>0$ and $f>u>0$. Hence
$\{f>u\}\cap\{W\ge c\}\subseteq\{fW\ge cu\}$. By Lemma~\ref{T15:lem:markov} applied to $V=fW\ge0$ and $s=cu>0$, and by
\eqref{T15:eq:id},
\[
  \Pc(fW\ge cu)\le \Ec[fW]/(cu)\le\pi(\A^c)/(cu).
\]
By the union bound,
\[
  \Pc(f>u)\le\pi(\A^c)/(cu)+\Pc(W<c). \qquad\qed
\]

\section{Proof of (iii): the lower tail from an $L^1$ bound}

Assume $\Ec|W-1|\le B$, and fix $c\in(0,1)$. On $\{W<c\}$, $|W-1|\ge 1-W>1-c>0$. Hence
$\{W<c\}\subseteq\{|W-1|\ge 1-c\}$, and Lemma~\ref{T15:lem:markov} with $V=|W-1|$ and $s=1-c>0$ gives
\[
  \Pc(W<c)\le\Pc(|W-1|\ge1-c)\le\frac{\Ec|W-1|}{1-c}\le\frac{B}{1-c}.
\]
Substituting this bound into (i) and (ii) gives, for every $c,u\in(0,1)$,
\[
  \Ec f\le\min\Big\{1,\ \frac{\pi(\A^c)}{c}+\frac{B}{1-c}\Big\},\qquad
  \Pc(f>u)\le\frac{\pi(\A^c)}{cu}+\frac{B}{1-c}. \qquad\qed
\]

\section{Proof of (iv): the lower tail must be islandwise}

Fix $c,u\in(0,1)$ and $a\in(0,cu/4]$. Note $a<1/4$.

\subsection{Proof of (iv-a)}
Set
\[
  \varepsilon:=\frac12\min\Big\{c,\ \frac{4(1-c)}3,\ \frac{4a}3\Big\}>0,\qquad
  h:=4-3\varepsilon,\qquad
  b:=\frac{a-3\varepsilon/4}{h}.
\]
\emph{Elementary bounds on the constants.} Since $\varepsilon\le 2a/3\le cu/6<1/6$, we have $h\in(7/2,4)$. Since
$3\varepsilon/4\le a/2<a$, we have $b>0$; and $b<a/h<a/3<1/12$. Moreover $\varepsilon\le c/2<c$, and
$\varepsilon\le 2(1-c)/3$ gives $3\varepsilon/4\le(1-c)/2$, hence
\begin{equation}\label{T15:eq:h4}
  h/4=1-3\varepsilon/4\ \ge\ (1+c)/2\ >\ c .
\end{equation}

\emph{The law.} Let $(W_r,f_r)$, $r=1,\dots,4$, be i.i.d.\ copies of a pair $(W,f)$ with the three-point law
\[
  \mathbb P\big((W,f)=(\varepsilon,1)\big)=\tfrac34,\qquad
  \mathbb P\big((W,f)=(h,1)\big)=b,\qquad
  \mathbb P\big((W,f)=(h,0)\big)=\tfrac14-b .
\]
The probabilities are positive (as $0<b<1/12$) and sum to $1$. Clearly $W\in\{\varepsilon,h\}\subset(0,\infty)$, so
$W>0$ strictly, and $f\in\{0,1\}$. Also $\mathbb P(W=h)=1/4$ and $\mathbb P(f=1)=3/4+b$.

\emph{The required moment properties.}
$\mathbb E W=\tfrac34\varepsilon+\tfrac14h=\tfrac34\varepsilon+1-\tfrac34\varepsilon=1$, and
$\mathbb E[Wf]=\tfrac34\varepsilon+hb=\tfrac34\varepsilon+a-\tfrac34\varepsilon=a$.

\emph{The pooled lower tail.} Let $K:=\#\{r:W_r=h\}\sim\mathrm{Bin}(4,1/4)$. Then
$\bar W=\tfrac14\big(Kh+(4-K)\varepsilon\big)$. If $K=0$, then $\bar W=\varepsilon<c$. If $K\ge1$, then
$\bar W\ge h/4>c$ by \eqref{T15:eq:h4}. Hence $\{\bar W<c\}=\{K=0\}$ and $\mathbb P(\bar W<c)=(3/4)^4=81/256$.

\emph{The two left-hand sides.} $\mathbb E f_1=\mathbb P(f_1=1)=\tfrac34+b>\tfrac34$. Since $f_1\in\{0,1\}$ and
$u<1$, $\{f_1>u\}=\{f_1=1\}$, so $\mathbb P(f_1>u)=\tfrac34+b>\tfrac34=\tfrac{192}{256}$.

\emph{The two right-hand sides.} Since $a\le cu/4$, $a/c\le u/4<1/4$ and $a/(cu)\le1/4$. Hence
\[
  \frac ac+\mathbb P(\bar W<c)<\frac{64}{256}+\frac{81}{256}=\frac{145}{256},\qquad
  \frac a{cu}+\mathbb P(\bar W<c)\le\frac{64}{256}+\frac{81}{256}=\frac{145}{256}.
\]
Since $145/256<192/256$, both
\[
  \mathbb E f_1> a/c+\mathbb P(\bar W<c)\qquad\text{and}\qquad \mathbb P(f_1>u)>a/(cu)+\mathbb P(\bar W<c)
\]
hold. As $\min\{1,x\}\le x$ for every real $x$, also $\mathbb E f_1>\min\{1,\ a/c+\mathbb P(\bar W<c)\}$. The pairs are
independent and identically distributed and satisfy $W_r>0$, $f_r\in\{0,1\}$, $\mathbb E W_r=1$ and
$\mathbb E[W_rf_r]=a$, which proves (iv-a). \qed

\subsection{Proof of (iv-b)}
Let $\Omega'$ be a finite set, $\pi'$ a probability on $\Omega'$, $\A'\subseteq\Omega'$ with $\pi'(\A'^c)=a$, and
$N\ge1$. Since $a\in(0,1)$, both $\pi'(\A'^c)=a>0$ and $\pi'(\A')=1-a>0$, so the conditional laws
\[
  \pi'_{\rm out}:=\pi'(\,\cdot\mid\A'^c),\qquad \pi'_{\rm in}:=\pi'(\,\cdot\mid\A')
\]
are well defined; $\pi'_{\rm out}$ is carried by $\A'^c$ and $\pi'_{\rm in}$ by $\A'$.

\emph{Construction.} Take the i.i.d.\ pairs $(W_r,f_r)$ of (iv-a). Independently across $r$, and given $(W_r,f_r)$,
draw $S_{r,1},\dots,S_{r,N}$ i.i.d.\ from $\pi'_{\rm out}$ if $f_r=1$ and from $\pi'_{\rm in}$ if $f_r=0$. (Formally:
let $(U_{r,i})$ and $(V_{r,i})$, $r\le4$, $i\le N$, be independent families of i.i.d.\ variables with laws
$\pi'_{\rm out}$ and $\pi'_{\rm in}$, independent of the pairs, and set $S_{r,i}:=U_{r,i}$ if $f_r=1$ and
$S_{r,i}:=V_{r,i}$ if $f_r=0$.) The four quadruples $\big(W_r,f_r,(U_{r,i})_i,(V_{r,i})_i\big)$ are i.i.d.: the
pairs $(W_r,f_r)$ are i.i.d., the families $(U_{r,i})$ and $(V_{r,i})$ are independent of each other and of the
pairs, and their laws do not depend on $r$. Since $(W_r,(S_{r,i})_i)$ is the same measurable function of each
quadruple, the four pairs $(W_r,(S_{r,i})_{i\le N})$ are i.i.d.

\emph{The outside fraction.} If $f_r=1$, every $S_{r,i}\in\A'^c$, so $\hat q_r(1_{\A'^c})=1$; if $f_r=0$, every
$S_{r,i}\in\A'$, so $\hat q_r(1_{\A'^c})=0$. Hence $f_r=\hat q_r(1_{\A'^c})$.

\emph{The identity.} Fix $g:\Omega'\to\mathbb R$ (bounded, as $\Omega'$ is finite). Since $(U_{r,i})$ and $(V_{r,i})$ are
independent of $(W_r,f_r)$,
\[
  \mathbb E\big[W_r\,\hat q_r(g)\,\big|\,W_r,f_r\big]
  =W_r\Big(f_r\,\pi'_{\rm out}(g)+(1-f_r)\,\pi'_{\rm in}(g)\Big),
\]
with $\pi'_{\rm out}(g)=\pi'(g1_{\A'^c})/a$ and $\pi'_{\rm in}(g)=\pi'(g1_{\A'})/(1-a)$. All variables involved are
bounded, so taking expectations and using $\mathbb E[W_rf_r]=a$ and $\mathbb E[W_r(1-f_r)]=\mathbb EW_r-\mathbb E[W_rf_r]=1-a$,
\[
  \mathbb E\big[W_r\,\hat q_r(g)\big]=a\cdot\frac{\pi'(g1_{\A'^c})}{a}+(1-a)\cdot\frac{\pi'(g1_{\A'})}{1-a}
  =\pi'(g1_{\A'^c})+\pi'(g1_{\A'})=\pi'(g).
\]
Multiplying by any constant $Z'>0$ gives $\mathbb E[(Z'W_r)\,\hat q_r(g)]=Z'\pi'(g)$, the identity of (H-T12) with
$\hat Z_r:=Z'W_r$, for every test function $g$. Each island has a single population; using it at every recording
position ($\hat q_{r,\tau}:=\hat q_r$ for all $\tau\in\mathcal T$) gives the identity at every position. This proves
(iv-b). \qed

\begin{T15remark}[Consistency with (i)--(iii)]
In the example of (iv-a), the islandwise bounds hold, as they must. Since $\varepsilon<c<h$,
$\mathbb P(W_1<c)=\mathbb P(W_1=\varepsilon)=3/4$, so $a/c+\mathbb P(W_1<c)\ge3/4+b=\mathbb E f_1$ because
$b<a/3<a/c$; and $a/(cu)+\mathbb P(W_1<c)\ge 3/4+b=\mathbb P(f_1>u)$ because $b<a/3<a/(cu)$. The pooled version fails
because of the event that island~1 has small evidence $W_1=\varepsilon$ and $f_1=1$ while another island has large
evidence $h$, which keeps $\bar W$ above $c$. Part (iv) is a non-implication about substituting the pooled lower
tail into (i)--(ii); it does not claim that the SMC final run, for a particular data set and calibration, realises
this law.
\end{T15remark}

\begin{T15remark}[Scope and use of hypotheses]
Parts (i)--(iii) bound the law of $f_r$ under $\Pc$; they give no bound on the outside fraction of a realised
population. Parts (i) and (ii) use only $W_r\ge0$, $0\le f_r\le1$ and the instance $f=1_{\A^c}$ of (H-T12) (in
fact only $\Ec[W_rf_r]\le\pi(\A^c)$). The first claim of part (iii), $\Pc(W_r<c)\le B/(1-c)$, uses only its own hypothesis
$\Ec|W_r-1|\le B$ (which also makes $W_r$ integrable) and Markov's inequality, and not (H-T12); its
``consequently'' clause substitutes this bound into (i) and (ii), and so rests on the same instance of (H-T12) as
they do. Part (iv) concerns explicitly constructed laws
and uses no property of the final run. Assumption EXP (or EXP$'$) is not used anywhere.
\end{T15remark}
\mdproofend
\mdsec{(d) Review record}
\mdsubsec{FINAL\_REPORT.md: statement changes}
\begin{mdblock}
\mdhead{2}{2. Statement changes and why}

The inequalities, constants and quantifiers of the theorem file are unchanged. The other changes:

\begin{mdbul}
\item \textbf{T12 dependency made explicit.} The theorem file says "Builds on T12". STATEMENT.md states T12(i) as the hypothesis (H-T12), as the run rules require. No statement-checker call was needed, because this follows the run rules and is not a new assumption.
\item \textbf{Pooled-tail sentence made precise.} The theorem file's sentence "a pooled-evidence lower tail does not substitute" became part (iv), a non-implication with an explicit example. It was strengthened during review:

\begin{mdbul}
\item v2: strictly positive evidence (Astra \#1);
\item v3: every outside mass a ∈ (0, cu/4], plus (iv-b), which realises the example with particle populations satisfying the T12 identity for every test function (suggested by the Opus reviewer);
\item v4: wording only.
\end{mdbul}
\item \textbf{Realised-population sentence kept as a scope remark.} The theorem file's sentence on realised populations is a scope remark: the bounds concern the law of the run, not a realised population.
\item \textbf{EXP / EXP′ not needed.} No decision from Decision Point 2 affects T15.
\end{mdbul}
\end{mdblock}

\mdsubsec{FINAL\_REPORT.md: review history}
\begin{mdblock}
\mdhead{2}{3. Review history}

{\small
\begin{longtable}{|>{\raggedright\arraybackslash}p{\dimexpr(\linewidth-10\tabcolsep-6\arrayrulewidth)*74/1000\relax}|>{\raggedright\arraybackslash}p{\dimexpr(\linewidth-10\tabcolsep-6\arrayrulewidth)*416/1000\relax}|>{\raggedright\arraybackslash}p{\dimexpr(\linewidth-10\tabcolsep-6\arrayrulewidth)*104/1000\relax}|>{\raggedright\arraybackslash}p{\dimexpr(\linewidth-10\tabcolsep-6\arrayrulewidth)*119/1000\relax}|>{\raggedright\arraybackslash}p{\dimexpr(\linewidth-10\tabcolsep-6\arrayrulewidth)*287/1000\relax}|}
\hline
\textbf{Stage} & \textbf{Reviewer} & \textbf{Version} & \textbf{Verdict} & \textbf{Findings} \\ \hline\endhead
A & Astra \#1 (session 01a16418-957c-7832-97e0-ccb8b8f10081), r1 & v1 & verified & 1 exposition (example used W = 0) → fixed in v2 \\ \hline
A & Astra \#1, r2 & v2 & verified & none \\ \hline
B & Astra \#2 (new session 01a1641c-660b-7822-a3a2-f4a9f9c67d7b), r1 & v2 & verified & none \\ \hline
B & Opus 5.5 xhigh (Claude Personal), r1 & v2 & verified & 5 exposition → fixed in v3 \\ \hline
B & Astra \#2, r2 (ruling on v3 + log) & v3 & verified & 2 exposition → fixed in v4 \\ \hline
B & Opus, r2 (ruling on v3 + log) & v3 & verified & 4 exposition (2 shared with Astra \#2) → fixed in v4 \\ \hline
B & Astra \#2, r3 & v4 & verified & none of any severity \\ \hline
B & Opus, r3 & v4 & verified & 1 optional wording note → aligned in STATEMENT §4 \\ \hline
\end{longtable}}

\begin{mdbul}
\item \textbf{Fatal or gap findings:} none, from any reviewer, at any round. CLAIMS\_OF\_FATAL\_FLAWS.md records this, with every exposition finding and the lead's answer.
\item \textbf{Rulings on the log:} both Stage B reviewers ruled on it and agree that nothing is outstanding.
\item \textbf{Astra judge:} not used.
\item \textbf{Lead's judgment of the findings:} all were accepted as correct. None changed parts (i)–(iii), which are word for word the same as in v1.
\end{mdbul}
\end{mdblock}

\mdsubsec{CLAIMS\_OF\_FATAL\_FLAWS.md (every entry, verbatim)}
\begin{mdblock}
\mdhead{1}{CLAIMS\_OF\_FATAL\_FLAWS — T15 (L-OCC)}

\mdhead{2}{Fatal or gap findings}

None. Eight reviews, and none raised a fatal or gap finding:

\begin{mdbul}
\item Astra \#1, round 1, on PROOF\_v1: verified; 1 exposition finding.
\item Astra \#1, round 2, on PROOF\_v2: verified; no new findings.
\item Stage B Astra \#2, round 1, on PROOF\_v2: verified; no findings.
\item Stage B Opus 5.5 xhigh, round 1, on PROOF\_v2: verified; 5 exposition findings.
\item Stage B Astra \#2, round 2, on PROOF\_v3 + this log: verified; 2 exposition findings. Agrees that no fatal or gap finding is outstanding. Rulings: A1-1, O-1, O-2, O-3, O-5 resolved; O-4 partially resolved.
\item Stage B Opus, round 2, on PROOF\_v3 + this log: verified; 4 exposition findings. Agrees that no fatal or gap finding is outstanding. Rulings: O-1 to O-5 resolved.
\item Stage B Astra \#2, round 3, on PROOF\_v4: verified; no findings of any severity; B2-1 to B2-4 resolved.
\item Stage B Opus, round 3, on PROOF\_v4: verified; B2-1 to B2-4 resolved; one optional exposition note (a leftover "full identity" in the STATEMENT §4 scope remark, which it judged accurate) was aligned afterwards (wording only).
\end{mdbul}

\mdhead{2}{Exposition findings and the lead's answers (kept for the record)}

{\small
\begin{longtable}{|>{\raggedright\arraybackslash}p{\dimexpr(\linewidth-10\tabcolsep-6\arrayrulewidth)*28/1000\relax}|>{\raggedright\arraybackslash}p{\dimexpr(\linewidth-10\tabcolsep-6\arrayrulewidth)*55/1000\relax}|>{\raggedright\arraybackslash}p{\dimexpr(\linewidth-10\tabcolsep-6\arrayrulewidth)*408/1000\relax}|>{\raggedright\arraybackslash}p{\dimexpr(\linewidth-10\tabcolsep-6\arrayrulewidth)*93/1000\relax}|>{\raggedright\arraybackslash}p{\dimexpr(\linewidth-10\tabcolsep-6\arrayrulewidth)*415/1000\relax}|}
\hline
\textbf{\#} & \textbf{Reviewer} & \textbf{Finding (short)} & \textbf{Answer} & \textbf{Change} \\ \hline\endhead
A1-1 & Astra \#1 r1 & The (iv) example has W = 0 with probability 3/4, but actual island evidence is strictly positive. Heading scope needs care. & Accepted. & v2: example uses ε > 0; scope sentence added to (iv). \\ \hline
O-1 & Opus B r1 & (H-T12) gives the identity for every test function g, but the example checks only g ∈ \{1, 1\_\{𝔄ᶜ\}\}. Opus suggested a populations-based extension. & Accepted; adopted as a strengthening. & v3: new (iv-b) realises the example with particle populations satisfying E[W q̂(g)] = π′(g) for all g, on any finite Ω′ (including the method's Ω, π) with π′(𝔄′ᶜ) = a. (iv) now holds for every a ∈ (0, cu/4]. \\ \hline
O-2 & Opus B r1 & Island i.i.d. given 𝓒 is presented as a method property, but it is not among the hypotheses. & Accepted. & v3: the claim about the method is removed; the example's i.i.d. property is described as only strengthening it. \\ \hline
O-3 & Opus B r1 & Z̄₀, d, m(∅), π₀ are undefined, and the positivity of Ẑ\_r is asserted, not derived. & Accepted. & v3: the symbols are defined briefly with a pointer to 00\_COMMON; positivity is derived from positive factors. The proof uses only W\_r ≥ 0. \\ \hline
O-4 & Opus B r1 & The Scope remark misdescribes the hypothesis use of (iii) and (iv). & Accepted. & v3: the remark is rewritten part by part. \\ \hline
O-5 & Opus B r1 & (iv) refutes the pooled form of (i) without the min\{1,·\} and does not state that the min form is refuted too. & Accepted. & v3: the one-line step min\{1,x\} ≤ x is added to (iv-a) and its proof. \\ \hline
B2-1 & Astra \#2 r2 = Opus r2 \#1 & The i.i.d. justification in (iv-b) cites (W\_r, U, V) and leaves out f\_r. & Accepted (both reviewers raised it independently). & v4: the argument now uses the quadruples (W\_r, f\_r, U\_r, V\_r). \\ \hline
B2-2 & Astra \#2 r2 = Opus r2 \#2 & "Part (iii) does not use (H-T12)" holds only for its first claim; the "consequently" clause rests on (i)–(ii). & Accepted. & v4: the note in STATEMENT §2 and the proof's final remark separate the two claims. \\ \hline
B2-3 & Opus r2 \#3 & (iv-b) gives the identity at one position, while "full identity" suggests all of 𝒯. & Accepted. & v4: "full identity" is replaced by "identity for every test function", plus a remark that reusing the population at every τ gives every position. \\ \hline
B2-4 & Opus r2 \#4 & "The method's Ω and π" example needs some set of posterior mass exactly a. & Accepted. & v4: the example states when such a set exists (smallest atom, if \#Ω ≥ 4/(cu)). \\ \hline
\end{longtable}}

Unified version after Stage B round 1: v3. After Stage B round 2 (wording only): STATEMENT.md (v4 bullet in §5) and proof/PROOF\_v4.tex/.pdf. (i)–(iii) are unchanged in statement and proof since v1.
\end{mdblock}

\chapter{T20 (T15), clauses (i)--(ii): local-search semantics}
\chaptersource{Final proof: \texttt{proofs/\allowbreak{}T20\_\allowbreak{}T15\_\allowbreak{}local\_\allowbreak{}search/\allowbreak{}proof/\allowbreak{}PROOF\_v3.tex}. Sources of parts (a), (b), (d): \texttt{theorems/\allowbreak{}T20\_\allowbreak{}T15\_\allowbreak{}local\_\allowbreak{}search.md}, \texttt{STATEMENT.md}, \texttt{FINAL\_REPORT.md}, \texttt{CLAIMS\_OF\_FATAL\_FLAWS.md}.}
\mdsec{(a) Original conjecture text (theorem file, verbatim)}
\begin{mdblock}
\mdhead{1}{T20 (T15): Local-search semantics and conditional repair}

\textbf{Status:} conjectured, not proven. Methodology SMC-BVS v3.0 (2026-10-19), SHA-256 93043758058a6d422cca6ca4b6be8f3a9e8149b6f6de2e1bab1ce31c44b8019c. Read with {\mdtt 00\_\allowbreak{}COMMON.\allowbreak{}md}.

\textbf{Category / rank:} deterministic component validity and statistical consistency. Clauses (i)–(ii): \textbf{core}; (iii) ORTH repair: \textbf{secondary}; (iv) COH extension: \textbf{optional}. v2's SE-only "every best move is a true move" clause is withdrawn.

\mdhead{2}{Definitions used}

The local search: from a centre S\_c, each scan scores every feasible flip and swap of the centre; best-improvement moves, at most ℓ moves, then one extra scan; {\mdtt local\_\allowbreak{}certified} flag; running all-scored HPM Ŝ; E\_n := \{S₀ is the strict global maximiser of F\} (00\_COMMON §§1–2; THEORY\_BRIEF §2).

\mdhead{2}{Assumptions}

EX-A and execution of the relevant scans for (i)–(ii); G, R, SE and BM-C (through T03) for the statistical reading; ORTH and BM-∞ for (iii); COH and BM-∞ for (iv).

\mdhead{2}{Precise statement}

(i) [certificate] {\mdtt local\_\allowbreak{}certified} certifies exactly that the returned model was itself scanned as a centre and had no strictly improving feasible flip or swap within the cap; equal-score neighbours are allowed; weak local maximality only. All scanned candidates, including rejected ones, enter the running HPM comparison; a model first discovered by the final extra scan can be the reported HPM without a certificate. (ii) [one neighbour] On E\_n, every completed scan from a feasible flip-or-swap neighbour of S₀ scores S₀; if a move remains the search moves to S₀; if the subsequent scan executes it certifies S₀. Under T03, P₀(E\_n) → 1, and a sufficient global margin is min\_\{S≠S₀\}\{F(S₀) − F(S)\} ≥ log[π(S₀)/(1−π(S₀))] ≥ ½ log n − ½ log log n − O\_\{P₀\}(1). (iii) [ORTH repair] Under G, R, ORTH, BM-∞: with probability → 1, simultaneously for every feasible starting centre with |S\_c Δ S₀| ≤ ℓ, each selected best improvement reduces |S Δ S₀| (by 1 for a true addition or noise deletion, by 2 for a true-for-noise swap); the search reaches S₀ within |S\_c Δ S₀| moves and the extra scan certifies it. Equal magnitudes are unnecessary. (iv) [COH extension] The same conclusion under G, R, COH, BM-∞, with the obligation of proving monotonicity of the selected best move uniformly over the centres with the integrated-variance score. (v) [scope] Posterior concentration does not imply the entry condition of (iii), nor bound an arbitrary centre's size by m\_n; the move limit can bind; a later search repeats an earlier ascent only if its centre and comparisons are unchanged. Under SE alone a signal-bearing proxy can win the first improvement from ∅ (correlation matrix [[1,0,.65],[0,1,.65],[.65,.65,1]], S₀ = \{1,2\}: eigenvalue floor 0.081 > 0, proxy population fit gain 1.69 vs 1 per true singleton).

\mdhead{2}{Dependencies}

Builds on: T03, T05, T11. Used by: T21, T24, T31.

\mdhead{2}{Numerical evidence (provisional)}

provisional evidence (Sol check J3): on the design above with θ = (1,1,0), best-improvement paths from ∅ at n = 64, 512, 2048 select the proxy first and recover S₀ afterwards; at n = 128 the path stops at a certified wrong local maximum although S₀ is globally best; one-neighbour scoring/certification semantics confirmed; a truth first hit on the final scan is scored but not certified. Files: judge/sol1/J3\_LOCAL\_SEARCH.md.

\mdhead{2}{Proof-route hint (optional)}

(i)–(ii): frozen-centre scan rule and the all-scored ledger (T11). (iii): the gain statistics G\_l do not depend on the centre under ORTH; one no-lucky event \{max\_l G\_l < x\_k for all k ≤ s₀+ℓ\} of probability → 1 separates true-coordinate gains from noise gains in the exact integrated odds. (iv): control residualised gains uniformly over admissible centres (deviation ≤ 4c\_n|T| log p/φ₀²(1+o(1)) from ORTH). Hardest step: best-move ordering after integrating variance. \emph{When proof.md is run on this file, its lead considers the statement both with and without these hints; the hints never reach the reviewers.}
\end{mdblock}

\mdsec{(b) STATEMENT.md (verbatim)}
\begin{mdblock}
\mdhead{1}{STATEMENT.md — T20 (T15), clauses (i)–(ii): local-search semantics}

Source: {\mdtt Theory\_\allowbreak{}architecture\_\allowbreak{}v3/\allowbreak{}theorems/\allowbreak{}T20\_\allowbreak{}T15\_\allowbreak{}local\_\allowbreak{}search.\allowbreak{}md} (clauses (i)–(ii) only; (iii)–(v) are out of scope). Method: SMC-BVS v3.0, {\mdtt Methodology/\allowbreak{}SMC\_\allowbreak{}BVS\_\allowbreak{}v3\_\allowbreak{}20261019/\allowbreak{}METHODOLOGY\_\allowbreak{}v3.\allowbreak{}md}, SHA-256 93043758058a6d422cca6ca4b6be8f3a9e8149b6f6de2e1bab1ce31c44b8019c (checked 2026-10-22). Statement analysis outcome: \textbf{(a) correct as stated}; no hypothesis was changed or added. Section 5 lists the readings used to make the informal wording precise.

\mdhead{2}{0. Original clauses (verbatim)}

\begin{mdquote}
(i) [certificate] {\mdtt local\_\allowbreak{}certified} certifies exactly that the returned model was itself scanned as a centre and had no strictly improving feasible flip or swap within the cap; equal-score neighbours are allowed; weak local maximality only. All scanned candidates, including rejected ones, enter the running HPM comparison; a model first discovered by the final extra scan can be the reported HPM without a certificate.

(ii) [one neighbour] On E\_n, every completed scan from a feasible flip-or-swap neighbour of S₀ scores S₀; if a move remains the search moves to S₀; if the subsequent scan executes it certifies S₀. Under T03, P₀(E\_n) → 1, and a sufficient global margin is min\_\{S≠S₀\}\{F(S₀) − F(S)\} ≥ log[π(S₀)/(1−π(S₀))] ≥ ½ log n − ½ log log n − O\_\{P₀\}(1).
\end{mdquote}

Assumptions named by the source for (i)–(ii): EX-A (exact scores and comparisons) and execution of the relevant scans; for the statistical reading, G, R, SE and BM-C or BM-∞ through T03. The source defines (its "Definitions used"): "E\_n := \{S₀ is the strict global maximiser of F\}".

Method text formalised in §2 (METHODOLOGY\_v3 §6.10, verbatim): "\textbf{Local search:} from the current best scored model (the centre), scan every feasible flip and swap (additions beyond K excluded) with the centre frozen during the scan; move to the best strict improvement (lexicographic ties); at most ℓ moves, then one further scan; stop earlier when a scan finds no improvement. Every evaluated model updates the global best. If the last scan still improves, set {\mdtt local\_\allowbreak{}limit} (no restart). {\mdtt local\_\allowbreak{}certified} only when the returned model is a scanned centre with no improving neighbour. The search never alters particles." and "\textbf{HPM:} the strict maximum of log ϖ(S) + log m(S) over every model ever scored (tables, pilot, calibration, final rounds, rejected proposals, independence proposals, local searches), kept as a running best that survives cache eviction. Only exactly equal scores invoke the tie rule: smallest |S|, then lexicographically smallest index set." Execution order (§6.1/§7): "… final local search; final checks (flags only); report" — nothing is scored after the final local search. ℓ = ⌈log₂(n + p + 1)⌉ (§3).

\mdhead{2}{1. Setting}

\begin{mdbul}
\item p ≥ 1 predictors, n ≥ 3, K = min(p, ⌊n/2⌋) ≥ 1, state space Ω = \{S ⊆ [p] : |S| ≤ K\}.
\item \textbf{Score.} F : Ω → ℝ. Clause (i) and the deterministic part of clause (ii) are proved for an \textbf{arbitrary} F : Ω → ℝ; only (ii-B)–(ii-C) use the method's score F(S) = log[ϖ(S) m(S)], π(S) = e\textasciicircum{}\{F(S)\}/Σ\_\{T∈Ω\} e\textasciicircum{}\{F(T)\} for S ∈ Ω, and the convention π(S) := 0 for S ∉ Ω. The method's score (00\_COMMON §1; METHODOLOGY\_v3 §5), on standardised data X ∈ ℝ\textasciicircum{}\{n×p\}, y ∈ ℝⁿ (y ≠ 0; columns and y centred, ‖x\_j‖² = ‖y‖² = n − 1): for S ∈ Ω with k = |S|, X\_S the columns in S, G\_S = X\_SᵀX\_S, z\_S = X\_Sᵀy, H\_S = (n+1)G\_S + I\_k, D\_S = det H\_S / det(G\_S + I\_k), Q\_S = yᵀy − n z\_SᵀH\_S⁻¹z\_S (for S = ∅: D\_∅ = 1, Q\_∅ = yᵀy); a₀ = b₀ = 0.01, ν = a₀ + n/2; m(S) = (2b₀)\textasciicircum{}\{a₀\} Γ(ν) / (π\textasciicircum{}\{n/2\} Γ(a₀)) · D\_S\textasciicircum{}\{−1/2\} (2b₀ + Q\_S)\textasciicircum{}\{−ν\}; w(k) = ½ / C(2p−1, k), W\_K = Σ\_\{j≤K\} C(p, j) w(j), ϖ(S) = w(|S|)/W\_K. (This is the marginal likelihood of the model y | β\_S, σ² \textasciitilde{} N(X\_Sβ\_S, σ²I), β\_S | σ² \textasciitilde{} N(0, nσ²(G\_S+I)⁻¹), σ² \textasciitilde{} IG(a₀, b₀), times the beta-binomial BB(1, p) prior capped at K; the proof uses only the displayed formulas.)
\item \textbf{EX-A.} All scores are exact real numbers and all comparisons are exact.
\item \textbf{HPM order ≻} (method §6.10): S ≻ T iff F(S) > F(T); or F(S) = F(T) and |S| < |T|; or F(S) = F(T), |S| = |T| and the increasingly sorted index list of S is lexicographically smaller than that of T. This is a strict total order on Ω. For a finite nonempty 𝒜 ⊆ Ω, best(𝒜) is its ≻-maximum. The running HPM after any step is best(𝒮) where 𝒮 is the set of all models scored so far in the run (tables, pilot, calibration, final rounds, accepted and rejected proposals including independence proposals, local scans).
\item \textbf{Move order ▷.} A fixed strict total order on Ω with F(S) > F(T) ⇒ S ▷ T (maximise F, exact ties broken by a fixed deterministic rule; the method says "lexicographic ties"). Every result holds for every such rule.
\item \textbf{Feasible flip-or-swap neighbourhood} of c ∈ Ω: 𝒩(c) = \{c ∪ \{j\} : j ∉ c, |c| < K\} ∪ \{c ∖ \{i\} : i ∈ c\} ∪ \{(c ∖ \{i\}) ∪ \{j\} : i ∈ c, j ∉ c\} ⊆ Ω. (All swaps, not only neighbour-graph swaps; additions beyond K excluded.) The relation N ∈ 𝒩(c) is symmetric on Ω.
\end{mdbul}

\mdhead{2}{2. The local search (method §6.10), as formalised}

Input: the finite nonempty set 𝒮⁻ ⊆ Ω of models scored in the run before the search starts; the move cap ℓ ≥ 0 (the method has ℓ = ⌈log₂(n+p+1)⌉).

\begin{Verbatim}[formatcom=\mdtt,fontsize=\scriptsize]
c_0 := best(𝒮⁻);  𝒮 := 𝒮⁻;  k := 0
loop:
   scan k (centre c_k frozen): 𝒮 := 𝒮 ∪ 𝒩(c_k)          # every neighbour scored; running HPM updated
   I_k := {N ∈ 𝒩(c_k) : F(N) > F(c_k)}                    # strict improvements over the centre
   if I_k = ∅:  stop; outcome NOIMP; m := k
   elif k < ℓ:  c_{k+1} := the ▷-maximum of I_k; k := k + 1  # a move
   else (k = ℓ): stop; outcome LIMIT; m := ℓ                 # the "one further scan" still improved
\end{Verbatim}

Outputs: centres 𝒞 = \{c\_0, …, c\_m\}; final scored set 𝒮\_end; \textbf{returned model R := best(𝒮\_end)} (the running HPM when the search ends); {\mdtt local\_\allowbreak{}limit} := 1 iff outcome LIMIT; \textbf{{\mdtt local\_\allowbreak{}certified} := 1 iff R ∈ 𝒞 and there is no N ∈ 𝒩(R) with F(N) > F(R)} (method: "{\mdtt local\_\allowbreak{}certified} only when the returned model is a scanned centre with no improving neighbour").

The run may contain several local searches (after the pilot, after each set of diagnostic islands, after the final run); between them the scored set may grow. In the method nothing is scored after the final local search, so for the final search R equals the reported HPM Ŝ.

\textbf{Execution hypothesis (EXEC).} Every scan that the procedure calls is completed, i.e. all of 𝒩(c\_k) is scored.

\mdhead{2}{3. Clause (i): certificate semantics (EX-A, EXEC; arbitrary F)}

For any execution of one local search:

\begin{mdbul}
\item \textbf{(i-a) Characterisation.} {\mdtt local\_\allowbreak{}certified} = 1 iff the outcome is NOIMP and R = c\_m. In that case F(R) = max\_\{𝒮\_end\} F and F(R) ≥ F(N) for every N ∈ 𝒩(R) (weak local maximality over the feasible flips and swaps; neighbours with F(N) = F(R) are permitted, and every such neighbour N satisfies R ≻ N, i.e. |N| > |R|, or |N| = |R| and N is lexicographically later). Moreover R ∈ 𝒞 holds iff R was a scan centre in any local search of the run up to and including this one.
\item \textbf{(i-b) Only weak local information.} The whole execution (centres, outcome, 𝒮\_end, R, both flags) is unchanged if F is replaced by any F′ : Ω → ℝ with F′ = F on 𝒮\_end, the F′-run using the same input 𝒮⁻, the same ℓ, the HPM order ≻ computed from F′, and the move order computed from F′ with the same fixed tie-break rule. Hence {\mdtt local\_\allowbreak{}certified} = 1 carries no information about F on Ω ∖ 𝒮\_end; in particular, when Ω ∖ 𝒮\_end ≠ ∅ it does not certify global maximality (when 𝒮\_end = Ω, R is the global maximiser of F trivially).
\item \textbf{(i-c) Ledger.} Every model scored by any scan of the search, whether or not it was selected as a move, belongs to 𝒮\_end, so R ⪰ N (R = N or R ≻ N) for each of them; in particular F(R) ≥ F(N). For the final search R = Ŝ.
\item \textbf{(i-d) Final extra scan.} On outcome LIMIT: R ∈ 𝒩(c\_ℓ), F(R) > F(c\_ℓ), R was first scored by scan ℓ (R ∉ the scored set before scan ℓ), R was never a scan centre of any local search of the run up to this one, and {\mdtt local\_\allowbreak{}certified} = 0. For the final search this R is the reported HPM Ŝ.
\item \textbf{(i-e) Generic case.} If exactly one model of 𝒮\_end attains max\_\{𝒮\_end\} F, then {\mdtt local\_\allowbreak{}certified} = 1 ⟺ outcome NOIMP ⟺ {\mdtt local\_\allowbreak{}limit} = 0.
\end{mdbul}

\mdhead{2}{4. Clause (ii): one neighbour of S₀}

Let S₀ ⊆ [p] be the true support and \textbf{E\_n := \{S₀ ∈ Ω and F(S₀) > F(S) for every S ∈ Ω ∖ \{S₀\}\}}.

\textbf{(ii-A) Deterministic (EX-A, EXEC; arbitrary F).} On E\_n, for every local search of the run and every scan k of it whose centre c\_k is a feasible flip-or-swap neighbour of S₀ (c\_k ∈ 𝒩(S₀)): 1. scan k scores S₀ (S₀ ∈ 𝒩(c\_k)); 2. S₀ ∈ I\_k and S₀ is the ▷-maximum of I\_k; if k < ℓ (a move remains), c\_\{k+1\} = S₀; 3. if k < ℓ, scan k+1 (with centre S₀) is called by the procedure; it finds I\_\{k+1\} = ∅, the search stops with outcome NOIMP, R = S₀ and {\mdtt local\_\allowbreak{}certified} = 1; 4. if k = ℓ (no move remains), the search stops with outcome LIMIT, R = S₀ and {\mdtt local\_\allowbreak{}certified} = 0. In both cases S₀ is the running HPM from the moment scan k scores it; for the final search Ŝ = S₀.

\textbf{(ii-B) Margin inequality (method's F; deterministic).} If S₀ ∈ Ω then 0 < π(S₀) < 1 and\\
\hspace*{1.5em}Δ\_n := min\_\{S∈Ω∖\{S₀\}\} [F(S₀) − F(S)] ≥ λ\_n := log[π(S₀)/(1 − π(S₀))]. Consequently \{π(S₀) > ½\} ⊆ E\_n.

\textbf{(ii-C) Statistical (method's F).} Regime: a sequence of problems indexed by n → ∞ (design fixed, P₀ = law of the Gaussian errors), r\_n := √(log n / n). Hypothesis \textbf{H03} (the conclusion of T03(i), conjectured there under R1 with G, R, SE and BM-C or BM-∞; it is used here only as a hypothesis): there is n₀ with S₀ ∈ Ω for all n ≥ n₀, and 1 − π(S₀) = O\_\{P₀\}(r\_n). Under H03: 1. P₀(E\_n) → 1; 2. (½ log n − ½ log log n − λ\_n)⁺ = O\_\{P₀\}(1) (over n ≥ n₀), i.e. Δ\_n ≥ λ\_n ≥ ½ log n − ½ log log n − O\_\{P₀\}(1); 3. under EX-A and EXEC, for any joint law of data and algorithm randomness whose data marginal is P₀, the probability that conclusions (ii-A)1–4 hold simultaneously for every scan of every local search of the run is at least P₀(E\_n), hence tends to 1.

\mdhead{2}{5. Readings used (the informal text made precise)}

\begin{mdnum}
\item[\mdlbl{1.}] "Returned model" = the running HPM when the search ends (R above); this matches the source's "a model first discovered by the final extra scan can be the reported HPM without a certificate". If one instead read it as the last centre c\_m: on outcome LIMIT, R ≠ c\_m (by (i-d), F(R) > F(c\_m)), but the flag is 0 under both readings (c\_m has a strictly improving neighbour); on outcome NOIMP, R = c\_m unless another model of 𝒮\_end has exactly the score F(c\_m) and precedes c\_m in the HPM order. So the two readings give different flag values only on NOIMP with such an exact tie.
\item[\mdlbl{2.}] "Scanned as a centre" = R ∈ 𝒞; (i-a) shows this is the same as having been a centre in any local search of the run up to and including the current one.
\item[\mdlbl{3.}] "Within the cap" = among feasible models, |S| ≤ K. (Reading it as the move cap ℓ gives nothing different: certification implies outcome NOIMP at some m ≤ ℓ.) The method's "{\mdtt local\_\allowbreak{}certified} only when …" is read as "iff", as the source's "certifies exactly that" requires.
\item[\mdlbl{4.}] "Weak local maximality only" is made precise by (i-a) (what is implied) and (i-b) (nothing about unscored models is implied).
\item[\mdlbl{5.}] "Can be the reported HPM without a certificate" is made precise by (i-d): on outcome LIMIT the reported model is such a model and is uncertified. Whether outcome LIMIT has positive probability for given data is not asserted. (A second, tie-only route to an uncertified reported model also exists: on NOIMP, if some scored model has exactly the score F(c\_m) and precedes c\_m in the HPM order, then R is the HPM-first of all such models (it may be a neighbour of c\_m or a model scored earlier), R ≠ c\_m, and both flags are 0; (i-a) covers it.)
\item[\mdlbl{6.}] "If the subsequent scan executes" = EXEC; in the procedure a scan always follows a move (a regular scan or the one further scan).
\item[\mdlbl{7.}] "Under T03" = hypothesis H03; "a sufficient global margin is … ≥ …" = the chain Δ\_n ≥ λ\_n ≥ ½ log n − ½ log log n − O\_\{P₀\}(1) of (ii-B)–(ii-C), with Δ\_n > 0 equivalent to E\_n when S₀ ∈ Ω.
\end{mdnum}

\mdhead{2}{6. Change log}

\begin{mdbul}
\item v3 → v3.1 (2026-10-22, after the Stage B ruling round; text of §1 and §5 only, proof unchanged): §5 item 5 tie-route sentence made precise (Opus N1 = Astra \#2 exposition finding); §1 list of HPM sources now names independence proposals (Opus N2).
\item v2 → v3 (2026-10-22, after the Stage B reviews; exposition only, no hypothesis added or changed): verbatim source text for E\_n and for METHODOLOGY\_v3 §6.10 added to §0; (i-a) states that equal-score neighbours of a certified R come after R in the HPM order; (i-b) says which move order the F′-run uses and qualifies "does not certify global maximality" by Ω ∖ 𝒮\_end ≠ ∅; (ii-A) closing sentence "from the moment scan k scores it"; (ii-C)3 names EX-A and EXEC; §5 items 3 and 5 extended.
\item v1 → v2 (2026-10-22, after Astra \#1 round 1; no hypothesis or claim of §§3–4 changed): §1 now writes out the method's score formulas (Q\_S, D\_S, m, ϖ) in full, so the statement is self-contained; §5 item 1 corrected (the two readings of "returned model" also differ on outcome LIMIT, where both give flag 0); §5 item 2 wording corrected ("up to and including the current one").
\end{mdbul}
\end{mdblock}

\mdsec{(c) Final proof, verbatim: body of PROOF\_v3.tex}
\noindent{\small\itshape The proof below is the body of \texttt{PROOF\_v3.tex} from \texttt{\textbackslash begin\{document\}} to \texttt{\textbackslash end\{document\}} without the title block. Its preamble definitions are local to this chapter; labels carry the prefix \texttt{T20:}.}\par\medskip
\mdproofbegin
% ---- preamble of PROOF_v3.tex (theorem environments carry the prefix T20) ----
\newtheorem{T20lemma}{Lemma}
\newtheorem{T20corollary}[T20lemma]{Corollary}
\theoremstyle{remark}
\newtheorem{T20remark}[T20lemma]{Remark}
\newcommand{\cS}{\mathcal{S}}
\newcommand{\cN}{\mathcal{N}}
\newcommand{\cC}{\mathcal{C}}
\newcommand{\best}{\operatorname{best}}
\newcommand{\NOIMP}{\textsf{NOIMP}}
\newcommand{\LIMIT}{\textsf{LIMIT}}
\newcommand{\cert}{\texttt{local\_certified}}
\newcommand{\llim}{\texttt{local\_limit}}
% ---- body of PROOF_v3.tex ----
All notation, the procedure and the claims are those of \texttt{STATEMENT.md}; section numbers \S1--\S4 below refer to it. We recall what is used.

\section{Setting and conventions}

$\Omega=\{S\subseteq[p]:|S|\le K\}$ with $p\ge1$, $n\ge3$, $K=\min(p,\lfloor n/2\rfloor)\ge1$. A score is a function $F:\Omega\to\mathbb{R}$. All comparisons are exact (EX-A).

\paragraph{Orders.} The HPM order $\succ$: $S\succ T$ iff $F(S)>F(T)$, or $F(S)=F(T)$ and $|S|<|T|$, or $F(S)=F(T)$, $|S|=|T|$ and the increasingly sorted index list of $S$ is lexicographically smaller than that of $T$. Distinct sets of equal size have distinct sorted lists, and lexicographic order on lists of equal length is a strict total order; hence $\succ$ is a strict total order on $\Omega$ (irreflexive, transitive, and any two distinct elements are comparable). We write $S\succeq T$ for ($S\succ T$ or $S=T$). For finite nonempty $\mathcal A\subseteq\Omega$, $\best(\mathcal A)$ is the $\succ$-maximum. Since $F(S)>F(T)\Rightarrow S\succ T$,
\begin{equation}\label{T20:eq:bestmax}
F(\best(\mathcal A))=\max_{\mathcal A}F .
\end{equation}

The move order: fix a strict total order $\tau$ on $\Omega$ (the tie-break rule; it does not depend on $F$). Put $S\rhd_F T$ iff $F(S)>F(T)$, or $F(S)=F(T)$ and $S\,\tau\,T$. This is a strict total order with $F(S)>F(T)\Rightarrow S\rhd_F T$. Every strict total order $\rhd$ on $\Omega$ with $F(S)>F(T)\Rightarrow S\rhd T$ (the move order of \S1) equals $\rhd_F$ for $\tau:=\rhd$; so nothing is lost by writing the move order as $\rhd_F$. Again, for finite nonempty $\mathcal A$ the $\rhd_F$-maximum of $\mathcal A$ has score $\max_{\mathcal A}F$.

\paragraph{Neighbourhoods.} For $c\in\Omega$,
$\cN(c)=\{c\cup\{j\}:j\notin c,\ |c|<K\}\cup\{c\setminus\{i\}:i\in c\}\cup\{(c\setminus\{i\})\cup\{j\}:i\in c,\ j\notin c\}$.
Every element of $\cN(c)$ has size at most $K$, so $\cN(c)\subseteq\Omega$; and $c\notin\cN(c)$ (a flip changes the size, a swap removes $i\in c$ and inserts $j\notin c$, $j\ne i$). Moreover $\cN(c)\ne\emptyset$: if $c=\emptyset$ the addition $\{1\}$ is feasible ($K\ge1$), and if $c\ne\emptyset$ any deletion is in $\cN(c)$. So maxima of $F$ over $\cN(c)$ are maxima over finite nonempty sets.

\begin{T20lemma}[symmetry]\label{T20:lem:sym}
For $c,S\in\Omega$: $S\in\cN(c)\iff c\in\cN(S)$.
\end{T20lemma}
\begin{proof}
By the symmetry of the roles it suffices to prove $\Rightarrow$. (a) $S=c\cup\{j\}$, $j\notin c$: then $j\in S$ and $c=S\setminus\{j\}$ is a deletion from $S$. (b) $S=c\setminus\{i\}$, $i\in c$: then $i\notin S$, $|S|=|c|-1<|c|\le K$, so the addition $S\cup\{i\}=c$ is feasible and lies in $\cN(S)$. (c) $S=(c\setminus\{i\})\cup\{j\}$ with $i\in c$, $j\notin c$: then $j\ne i$, $j\in S$, $i\notin S$, and $c=(S\setminus\{j\})\cup\{i\}$ is a swap of $S$.
\end{proof}

\paragraph{The procedure (\S2).} Input: finite nonempty $\cS^-\subseteq\Omega$ (all models scored in the run before the search) and $\ell\ge0$. Put $c_0=\best(\cS^-)$, $\cS_{<0}=\cS^-$. For $k=0,1,\dots$: scan $k$ scores $\cN(c_k)$, so the scored set becomes $\cS_{<k+1}:=\cS_{<k}\cup\cN(c_k)$; let $I_k=\{N\in\cN(c_k):F(N)>F(c_k)\}$. If $I_k=\emptyset$, stop with outcome \NOIMP\ and $m=k$. Else, if $k<\ell$, let $c_{k+1}$ be the $\rhd_F$-maximum of $I_k$ and continue with $k+1$. Else ($k=\ell$) stop with outcome \LIMIT\ and $m=\ell$. Nothing else is scored during the search. Outputs: $\cC=\{c_0,\dots,c_m\}$, $\cS_{\rm end}=\cS_{<m+1}$, $R=\best(\cS_{\rm end})$, $\llim=1$ iff \LIMIT, and $\cert=1$ iff $R\in\cC$ and no $N\in\cN(R)$ has $F(N)>F(R)$. (EXEC: each scan scores all of $\cN(c_k)$.)

\section{The basic invariant}

\begin{T20lemma}\label{T20:lem:inv}
The procedure stops with $m\le\ell$. For every $k\le m$:
\begin{enumerate}[label=(\alph*)]
\item $c_k\in\cS_{<k}$ (each centre has been scored before its scan);
\item $F(c_k)=\max_{\cS_{<k}}F$;
\item if $k<m$ then $F(c_{k+1})=\max_{\cN(c_k)}F>F(c_k)$; in particular $F(c_0)<F(c_1)<\dots<F(c_m)$ and the centres are distinct.
\end{enumerate}
Moreover:
\begin{enumerate}[label=(\alph*),resume]
\item on \NOIMP, $\max_{\cS_{\rm end}}F=F(c_m)$;
\item on \LIMIT, $\max_{\cS_{\rm end}}F=\max_{\cN(c_\ell)}F>F(c_\ell)=\max_{\cS_{<\ell}}F$, so every maximiser of $F$ on $\cS_{\rm end}$ lies in $\cN(c_\ell)\setminus\cS_{<\ell}$.
\end{enumerate}
\end{T20lemma}
\begin{proof}
A move happens only from a scan with index $k<\ell$, and each move raises the index by one; so at most $\ell$ moves occur and the procedure stops at some $m\le\ell$ (each scan is a finite computation).

(a) $c_0=\best(\cS^-)\in\cS^-=\cS_{<0}$; and for $k<m$, $c_{k+1}\in I_k\subseteq\cN(c_k)\subseteq\cS_{<k+1}$.

(b), (c) by induction. For $k=0$, (b) is \eqref{T20:eq:bestmax}. Assume (b) for some $k<m$. Since $k<m$, scan $k$ led to a move, so $I_k\ne\emptyset$. Let $M=\max_{I_k}F$. Elements of $\cN(c_k)\setminus I_k$ have $F\le F(c_k)<M$, so $M=\max_{\cN(c_k)}F$. The $\rhd_F$-maximum $c_{k+1}$ of $I_k$ has $F(c_{k+1})=M$. Hence $F(c_{k+1})=\max_{\cN(c_k)}F>F(c_k)$, which is (c), and
$\max_{\cS_{<k+1}}F=\max\{\max_{\cS_{<k}}F,\max_{\cN(c_k)}F\}=\max\{F(c_k),M\}=M=F(c_{k+1})$, which is (b) for $k+1$.

(d) On \NOIMP, $I_m=\emptyset$, i.e.\ $F(N)\le F(c_m)$ for all $N\in\cN(c_m)$; with (b), $\max_{\cS_{\rm end}}F=\max\{\max_{\cS_{<m}}F,\max_{\cN(c_m)}F\}=F(c_m)$.

(e) On \LIMIT, $m=\ell$ and $I_\ell\ne\emptyset$, so $\max_{\cN(c_\ell)}F>F(c_\ell)=\max_{\cS_{<\ell}}F$ by (b); hence $\max_{\cS_{\rm end}}F=\max_{\cN(c_\ell)}F$, and any maximiser $T$ on $\cS_{\rm end}$ has $F(T)>\max_{\cS_{<\ell}}F$, so $T\notin\cS_{<\ell}$ and $T\in\cN(c_\ell)$.
\end{proof}

\begin{T20lemma}[earlier scores]\label{T20:lem:earlier}
If $R\in\cS^-$ then $R=c_0$.
\end{T20lemma}
\begin{proof}
$c_0=\best(\cS^-)$ and $R\in\cS^-$ give $c_0\succeq R$. $R=\best(\cS_{\rm end})$ and $c_0\in\cS^-\subseteq\cS_{\rm end}$ give $R\succeq c_0$. If $R\ne c_0$ we would have $c_0\succ R\succ c_0$, hence $c_0\succ c_0$ by transitivity, contradicting irreflexivity.
\end{proof}

\section{Clause (i)}

\paragraph{(i-a).}
Suppose $\cert=1$. Then $R=c_j$ for some $j\le m$, and $\cN(c_j)$ contains no $N$ with $F(N)>F(c_j)$, i.e.\ $I_j=\emptyset$. If $j<m$ the procedure moved after scan $j$, which requires $I_j\ne\emptyset$; so $j=m$. On \LIMIT, $I_m=I_\ell\ne\emptyset$; so the outcome is \NOIMP. Conversely, if the outcome is \NOIMP\ and $R=c_m$, then $R\in\cC$ and $I_m=\emptyset$ says that no $N\in\cN(R)$ has $F(N)>F(R)$, so $\cert=1$.

In that case $F(R)=F(c_m)=\max_{\cS_{\rm end}}F$ by Lemma~\ref{T20:lem:inv}(d), and $F(R)\ge F(N)$ for all $N\in\cN(R)$ because $I_m=\emptyset$. Equality $F(N)=F(R)$ is not excluded by the definition (only strict improvements are), so equal-score neighbours are compatible with $\cert=1$; for such $N$, since $N\in\cS_{\rm end}$ and $R=\best(\cS_{\rm end})$, the tie rule gives $R\succ N$.

Finally, $R\in\cC$ trivially implies that $R$ was a scan centre in some local search of the run up to this one. Conversely, suppose $R$ was a scan centre of an earlier local search. By Lemma~\ref{T20:lem:inv}(a) applied to that search, $R$ was scored then, so $R\in\cS^-$; by Lemma~\ref{T20:lem:earlier}, $R=c_0\in\cC$.

\paragraph{(i-b).}
Let $F':\Omega\to\mathbb{R}$ with $F'=F$ on $\cS_{\rm end}$, and run the procedure with $F'$ in place of $F$ (same $\cS^-$, $\ell$, tie rules $\succ$ via sizes and indices, and $\tau$). Write primes for the $F'$-run. We show by induction on $k$ that the two runs coincide step by step. The restriction of $\succ$ to a set $\mathcal A$ depends on $F$ only through $F|_{\mathcal A}$ (and on sizes and indices); since $\cS^-\subseteq\cS_{\rm end}$, $c_0'=\best'(\cS^-)=\best(\cS^-)=c_0$ and $\cS'_{<0}=\cS_{<0}$. Suppose $c_j'=c_j$ and $\cS'_{<j}=\cS_{<j}$ for all $j\le k$, where $k\le m$. Scan $k$ scores the same set $\cN(c_k)$ (it depends only on $c_k$ and $K$), so $\cS'_{<k+1}=\cS_{<k+1}$. Since $\{c_k\}\cup\cN(c_k)\subseteq\cS_{\rm end}$, $F'$ and $F$ agree on it, so $I_k'=I_k$; hence the stop/move decision at scan $k$ is the same (it depends on whether $I_k$ is empty and on $k$ versus $\ell$). If a move occurs, $\rhd_{F'}$ and $\rhd_F$ agree on $I_k\subseteq\cS_{\rm end}$ (they are defined by the values of $F'$, resp.\ $F$, there and the fixed $\tau$), so $c'_{k+1}=c_{k+1}$. Hence $m'=m$, the outcomes coincide, $\cC'=\cC$, $\cS'_{\rm end}=\cS_{\rm end}$, and $R'=\best'(\cS_{\rm end})=R$ because $\succ'$ and $\succ$ agree on $\cS_{\rm end}$. The flag $\llim$ coincides. For $\cert$: if $R\notin\cC$ both flags are $0$; if $R\in\cC$ then $\cN(R)\subseteq\cS_{\rm end}$ (it was scanned), so the condition ``no $N\in\cN(R)$ with $F'(N)>F'(R)$'' coincides with the one for $F$.

Consequently the value of $\cert$ is the same for every assignment of scores to the models of $\Omega\setminus\cS_{\rm end}$. In particular, if $\Omega\setminus\cS_{\rm end}\ne\emptyset$, the choice $F'=F$ on $\cS_{\rm end}$ and $F'=F(R)+1$ elsewhere leaves $\cert$ unchanged while $R$ is not a global maximiser of $F'$. So, as a property of the procedure, $\cert=1$ certifies nothing beyond the facts in (i-a) about scored models; when some model is unscored it is not a certificate of global maximality. (If $\cS_{\rm end}=\Omega$, then $R=\best(\Omega)$ is a global maximiser of $F$ by \eqref{T20:eq:bestmax}, whatever the flag.) (No claim is made here about which score functions arise from data.)

\paragraph{(i-c).}
Every model scored by scan $k$ is in $\cN(c_k)\subseteq\cS_{\rm end}$, whether or not it is $c_{k+1}$. As $R=\best(\cS_{\rm end})$, $R\succeq N$ and, by \eqref{T20:eq:bestmax}, $F(R)\ge F(N)$ for each such $N$. In the method the reported HPM $\hat S$ is the running best over every model scored in the run, and nothing is scored after the final local search (STATEMENT.md \S2); so for the final search $\hat S=\best(\cS_{\rm end})=R$.

\paragraph{(i-d).}
On \LIMIT, $R$ maximises $F$ on $\cS_{\rm end}$ by \eqref{T20:eq:bestmax}; Lemma~\ref{T20:lem:inv}(e) gives $R\in\cN(c_\ell)\setminus\cS_{<\ell}$ and $F(R)>F(c_\ell)$. Thus $R$ was first scored by scan $\ell$. Every centre $c_k$ of this search lies in $\cS_{<k}\subseteq\cS_{<\ell}$ (Lemma~\ref{T20:lem:inv}(a)), and every centre of an earlier local search was scored in that search, hence lies in $\cS^-\subseteq\cS_{<\ell}$ (the scored set is cumulative: it records every model ever scored in the run). So $R$ was never a scan centre of any local search of the run up to this one; in particular $R\notin\cC$ and $\cert=0$. For the final search, $\hat S=R$ by (i-c).

\paragraph{(i-e).}
Assume exactly one model of $\cS_{\rm end}$ attains $\max_{\cS_{\rm end}}F$. On \NOIMP, $c_m\in\cS_{\rm end}$ attains it (Lemma~\ref{T20:lem:inv}(d)) and so does $R$ by \eqref{T20:eq:bestmax}; by uniqueness $R=c_m$, and (i-a) gives $\cert=1$. On \LIMIT, $\cert=0$ by (i-d). Since $\llim=1$ exactly on \LIMIT, $\cert=1\iff\NOIMP\iff\llim=0$. \qed

\section{Clause (ii), deterministic part (ii-A)}

Assume $E_n$: $S_0\in\Omega$ and $F(S_0)>F(S)$ for all $S\in\Omega\setminus\{S_0\}$. Consider any local search of the run and any scan $k$ of it with $c_k\in\cN(S_0)$.

\emph{1.} By Lemma~\ref{T20:lem:sym}, $S_0\in\cN(c_k)$, and by EXEC scan $k$ scores it.

\emph{2.} $c_k\ne S_0$ (as $S_0\notin\cN(S_0)$), so $F(S_0)>F(c_k)$ and $S_0\in I_k$. For $N\in I_k\setminus\{S_0\}$, $F(S_0)>F(N)$, hence $S_0\rhd_F N$. So $S_0$ is the $\rhd_F$-maximum of $I_k$, and if $k<\ell$ the procedure sets $c_{k+1}=S_0$.

\emph{3.} If $k<\ell$, the procedure continues to scan $k+1$ with centre $S_0$ (after any move the loop performs another scan: a regular one if $k+1<\ell$, the one further scan if $k+1=\ell$). Every $N\in\cN(S_0)$ differs from $S_0$, so $F(N)<F(S_0)$ and $I_{k+1}=\emptyset$: the search stops with \NOIMP, $m=k+1$, $c_m=S_0$. Since $S_0\in\cS_{\rm end}$ and $F(S_0)>F(T)$ for every other $T$, $R=\best(\cS_{\rm end})=S_0=c_m$, and (i-a) gives $\cert=1$.

\emph{4.} If $k=\ell$, then $I_\ell\ni S_0$ is nonempty and the search stops with \LIMIT; $R=S_0$ as in 3; $\cert=0$ by (i-d) (and $\llim=1$).

From the moment scan $k$ scores $S_0$, it is in the scored set and strictly beats every other model, so the running HPM, being the $\succ$-maximum of a scored set containing $S_0$, equals $S_0$ for the rest of the run; for the final search, $\hat S=S_0$. \qed

\begin{T20corollary}\label{T20:cor:already}
On $E_n$, if $S_0\in\cS^-$ at the start of some local search, that search has $c_0=S_0$, stops after scan $0$ with \NOIMP, and returns $R=S_0$ with $\cert=1$.
\end{T20corollary}
\begin{proof}
$c_0=\best(\cS^-)=S_0$ since $S_0$ has the strictly largest score; $I_0=\emptyset$ as in step 3; $R=S_0=c_0$ and (i-a) applies.
\end{proof}

\section{Clause (ii), margin and statistical parts}

From here on $F(S)=\log\varpi(S)+\log m(S)$ is the method's score, with $\varpi$, $m$, $G_S$, $z_S$, $H_S$, $D_S$ and $Q_S=y^\top y-n\,z_S^\top H_S^{-1}z_S$ exactly as defined in STATEMENT.md \S1, $\pi(S)=e^{F(S)}/\sum_{T\in\Omega}e^{F(T)}$ for $S\in\Omega$ and $\pi(S):=0$ for $S\notin\Omega$.

\begin{T20lemma}\label{T20:lem:finite}
For every $S\in\Omega$, $0<\varpi(S)m(S)<\infty$. Hence $F$ is real-valued on $\Omega$ and $0<\pi(S)<1$ for every $S\in\Omega$.
\end{T20lemma}
\begin{proof}
$\varpi(S)=w(|S|)/W_K$ with $w(k)=\tfrac12/\binom{2p-1}{k}>0$ for $0\le k\le K\le p\le 2p-1$, and $W_K$ a finite positive sum. In $m(S)=\frac{(2b_0)^{a_0}\Gamma(\nu)}{\pi^{n/2}\Gamma(a_0)}D_S^{-1/2}(2b_0+Q_S)^{-\nu}$ the constant is positive and finite ($a_0=b_0=0.01$, $\nu=a_0+n/2$). For $S=\emptyset$ the conventions are $D_\emptyset=1$ and $Q_\emptyset=y^\top y$. For $k=|S|\ge1$: $G_S=X_S^\top X_S\succeq0$, so $G_S+I$ and $H_S=(n+1)G_S+I$ are positive definite and $D_S=\det H_S/\det(G_S+I)\in(0,\infty)$. Since $k\le K\le n$, take a thin SVD $X_S=U\Sigma V^\top$ with $U\in\mathbb{R}^{n\times k}$, $U^\top U=I_k$, $V$ orthogonal and $\Sigma=\mathrm{diag}(\sigma_1,\dots,\sigma_k)\ge0$. Then $H_S^{-1}=V((n+1)\Sigma^2+I)^{-1}V^\top$ and
$X_SH_S^{-1}X_S^\top=U\,\mathrm{diag}\big(\sigma_i^2/((n+1)\sigma_i^2+1)\big)U^\top\preceq\tfrac1{n+1}UU^\top\preceq\tfrac1{n+1}I_n$.
With $z_S=X_S^\top y$, $n z_S^\top H_S^{-1}z_S=n\,y^\top X_SH_S^{-1}X_S^\top y\le\frac{n}{n+1}y^\top y$, so, by the definition $Q_S=y^\top y-n z_S^\top H_S^{-1}z_S$, $Q_S\ge y^\top y/(n+1)\ge0$ and $2b_0+Q_S>0$. Hence $m(S)\in(0,\infty)$. Finally $\emptyset$ and $\{1\}$ belong to $\Omega$ ($K\ge1$), so $|\Omega|\ge2$, and each $\pi(S)$ is a positive term of a sum of at least two positive terms divided by that sum: $0<\pi(S)<1$.
\end{proof}

\paragraph{(ii-B).}
Let $S_0\in\Omega$. By Lemma~\ref{T20:lem:finite}, $\lambda_n=\log[\pi(S_0)/(1-\pi(S_0))]$ is a real number and $\Omega\setminus\{S_0\}$ is nonempty and finite; choose $S^*\in\Omega\setminus\{S_0\}$ with $F(S^*)=\max_{S\ne S_0}F(S)$, so $\Delta_n=F(S_0)-F(S^*)$. Then
\[
1-\pi(S_0)=\sum_{S\in\Omega\setminus\{S_0\}}\pi(S)\ \ge\ \pi(S^*),\qquad\text{so}\qquad
\frac{\pi(S_0)}{1-\pi(S_0)}\le\frac{\pi(S_0)}{\pi(S^*)}=e^{F(S_0)-F(S^*)}=e^{\Delta_n}.
\]
Taking logarithms, $\lambda_n\le\Delta_n$. If $\pi(S_0)>\frac12$ then $S_0\in\Omega$ (by the convention $\pi=0$ off $\Omega$) and $\lambda_n>0$, so $\Delta_n>0$, i.e.\ $F(S_0)>F(S)$ for every $S\in\Omega\setminus\{S_0\}$: $E_n$ holds. \qed

\paragraph{(ii-C).}
We use the standard definitions: real random variables $Y_n$ are $O_{P_0}(1)$ if for every $\varepsilon>0$ there are $M<\infty$ and $N$ with $P_0(|Y_n|>M)<\varepsilon$ for all $n\ge N$; and $X_n=O_{P_0}(a_n)$ means $X_n/a_n=O_{P_0}(1)$. (If one prefers ``for all $n$'' in place of ``for all $n\ge N$'', the two agree for a.s.\ finite variables, by enlarging $M$ to cover the finitely many $n<N$.) Note $r_n=\sqrt{\log n/n}\to0$ and, for $n\ge3$, $\log\log n$ is defined and $\log r_n=\tfrac12\log\log n-\tfrac12\log n$.

Assume H03 and fix $\varepsilon>0$. There are $C\in(0,\infty)$ and $N_1$ (enlarging $C$ if necessary, we may take $C\ge1$, so $\log(2C)>0$) with $P_0(1-\pi(S_0)>Cr_n)<\varepsilon$ for $n\ge N_1$. Let $N_2\ge\max(n_0,N_1,3)$ be such that $Cr_n<\frac12$ for $n\ge N_2$. Fix $n\ge N_2$ and work on the event $B_n=\{1-\pi(S_0)\le Cr_n\}$, which has $P_0(B_n)>1-\varepsilon$. On $B_n$, $\pi(S_0)>\frac12$; we also know $1-\pi(S_0)>0$ (Lemma~\ref{T20:lem:finite}, $S_0\in\Omega$ as $n\ge n_0$).

\emph{1.} On $B_n$, $\pi(S_0)>\frac12$, so $E_n$ holds by (ii-B). Hence $P_0(E_n^c)\le P_0(B_n^c)<\varepsilon$ for all $n\ge N_2$. As $\varepsilon$ was arbitrary, $P_0(E_n)\to1$.

\emph{2.} On $B_n$ (with $\pi(S_0)>\frac12$),
\[
\lambda_n=\log\pi(S_0)-\log(1-\pi(S_0))\ \ge\ -\log2-\log C-\log r_n=\tfrac12\log n-\tfrac12\log\log n-\log(2C).
\]
Let $Y_n=(\tfrac12\log n-\tfrac12\log\log n-\lambda_n)^+$, a finite random variable for $n\ge n_0$. On $B_n$, $Y_n\le\log(2C)=:M$. Hence $P_0(Y_n>M)<\varepsilon$ for all $n\ge N_2$, i.e.\ $Y_n=O_{P_0}(1)$. Together with (ii-B), for $n\ge n_0$,
$\Delta_n\ge\lambda_n\ge\tfrac12\log n-\tfrac12\log\log n-Y_n$ with $Y_n=O_{P_0}(1)$.

\emph{3.} Assume EX-A and EXEC, as in (ii-A). The conclusions (ii-A)1--4 were derived on $E_n$ for every local search, every scan with a centre in $\cN(S_0)$, every input $\cS^-$ and every $\ell$; they use no property of the algorithm's randomness. Hence, for any joint law of data and algorithm randomness whose data marginal is $P_0$, the set $A_n$ of outcomes on which these conclusions hold for every scan of every local search of the run contains the event $E_n$, which depends on the data only. Therefore $P(A_n)\ge P_0(E_n)\to1$ (read $P(A_n)$ as an inner probability if the measurability of $A_n$ is not established; $E_n$ is measurable, being defined by finitely many strict inequalities between continuous functions of $y$). \qed
\mdproofend
\mdsec{(d) Review record}
\mdsubsec{FINAL\_REPORT.md: statement changes}
\begin{mdblock}
\mdhead{2}{2. Changes to the statement}

No hypothesis was changed or added. The step 1 outcome was (a), correct as stated.

To make the informal clauses precise, STATEMENT.md:

\begin{mdbul}
\item formalises the method's §6.10 search as a short algorithm (§2), with the source text quoted verbatim in §0;
\item uses T03 only through the explicit hypothesis H03;
\item names the "scans execute" assumption EXEC;
\item records seven readings in §5. Two examples:

\begin{mdbul}
\item "returned model" means the running HPM at the end of the search;
\item "can be the reported HPM without a certificate" is proved as a conditional ("on LIMIT, …"). No claim is made that LIMIT actually happens for real data.
\end{mdbul}
\end{mdbul}

Reviewer-driven text corrections (change log in STATEMENT.md §6):

\begin{mdbul}
\item v1 → v2:

\begin{mdbul}
\item an explanatory remark in §5 was wrong and was fixed (A1-1);
\item the score formulas were written out in full (A1-2).
\end{mdbul}
\item v2 → v3:

\begin{mdbul}
\item (i-b)'s "does not certify global maximality" is now qualified by "when some model is unscored". A reviewer noted it was false if every model had been scored;
\item other edits are clarifications only.
\end{mdbul}
\item v3 → v3.1: readings and setting wording only.
\end{mdbul}
\end{mdblock}

\mdsubsec{FINAL\_REPORT.md: review history}
\begin{mdblock}
\mdhead{2}{3. Review history}

\textbf{Stage A: lead + Astra \#1} (GPT-6 Astra, high, session 01a1641e-caed-7313-8c23-8d32e3a413f3).

\begin{mdbul}
\item Round 1 on v1: \emph{not verified}.

\begin{mdbul}
\item One "fatal" finding, in an explanatory reading note rather than a claim. The reviewer said it does not affect §§3–4.
\item One gap: Q\_S was not defined in the packet.
\item One exposition finding.
\item All three were accepted and fixed.
\end{mdbul}
\item Round 2 on v2: \emph{verified}. Astra agreed with both answers.
\end{mdbul}

\textbf{Stage B: two fresh independent reviews of v2, started together.}

\begin{mdbul}
\item \textbf{Astra \#2} (new session 01a16424-4661-7c72-9494-a9613c43944d): \emph{verified}. No fatal or gap findings; one cosmetic remark.
\item \textbf{Opus 5.5 xhigh} (native Agent, read-only Plan type, Claude Work account): \emph{verified}. Ten exposition or scope notes (F1–F10); seven applied in v3.
\item \textbf{Ruling round on v3 and the claims log:} both \emph{verified}, agreeing with every answer.

\begin{mdbul}
\item Both raised the same remaining exposition point, fixed in STATEMENT v3.1.
\item One harmless note (N3, on measurability) was left unapplied so the reviewed proof text stays final.
\end{mdbul}
\end{mdbul}

No disputes; the Astra judge was not used. The full log, with every fatal and gap finding verbatim and the lead's answers, is in {\mdtt CLAIMS\_\allowbreak{}OF\_\allowbreak{}FATAL\_\allowbreak{}FLAWS.\allowbreak{}md}. All replies are in {\mdtt review/\allowbreak{}}.

\textbf{Numerical sanity check.} A Sonnet assistant ran a single-seed local check ({\mdtt helpers/\allowbreak{}numcheck/\allowbreak{}NUMCHECK.\allowbreak{}md}).

\begin{mdbul}
\item Scope: 72 small problems (p ≤ 8), where the whole model space can be scored, and 8,640 random search runs.
\item Result: 0 failures across all checks of (i)–(ii), including tie canaries.
\item This is a sanity check, not part of the proof.
\end{mdbul}
\end{mdblock}

\mdsubsec{CLAIMS\_OF\_FATAL\_FLAWS.md (every entry, verbatim)}
\begin{mdblock}
\mdhead{1}{CLAIMS\_OF\_FATAL\_FLAWS.md — T20 (T15) (i)–(ii)}

Every fatal or gap finding, verbatim, with the lead's answer.

\mdhead{2}{Stage A — Astra \#1 (session 01a1641e-caed-7313-8c23-8d32e3a413f3), round 1, on PROOF\_v1 + STATEMENT v1}

\mdhead{3}{A1-1 (fatal, per the reviewer: "it does not invalidate the formal conclusions in §§3–4")}

\begin{mdquote}
\begin{mdnum}
\item[\mdlbl{1.}] \textbf{Location:} STATEMENT.md §5, item 1, beginning “If one instead reads it as the last centre…”.

\textbf{Severity:} fatal — this explanatory claim is false; it does not invalidate the formal conclusions in §§3–4.

\textbf{Claim:} The running HPM \(R\) and the last centre \(c_m\) do not differ only in the presence of score ties. Unique attainment of the largest scored value does not ensure that they coincide.

\textbf{Evidence:} Take \(p=1\), \(n=3\), hence \(K=1\), with 
\[
\mathcal S^-=\{\varnothing\},\qquad \ell=0,\qquad
F(\varnothing)=0,\quad F(\{1\})=1.
\]
 Scan \(0\) scores \(\{1\}\) and stops with LIMIT. The last centre is \(c_m=\varnothing\), whereas \(R=\{1\}\). There are no score ties, and the scored maximum is unique. More generally, (i-d) itself establishes \(F(R)>F(c_m)\) on every LIMIT outcome. Thus the explanatory sentence contradicts a formal claim that the proof correctly establishes.
\end{mdnum}
\end{mdquote}

\textbf{Lead's answer: accepted and fixed.} The sentence in the readings section was wrong: on LIMIT, R ≠ c\_m always. It was an explanatory remark, not one of the claims of §§3–4, and the proof never used it. STATEMENT.md §5 item 1 now reads: on LIMIT, R ≠ c\_m but the flag is 0 under both readings (c\_m has a strictly improving neighbour); on NOIMP, R = c\_m unless another scored model has exactly the score F(c\_m) and precedes c\_m in the HPM order; so the two readings give different flag values only on NOIMP with such an exact tie. (Check: under the "last centre" reading the flag is 1 iff NOIMP; under the adopted reading it is 1 iff NOIMP and R = c\_m, by (i-a).)

\mdhead{3}{A1-2 (gap)}

\begin{mdquote}
\begin{mdnum}
\item[\mdlbl{2.}] \textbf{Location:} PROOF\_v1.tex, “Clause (ii), margin and statistical parts,” lemma labelled {\mdtt lem:finite}, proof for nonempty \(S\).

\textbf{Severity:} gap.

\textbf{Claim:} The asserted lower bound on \(Q_S\) is not justified by definitions supplied in the permitted files.

\textbf{Evidence:} The proof establishes 
\[
n z_S^\top H_S^{-1}z_S\le \frac{n}{n+1}y^\top y
\]
 and immediately concludes 
\[
Q_S\ge \frac{y^\top y}{n+1}\ge0.
\]
 However, neither supplied text defines \(Q_S\) for nonempty \(S\), or states an identity connecting it to the preceding quadratic form. Only \(Q_\varnothing\) is defined. Consequently this implication, and therefore this particular derivation of positivity of the displayed marginal likelihood, cannot be verified from the permitted files. The identification of those detailed formulas with the method’s marginal likelihood also relies on an external source whose contents are not supplied.

This is an auxiliary derivational gap: the stated setting already assumes a real-valued score, and the subsequent posterior-odds calculation is correct.
\end{mdnum}
\end{mdquote}

\textbf{Lead's answer: accepted and fixed.} The packet lacked the definition. STATEMENT.md §1 now writes out the method's score in full, copied from 00\_COMMON §1 and METHODOLOGY\_v3 §5: G\_S, z\_S, H\_S, D\_S, Q\_S = yᵀy − n z\_SᵀH\_S⁻¹z\_S, ν, m(S), w, W\_K, ϖ. It also says the proof uses only these displayed formulas. PROOF\_v2's Lemma (finite) now cites that definition explicitly at the step Q\_S ≥ yᵀy/(n+1). The identification with a Bayesian marginal likelihood is not used anywhere in the proof; the formulas are the definition of the score.

Astra \#1 round 2 (same session, on PROOF\_v2 + STATEMENT v2 + this log): agrees with both answers ("I agree with the answer and confirm that the fix works" for A1-1 and A1-2); no new findings; verdict \textbf{verified} (review/astra1\_r2.md). End of Stage A; proof candidate = PROOF\_v2.

\mdhead{2}{Stage B — Astra \#2 (fresh session 01a16424-4661-7c72-9494-a9613c43944d) and Opus 5.5 xhigh reviewer, round 1, on PROOF\_v2 + STATEMENT v2 (change log removed from the packet copy)}

\textbf{No fatal or gap findings.} Astra \#2: verdict verified; one remark (the inner-probability parenthetical in (ii-C)3 is unnecessary; kept, harmless). Opus: verdict verified; findings F1–F10, all exposition or scope (review/astra2\_r1.md, review/opus\_r1.md).

Lead's handling of the exposition/scope findings (unified revision PROOF\_v3 + STATEMENT v3; no hypothesis or claim content changed):

\begin{mdbul}
\item F1 accepted: (ii-C)3 and proof step 3 now name EX-A and EXEC.
\item F2 accepted: (i-b) now says the F′-run uses ≻ and the move order computed from F′ with the same fixed tie-break rule (the proof already did this via τ).
\item F3 accepted: (i-b) "does not certify global maximality" now qualified by Ω ∖ 𝒮\_end ≠ ∅, with the trivial 𝒮\_end = Ω case noted (statement and proof).
\item F4 accepted: (i-a) now states that every equal-score neighbour N of a certified R satisfies R ≻ N.
\item F5 accepted: proof §1 shows 𝒩(c) ≠ ∅.
\item F6 accepted: proof (ii-C) takes C ≥ 1 (so log(2C) > 0).
\item F7 accepted: "from the moment scan k scores S₀" (statement and proof).
\item F8 noted (scope): (i-c) holds by construction of the formal model; the link to the method is now quoted verbatim in STATEMENT §0 (METHODOLOGY\_v3 §6.10 and §6.1).
\item F9 partly addressed: the source's E\_n definition and the method's §6.10 text are now quoted verbatim in STATEMENT §0; H03 remains an explicit hypothesis (T03 is a conjecture being proved in a sibling thread), as the statement says.
\item F10 accepted: §5 item 5 now mentions the tie-only route (NOIMP with an equal-score neighbour preceding c\_m), covered by (i-a).
\end{mdbul}

\mdhead{3}{Stage B ruling round (on PROOF\_v3 + STATEMENT v3 + this log)}

\begin{mdbul}
\item \textbf{Opus reviewer (r2): verified.} Agrees with every answer in the log; each change works. New findings N1–N3, exposition only (review/opus\_r2.md).
\item \textbf{Astra \#2 (r2, same session): verified.} Agrees with every answer (table of rulings); one exposition finding, the same point as Opus N1 (review/astra2\_r2.md).
\item Lead: N1 (= Astra \#2's point) and N2 accepted, fixed in STATEMENT v3.1 (§5 item 5 and §1 text only; proof unchanged). N3 (measurability of E\_n via a.s. continuity of standardisation) noted; harmless, both reviewers agree; not applied, to keep the reviewed proof text final.
\item No fatal or gap finding is outstanding. No dispute; Astra judge not used.
\end{mdbul}
\end{mdblock}

\chapter{T22 (T16), parts (i), (ii), (iv): aggregate start bounds and endpoint propagation}
\chaptersource{Final proof: \texttt{proofs/\allowbreak{}T22\_\allowbreak{}T16\_\allowbreak{}start\_\allowbreak{}aggregation/\allowbreak{}proof/\allowbreak{}PROOF\_v3.tex}. Sources of parts (a), (b), (d): \texttt{theorems/\allowbreak{}T22\_\allowbreak{}T16\_\allowbreak{}start\_\allowbreak{}aggregation.md}, \texttt{STATEMENT.md}, \texttt{FINAL\_REPORT.md}, \texttt{CLAIMS\_OF\_FATAL\_FLAWS.md}.}
\mdsec{(a) Original conjecture text (theorem file, verbatim)}
\begin{mdblock}
\mdhead{1}{T22 (T16): Aggregate start bounds and conditional concentration for the actual blocks}

\textbf{Status:} conjectured, not proven. Methodology SMC-BVS v3.0 (2026-10-19), SHA-256 93043758058a6d422cca6ca4b6be8f3a9e8149b6f6de2e1bab1ce31c44b8019c. Read with {\mdtt 00\_\allowbreak{}COMMON.\allowbreak{}md}.

\textbf{Category / rank:} deterministic start analysis (i)–(ii), (iv): \textbf{core} (lemma); statistical concentration (iii), (v): \textbf{secondary} (theorem).

\mdhead{2}{Definitions used}

Tables r\_b, γ₀, π₀, w(k) (00\_COMMON §2); AGG's q\_n, e\_n for a feasible reference S° (s = |S°|, U\_b* = S°∩C\_b) with the stated conventions; u\_n := (p−s)⁻¹Σ\_\{l∉S°\} max\_\{S°∩C\_\{b(l)\}⊆T⊆C\_\{b(l)\}∖l\} r\_\{b(l)\}(T∪l)/r\_\{b(l)\}(T); uncapped weight γ₀\textasciicircum{}\{unc\}(S) = w(|S|)∏\_b r\_b(S∩C\_b); δ\_n\textasciicircum{}\{comp\} = sup\_\{S⊉S°\}γ₀\textasciicircum{}\{unc\}(S)/γ₀\textasciicircum{}\{unc\}(S∪S°); r\_k (T10); R\_i = 1/q\_i(S₀) − 1 for i ∈ \{0,1\}; leakage quantities E\_\{T,l\}, v\_T; inventory codes VIS, NL\_avg(κ), EXP, BO.

\mdhead{2}{Assumptions}

EX-A for (i), (ii), (iv); G, R, within-block SE, VIS, NL\_avg(κ) and EXP for (iii); T03's hypotheses for the full-endpoint conclusion in (iv); BO, SE, bounded signal energy and BM-∞ for (v).

\mdhead{2}{Precise statement}

(i) [deterministic aggregation] If q\_n < 1: π₀(S°) ≥ (1−q\_n)\textasciicircum{}\{s+1\}/(1+e\_n)\textasciicircum{}s and 1 − π₀(S°) ≤ (s+1)q\_n/(1−q\_n) + se\_n. Moreover q\_n ≤ u\_n, so u\_n < 1 permits its substitution. A further candidate bound is π₀(S°) ≥ (1−u\_n)\textasciicircum{}\{s+1\}/[1 + (2\textasciicircum{}s − 1)δ\_n\textasciicircum{}\{comp\}]. Conventions: empty maxima are 0; q = u = 0 if s = p; e = δ = 0 if s = 0. (ii) [reference-to-truth odds] π₀(S₀) = π₀(S°)·[w(s₀)/w(s)]·∏\_b r\_b(S₀∩C\_b)/r\_b(S°∩C\_b) exactly; if S° ⊆ S₀, successive feasible additions give π₀(S₀)/π₀(S°) ≥ ∏\_\{j=0\}\textasciicircum{}\{s₀−s−1\}r\_\{s+j\}. Normalised reference mass remains essential. (iii) [statistical, with EXP] For fixed κ > 0: q\_n ≤ u\_n = O\_\{P₀\}(2\textasciicircum{}B√(log n)n\textasciicircum{}\{−κ\}) when EXP's remainder is O\_\{P₀\}(1) (exponent form 2\textasciicircum{}B√(log n)n\textasciicircum{}\{−κ+o\_\{P₀\}(1)\} with an o\_\{P₀\}(log n) remainder); EXP's visibility lower bound with M\_n\textasciicircum{}V → ∞ and B ≤ 16 gives e\_n = o\_\{P₀\}(n\textasciicircum{}\{−m\}) for every fixed m; hence 1 − π₀(S₀) = O\_\{P₀\}(2\textasciicircum{}B√(log n)n\textasciicircum{}\{−κ\} + e\_n). Exact zero leakage permits κ = ½; NL\_η supplies the admissible κ of the inventory. Neither substitutes for EXP. (iv) [endpoint propagation] For 0 ≤ φ ≤ 1: 1/q\_φ(S₀) − 1 ≤ R₀\textasciicircum{}\{1−φ\}R₁\textasciicircum{}φ ≤ max(R₀, R₁); hence inf\_\{φ∈[0,1]\}q\_φ(S₀) ≥ min\{π₀(S₀), π(S₀)\}. With T03 and (iii) this gives uniform concentration on [0,1]. (v) [BO] Under BO, completing the true variables in a block removes its signal leakage into the remaining coordinates of that block; with SE, bounded signal energy and BM-∞ the projection argument supplies VIS; the score conclusion still requires EXP. No over-tempered clause is retained. Start concentration implies neither schedule completion nor evidence accuracy, nor that every initial particle equals S₀.

\mdhead{2}{Dependencies}

Builds on: T01, T03, T10. Used by: T21, T26, T27.

\mdhead{2}{Numerical evidence (provisional)}

provisional evidence (Sol check J5): 468 table cases (adversarial and seeded), 376 applicable bounds, 0 failures, including empty, full and binding-cap references; actual partition at n = 64, p = 9, B = 3; observed EXP remainder range [−0.10, 0.17] on 20 additions (cannot establish its stochastic order); two-branch bound verified algebraically. Files: judge/sol1/J5\_COMPOSITE\_START.md, J5\_tables.json.

\mdhead{2}{Proof-route hint (optional)}

Dominate the completed-table generating function coefficientwise by exp\{z(p−s)q\_n\}; prior ratio w(s+j)/w(s) ≤ (s+j)!/(s!(p−s)\textasciicircum{}j); Σ\_jC(s+j, j)q\_n\textasciicircum{}j = (1−q\_n)\textasciicircum{}\{−(s+1)\}; charge incomplete configurations to their completions (≤ e\_n per missing coordinate, (1+e\_n)\textasciicircum{}s); the uncapped sum dominates the capped normaliser. (iii): truncate |Z| ≤ h\_n uniformly over the ≤ p2\textasciicircum{}B table comparisons; a\_l ≤ C(n+1)\textasciicircum{}\{−1/2\}exp\{(δ+σ*Z)²/(2v\_T)\} by EXP; the two-branch bound gives E[·; |Z| ≤ h\_n] ≤ (2h\_n/√(2π))n\textasciicircum{}\{−1/2\}e\textasciicircum{}\{E\_\{T,l\}\}; linearity and Markov, no independence. (iv): Hölder on the total odds Σ\_\{S≠S₀\}[γ₀(S)/γ₀(S₀)]\textasciicircum{}\{1−φ\}[γ₁(S)/γ₁(S₀)]\textasciicircum{}φ ≤ R₀\textasciicircum{}\{1−φ\}R₁\textasciicircum{}φ. Hardest step: deriving EXP from primitive assumptions (the uniform block-score expansion over O(p2\textasciicircum{}B) entries). \emph{When proof.md is run on this file, its lead considers the statement both with and without these hints; the hints never reach the reviewers.}
\end{mdblock}

\mdsec{(b) STATEMENT.md (verbatim)}
\begin{mdblock}
\mdhead{1}{STATEMENT — T22 (T16), parts (i), (ii), (iv)}

Source: {\mdtt Theory\_\allowbreak{}architecture\_\allowbreak{}v3/\allowbreak{}theorems/\allowbreak{}T22\_\allowbreak{}T16\_\allowbreak{}start\_\allowbreak{}aggregation.\allowbreak{}md} (SMC-BVS v3.0, method SHA-256 93043758…8019c). Only parts (i), (ii), (iv) are stated here; parts (iii) and (v) are out of scope for this run. Parts (i), (ii), (iv-a), (iv-b) are deterministic: they hold for every realisation of the data; (iv-c) is a conditional probabilistic consequence of (iv-b). Exact arithmetic is assumed (idealisation EX-A): π₀, q\_φ and all tables are the exact mathematical objects defined below.

Status: statement approved with corrections by Astra (step 1b), 2026-10-22; see "Changes to the statement" at the end.

\mdhead{2}{1. Setting and definitions}

\textbf{Data and model space.} Integers n ≥ 3 and p ≥ 1. X ∈ ℝ\textasciicircum{}\{n×p\} (columns x\_1, …, x\_p) and y ∈ ℝⁿ are fixed (in the method they are the standardised design and response, ‖x\_j‖² = ‖y‖² = n − 1; the statements below use only that they are real). K = min(p, ⌊n/2⌋). Ω = \{S ⊆ [p] : |S| ≤ K\} ("feasible models"). For S ⊆ [p] (any size), with X\_S the columns in S:

\begin{mdbul}
\item G\_S = X\_SᵀX\_S, z\_S = X\_Sᵀy, H\_S = (n+1)G\_S + I\_\{|S|\}, D\_S = det H\_S / det(G\_S + I\_\{|S|\}),
\item t\_S = n z\_SᵀH\_S⁻¹z\_S, Q\_S = yᵀy − t\_S, ν = a₀ + n/2 with a₀ = b₀ = 0.01,
\item m(S) = c\_n · D\_S\textasciicircum{}\{−1/2\}(2b₀ + Q\_S)\textasciicircum{}\{−ν\}, where c\_n = (2b₀)\textasciicircum{}\{a₀\}Γ(ν)/(π\textasciicircum{}\{n/2\}Γ(a₀)) > 0 (for S = ∅: D\_∅ = 1, Q\_∅ = yᵀy),
\item BF(S) = m(S)/m(∅).
\item Size prior: w(k) = ½ / C(2p−1, k) for 0 ≤ k ≤ p; W\_K = Σ\_\{k≤K\} C(p,k)w(k); ϖ(S) = w(|S|)/W\_K for S ∈ Ω.
\item Posterior: π(S) = ϖ(S)m(S)/Z for S ∈ Ω, Z = Σ\_\{S∈Ω\} ϖ(S)m(S).
\end{mdbul}

\textbf{Blocks and tables.} C\_1, …, C\_J is a partition of [p] into nonempty blocks with |C\_b| ≤ B ≤ K (in the method: the deterministic complete-linkage/cut/next-fit partition with B = min(⌊√n⌋, ⌊√p⌋, 16); parts (i) and (iv) use only that it is a partition and that tables are positive). b(l) denotes the block containing l. Tables: r\_b(U) = BF(U) for U ⊆ C\_b (so r\_b(∅) = 1 and r\_b(U) > 0).

\textbf{Start law and bridge.} For S ∈ Ω: γ₀(S) = ϖ(S)∏\_b r\_b(S ∩ C\_b), Z̄₀ = Σ\_\{S∈Ω\} γ₀(S), π₀(S) = γ₀(S)/Z̄₀. Uncapped weight, for every S ⊆ [p]: γ₀\textasciicircum{}\{unc\}(S) = w(|S|)∏\_b r\_b(S ∩ C\_b). Discrepancy d(S) = log BF(S) − Σ\_b log r\_b(S ∩ C\_b). For φ ∈ [0,1] and S ∈ Ω: γ\_φ(S) = γ₀(S)e\textasciicircum{}\{φ d(S)\}, Z̄\_φ = Σ\_\{S∈Ω\} γ\_φ(S), q\_φ(S) = γ\_φ(S)/Z̄\_φ. (Thus q₀ = π₀; that q₁ = π is a fact to be checked in the proof.) Masses of sets outside Ω are 0.

\textbf{Reference quantities (AGG).} Fix a reference S° ∈ Ω, s = |S°|, U\_b* = S° ∩ C\_b.

\begin{mdbul}
\item For each block b and integer j ≥ 0: a\_\{bj\} = Σ\_\{U : U\_b* ⊆ U ⊆ C\_b, |U ∖ U\_b*| = j\} r\_b(U)/r\_b(U\_b*) (so a\_\{b0\} = 1, and a\_\{bj\} = 0 for j > |C\_b ∖ U\_b*|).
\item D\_b = max\_\{1 ≤ j ≤ |C\_b∖U\_b*|\} (j! a\_\{bj\})\textasciicircum{}\{1/j\}, with D\_b = 0 when C\_b = U\_b* (empty maximum).
\item q = (p−s)⁻¹ Σ\_b D\_b if s < p; q = 0 if s = p.
\item e = 2p · max\{ [r\_b(U)/r\_b(U ∪ U\_b*)]\textasciicircum{}\{1/a\} : b, U ⊆ C\_b, a := |U\_b* ∖ U| > 0 \}; e = 0 when the set is empty (in particular when s = 0).
\item u = (p−s)⁻¹ Σ\_\{l ∉ S°\} max\{ r\_\{b(l)\}(T ∪ \{l\})/r\_\{b(l)\}(T) : U*\_\{b(l)\} ⊆ T ⊆ C\_\{b(l)\} ∖ \{l\} \} if s < p; u = 0 if s = p.
\item δ\textasciicircum{}\{comp\} = sup\{ γ₀\textasciicircum{}\{unc\}(S)/γ₀\textasciicircum{}\{unc\}(S ∪ S°) : S ∈ Ω, S ⊉ S° \}, with δ\textasciicircum{}\{comp\} = 0 when s = 0. (The proof may use this sup over feasible S; the bound then also holds for the larger sup over all S ⊆ [p] with S ⊉ S°.)
\end{mdbul}

\textbf{Addition floors.} For 0 ≤ k ≤ p − 1: r\_k = (k+1)/((2p−k−1)√(n+1)).

\textbf{Truth.} S₀ ⊆ [p] is a fixed set (the support of the true coefficient vector), s₀ = |S₀|. For i ∈ \{0, 1\}, when S₀ ∈ Ω: R\_i = 1/q\_i(S₀) − 1.

\mdhead{2}{2. Claims}

\textbf{(i) Deterministic aggregation.} For every partition and positive tables as above and every S° ∈ Ω:

\begin{mdbul}
\item (i-a) If q < 1, then π₀(S°) ≥ (1−q)\textasciicircum{}\{s+1\}/(1+e)\textasciicircum{}s and 1 − π₀(S°) ≤ (s+1)q/(1−q) + s·e.
\item (i-b) q ≤ u. Hence, if u < 1, both bounds of (i-a) hold with u in place of q.
\item (i-c) If u < 1, then π₀(S°) ≥ (1−u)\textasciicircum{}\{s+1\}/[1 + (2\textasciicircum{}s − 1)δ\textasciicircum{}\{comp\}].
\end{mdbul}

\textbf{(ii) Reference-to-truth odds.} Assume S₀ ∈ Ω (i.e. s₀ ≤ K).

\begin{mdbul}
\item (ii-a) π₀(S₀) = π₀(S°) · [w(s₀)/w(s)] · ∏\_b r\_b(S₀ ∩ C\_b)/r\_b(S° ∩ C\_b), exactly.
\item (ii-b) If moreover S° ⊆ S₀, then π₀(S₀)/π₀(S°) ≥ ∏\_\{j=0\}\textasciicircum{}\{s₀−s−1\} r\_\{s+j\} (empty product = 1).
\item Remark (interpretive, not a separate mathematical claim): (ii) controls π₀(S₀) only relative to π₀(S°); a lower bound on the normalised reference mass π₀(S°) (e.g. from (i)) is still required.
\end{mdbul}

\textbf{(iv) Endpoint propagation.} Assume S₀ ∈ Ω.

\begin{mdbul}
\item (iv-a) For every φ ∈ [0,1]: 1/q\_φ(S₀) − 1 ≤ R₀\textasciicircum{}\{1−φ\}R₁\textasciicircum{}\{φ\} ≤ max(R₀, R₁), with the convention 0⁰ = 1.
\item (iv-b) Hence inf\_\{φ∈[0,1]\} q\_φ(S₀) ≥ min\{π₀(S₀), π(S₀)\}; equivalently sup\_\{φ∈[0,1]\} [1 − q\_φ(S₀)] ≤ max\{1 − π₀(S₀), 1 − π(S₀)\}.
\item (iv-c) Consequence (conditional). Consider a sequence of data sets indexed by n → ∞ under a probability P₀. If P₀(S₀ ∈ Ω) → 1, 1 − π(S₀) → 0 in P₀-probability, and 1 − π₀(S₀) → 0 in P₀-probability, then sup\_\{φ∈[0,1]\} [1 − q\_φ(S₀)] → 0 in P₀-probability (uniform concentration of the bridge on [0,1]). The two convergence hypotheses are the conclusions of T03 (posterior concentration) and of part (iii) of T22; here they are explicit hypotheses, not results of this run.
\end{mdbul}

\mdhead{2}{Changes to the statement}

Statement check by Astra (session 01a16418-acfd-7cb3-8783-b91d3752db30, reply in review/astra\_01\_statement\_check.log): \textbf{APPROVED WITH CORRECTIONS}, 2026-10-22. The corrections are applied below.

Approved changes:

\begin{mdbul}
\item C1 (added hypothesis): (ii) and (iv) assume S₀ ∈ Ω (s₀ ≤ K). Reason: π₀, q\_φ and π are distributions on Ω; if s₀ > K then π₀(S₀) = 0 while the right side of (ii-a) is positive, so (ii-a) is false (instance: n = 3, p = 2, x₁ = x₂ = y = (−1, 0, 1)ᵀ, singleton blocks, S₀ = \{1,2\}, S° = ∅), and R\_i is undefined in (iv). The method's regime R assumes "eventually s₀ ≤ K", so nothing is lost asymptotically. Astra: necessary and minimal (not needed for (iv-b), which is trivial outside Ω).
\item C6 (generalisation, recorded at Astra's request): the source's setting is the standardised design and response and the method's complete-linkage/cut/next-fit partition. STATEMENT.md quantifies over all real X, y and all partitions of [p] into blocks with |C\_b| ≤ K (and, for (i), all positive tables). This is a strictly more general statement; the source setting is a special case.
\item C7 (scope reading): (i-c) is stated under u < 1, the condition the source attaches to substituting u ("u\_n < 1 permits its substitution"); the source's "candidate bound" sentence does not repeat it.
\end{mdbul}

Accepted readings:

\begin{mdbul}
\item C2: the final sentence of the source's (iv), "With T03 and (iii) this gives uniform concentration on [0,1]", is stated as (iv-c) with the conclusions of T03 and (iii) as explicit hypotheses. (Astra notes that the hypothesis P₀(S₀ ∈ Ω) → 1 is redundant given either concentration hypothesis; it is kept, harmlessly.) (iv-c) is a conditional probabilistic consequence of the deterministic (iv-b); "all three parts are deterministic" refers to (i), (ii), (iv-a), (iv-b).
\item C3: δ\textasciicircum{}\{comp\}'s supremum is taken over feasible S ⊉ S° (a smaller sup, hence a stronger result that implies the reading over all S ⊆ [p]).
\item C4: "Normalised reference mass remains essential" is recorded as an interpretive remark, not a claim.
\item C5 (provenance, corrected): r\_k = (k+1)/((2p−k−1)√(n+1)) is copied verbatim from {\mdtt theorems/\allowbreak{}T10\_\allowbreak{}T11\_\allowbreak{}addition\_\allowbreak{}hazard.\allowbreak{}md}, "Definitions used" ("r\_k := (k+1)/((2p−k−1)√(n+1))"); the lead checked this. Astra objected only because T10 was not in its packet; the formula is used here as an explicit definition. All other definitions are transcriptions of the source and 00\_COMMON.md, confirmed by Astra.
\item No EXP / EXP′ hypothesis is needed for (i), (ii), (iv) (confirmed by Astra).
\end{mdbul}
\end{mdblock}

\mdsec{(c) Final proof, verbatim: body of PROOF\_v3.tex}
\noindent{\small\itshape The proof below is the body of \texttt{PROOF\_v3.tex} from \texttt{\textbackslash begin\{document\}} to \texttt{\textbackslash end\{document\}} without the title block. Its preamble definitions are local to this chapter; labels carry the prefix \texttt{T22:}.}\par\medskip
\mdproofbegin
% ---- preamble of PROOF_v3.tex (theorem environments carry the prefix T22) ----
\newtheorem{T22lemma}{Lemma}
\newtheorem{T22proposition}{Proposition}
\theoremstyle{remark}
\newtheorem{T22remark}{Remark}
\newcommand{\unc}{\mathrm{unc}}
\newcommand{\comp}{\mathrm{comp}}
\newcommand{\gu}{\gamma_0^{\unc}}
% ---- body of PROOF_v3.tex ----
\noindent All notation is that of \texttt{STATEMENT.md}. Every claim except (iv-c) is deterministic (and (iv-c) is a direct consequence of (iv-b)): the data $(X,y)$, the partition $C_1,\dots,C_J$ and the reference $S^\circ\in\Omega$ are fixed, and all quantities are exact. We write $S+l:=S\cup\{l\}$, $\varrho:=n+1$, and, for $S\subseteq[p]$, $M_S:=X_SX_S^\top\in\mathbb{R}^{n\times n}$ (with $M_\emptyset=0$). Throughout, $A\succeq B$ is the Loewner order. The letter $e$ always denotes the AGG quantity of the statement (never Euler's number), $e^{k}$ is its $k$-th power, and exponentials are written $\exp(\cdot)$. We use $x^0:=1$ for every $x\ge0$ (in particular $e^0=1$, and $D_b^0=1$ when $D_b=0$).

\section{Preliminaries on the score}

\begin{T22lemma}[Positivity and a closed form]\label{T22:lem:pos}
For every $S\subseteq[p]$,
\[
t_S=\frac{n}{\varrho}\,y^\top\bigl(I_n-(I_n+\varrho M_S)^{-1}\bigr)y,\qquad D_S=\frac{\det(I_n+\varrho M_S)}{\det(I_n+M_S)} .
\]
Consequently $0\le t_S\le \tfrac{n}{n+1}y^\top y$, $Q_S\ge y^\top y/(n+1)\ge 0$, $D_S\ge 1$, and $m(S)\in(0,\infty)$. In particular every table entry $r_b(U)$, every $w(k)$ ($0\le k\le p$), every $\gamma_0(S)$, $\gu(S)$, $\gamma_\varphi(S)$ is strictly positive and finite.
\end{T22lemma}

\begin{proof}
For $S=\emptyset$ both formulas give $t_\emptyset=0$, $D_\emptyset=1$, matching the definitions. Let $S\neq\emptyset$, $k=|S|$, $A=X_S$. The matrices $I_k+\varrho A^\top A$ and $I_n+\varrho AA^\top$ are positive definite, and $(I_n+\varrho AA^\top)A=A(I_k+\varrho A^\top A)$, so $A(I_k+\varrho A^\top A)^{-1}=(I_n+\varrho AA^\top)^{-1}A$. Hence, since $H_S=\varrho A^\top A+I_k$,
\[
AH_S^{-1}A^\top=(I_n+\varrho AA^\top)^{-1}AA^\top=\tfrac1\varrho(I_n+\varrho AA^\top)^{-1}\bigl((I_n+\varrho AA^\top)-I_n\bigr)=\tfrac1\varrho\bigl(I_n-(I_n+\varrho M_S)^{-1}\bigr),
\]
and $t_S=n\,y^\top AH_S^{-1}A^\top y$ gives the formula for $t_S$. Since $0\prec(I_n+\varrho M_S)^{-1}\preceq I_n$, the matrix $I_n-(I_n+\varrho M_S)^{-1}$ has eigenvalues in $[0,1)$, so $0\le t_S\le \frac{n}{\varrho} y^\top y$ and $Q_S=y^\top y-t_S\ge y^\top y/\varrho\ge 0$. By Sylvester's determinant identity $\det(I_k+\tau A^\top A)=\det(I_n+\tau AA^\top)$ for every real $\tau$, applied with $\tau=\varrho$ and $\tau=1$, we get the formula for $D_S$; if $\lambda_1,\dots,\lambda_n\ge0$ are the eigenvalues of $M_S\succeq0$, then $D_S=\prod_{i=1}^n(1+\varrho\lambda_i)/(1+\lambda_i)\ge1$ because $\varrho\ge1$. Finally $2b_0+Q_S\ge 2b_0>0$, $c_n>0$, so $m(S)=c_nD_S^{-1/2}(2b_0+Q_S)^{-\nu}\in(0,\infty)$. Positivity of the remaining objects follows: $w(k)=\tfrac12/\binom{2p-1}{k}>0$ for $0\le k\le p\le 2p-1$, $W_K\ge w(0)>0$, $r_b(U)=m(U)/m(\emptyset)>0$, and $\exp(\varphi d(S))\in(0,\infty)$.
\end{proof}

\begin{T22lemma}[One addition]\label{T22:lem:add}
For every $S\subseteq[p]$ and $l\notin S$: $Q_{S+l}\le Q_S$, $D_{S+l}/D_S\le n+1$, and hence
\[
\frac{m(S+l)}{m(S)}\ \ge\ (n+1)^{-1/2}.
\]
\end{T22lemma}

\begin{proof}
The matrix $M_S$ does not depend on the order of the columns, and $M_{S+l}=M_S+x_lx_l^\top\succeq M_S$. Then $I_n+\varrho M_{S+l}\succeq I_n+\varrho M_S\succ0$, and inversion reverses the Loewner order on positive definite matrices, so $(I_n+\varrho M_{S+l})^{-1}\preceq (I_n+\varrho M_S)^{-1}$; by Lemma~\ref{T22:lem:pos}, $t_{S+l}\ge t_S$, i.e.\ $Q_{S+l}\le Q_S$.

For the determinant ratio, the matrix determinant lemma $\det(P+xx^\top)=\det P\,(1+x^\top P^{-1}x)$ ($P\succ0$) and Lemma~\ref{T22:lem:pos} give
\[
\frac{D_{S+l}}{D_S}=\frac{\det(I_n+\varrho M_S+\varrho\,x_lx_l^\top)/\det(I_n+\varrho M_S)}{\det(I_n+M_S+x_lx_l^\top)/\det(I_n+M_S)}=\frac{1+\varrho\,\xi}{1+\zeta},
\]
with $\xi=x_l^\top(I_n+\varrho M_S)^{-1}x_l$ and $\zeta=x_l^\top(I_n+M_S)^{-1}x_l$. Since $I_n+\varrho M_S\succeq I_n+M_S\succ0$, $0\le\xi\le\zeta$. Thus $1+\varrho\xi\le 1+\varrho\zeta\le \varrho(1+\zeta)$ (as $\varrho\ge1$), i.e.\ $D_{S+l}/D_S\le \varrho$.

Therefore, with $\nu>0$ and $2b_0+Q_{S+l}\le 2b_0+Q_S$ (both positive),
\[
\frac{m(S+l)}{m(S)}=\Bigl(\frac{D_S}{D_{S+l}}\Bigr)^{1/2}\Bigl(\frac{2b_0+Q_S}{2b_0+Q_{S+l}}\Bigr)^{\nu}\ \ge\ \varrho^{-1/2}\cdot 1 .\qedhere
\]
\end{proof}

\begin{T22lemma}[Prior ratios]\label{T22:lem:prior}
For $0\le k\le p-1$: $w(k+1)/w(k)=(k+1)/(2p-k-1)$, and $w(k)/w(k+1)\le 2p-1<2p$. If $0\le s<p$ and $1\le j\le p-s$, then
\[
\frac{w(s+j)}{w(s)}\ \le\ \frac{(s+j)!}{s!\,(p-s)^j}.
\]
\end{T22lemma}

\begin{proof}
With $N=2p-1$, $w(k+1)/w(k)=\binom{N}{k}/\binom{N}{k+1}=(k+1)/(N-k)=(k+1)/(2p-k-1)$; its inverse is $(2p-k-1)/(k+1)\le 2p-1$. For the product, $w(s+j)/w(s)=\prod_{i=0}^{j-1}(s+i+1)/(2p-s-i-1)$, and for $0\le i\le j-1\le p-s-1$ we have $2p-s-i-1\ge 2p-s-(p-s-1)-1=p\ge p-s>0$. Hence the product is at most $\prod_{i=0}^{j-1}(s+i+1)/(p-s)=(s+j)!/(s!(p-s)^j)$.
\end{proof}

\section{Part (i)}

Because $\gamma_0(S)=\gu(S)/W_K$ for $S\in\Omega$ and $S^\circ\in\Omega$,
\begin{equation}\label{T22:eq:pi0ratio}
\pi_0(S^\circ)=\frac{\gu(S^\circ)}{\sum_{S\in\Omega}\gu(S)} .
\end{equation}

\paragraph{Completion bijection.} Let $\mathcal T=\{T: S^\circ\subseteq T\subseteq[p]\}$. The map $\Phi(S)=(S\cup S^\circ,\ S^\circ\setminus S)$ is a bijection from the power set of $[p]$ onto $\mathcal T\times 2^{S^\circ}$, with inverse $(T,M)\mapsto T\setminus M$. Indeed, $(S\cup S^\circ)\setminus(S^\circ\setminus S)=S$; and for $(T,M)\in\mathcal T\times2^{S^\circ}$, $(T\setminus M)\cup S^\circ=T$ (since $M\subseteq S^\circ\subseteq T$) and $S^\circ\setminus(T\setminus M)=S^\circ\cap M=M$. Note $|T\setminus M|=|T|-|M|$.

\begin{T22lemma}[Completed sum]\label{T22:lem:completed}
If $q<1$, then $\sum_{T\in\mathcal T}\gu(T)\le (1-q)^{-(s+1)}\,\gu(S^\circ)$.
\end{T22lemma}

\begin{proof}
If $s=p$, then $\mathcal T=\{S^\circ\}$ and the claim is $\gu(S^\circ)\le\gu(S^\circ)$ (here $q=0$). Let $s<p$. For $T\in\mathcal T$ put $V_b=(T\cap C_b)\setminus U_b^*\subseteq C_b\setminus U_b^*$. Since the blocks partition $[p]$, $T\mapsto(V_b)_b$ is a bijection from $\mathcal T$ onto $\prod_b 2^{C_b\setminus U_b^*}$, with $T\cap C_b=U_b^*\cup V_b$ and $|T|=s+\sum_b|V_b|$. Hence
\[
\frac{\gu(T)}{\gu(S^\circ)}=\frac{w(|T|)}{w(s)}\prod_b\frac{r_b(U_b^*\cup V_b)}{r_b(U_b^*)} ,
\]
and grouping the $T$ by $j=|T|-s$ and expanding the product of sums,
\[
\sum_{T\in\mathcal T}\frac{\gu(T)}{\gu(S^\circ)}=\sum_{j=0}^{p-s}\frac{w(s+j)}{w(s)}\,c_j,\qquad c_j:=\sum_{(j_b)_b:\ \sum_b j_b=j}\ \prod_b a_{b j_b},
\]
where $j_b$ ranges over nonnegative integers ($a_{bj_b}=0$ when $j_b>|C_b\setminus U_b^*|$).

\emph{Coefficient bound.} For every $b$ and $j\ge0$, $a_{bj}\le D_b^{\,j}/j!$: for $j=0$ both sides equal $1$ ($D_b^0=1$); for $1\le j\le|C_b\setminus U_b^*|$ it is the definition of $D_b$ as a maximum; for larger $j$, $a_{bj}=0$. By the multinomial theorem,
\[
c_j\le\sum_{\sum_b j_b=j}\ \prod_b\frac{D_b^{\,j_b}}{j_b!}=\frac{(\sum_b D_b)^j}{j!}=\frac{((p-s)q)^j}{j!}.
\]

\emph{Prior bound.} By Lemma~\ref{T22:lem:prior} (and trivially for $j=0$), $w(s+j)/w(s)\le (s+j)!/(s!(p-s)^j)$ for $0\le j\le p-s$. Therefore
\[
\sum_{T\in\mathcal T}\frac{\gu(T)}{\gu(S^\circ)}\le\sum_{j=0}^{p-s}\frac{(s+j)!}{s!\,(p-s)^j}\cdot\frac{(p-s)^jq^j}{j!}=\sum_{j=0}^{p-s}\binom{s+j}{j}q^j\le\sum_{j=0}^{\infty}\binom{s+j}{j}q^j=(1-q)^{-(s+1)},
\]
using $q\ge0$ and the negative binomial series $\sum_{j\ge0}\binom{s+j}{j}x^j=(1-x)^{-(s+1)}$ for $|x|<1$ (obtained by differentiating the geometric series $s$ times and dividing by $s!$).
\end{proof}

\begin{T22lemma}[Charging incomplete sets to completions]\label{T22:lem:charge}
For every $S\subseteq[p]$ with $\Phi(S)=(T,M)$: $\gu(S)\le e^{|M|}\,\gu(T)$.
\end{T22lemma}

\begin{proof}
If $M=\emptyset$ then $S=T$. Let $M\neq\emptyset$; then $s\ge1$, the set over which $e$ is a maximum is nonempty (it contains $(b,\emptyset)$ for any $b$ with $U_b^*\neq\emptyset$), and $\epsilon:=e/(2p)>0$ satisfies $r_b(U)/r_b(U\cup U_b^*)\le\epsilon^{a}$ whenever $U\subseteq C_b$ and $a=|U_b^*\setminus U|>0$. Fix $b$ and put $U=S\cap C_b$. Since $T=S\cup S^\circ$, $T\cap C_b=U\cup U_b^*$, and $a_b:=|U_b^*\setminus U|=|M\cap C_b|$. If $a_b>0$, $r_b(S\cap C_b)/r_b(T\cap C_b)\le\epsilon^{a_b}$; if $a_b=0$ then $U\supseteq U_b^*$ and the ratio is $1=\epsilon^0$. Multiplying over $b$ (the $M\cap C_b$ partition $M$): $\prod_b r_b(S\cap C_b)/r_b(T\cap C_b)\le\epsilon^{|M|}$. For the prior, $|S|=|T|-|M|$ and, by Lemma~\ref{T22:lem:prior} applied $|M|$ times (each intermediate size $k$ satisfies $k\le p-1$), $w(|S|)/w(|T|)\le(2p)^{|M|}$. Hence $\gu(S)/\gu(T)\le(2p\epsilon)^{|M|}=e^{|M|}$.
\end{proof}

\begin{T22proposition}[(i-a)]
If $q<1$, then $\pi_0(S^\circ)\ge(1-q)^{s+1}/(1+e)^s$ and $1-\pi_0(S^\circ)\le (s+1)q/(1-q)+s\,e$.
\end{T22proposition}

\begin{proof}
All terms are positive, so dropping the feasibility restriction and using the bijection $\Phi$, then Lemmas~\ref{T22:lem:charge} and \ref{T22:lem:completed},
\[
\begin{aligned}
\sum_{S\in\Omega}\gu(S)&\le\sum_{S\subseteq[p]}\gu(S)=\sum_{T\in\mathcal T}\sum_{M\subseteq S^\circ}\gu(T\setminus M)\\
&\le\sum_{T\in\mathcal T}\gu(T)\sum_{M\subseteq S^\circ}e^{|M|}=(1+e)^s\sum_{T\in\mathcal T}\gu(T)\\
&\le\frac{(1+e)^s}{(1-q)^{s+1}}\,\gu(S^\circ).
\end{aligned}
\]
With \eqref{T22:eq:pi0ratio} this is the first bound. (The completions $T$ may have $|T|>K$; their uncapped weights are used only in this upper bound.) For the second, put $x=(1-q)^{s+1}\in(0,1]$ and $z=(1+e)^{-s}\in(0,1]$. Then $1-\pi_0(S^\circ)\le 1-xz=(1-x)+x(1-z)\le(1-x)+(1-z)$. Bernoulli's inequality gives $1-x\le(s+1)q\le (s+1)q/(1-q)$. The function $g(x)=1-(1+x)^{-s}$ satisfies $g(0)=0$ and $g'(x)=s(1+x)^{-s-1}\le s$ for $x\ge0$, so $1-z=g(e)\le se$.
\end{proof}

\begin{T22proposition}[(i-b)]
$q\le u$. Hence, if $u<1$, both bounds of (i-a) hold with $u$ in place of $q$.
\end{T22proposition}

\begin{proof}
If $s=p$, $q=u=0$. Let $s<p$. For $l\notin S^\circ$, with $b=b(l)$, put $\mu_l=\max\{r_b(T+l)/r_b(T):U_b^*\subseteq T\subseteq C_b\setminus\{l\}\}$; the index set is nonempty because $l\notin U_b^*$ (so $T=U_b^*$ is admissible), and $u=(p-s)^{-1}\sum_{l\notin S^\circ}\mu_l$. Fix $b$ and $V\subseteq C_b\setminus U_b^*$ with $|V|=j\ge1$, and list $V=\{v_1,\dots,v_j\}$ in any order. With $T_i=U_b^*\cup\{v_1,\dots,v_i\}$ we have $U_b^*\subseteq T_{i-1}\subseteq C_b\setminus\{v_i\}$, so telescoping,
\[
\frac{r_b(U_b^*\cup V)}{r_b(U_b^*)}=\prod_{i=1}^{j}\frac{r_b(T_{i-1}+v_i)}{r_b(T_{i-1})}\le\prod_{l\in V}\mu_l .
\]
Summing over $V$, $j!\,a_{bj}\le j!\sum_{V\subseteq C_b\setminus U_b^*,\,|V|=j}\prod_{l\in V}\mu_l\le\bigl(\sum_{l\in C_b\setminus U_b^*}\mu_l\bigr)^j$, because expanding the $j$-th power of a sum of nonnegative terms produces each product over a $j$-subset exactly $j!$ times (once per ordering), plus further nonnegative terms. Taking $j$-th roots and the maximum over $1\le j\le|C_b\setminus U_b^*|$, $D_b\le\sum_{l\in C_b\setminus U_b^*}\mu_l$ (also when $C_b=U_b^*$, both sides being $0$). The sets $C_b\setminus U_b^*$ partition $[p]\setminus S^\circ$, so $(p-s)q=\sum_bD_b\le\sum_{l\notin S^\circ}\mu_l=(p-s)u$.

If $u<1$ then $q\le u<1$, so (i-a) applies; since $x\mapsto(1-x)^{s+1}$ is nonincreasing and $x\mapsto(s+1)x/(1-x)$ is nondecreasing on $[0,1)$, the bounds with $q$ imply those with $u$.
\end{proof}

\begin{T22proposition}[(i-c)]
If $u<1$, then $\pi_0(S^\circ)\ge(1-u)^{s+1}/[1+(2^s-1)\delta^{\comp}]$.
\end{T22proposition}

\begin{proof}
By the bijection $\Phi$, $\sum_{S\in\Omega}\gu(S)=\sum_{T\in\mathcal T}\sum_{M\subseteq S^\circ}\mathbf 1\{T\setminus M\in\Omega\}\gu(T\setminus M)$. Fix $T\in\mathcal T$. The term $M=\emptyset$ is at most $\gu(T)$. For $\emptyset\ne M\subseteq S^\circ$ with $S:=T\setminus M\in\Omega$: $S\not\supseteq S^\circ$ (it misses $M\neq\emptyset$) and $S\cup S^\circ=T$, so $\gu(S)\le\delta^{\comp}\gu(T)$ by definition of $\delta^{\comp}$ (a maximum over a finite set). There are at most $2^s-1$ such $M$ (the nonempty subsets of $S^\circ$, some of which may be excluded by feasibility; none when $s=0$). Hence $\sum_{S\in\Omega}\gu(S)\le[1+(2^s-1)\delta^{\comp}]\sum_{T\in\mathcal T}\gu(T)$. Since $q\le u<1$ (part (i-b)), Lemma~\ref{T22:lem:completed} and monotonicity give $\sum_{T}\gu(T)\le(1-q)^{-(s+1)}\gu(S^\circ)\le(1-u)^{-(s+1)}\gu(S^\circ)$, and \eqref{T22:eq:pi0ratio} concludes. The argument used the supremum over feasible $S$ only, so the bound also holds with the (larger) supremum over all $S\subseteq[p]$, $S\not\supseteq S^\circ$.
\end{proof}

\section{Part (ii)}

Assume $S_0\in\Omega$.

\begin{T22proposition}[(ii-a)]
$\pi_0(S_0)=\pi_0(S^\circ)\,[w(s_0)/w(s)]\,\prod_b r_b(S_0\cap C_b)/r_b(S^\circ\cap C_b)$.
\end{T22proposition}

\begin{proof}
For $S\in\Omega$, $\pi_0(S)=w(|S|)\prod_b r_b(S\cap C_b)/(W_K\bar Z_0)$, and $\pi_0(S^\circ)>0$ (Lemma~\ref{T22:lem:pos}). Divide the expression for $S_0$ by that for $S^\circ$.
\end{proof}

\begin{T22lemma}[One feasible addition under $\pi_0$]\label{T22:lem:pi0add}
If $S\in\Omega$, $k=|S|<K$ and $l\notin S$, then $S+l\in\Omega$ and $\pi_0(S+l)\ge r_k\,\pi_0(S)$.
\end{T22lemma}

\begin{proof}
$|S+l|=k+1\le K$. Let $b=b(l)$ and $U=S\cap C_b$; only block $b$ changes, $(S+l)\cap C_b=U+l$. By (ii-a)'s computation and Lemmas~\ref{T22:lem:prior} and \ref{T22:lem:add} (note $k\le K-1\le p-1$),
\[
\frac{\pi_0(S+l)}{\pi_0(S)}=\frac{w(k+1)}{w(k)}\cdot\frac{r_b(U+l)}{r_b(U)}=\frac{k+1}{2p-k-1}\cdot\frac{m(U+l)}{m(U)}\ge\frac{k+1}{(2p-k-1)\sqrt{n+1}}=r_k.\qedhere
\]
\end{proof}

\begin{T22proposition}[(ii-b)]
If $S^\circ\subseteq S_0$, then $\pi_0(S_0)/\pi_0(S^\circ)\ge\prod_{j=0}^{s_0-s-1}r_{s+j}$.
\end{T22proposition}

\begin{proof}
If $s_0=s$ then $S_0=S^\circ$ and both sides equal $1$. Otherwise list $S_0\setminus S^\circ=\{l_1,\dots,l_{s_0-s}\}$ and put $S^{(i)}=S^\circ\cup\{l_1,\dots,l_i\}$, so $|S^{(i)}|=s+i$. For $0\le i\le s_0-s-1$, $|S^{(i)}|=s+i<s_0\le K$, so Lemma~\ref{T22:lem:pi0add} gives $\pi_0(S^{(i+1)})\ge r_{s+i}\pi_0(S^{(i)})$. Multiplying these $s_0-s$ inequalities (all terms positive) gives the claim, since $S^{(0)}=S^\circ$ and $S^{(s_0-s)}=S_0$.
\end{proof}

\begin{T22remark}
(ii) bounds $\pi_0(S_0)$ only relative to $\pi_0(S^\circ)$; an absolute bound needs a lower bound on the normalised mass $\pi_0(S^\circ)$, such as part (i).
\end{T22remark}

\section{Part (iv)}

Assume $S_0\in\Omega$.

\begin{T22lemma}[Endpoints and geometric interpolation]\label{T22:lem:bridge}
$q_1=\pi$, and for every $\varphi\in[0,1]$ and $S\in\Omega$, $\gamma_\varphi(S)=\gamma_0(S)^{1-\varphi}\gamma_1(S)^{\varphi}$.
\end{T22lemma}

\begin{proof}
By Lemma~\ref{T22:lem:pos}, $d(S)$ is finite and
\[
\gamma_1(S)=\gamma_0(S)\exp(d(S))=\varpi(S)\prod_br_b(S\cap C_b)\cdot\frac{\mathrm{BF}(S)}{\prod_br_b(S\cap C_b)}=\frac{\varpi(S)m(S)}{m(\emptyset)} .
\] Hence $q_1(S)=\varpi(S)m(S)/\sum_{S'\in\Omega}\varpi(S')m(S')=\pi(S)$. Also $\gamma_0\exp(\varphi d)=\gamma_0^{1-\varphi}(\gamma_0\exp(d))^{\varphi}$.
\end{proof}

\begin{T22proposition}[(iv-a), (iv-b)]
For every $\varphi\in[0,1]$, $1/q_\varphi(S_0)-1\le R_0^{1-\varphi}R_1^{\varphi}\le\max(R_0,R_1)$ (with $0^0=1$). Hence $\inf_{\varphi\in[0,1]}q_\varphi(S_0)\ge\min\{\pi_0(S_0),\pi(S_0)\}$, equivalently $\sup_{\varphi\in[0,1]}[1-q_\varphi(S_0)]\le\max\{1-\pi_0(S_0),1-\pi(S_0)\}$.
\end{T22proposition}

\begin{proof}
Since $\gamma_\varphi>0$ on $\Omega$, $1/q_\varphi(S_0)-1=\sum_{S\in\Omega\setminus\{S_0\}}\gamma_\varphi(S)/\gamma_\varphi(S_0)$. With $\alpha_S=\gamma_0(S)/\gamma_0(S_0)$ and $\beta_S=\gamma_1(S)/\gamma_1(S_0)$ (positive), Lemma~\ref{T22:lem:bridge} gives $\gamma_\varphi(S)/\gamma_\varphi(S_0)=\alpha_S^{1-\varphi}\beta_S^{\varphi}$, and $R_0=\sum_{S\ne S_0}\alpha_S$, $R_1=\sum_{S\neq S_0}\beta_S$ (sums over $\Omega\setminus\{S_0\}$), using $q_0=\pi_0$, $q_1=\pi$.

Since $n\ge3$ and $p\ge1$, $K\ge1$, so $\emptyset$ and $\{1\}$ both lie in $\Omega$ and $\Omega\setminus\{S_0\}\neq\emptyset$; hence $R_0,R_1>0$ (the convention $0^0=1$ in the statement is never needed). For $\varphi\in\{0,1\}$ the first inequality is an equality. For $\varphi\in(0,1)$, Hölder's inequality with conjugate exponents $1/(1-\varphi)$ and $1/\varphi$ gives
\[
\sum_{S\ne S_0}\alpha_S^{1-\varphi}\beta_S^{\varphi}\le\Bigl(\sum_{S\neq S_0}\alpha_S\Bigr)^{1-\varphi}\Bigl(\sum_{S\ne S_0}\beta_S\Bigr)^{\varphi}=R_0^{1-\varphi}R_1^{\varphi}.
\]
Finally $R_0^{1-\varphi}R_1^{\varphi}\le\max(R_0,R_1)^{1-\varphi}\max(R_0,R_1)^{\varphi}=\max(R_0,R_1)$.

For (iv-b): $q_\varphi(S_0)=1/(1+[1/q_\varphi(S_0)-1])\ge1/(1+\max(R_0,R_1))=\min\{1/(1+R_0),1/(1+R_1)\}=\min\{q_0(S_0),q_1(S_0)\}=\min\{\pi_0(S_0),\pi(S_0)\}$, for every $\varphi\in[0,1]$. Taking the infimum gives the claim; the supremum form is the same statement.
\end{proof}

\begin{T22proposition}[(iv-c)]
Along a sequence of data sets with $P_0(S_0\in\Omega)\to1$, $1-\pi(S_0)\to0$ and $1-\pi_0(S_0)\to0$ in $P_0$-probability, $\sup_{\varphi\in[0,1]}[1-q_\varphi(S_0)]\to0$ in $P_0$-probability.
\end{T22proposition}

\begin{proof}
(Masses of $S_0$ are $0$ when $S_0\notin\Omega$.) Let $\Sigma_n=\sup_{\varphi\in[0,1]}[1-q_\varphi(S_0)]$. On $E_n=\{S_0\in\Omega\}$, (iv-b) gives $\Sigma_n\le\max\{1-\pi_0(S_0),1-\pi(S_0)\}\le[1-\pi_0(S_0)]+[1-\pi(S_0)]$. For $\eta>0$,
\[
P_0(\Sigma_n>\eta)\le P_0(E_n^c)+P_0\bigl(1-\pi_0(S_0)>\eta/2\bigr)+P_0\bigl(1-\pi(S_0)>\eta/2\bigr)\to0 .
\]
Measurability: $\varphi\mapsto q_\varphi(S_0)$ is continuous on $[0,1]$ (a ratio of finite sums of positive functions $\gamma_0(S)\exp(\varphi d(S))$), so $\Sigma_n$ equals the supremum over rational $\varphi$ and is measurable whenever each $q_\varphi(S_0)$ is; in any case the displayed bound holds with outer probability.
\end{proof}

\begin{T22remark}[Scope]
No assumption on the noise, the design, or the score expansion (EXP or EXP$'$) was used: (i), (ii) and (iv) are identities and inequalities valid for every data set, every partition with positive tables (for (ii-b): tables equal to Bayes factors), and every feasible reference. Statistical content enters only through (iv-c)'s hypotheses.
\end{T22remark}
\mdproofend
\mdsec{(d) Review record}
\mdsubsec{FINAL\_REPORT.md: statement changes}
\begin{mdblock}
\mdhead{2}{Changes to the statement}

All changes were approved by Astra \#1 as statement checker (step 1b), with the verdict "APPROVED WITH CORRECTIONS". Details are in STATEMENT.md, "Changes to the statement".

\begin{mdbul}
\item \textbf{C1, added hypothesis:} (ii) and (iv) assume S₀ is a feasible model, |S₀| ≤ K = min(p, ⌊n/2⌋).

\begin{mdbul}
\item Without it, the identity in (ii) is false: π₀(S₀) = 0 while the right side is positive.
\item Astra confirmed this with a concrete instance (n = 3, p = 2, S₀ = \{1,2\}).
\item Asymptotically nothing is lost: the regime R already assumes "eventually s₀ ≤ K".
\end{mdbul}
\item \textbf{C6, generalisation:} the proof holds for any real X, y and any block partition, and (i) holds for any positive tables. The method's setting is a special case.
\item \textbf{C7, scope:} (i-c) is stated under u < 1.
\item \textbf{Readings:}

\begin{mdbul}
\item "With T03 and (iii) …" in (iv) became a conditional corollary, (iv-c). The conclusions of T03 and of part (iii) enter only as explicit hypotheses; nothing about T03 or (iii) is proved here.
\item δ\textasciicircum{}comp's supremum is taken over feasible models. This makes the bound stronger and implies the other reading.
\item "Normalised reference mass remains essential" is treated as a remark, not a claim.
\end{mdbul}
\item \textbf{EXP / EXP′ is not needed} for (i), (ii) or (iv). Astra confirmed this. Part (iii), which needs EXP′, was left alone as instructed, and so was part (v).
\item \textbf{The proof-route hint} never entered STATEMENT.md or any packet.
\end{mdbul}
\end{mdblock}

\mdsubsec{FINAL\_REPORT.md: review history}
\begin{mdblock}
\mdhead{2}{Review history}

\mdhead{3}{Stage A: lead and Astra \#1 (session 01a16418-acfd-7cb3-8783-b91d3752db30)}

\begin{mdbul}
\item \textbf{Statement check:} approved with corrections ({\mdtt review/\allowbreak{}astra\_\allowbreak{}01\_\allowbreak{}statement\_\allowbreak{}check.\allowbreak{}log}).
\item \textbf{Review of PROOF\_v1:} VERDICT verified, one exposition finding ({\mdtt review/\allowbreak{}astra\_\allowbreak{}02\_\allowbreak{}review\_\allowbreak{}A1\_\allowbreak{}v1.\allowbreak{}log}).

\begin{mdbul}
\item The finding: "2\textasciicircum{}s − 1 such M" should read "at most 2\textasciicircum{}s − 1".
\item I accepted it and fixed it in PROOF\_v2, the proof candidate.
\end{mdbul}
\end{mdbul}

\mdhead{3}{Stage B: two independent reviews of PROOF\_v2, started together with fresh context and given only the packet}

\begin{mdbul}
\item \textbf{Astra \#2} (new session 01a1641f-ea40-7570-bde0-1de23647a8bc): no findings, VERDICT verified ({\mdtt review/\allowbreak{}astra2\_\allowbreak{}B\_\allowbreak{}review\_\allowbreak{}v2.\allowbreak{}log}).
\item \textbf{Opus 5.5, xhigh:} five exposition findings, VERDICT verified ({\mdtt review/\allowbreak{}opus\_\allowbreak{}B\_\allowbreak{}review\_\allowbreak{}v2.\allowbreak{}md}).

\begin{mdbul}
\item The findings were: a display overflowing the margin, an unstated step for D\_S ≥ 1, symbol clashes, an impossible case Ω = \{S₀\}, and the 0⁰ convention.
\item I accepted all five and fixed them in the unified PROOF\_v3 ({\mdtt review/\allowbreak{}CHANGES\_\allowbreak{}v2\_\allowbreak{}to\_\allowbreak{}v3.\allowbreak{}md}, {\mdtt review/\allowbreak{}REVIEW\_\allowbreak{}PROCESSING.\allowbreak{}md}).
\end{mdbul}
\item \textbf{Round 2:} PROOF\_v3 and the claims log went back to both reviewers.

\begin{mdbul}
\item Astra \#2: agrees with the log, VERDICT verified ({\mdtt review/\allowbreak{}astra2\_\allowbreak{}B\_\allowbreak{}round2\_\allowbreak{}v3.\allowbreak{}log}).
\item Opus: agrees with the log, confirms all five fixes, VERDICT verified ({\mdtt review/\allowbreak{}opus\_\allowbreak{}B\_\allowbreak{}round2\_\allowbreak{}v3.\allowbreak{}md}).
\item Opus raised one new cosmetic point, N1: the letter x is reused as a dummy variable in one paragraph of (i-a).
\item I accepted N1 but deferred it, so that the final file is exactly the version every reviewer verified.
\end{mdbul}
\item \textbf{CLAIMS\_OF\_FATAL\_FLAWS.md:} no fatal or gap finding in any review. Both Stage B reviewers agreed with the log.
\item \textbf{Disputes:} none, and the Astra judge was not used.
\end{mdbul}

\mdhead{3}{Numerical sanity check (not a review, not evidence of proof)}

\begin{mdbul}
\item A Sonnet assistant wrote an independent brute-force implementation from STATEMENT.md and ran it with one seed, locally ({\mdtt helpers/\allowbreak{}numcheck/\allowbreak{}REPORT.\allowbreak{}md}).
\item Scale: 10,000 small instances (p ≤ 9, the cap often binding, adversarial tables for (i)). In total about 1.1M checks of part (i), 3.1M of (ii-a) and 0.54M φ-points of (iv).
\item 0 failures.
\item All 4 planted false variants were detected.
\end{mdbul}
\end{mdblock}

\mdsubsec{CLAIMS\_OF\_FATAL\_FLAWS.md (every entry, verbatim)}
\begin{mdblock}
\mdhead{1}{CLAIMS\_OF\_FATAL\_FLAWS.md — T22 (T16) parts (i), (ii), (iv)}

Every fatal or gap finding from every review, verbatim, with the lead's answer.

\mdhead{2}{Stage A (Astra \#1, review of PROOF\_v1; review/astra\_02\_review\_A1\_v1.log)}

No fatal or gap findings (one exposition finding; VERDICT: verified).

\mdhead{2}{Stage B (Astra \#2 and Opus xhigh, reviews of PROOF\_v2; review/astra2\_B\_review\_v2.log, review/opus\_B\_review\_v2.md)}

No fatal or gap findings. Astra \#2: no findings, VERDICT: verified. Opus: five exposition findings, VERDICT: verified.

The log is therefore empty of fatal and gap claims. Exposition findings and their handling are in review/REVIEW\_PROCESSING.md.

\mdhead{2}{Stage B round 2 — rulings on this log (PROOF\_v3)}

\begin{mdbul}
\item Astra \#2 (review/astra2\_B\_round2\_v3.log): "I agree with the log concerning my review: Round 1 raised no fatal or gap finding." No findings in PROOF\_v3. VERDICT: verified.
\item Opus xhigh (review/opus\_B\_round2\_v3.md): "The log records no fatal or gap claims from any review. I agree with that ruling." All five exposition fixes confirmed. One new exposition finding (N1: dummy variable x reused in the (i-a) paragraph); accepted by the lead, deferred (cosmetic; PROOF\_v3 kept as the verified final). VERDICT: verified. No dispute; the Astra judge was not needed.
\end{mdbul}
\end{mdblock}

\chapter{T23(ii): pathwise Rao--Blackwell PIP bound and its basin consequence}
\chaptersource{Final proof: \texttt{proofs/\allowbreak{}T23\_\allowbreak{}T18\_\allowbreak{}basin\_\allowbreak{}prefix/\allowbreak{}proof/\allowbreak{}PROOF\_v5.tex}. Sources of parts (a), (b), (d): \texttt{theorems/\allowbreak{}T23\_\allowbreak{}T18\_\allowbreak{}basin\_\allowbreak{}prefix.md}, \texttt{STATEMENT.md}, \texttt{FINAL\_REPORT.md}, \texttt{CLAIMS\_OF\_FATAL\_FLAWS.md}.}
\mdsec{(a) Original conjecture text (theorem file, verbatim)}
\begin{mdblock}
\mdhead{1}{T23 (T18): Basin-to-output guarantee with recording-prefix control}

\textbf{Status:} conjectured, not proven. Methodology SMC-BVS v3.0 (2026-10-19), SHA-256 93043758058a6d422cca6ca4b6be8f3a9e8149b6f6de2e1bab1ce31c44b8019c. Read with {\mdtt 00\_\allowbreak{}COMMON.\allowbreak{}md}.

\textbf{Category / rank:} conditional algorithmic convergence; abstract theorem (i)–(iv): \textbf{core}; conditional instances (v): \textbf{secondary}.

\mdhead{2}{Definitions used}

Final-run terminal rounds (L burn-in + 4 measured) at φ = 1; f\_H(S₀ᶜ) = fraction of an island's particles outside S₀ after H terminal rounds; B\_\{i,h\} = indicator that particle i is outside S₀ at some point during measured round h (including its start); B̄\_r = (4N\_e)⁻¹Σ\_\{h=1\}\textasciicircum{}4Σ\_iB\_\{i,h\}; PIP̂\_\{rj\}, primary average; u\_π = δ₁/π(S₀); ρ\_* = min\{1,(B+3)u\_π\}; BASIN(𝓡, r₀, a, η, ρ, δ) (inventory); 𝔓 = \{S̃₀, S₀\}.

\mdhead{2}{Assumptions}

EX-A, EX-S, CAL, BASIN; T03's hypotheses supply statistical truth stability; the instances in (v) need their own exit, gain and initial-population hypotheses.

\mdhead{2}{Precise statement}

(i) For one island, after H terminal rounds: E\_𝓒 f\_H(S₀ᶜ) ≤ δ + (1−a)\textasciicircum{}\{⌊H/r₀⌋\} + H(η + ρ). Conditioning on the realised initial population replaces δ by its actual outside-basin fraction. (ii) [pathwise RB bound] max\_j|PIP̂\_\{rj\} − 1\{j ∈ S₀\}| ≤ B̄\_r + max\_j|o\_j(S₀,1) − 1\{j ∈ S₀\}| ≤ B̄\_r + u\_π. Hence for the primary average, E\_𝓒 max\_j|PIP̂\_j\textasciicircum{}\{primary\} − 1\{j∈S₀\}| ≤ δ + (1−a)\textasciicircum{}\{⌊L/r₀⌋\} + (L+4)(η+ρ) + u\_π =: 𝓑\_n; no particle independence and no union over p coordinates is needed; the conditional MPM error probability is ≤ 2𝓑\_n. (iii) [truth stability] For the actual hybrid round, ρ ≤ ρ\_*: the Gibbs first-departure probabilities sum over distinct one-flip neighbours, and each of the B+2 other labels contributes ≤ u\_π by reversibility. (iv) [initial occupancy] T15 supplies δ ≤ π(𝓡ᶜ)/c + P\_𝓒(Ẑ\_r/Z < c), 0 < c < 1, with a per-island lower tail. (v) [conditional instances] For either single-flip form of ONE, if the actual one-round first-exit probability from 𝔓 is uniformly ≤ η\_𝔓, reflection and T18 give the candidate entrance bound a ≥ π(S₀|𝔓) − η\_𝔓; the pair is useful when (L+4)η\_𝔓 → 0 and the initial-population condition holds. For a proxy swap, the natural basin \{proxy state, S₀, S₀ ∪ proxy\} needs verified entrance and exit bounds, or T24's 𝓡₁. ONE and TL\_avg alone establish neither basin claim.

\mdhead{2}{Dependencies}

Builds on: T03, T09, T12, T15, T18. Used by: T24, T26, T31. T25 may supply the per-island lower tail.

\mdhead{2}{Numerical evidence (provisional)}

provisional evidence (Sol check J4): labelled-round enumeration (n = 8, p = 4) confirms the first-exit partition (residual 2.2·10⁻¹⁶) and that end-of-round fractions alone do not give the pathwise RB bound (max visit-minus-endpoint 0.578). provisional evidence (Sol check J6): occupancy identity. Files: judge/sol2/J4\_REPAIR\_BASIN.md.

\mdhead{2}{Proof-route hint (optional)}

Blockwise entrance trials over ⌊H/r₀⌋ blocks; separate basin exits from truth departures; bound measured rewards by whole-round bad-path indicators (a pre-coordinate fraction is bounded in expectation by its round-start fraction plus truth's prefix-departure probability); T12 and T15 for the initial occupancy. Hardest step: deriving a useful basin for the actual mixed-label kernel. \emph{When proof.md is run on this file, its lead considers the statement both with and without these hints; the hints never reach the reviewers.}
\end{mdblock}

\mdsec{(b) STATEMENT.md (verbatim)}
\begin{mdblock}
\mdhead{1}{STATEMENT — T23(ii): pathwise Rao–Blackwell PIP bound and its basin consequence}

Source: {\mdtt theorems/\allowbreak{}T23\_\allowbreak{}T18\_\allowbreak{}basin\_\allowbreak{}prefix.\allowbreak{}md}, part (ii) only (SMC-BVS v3.0, METHODOLOGY\_v3.md SHA-256 93043758058a6d422cca6ca4b6be8f3a9e8149b6f6de2e1bab1ce31c44b8019c), read with {\mdtt theorems/\allowbreak{}00\_\allowbreak{}COMMON.\allowbreak{}md}. Parts (i), (iii), (iv), (v) are not claimed here.

Statement analysis outcome: \textbf{(a) correct as stated.} No assumption is changed or added. EXP / EXP′ is not needed. No other v3 conjecture is used as a hypothesis: the expectation bound is proved from the BASIN conditions directly (it does not cite part (i)). Section 5 lists the readings fixed here.

\mdhead{2}{1. Setting (compact; the method file is authoritative)}

\begin{mdbul}
\item \textbf{Models and target.} Ω = \{S ⊆ [p] : |S| ≤ K\}, K = min(p, ⌊n/2⌋). Each S ∈ Ω has weight γ₁(S) = ϖ(S)m(S) > 0; Z = Σ\_Ω γ₁ and π = γ₁/Z on Ω. The data (X, y) are fixed throughout.
\item \textbf{Gibbs conditional.} For S ∈ Ω and j ∈ [p], o\_j(S, 1) is the inclusion probability that the Gibbs step for coordinate j uses at state S (at φ = 1):

\begin{mdbul}
\item if j ∈ S, or j ∉ S and |S| < K: o\_j(S,1) = γ₁(S ∪ \{j\}) / [γ₁(S ∪ \{j\}) + γ₁(S ∖ \{j\})];
\item if j ∉ S and |S| = K (addition excluded): o\_j(S,1) = 0.
\end{mdbul}
\item \textbf{Truth and stability constant.} S₀ ⊆ [p] is a fixed set (the true support). δ₁ = 1 − π(S₀), u\_π = δ₁/π(S₀) (u\_π = +∞ if π(S₀) = 0).
\item \textbf{Round at φ = 1.} One round of a particle is m = p + B + 2 labelled elementary steps (p Gibbs coordinate steps, one uniform swap, B neighbour swaps, one independence proposal from π₀), applied in a uniformly random order drawn independently of the particle's state, afresh for every particle and every round. Each Gibbs step for coordinate j records o\_j(S,1) at the state S held immediately before its draw. A round started at x produces a path x = X⁽⁰⁾, X⁽¹⁾, …, X⁽ᵐ⁾ (X⁽ᵏ⁾ = state after k steps). K₁ is the one-round transition kernel (law of X⁽ᵐ⁾ given X⁽⁰⁾ = x); K₁\textasciicircum{}r its r-fold power.
\item \textbf{Calibration.} 𝓒 is the frozen calibration output (grid, T, L, N\_e, all settings), measurable with respect to streams 0–9 and the data; we fix a 𝓒 for which the final run is executed (successful grid construction). P\_𝓒, E\_𝓒 are probability and expectation over the final-run randomness (streams 10–13) given 𝓒.
\item \textbf{Final run, island r ∈ \{1,2,3,4\}.} N = N\_e particles; T bridge stages (reweight, multinomial resampling, one round per stage); then the \textbf{terminal rounds}: L burn-in rounds and 4 measured rounds, all with the kernel K₁ and with no reweighting or resampling. For particle i, X\_\{i,s\} denotes its state after s terminal rounds (s = 0: after the stage-T resampling and the stage-T round). Measured round h ∈ \{1,2,3,4\} is terminal round L + h; it starts at X\_\{i,L+h−1\}.
\item \textbf{Estimators.} For island r, PIP̂\_\{rj\} = (4N)⁻¹ Σ\_\{h=1\}\textasciicircum{}\{4\} Σ\_\{i=1\}\textasciicircum{}\{N\} o\_j(S\_\{i,h,j\}, 1), where S\_\{i,h,j\} is particle i's state immediately before its Gibbs step for j in measured round h. Primary PIP̂\_j = ¼ Σ\_r PIP̂\_\{rj\}. MPM̂ = \{j : PIP̂\_j > ½\}.
\item \textbf{Bad-round indicators.} B\_\{i,h\} = 1 if particle i's path in measured round h, X⁽⁰⁾, …, X⁽ᵐ⁾ (including its start X⁽⁰⁾ = X\_\{i,L+h−1\}), contains a state ≠ S₀; else 0. B̄\_r = (4N)⁻¹ Σ\_\{h=1\}\textasciicircum{}\{4\} Σ\_\{i=1\}\textasciicircum{}\{N\} B\_\{i,h\}.
\end{mdbul}

\mdhead{2}{2. Hypotheses}

\begin{mdbul}
\item \textbf{EX-A} (exact scores and probabilities, so o\_j and π₀ draws are exact), \textbf{EX-S} (independent ideal uniform streams with the method's stream separation), \textbf{CAL} (𝓒 uses streams 0–9 only; streams 10–13 are fresh).
\item \textbf{BASIN(𝓡, r₀, a, η, ρ, δ)}, given 𝓒: 𝓡 ⊆ Ω with S₀ ∈ 𝓡; an integer r₀ ≥ 1; numbers a, η, ρ ∈ [0,1], δ ≥ 0, such that

\begin{mdbul}
\item (entrance) K₁\textasciicircum{}\{r₀\}(x, \{S₀\}) ≥ a for every x ∈ 𝓡;
\item (basin exit) for every x ∈ 𝓡, the probability that a round started at x has some state outside 𝓡 on its path X⁽¹⁾, …, X⁽ᵐ⁾ is ≤ η;
\item (truth departure) the probability that a round started at S₀ has some state ≠ S₀ on its path X⁽¹⁾, …, X⁽ᵐ⁾ is ≤ ρ;
\item (initial occupancy) for every island r, E\_𝓒[ N⁻¹ Σ\_i 1\{X\_\{i,0\} ∉ 𝓡\} ] ≤ δ.
\end{mdbul}
\end{mdbul}

Part (ii)'s first inequality (pathwise) uses only the setting and EX-A; the second uses S₀ ∈ Ω; the expectation and MPM bounds use EX-S, CAL and BASIN.

\mdhead{2}{3. Conclusions}

(ii-a) \textbf{Pathwise RB bound.} On every realisation of the final run, for each island r,\\
\hspace*{1.5em}max\_j |PIP̂\_\{rj\} − 1\{j ∈ S₀\}| ≤ B̄\_r + max\_j |o\_j(S₀,1) − 1\{j ∈ S₀\}|.\\
\hspace*{1.5em}If S₀ ∉ Ω, the method does not define o\_j(S₀, 1); it may then be given any real value, and the inequality holds for every such choice, because every path state lies in Ω and so B̄\_r = 1.

(ii-b) \textbf{Truth term.} If S₀ ∈ Ω, then max\_j |o\_j(S₀,1) − 1\{j ∈ S₀\}| ≤ u\_π. (If S₀ ∉ Ω, u\_π = +∞ and the bound is trivial.)

(ii-c) \textbf{Expectation bound for the primary PIP.} Under BASIN,\\
\hspace*{1.5em}E\_𝓒 max\_j |PIP̂\_j − 1\{j ∈ S₀\}| ≤ δ + (1−a)\textasciicircum{}\{⌊L/r₀⌋\} + (L+4)(η+ρ) + u\_π =: 𝓑\_n.

(ii-d) \textbf{MPM.} P\_𝓒(MPM̂ ≠ S₀) ≤ 2𝓑\_n.

The bounds use no independence between particles and no union bound over the p coordinates.

\mdhead{2}{4. Scope}

\begin{mdbul}
\item All conclusions are given the data and a fixed successful 𝓒 (calibration-conditional). They say nothing about whether BASIN holds; statistical instances are parts (iii)–(v), not claimed here.
\end{mdbul}

\mdhead{2}{5. Readings fixed here (precisions, not changes)}

\begin{mdnum}
\item[\mdlbl{1.}] \textbf{ρ is the whole-round departure probability from S₀.} The theorem file says "truth prefix-departure ≤ ρ" and, in (iii), bounds ρ by summing first-departure probabilities over all p + B + 2 labels of a round. B\_\{i,h\} counts a departure anywhere in the round. The whole-round reading is the one consistent with both.
\item[\mdlbl{2.}] \textbf{η is the one-round first-exit probability from 𝓡,} maximised over starts in 𝓡 (path version). The proof uses only its consequence K₁(x, 𝓡ᶜ) ≤ η for x ∈ 𝓡, so the end-of-round version would also suffice.
\item[\mdlbl{3.}] \textbf{Terminal rounds have no resampling} (the method resamples after every reweighting; there is no reweighting at φ = 1 after stage T). Particle identities persist through the L + 4 terminal rounds.
\item[\mdlbl{4.}] \textbf{δ is per island} (islands are i.i.d. given 𝓒 under EX-S and CAL, so a bound for one island holds for all).
\item[\mdlbl{5.}] "MPM error" means the event MPM̂ ≠ S₀.
\end{mdnum}
\end{mdblock}

\mdsec{(c) Final proof, verbatim: body of PROOF\_v5.tex}
\noindent{\small\itshape The proof below is the body of \texttt{PROOF\_v5.tex} from \texttt{\textbackslash begin\{document\}} to \texttt{\textbackslash end\{document\}} without the title block. Its preamble definitions are local to this chapter; labels carry the prefix \texttt{T23:}.}\par\medskip
\mdproofbegin
% ---- preamble of PROOF_v5.tex (theorem environments carry the prefix T23) ----
\newtheorem{T23theorem}{Theorem}
\newtheorem{T23lemma}[T23theorem]{Lemma}
\theoremstyle{remark}
\newtheorem{T23remark}[T23theorem]{Remark}
\newcommand{\one}{\mathbf 1}
\newcommand{\R}{\mathcal R}
\newcommand{\C}{\mathcal C}
\newcommand{\PIP}{\widehat{\mathrm{PIP}}}
\newcommand{\Pc}{P_{\C}}
\newcommand{\Ec}{E_{\C}}
% ---- body of PROOF_v5.tex ----
\noindent All notation, the setting and the hypotheses are those of \texttt{STATEMENT.md}.
We recall the objects used. The data are fixed throughout.

\section{Setting}

\paragraph{Models.} $\Omega=\{S\subseteq[p]:|S|\le K\}$; $\gamma_1(S)>0$ for $S\in\Omega$; $\pi=\gamma_1/Z$ on $\Omega$.
For $S\in\Omega$ and $j\in[p]$ the Gibbs inclusion probability at $\varphi=1$ is
\[
o_j(S)=\begin{cases}
\dfrac{\gamma_1(S\cup\{j\})}{\gamma_1(S\cup\{j\})+\gamma_1(S\setminus\{j\})}, & j\in S,\ \text{or } j\notin S \text{ and } |S|<K,\\[8pt]
0, & j\notin S \text{ and } |S|=K .
\end{cases}
\]
(We write $o_j(S)$ for $o_j(S,1)$.) In both cases $o_j(S)\in[0,1]$, and both sets in the first case belong to $\Omega$.
$S_0\subseteq[p]$ is fixed. We extend $\pi$ by $\pi(S)=0$ for $S\notin\Omega$; then $\delta_1=1-\pi(S_0)$, $u_\pi=\delta_1/\pi(S_0)$, with $u_\pi=+\infty$ when $\pi(S_0)=0$, i.e.\ exactly when $S_0\notin\Omega$. When $S_0\notin\Omega$ the quantity $o_j(S_0)$ is not defined by the method; in that case we give it an arbitrary real value $c_j^\ast$ (the results below hold for every choice).

\paragraph{Rounds.} A round started at $x\in\Omega$ applies the $m=p+B+2$ labelled elementary steps in a uniformly random order and produces a path $x=X^{(0)},X^{(1)},\dots,X^{(m)}$, all in $\Omega$. The Gibbs step for coordinate $j$ occurs exactly once in the round; if it is the $k$-th step ($1\le k\le m$), it records $o_j(X^{(k-1)})$. $K_1(x,\cdot)$ is the law of $X^{(m)}$ and $K_1^{r}$ its $r$-fold power. For $x\in\Omega$ and $A\subseteq\Omega$ write
\[
\mathrm{Path}_x(A^c)=\Pr\bigl(X^{(k)}\notin A\text{ for some }1\le k\le m \bigm| X^{(0)}=x\bigr),
\]
the probability that a round started at $x$ visits $A^c$.

\paragraph{Terminal rounds of island $r$.} $N=N_e$ particles. $X_{i,s}$ is particle $i$'s state after $s$ terminal rounds ($s=0$: after the stage-$T$ resampling and the stage-$T$ round). The terminal rounds are $L$ burn-in and $4$ measured rounds, all with kernel $K_1$, with no reweighting and no resampling. Measured round $h\in\{1,2,3,4\}$ is terminal round $L+h$; its path for particle $i$ is $X^{(0)}_{i,h}=X_{i,L+h-1},X^{(1)}_{i,h},\dots,X^{(m)}_{i,h}=X_{i,L+h}$. Let $k_{i,h}(j)$ be the position of the Gibbs step for $j$ in that round and $S_{i,h,j}=X^{(k_{i,h}(j)-1)}_{i,h}$ the state immediately before it. Then
\begin{gather*}
\PIP_{rj}=\frac1{4N}\sum_{h=1}^4\sum_{i=1}^N o_j(S_{i,h,j}),\\
B_{i,h}=\one\bigl\{X^{(k)}_{i,h}\neq S_0\text{ for some }0\le k\le m\bigr\},\qquad
\overline B_r=\frac1{4N}\sum_{h=1}^4\sum_{i=1}^N B_{i,h}.
\end{gather*}
These objects belong to island $r$; we suppress the index $r$ in $X_{i,s}$, $S_{i,h,j}$, $B_{i,h}$, $k_{i,h}(j)$ and $\mathcal F^i_s$ below, and keep it in $\PIP_{rj}$ and $\overline B_r$. The primary estimator is $\PIP_j=\tfrac14\sum_{r=1}^4\PIP_{rj}$ and $\widehat{\mathrm{MPM}}=\{j:\PIP_j>\tfrac12\}$.

\paragraph{BASIN$(\R,r_0,a,\eta,\rho,\delta)$.} $\R\subseteq\Omega$, $S_0\in\R$, $r_0\ge1$ an integer, $a,\eta,\rho\in[0,1]$, $\delta\ge0$, and
\begin{enumerate}[label=(B\arabic*)]
\item $K_1^{r_0}(x,\{S_0\})\ge a$ for every $x\in\R$;
\item $\mathrm{Path}_x(\R^c)\le\eta$ for every $x\in\R$;
\item $\mathrm{Path}_{S_0}(\{S_0\}^c)\le\rho$;
\item $\Ec\bigl[N^{-1}\sum_{i}\one\{X_{i,0}\notin\R\}\bigr]\le\delta$ for every island $r$.
\end{enumerate}
Since the event ``$X^{(m)}\notin A$'' is contained in the event ``$X^{(k)}\notin A$ for some $1\le k\le m$'',
\begin{equation}\label{T23:eq:endpoint}
K_1(x,\R^c)\le\eta\ (x\in\R),\qquad K_1(S_0,\{S_0\}^c)\le\rho .
\end{equation}

\section{The Markov structure of the terminal rounds}

Throughout, $\C$ is fixed (a successful calibration) and all probabilities are $\Pc$. Fix an island $r$.

\paragraph{Elementary steps.} For a label $e$ (one of the $p$ Gibbs coordinates, the uniform swap, the $B$ neighbour swaps, the independence proposal) let $K_e(x,\cdot)$ be the law of the state after step $e$ is applied at state $x$, as prescribed by the method at $\varphi=1$. By \texttt{STATEMENT.md} \S1, every terminal round uses the same round, with the same kernel $K_1$, fixed by the data and $\C$. So the $K_e$ do not depend on the round, the particle or the particle population. By the definition of a round, a round started at $x_0$ has path law
\begin{equation}\label{T23:eq:pathlaw}
\Pr\bigl(X^{(1)}=x_1,\dots,X^{(m)}=x_m\bigm|X^{(0)}=x_0\bigr)=\frac1{m!}\sum_{\sigma}\prod_{t=1}^m K_{\sigma(t)}(x_{t-1},\{x_t\}),
\end{equation}
where $\sigma$ runs over the orderings of the $m$ labels and $\sigma(t)$ is the label at position $t$.

\paragraph{The form of EX-S used.} Each random choice of the algorithm (an order of the labels, or the outcome of one elementary step) is made from uniform variates. EX-S (ideal independent uniform streams), together with CAL (island $r$'s streams are among $10$--$13$, independent of streams $0$--$9$ and hence of $\C$), is used in the following form:
\begin{itemize}
\item[(F)] For each random choice $c$ of island $r$'s final run, let $\mathcal D_c$ be the $\sigma$-field generated by all variates of island $r$ used by the random choices made before $c$. Conditionally on $\mathcal D_c$, the fresh variates from which $c$ is made (a sequence, if $c$ uses a random number of them) are i.i.d.\ uniform. Hence, given $\mathcal D_c$, the choice $c$ has the law that the method prescribes for it given the current state.
\end{itemize}
Concretely: the order of a round is uniform over the $m!$ orderings given $\mathcal D_c$ (the method draws it ``independently of the particle's state, afresh for every particle and every round''); and the outcome of step $e$ applied at the current state $x$ has law $K_e(x,\cdot)$ given $\mathcal D_c$. Under EX-A these laws are realised exactly. Variates are attributed to the choice that uses them, whenever they were generated. So (F) does not depend on whether particles are processed round by round across the population or one particle at a time, on how the variates are split into sub-streams, or on whether variates are generated in batches ahead of use. It is the idealisation EX-S itself, stated for the choices made by the algorithm.

\paragraph{Per-particle $\sigma$-fields.} For a particle $i$ and $0\le s\le L+3$, let $\mathcal F^i_s$ be the $\sigma$-field generated by all variates of island $r$ used by the random choices made before particle $i$'s terminal round $s+1$ begins. Then $\mathcal F^i_s\subseteq\mathcal F^i_{s+1}$, because every choice made before round $s+1$ of particle $i$ begins is also made before its round $s+2$ begins. Also $X_{i,0},\dots,X_{i,s}$ are $\mathcal F^i_s$-measurable, because they are computed from the initial draws, the bridge stages and particle $i$'s own terminal rounds $1,\dots,s$, all of which precede its round $s+1$.

\begin{T23lemma}\label{T23:lem:markov}
Assume EX-A, EX-S (in the form (F)) and CAL. For every particle $i$, every $0\le s\le L+3$ and every set $\Gamma\subseteq\Omega^{m+1}$ of paths,
\[
\Pc\bigl((Z^{(0)}_{i,s+1},\dots,Z^{(m)}_{i,s+1})\in\Gamma\bigm|\mathcal F^i_s\bigr)=\Pr\bigl((X^{(0)},\dots,X^{(m)})\in\Gamma\bigm|X^{(0)}=x\bigr)\Big|_{x=X_{i,s}}\quad\text{a.s.},
\]
where $(Z^{(0)}_{i,s+1},\dots,Z^{(m)}_{i,s+1})$ is particle $i$'s path in terminal round $s+1$, so $Z^{(0)}_{i,s+1}=X_{i,s}$ (no resampling between terminal rounds). In the notation of \S1, the path of measured round $h$ is $X^{(k)}_{i,h}=Z^{(k)}_{i,L+h}$. In particular, $\Pc(X_{i,s+1}\in A\mid\mathcal F^i_s)=K_1(X_{i,s},A)$. Moreover, for $0\le s<s+q\le L+4$ and $A\subseteq\Omega$, $\Pc(X_{i,s+q}\in A\mid\mathcal F^i_s)=K_1^{q}(X_{i,s},A)$.
\end{T23lemma}

\begin{proof}
Write $Y_t=Z^{(t)}_{i,s+1}$ for $0\le t\le m$, and let $\Sigma$ be the order drawn for this round. Let $\mathcal D_0$ be the $\sigma$-field of all variates used by the random choices made before $\Sigma$ is chosen, and, for $1\le t\le m$, $\mathcal D_t$ the $\sigma$-field of all variates used by the random choices made before the outcome of the $t$-th step of this round is chosen. Then $\mathcal F^i_s\subseteq\mathcal D_0\subseteq\mathcal D_1\subseteq\dots\subseteq\mathcal D_m$. $Y_0=X_{i,s}$ is $\mathcal D_0$-measurable, and $\Sigma,Y_0,\dots,Y_{t-1}$ are $\mathcal D_t$-measurable. By (F), $\Pc(\Sigma=\sigma\mid\mathcal D_0)=1/m!$ for every ordering $\sigma$, and $\Pc(Y_t=y\mid\mathcal D_t)=K_{\Sigma(t)}(Y_{t-1},\{y\})$ for $1\le t\le m$. Fix $x_0,\dots,x_m\in\Omega$ and an ordering $\sigma$. Conditioning successively on $\mathcal D_m,\mathcal D_{m-1},\dots,\mathcal D_1,\mathcal D_0$ (tower property; all spaces are finite) gives
\[
\Pc\bigl(\Sigma=\sigma,\ Y_0=x_0,\dots,Y_m=x_m\bigm|\mathcal D_0\bigr)=\one\{X_{i,s}=x_0\}\,\frac1{m!}\prod_{t=1}^mK_{\sigma(t)}(x_{t-1},\{x_t\}).
\]
Indeed, at stage $t$ the factor $\one\{\Sigma=\sigma,Y_0=x_0,\dots,Y_{t-1}=x_{t-1}\}$ is $\mathcal D_t$-measurable, and on that event $\Pc(Y_t=x_t\mid\mathcal D_t)=K_{\sigma(t)}(x_{t-1},\{x_t\})$, a constant. Sum over $\sigma$ and condition on $\mathcal F^i_s\subseteq\mathcal D_0$. The right side is $\mathcal F^i_s$-measurable, so by \eqref{T23:eq:pathlaw} the conditional law of the path given $\mathcal F^i_s$ is the round path law started at $X_{i,s}$. Summing over the paths in $\Gamma$ gives the first claim. The second claim is the case $\Gamma=\{\text{paths with }x_m\in A\}$.

For the $q$-round claim, induct on $q$. The case $q=1$ is the second claim. If it holds for $q$, then, since $\mathcal F^i_s\subseteq\mathcal F^i_{s+q}$ and $X_{i,s+q}$ is $\mathcal F^i_{s+q}$-measurable (here $s+q\le L+3$),
\[
\Pc(X_{i,s+q+1}\in A\mid\mathcal F^i_s)=\Ec\bigl[K_1(X_{i,s+q},A)\bigm|\mathcal F^i_s\bigr]=\sum_{y}K_1^{q}(X_{i,s},\{y\})K_1(y,A)=K_1^{q+1}(X_{i,s},A).\qedhere
\]
\end{proof}

\section{The pathwise bound}

\begin{T23lemma}[pathwise RB bound]\label{T23:lem:pathwise}
Assume EX-A (the recorded value is the exact $o_j$ at the pre-draw state). On every realisation, for every island $r$ and every $j\in[p]$,
\[
|\PIP_{rj}-\one\{j\in S_0\}|\le \overline B_r+(1-\overline B_r)\,|o_j(S_0)-\one\{j\in S_0\}|
\le \overline B_r+\max_{l\in[p]}|o_l(S_0)-\one\{l\in S_0\}|.
\]
\end{T23lemma}

\begin{proof}
If $S_0\notin\Omega$, then no path state (all lie in $\Omega$) equals $S_0$, so $B_{i,h}=1$ for all $(i,h)$ and $\overline B_r=1$; since $o_j\in[0,1]$, the left side is $\le1\le1+|c_l^\ast-\one\{l\in S_0\}|$ for any value $c_l^\ast$, and both inequalities hold. Assume now $S_0\in\Omega$.
Write $c_j=\one\{j\in S_0\}$. By the triangle inequality,
\[
|\PIP_{rj}-c_j|=\Bigl|\frac1{4N}\sum_{h,i}\bigl(o_j(S_{i,h,j})-c_j\bigr)\Bigr|\le\frac1{4N}\sum_{h,i}|o_j(S_{i,h,j})-c_j| .
\]
Fix $(i,h)$. The state $S_{i,h,j}=X^{(k_{i,h}(j)-1)}_{i,h}$ is one of the states $X^{(0)}_{i,h},\dots,X^{(m-1)}_{i,h}$ of the path of measured round $h$.
If $B_{i,h}=0$, every state of that path equals $S_0$ (so $S_0\in\Omega$), hence $S_{i,h,j}=S_0$ and $|o_j(S_{i,h,j})-c_j|=|o_j(S_0)-c_j|$.
If $B_{i,h}=1$, then $|o_j(S_{i,h,j})-c_j|\le1$ because $o_j\in[0,1]$ and $c_j\in\{0,1\}$.
Thus $|o_j(S_{i,h,j})-c_j|\le B_{i,h}+(1-B_{i,h})|o_j(S_0)-c_j|$. Averaging over $(i,h)$ gives the first inequality, and $(1-\overline B_r)\in[0,1]$ gives the second. The bound holds for each $j$, hence for $\max_j$.
\end{proof}

\begin{T23lemma}[truth term]\label{T23:lem:truth}
If $S_0\in\Omega$, then $\max_{j\in[p]}|o_j(S_0)-\one\{j\in S_0\}|\le u_\pi$.
\end{T23lemma}

\begin{proof}
$\pi(S_0)>0$ since $\gamma_1>0$ on $\Omega$. Let $j\in S_0$. Then $S_0\setminus\{j\}\in\Omega$, $S_0\setminus\{j\}\neq S_0$, and
\[
1-o_j(S_0)=\frac{\gamma_1(S_0\setminus\{j\})}{\gamma_1(S_0)+\gamma_1(S_0\setminus\{j\})}\le\frac{\gamma_1(S_0\setminus\{j\})}{\gamma_1(S_0)}=\frac{\pi(S_0\setminus\{j\})}{\pi(S_0)}\le\frac{1-\pi(S_0)}{\pi(S_0)}=u_\pi .
\]
Let $j\notin S_0$. If $|S_0|=K$, then $o_j(S_0)=0$ and the bound is trivial. If $|S_0|<K$, then $S_0\cup\{j\}\in\Omega$, $S_0\cup\{j\}\ne S_0$, and
\[
o_j(S_0)=\frac{\gamma_1(S_0\cup\{j\})}{\gamma_1(S_0\cup\{j\})+\gamma_1(S_0)}\le\frac{\pi(S_0\cup\{j\})}{\pi(S_0)}\le\frac{1-\pi(S_0)}{\pi(S_0)}=u_\pi .
\]
In both displays we used $\pi(S')\le\pi(\Omega\setminus\{S_0\})=1-\pi(S_0)$ for $S'\in\Omega\setminus\{S_0\}$.
\end{proof}

\section{Expected bad-round fraction}

Fix an island $r$ and a particle $i$; write $X_s=X_{i,s}$ and $\mathcal F_s=\mathcal F^i_s$ ($0\le s\le L+3$), and $\Pc$ for the conditional law given $\C$. All conditioning below is on $\mathcal F_s$ with $s\le L+3$, as required by Lemma~\ref{T23:lem:markov}.
Put $\bar\eta=\max_{x\in\R}K_1(x,\R^c)$ and $\bar\rho=K_1(S_0,\{S_0\}^c)$; by \eqref{T23:eq:endpoint}, $\bar\eta\le\eta$ and $\bar\rho\le\rho$.

\begin{T23lemma}[occupancy at round ends]\label{T23:lem:occ}
Under EX-A, EX-S, CAL and (B1), for every integer $0\le H\le L+4$,
\[
\Pc(X_H\neq S_0)\le\Pc(X_0\notin\R)+(1-a)^{\lfloor H/r_0\rfloor}+H(\bar\eta+\bar\rho).
\]
\end{T23lemma}

\begin{proof}
Let $k=\lfloor H/r_0\rfloor$, so $kr_0\le H$. Define the events
\begin{align*}
E_1&=\{X_0\notin\R\},\\
E_2&=\{X_s\in\R,\ X_{s+1}\notin\R\ \text{for some }0\le s\le H-1\},\\
E_3&=\{X_s=S_0,\ X_{s+1}\neq S_0\ \text{for some }0\le s\le H-1\},\\
E_4&=\{X_{jr_0}\in\R\ \text{for }0\le j\le k-1\}\cap\{X_{jr_0}\neq S_0\ \text{for }1\le j\le k\}.
\end{align*}
(For $k=0$, $E_4$ is the whole space.)

\emph{Inclusion: $\{X_H\neq S_0\}\subseteq E_1\cup E_2\cup E_3\cup E_4$.}
Suppose $X_H\ne S_0$ and none of $E_1,E_2,E_3$ occurs. Since $E_1$ fails, $X_0\in\R$; since $E_2$ fails, induction on $s$ gives $X_s\in\R$ for all $0\le s\le H$. If $X_s=S_0$ for some $s\le H$, then, since $E_3$ fails, induction gives $X_{s'}=S_0$ for all $s\le s'\le H$, contradicting $X_H\neq S_0$. So $X_s\in\R\setminus\{S_0\}$ for all $0\le s\le H$; as $jr_0\le H$ for $0\le j\le k$, $E_4$ occurs.

\emph{Bound on $E_2$ and $E_3$.} By the union bound and Lemma~\ref{T23:lem:markov},
\begin{align*}
\Pc(E_2)&\le\sum_{s=0}^{H-1}\Pc(X_s\in\R,\,X_{s+1}\notin\R)\\&=\sum_{s=0}^{H-1}\Ec\bigl[\one\{X_s\in\R\}K_1(X_s,\R^c)\bigr]\le H\bar\eta,\\
\Pc(E_3)&\le\sum_{s=0}^{H-1}\Pc(X_s=S_0,\,X_{s+1}\ne S_0)\\&=\sum_{s=0}^{H-1}\Ec\bigl[\one\{X_s=S_0\}K_1(S_0,\{S_0\}^c)\bigr]\le H\bar\rho .
\end{align*}

\emph{Bound on $E_4$.} For $1\le l\le k$ let $G_l=\{X_{jr_0}\in\R\text{ for }0\le j\le l-1\}\cap\{X_{jr_0}\ne S_0\text{ for }1\le j\le l\}$, and $G_0$ the whole space; so $E_4=G_k$. For $1\le l\le k$, $G_l=G_{l-1}\cap\{X_{(l-1)r_0}\in\R\}\cap\{X_{lr_0}\neq S_0\}$, and $G_{l-1}\cap\{X_{(l-1)r_0}\in\R\}\in\mathcal F_{(l-1)r_0}$. Since $(l-1)r_0\le H-r_0\le L+3$ and $lr_0\le H\le L+4$, the $q$-round part of Lemma~\ref{T23:lem:markov} with $q=r_0$ gives $\Pc(X_{lr_0}=S_0\mid\mathcal F_{(l-1)r_0})=K_1^{r_0}(X_{(l-1)r_0},\{S_0\})$, which is $\ge a$ on $\{X_{(l-1)r_0}\in\R\}$ by (B1). Hence
\[
\Pc(G_l)=\Ec\Bigl[\one_{G_{l-1}\cap\{X_{(l-1)r_0}\in\R\}}\bigl(1-K_1^{r_0}(X_{(l-1)r_0},\{S_0\})\bigr)\Bigr]\le(1-a)\Pc(G_{l-1}),
\]
and $\Pc(E_4)=\Pc(G_k)\le(1-a)^k$.

Summing the four bounds proves the lemma.
\end{proof}

\begin{T23lemma}[one measured round]\label{T23:lem:round}
Under EX-A, EX-S, CAL and (B1)--(B3), for every particle $i$ and $h\in\{1,2,3,4\}$,
\[
\Pc(B_{i,h}=1)\le\Pc(X_{i,0}\notin\R)+(1-a)^{\lfloor L/r_0\rfloor}+(L+3)\eta+(L+4)\rho .
\]
\end{T23lemma}

\begin{proof}
Let $H=L+h-1$, the number of terminal rounds before measured round $h$; $X^{(0)}_{i,h}=X_{i,H}$. If $B_{i,h}=1$ then either $X_{i,H}\ne S_0$, or $X_{i,H}=S_0$ and some state $X^{(k)}_{i,h}$, $1\le k\le m$, differs from $S_0$. By Lemma~\ref{T23:lem:markov} with $s=H$ (the path of measured round $h$ is $Z^{(\cdot)}_{i,H+1}$) and (B3), the second event has probability $\Ec[\one\{X_{i,H}=S_0\}\mathrm{Path}_{S_0}(\{S_0\}^c)]\le\rho$. By Lemma~\ref{T23:lem:occ},
\[
\Pc(B_{i,h}=1)\le\Pc(X_{i,0}\notin\R)+(1-a)^{\lfloor H/r_0\rfloor}+H(\bar\eta+\bar\rho)+\rho .
\]
Now $H\ge L$ and $0\le1-a\le1$ give $(1-a)^{\lfloor H/r_0\rfloor}\le(1-a)^{\lfloor L/r_0\rfloor}$; and $H\le L+3$, $\bar\eta\le\eta$, $\bar\rho\le\rho$ give $H(\bar\eta+\bar\rho)+\rho\le(L+3)\eta+(L+4)\rho$.
\end{proof}

\section{The theorem}

\begin{T23theorem}[T23(ii)]\label{T23:thm:main}
\begin{enumerate}[label=(\alph*)]
\item Under EX-A, on every realisation, for every island $r$: $\displaystyle\max_j|\PIP_{rj}-\one\{j\in S_0\}|\le\overline B_r+\max_j|o_j(S_0)-\one\{j\in S_0\}|$. If $S_0\notin\Omega$, the method does not define $o_j(S_0,1)$; it may then be given any real value, and the inequality holds for every such choice, because every path state lies in $\Omega$ and so $\overline B_r=1$.
\item If $S_0\in\Omega$: $\max_j|o_j(S_0)-\one\{j\in S_0\}|\le u_\pi$.
\item Under EX-A, EX-S, CAL and BASIN$(\R,r_0,a,\eta,\rho,\delta)$,
\[
\Ec\max_j|\PIP_j-\one\{j\in S_0\}|\le\delta+(1-a)^{\lfloor L/r_0\rfloor}+(L+4)(\eta+\rho)+u_\pi=:\mathcal B_n .
\]
\item Under the same hypotheses, $\Pc(\widehat{\mathrm{MPM}}\neq S_0)\le2\mathcal B_n$.
\end{enumerate}
\end{T23theorem}

\begin{proof}
(a) is Lemma~\ref{T23:lem:pathwise} and (b) is Lemma~\ref{T23:lem:truth}.

(c) BASIN gives $S_0\in\R\subseteq\Omega$, so (a) and (b) give, for each island $r$, on every realisation,
$\max_j|\PIP_{rj}-\one\{j\in S_0\}|\le\overline B_r+u_\pi$. For the primary estimator, by the triangle inequality,
\begin{align*}
\max_j|\PIP_j-\one\{j\in S_0\}|&=\max_j\Bigl|\frac14\sum_{r=1}^4\bigl(\PIP_{rj}-\one\{j\in S_0\}\bigr)\Bigr|\\
&\le\frac14\sum_{r=1}^4\max_j|\PIP_{rj}-\one\{j\in S_0\}|\\
&\le\frac14\sum_{r=1}^4\overline B_r+u_\pi .
\end{align*}
For each island, linearity of expectation and Lemma~\ref{T23:lem:round} give
\[
\Ec\overline B_r=\frac1{4N}\sum_{h=1}^4\sum_{i=1}^N\Pc(B_{i,h}=1)
\le\Ec\Bigl[\frac1N\sum_{i=1}^N\one\{X_{i,0}\notin\R\}\Bigr]+(1-a)^{\lfloor L/r_0\rfloor}+(L+3)\eta+(L+4)\rho,
\]
which is $\le\delta+(1-a)^{\lfloor L/r_0\rfloor}+(L+4)(\eta+\rho)$ by (B4). Taking $\Ec$ in the previous display proves (c). (When $u_\pi=+\infty$ the claim is trivial; under BASIN, $u_\pi<\infty$.)

(d) If $\widehat{\mathrm{MPM}}\neq S_0$, there is $j$ with either $j\in S_0$ and $\PIP_j\le\frac12$, so $|\PIP_j-1|\ge\frac12$; or $j\notin S_0$ and $\PIP_j>\frac12$, so $|\PIP_j-0|>\frac12$. Hence $\{\widehat{\mathrm{MPM}}\ne S_0\}\subseteq\{\max_j|\PIP_j-\one\{j\in S_0\}|\ge\frac12\}$, and Markov's inequality and (c) give $\Pc(\widehat{\mathrm{MPM}}\ne S_0)\le2\mathcal B_n$.
\end{proof}

\begin{T23remark}
No independence between particles is used: Lemma~\ref{T23:lem:round} is a statement about one particle, and (c) uses only linearity of expectation over particles. No union bound over coordinates is used: the bad-round indicators $B_{i,h}$ do not depend on $j$, so one bound serves all $p$ coordinates at once.
\end{T23remark}

\begin{T23remark}
Lemma~\ref{T23:lem:occ} uses only (B1), with the end-of-round quantities $\bar\eta,\bar\rho$; (B2)--(B3) enter Lemma~\ref{T23:lem:round} through $\bar\eta\le\eta$ and $\bar\rho\le\rho$ \eqref{T23:eq:endpoint}. The path form of (B3) is needed once, in Lemma~\ref{T23:lem:round}, because the Gibbs recordings of a measured round are taken at intermediate states of the round, not at its end.
\end{T23remark}
\mdproofend
\mdsec{(d) Review record}
\mdsubsec{FINAL\_REPORT.md: statement changes}
\begin{mdblock}
\mdhead{2}{Statement changes}

None. Step-1 outcome: (a), correct as stated.

\begin{mdbul}
\item EXP/EXP′ is not needed.
\item No sibling conjecture is used as a hypothesis; the occupancy step is proved directly from BASIN, without citing part (i).
\item STATEMENT.md §5 records five readings. These are precisions, not changes: 1. ρ is the whole-round departure probability from S₀. 2. η is the one-round first-exit probability from the basin 𝓡. 3. There is no resampling in the L + 4 terminal rounds. 4. δ applies to every island. 5. "MPM error" means the event MPM̂ ≠ S₀.
\item The proof states EX-S in the form it uses, (F): each random choice of the algorithm, given everything used by earlier choices, has the law the method prescribes. Both Stage B reviewers ruled (F) a faithful reading of EX-S and CAL, not an added assumption.
\end{mdbul}
\end{mdblock}

\mdsubsec{FINAL\_REPORT.md: review history}
\begin{mdblock}
\mdhead{2}{Review history}

\textbf{Stage A (lead + Astra \#1).} Astra \#1: gpt-6-astra, reasoning high, session 01a16419-ed12-7492-9b6f-8736751fe274.

\begin{mdbul}
\item PROOF\_v1 → "No findings … VERDICT: verified".
\end{mdbul}

\textbf{Stage B, round 1} (independent, started together; each reviewer had only STATEMENT.md, the proof and the report format).

\begin{mdbul}
\item \textbf{Astra \#2} (new session 01a1641c-c05f-7d81-9ac6-6bcc512573e3, high): no findings, verified.
\item \textbf{Opus 5.5 reviewer} (effort xhigh, read-only Plan agent, own account): verified, with 6 exposition findings, all accepted → PROOF\_v2. The findings:

\begin{mdbul}
\item Lemma 1's stopping-time step was not written out.
\item The kernel being fixed given 𝓒 was not cited to the statement.
\item The case S₀ ∉ Ω needed a convention.
\item Lemma 4's range of H was not bounded.
\item Lemma 4's hypothesis list was not minimal.
\item EX-A and the island index were not stated explicitly.
\end{mdbul}
\end{mdbul}

\textbf{Round 2 (rulings on v2).}

\begin{mdbul}
\item Both verified and agreed with all six changes.
\item They raised 7 more exposition points, all accepted → PROOF\_v3. The main ones:

\begin{mdbul}
\item Lemma 1 cited method details that STATEMENT.md does not contain (both reviewers).
\item Lemma 1 assumed the population is processed round by round.
\end{mdbul}
\item The fix rewrote Lemma 1 with per-particle σ-fields and a step-by-step conditioning proof.
\end{mdbul}

\textbf{Round 3 (rulings on v3).}

\begin{mdbul}
\item Astra \#2: no findings, verified; it ruled (F) faithful.
\item Opus: verified, with 2 exposition points (define the σ-fields by variates \emph{used}, not \emph{drawn}; a clash of path notation). Both accepted → PROOF\_v4, wording only.
\end{mdbul}

\textbf{Round 4 (confirmation of v4).} Both verified.

\begin{mdbul}
\item The Opus reviewer added one cosmetic note: the path symbol Z reuses the evidence symbol Z. It said this "needs no action". It is deferred, not changed, so the final proof is exactly the text both reviewers confirmed.
\end{mdbul}

\textbf{Overall.}

\begin{mdbul}
\item No fatal or gap finding in any round, from any reviewer (CLAIMS\_OF\_FATAL\_FLAWS.md).
\item No dispute, so no Astra judge.
\item The lead holds the proof correct. Astra \#1 verified v1; Astra \#2 and the Opus reviewer verified v2, v3 and v4. The v1→v4 changes are exposition plus a more careful proof of Lemma 1. Astra \#1 was not re-consulted after Stage A.
\end{mdbul}
\end{mdblock}

\mdsubsec{CLAIMS\_OF\_FATAL\_FLAWS.md (every entry, verbatim)}
\begin{mdblock}
\mdhead{1}{CLAIMS\_OF\_FATAL\_FLAWS — T23(ii)}

Fatal or gap findings, verbatim, with the lead's answer.

\mdhead{2}{Stage A (Astra \#1, session 01a16419-ed12-7492-9b6f-8736751fe274)}

None. Reply: "No findings." / "VERDICT: verified" (review/r1\_ASTRA\_REPLY.md).

\mdhead{2}{Stage B, round 1}

\begin{mdbul}
\item Astra \#2 (session 01a1641c-c05f-7d81-9ac6-6bcc512573e3): none. Reply: "No findings." / "VERDICT: verified" (review/B\_astra2\_r1\_REPLY.md).
\item Opus 5.5 xhigh reviewer: none. Its six findings are all graded "exposition" by the reviewer. Its summary says "I found no fatal error and no gap" and ends "VERDICT: verified" (review/B\_opus\_r1\_REPLY.md).
\end{mdbul}

There are no fatal or gap findings to log. The exposition findings and the changes made for them are in review/B\_CHANGES\_v1\_to\_v2.md; the unified proof is proof/PROOF\_v2.tex.

\mdhead{2}{Stage B, round 2 (rulings on PROOF\_v2)}

\begin{mdbul}
\item Astra \#2: "I found no fatal error or gap, and one exposition issue"; agrees with all six changes; "VERDICT: verified" (review/B\_astra2\_r2\_REPLY.md).
\item Opus reviewer: six new findings, all exposition; "I add no fatal or gap entry for round 2"; agrees with the log entry describing its round-1 report; "VERDICT: verified" (review/B\_opus\_r2\_REPLY.md).
\end{mdbul}

Still no fatal or gap findings to log. The round-2 exposition findings and the changes made for them are in review/B\_CHANGES\_v2\_to\_v3.md; the unified proof is proof/PROOF\_v3.tex.

\mdhead{2}{Stage B, round 3 (rulings on PROOF\_v3)}

\begin{mdbul}
\item Astra \#2: "No findings."; agrees with every v2→v3 change; finds (F) "a faithful reading of EX-S and CAL, rather than an added assumption"; "VERDICT: verified" (review/B\_astra2\_r3\_REPLY.md).
\item Opus reviewer: two findings, both exposition; "I add no fatal or gap entry"; finds (F) a faithful reading provided "drawn before" is read as "used by earlier choices"; "VERDICT: verified" (review/B\_opus\_r3\_REPLY.md).
\end{mdbul}

No fatal or gap finding was raised in any round. The two round-3 exposition points are fixed in proof/PROOF\_v4.tex (review/B\_CHANGES\_v3\_to\_v4.md).

\mdhead{2}{Stage B, round 4 (confirmation of PROOF\_v4)}

\begin{mdbul}
\item Astra \#2: "No findings."; the changes "preserve every argument, hypothesis, conclusion, and constant"; "VERDICT: verified" (review/B\_astra2\_r4\_REPLY.md).
\item Opus reviewer: one cosmetic exposition note (the symbol Z for terminal-round paths reuses the evidence symbol Z), with "This needs no action for the verdict"; "VERDICT: verified" (review/B\_opus\_r4\_REPLY.md). Lead: deferred, not changed, so that the final proof is exactly the text both reviewers confirmed.
\end{mdbul}

Final: no fatal or gap finding in any round (Stage A or B); no dispute; no Astra judge.
\end{mdblock}

\chapter{T31 (T22), part (i): Markov coverage of \texorpdfstring{$0.05\,\widehat Z$}{0.05 Zhat} and what it certifies}
\chaptersource{Final proof: \texttt{proofs/\allowbreak{}T31\_\allowbreak{}T22\_\allowbreak{}certificates/\allowbreak{}proof/\allowbreak{}PROOF\_v3.tex}. Sources of parts (a), (b), (d): \texttt{theorems/\allowbreak{}T31\_\allowbreak{}T22\_\allowbreak{}certificates.md}, \texttt{STATEMENT.md}, \texttt{FINAL\_REPORT.md}, \texttt{CLAIMS\_OF\_FATAL\_FLAWS.md}.}
\mdsec{(a) Original conjecture text (theorem file, verbatim)}
\begin{mdblock}
\mdhead{1}{T31 (T22): What the reported checks and certificates can establish}

\textbf{Status:} conjectured, not proven. Methodology SMC-BVS v3.0 (2026-10-19), SHA-256 93043758058a6d422cca6ca4b6be8f3a9e8149b6f6de2e1bab1ce31c44b8019c. Read with {\mdtt 00\_\allowbreak{}COMMON.\allowbreak{}md}.

\textbf{Category / rank:} validity and certification limits; \textbf{core} (proposition).

\mdhead{2}{Definitions used}

Reported quantities: Ẑ, the bound 0.05Ẑ, island SD of log Ẑ\_r (E2), E1, E3, low\_empirical\_ESS, local\_certified, global\_certified; Z\_W := Zπ(W) (distinct from the comparator mean r\_WZ); Z\_V := Σ\_\{S∈V\}ϖ(S)m(S); Q\_f; thresholds t fixed before final randomness; weight comparability ratio C\_eq.

\mdhead{2}{Assumptions}

EX-A, EX-S, CAL; MISS-E/CPL-E for (iii)–(iv); the usability, threshold-measurability or external-bound assumptions stated in the clauses.

\mdhead{2}{Precise statement}

(i) Under T12: P\_𝓒(Z ≥ αẐ) ≥ 1 − α and E\_𝓒 log Ẑ ≤ log Z; so 0.05Ẑ has coverage at least 95\% under the ideal final-run assumptions; it is not a relative-error interval; conditioning afterwards on final flags does not preserve this coverage automatically. (ii) [estimator-class example, not a regression construction] For 0 < a, ε < 1, an unbiased nonnegative estimator equal to aZ with probability 1−ε and to Z\{1 − (1−ε)a\}/ε otherwise has four independent copies all equal to aZ with probability (1−ε)⁴: zero log dispersion with arbitrarily severe underestimation. Independence and agreement give no uniformly informative multiplicative-accuracy certificate over this class. (iii) Under T28, diagnostics measurable from the evidence coupling have the comparator's law up to h\_n\textasciicircum{}E; score/search diagnostics need T29's information control; passing transfers only if the comparator passes with high probability; calibration flags are fixed given 𝓒, recomputed final E2 is not. (iv) [upper bounds] If U uses only coupled information and P\_\{ref,𝓒\}(U\textasciicircum{}W ≤ CZ\_W) ≥ 1 − α′ with CZ\_W < Z, then P\_𝓒(Z ≤ U) ≤ h\_n\textasciicircum{}E + α′ on that loss instance; not an unconditional impossibility theorem. (v) [discovery after selection] For a threshold t fixed before final randomness: P\_𝓒\{∃S : π₀(S) ≥ t, S missed by the iid draws\} ≤ min\{1, ⌊1/t⌋e\textasciicircum{}\{−Q\_ft\}\}; and with probability ≥ 1 − δ every model with Q\_fπ₀(S) ≥ log[1/(δπ₀(S))] has been scored (usable after selection). Substituting π₀(Ŝ) into a fixed-model miss formula is invalid as a competitor certificate; neither statement shows all better-posterior models have large start mass. (vi) local\_certified has exactly T20(i)'s meaning; global\_certified is restricted to completed validated enumeration; a genuine external U ≥ Z would certify the best scored model globally whenever U − Z\_V < max\_\{S∈V\}ϖ(S)m(S) (v3 supplies no such U). (vii) Restricted ESS controls only T13's restricted moment; grid equality across rebuilds certifies neither omitted mass nor mixing; if M models have unnormalised weight ratio ≤ C\_eq at a temperature, V captures at most C\_eq|V|/M of their mass; "near-equivalent" without a quantitative condition supplies no bound.

\mdhead{2}{Dependencies}

Builds on: T12, T13, T11, T20, T28, T29 (and T05–T07, T27 for interpretation). Used by: nothing (terminal interpretation).

\mdhead{2}{Numerical evidence (provisional)}

provisional evidence (Sol check J7): exact finite law p = (0.4, 0.3, 0.2, 0.1), 3 draws: fixed-model miss 0.216, all-above-0.25 miss 0.532, selected-threshold competitor miss 0.54 vs naive 0.545 (conditional on the selected label: 1 vs 0.74); fixed-threshold union bound 0.559. provisional evidence (Sol check J8): information-transfer corrections.

\mdhead{2}{Proof-route hint (optional)}

Markov and Jensen; the explicit rare-tail law; information-preserving coupling (T17); at most ⌊1/t⌋ atoms of mass ≥ t for the union bound; allocate error δπ₀(S) per atom for the simultaneous statement; T13 for (vii). Hardest step: none in the algebra; a verified production loss instance is the difficult premise. \emph{When proof.md is run on this file, its lead considers the statement both with and without these hints; the hints never reach the reviewers.}
\end{mdblock}

\mdsec{(b) STATEMENT.md (verbatim)}
\begin{mdblock}
\mdhead{1}{STATEMENT — T31 (T22), part (i) only}

Method: SMC-BVS v3.0 (METHODOLOGY\_v3.md, SHA-256 93043758058a6d422cca6ca4b6be8f3a9e8149b6f6de2e1bab1ce31c44b8019c).

\mdhead{2}{Source text (verbatim, clause (i) of the theorem file)}

\begin{mdquote}
(i) Under T12: P\_𝓒(Z ≥ αẐ) ≥ 1 − α and E\_𝓒 log Ẑ ≤ log Z; so 0.05Ẑ has coverage at least 95\% under the ideal final-run assumptions; it is not a relative-error interval; conditioning afterwards on final flags does not preserve this coverage automatically.
\end{mdquote}

Assumptions listed for the whole file: EX-A, EX-S, CAL. Clause (i) uses T12 (a sibling conjecture), which enters below only as the explicit hypothesis (H-T12).

\mdhead{2}{1. Setting and definitions}

Fixed data (X, y), already standardised as in the method (§4): n ≥ 3, ‖y‖² = n − 1 > 0, p ≥ 1 retained columns. All objects below are deterministic functions of the data unless stated otherwise.

\begin{mdbul}
\item State space Ω = \{S ⊆ [p] : |S| ≤ K\}, K = min(p, ⌊n/2⌋) (finite).
\item G\_S = X\_SᵀX\_S, z\_S = X\_Sᵀy, H\_S = (n+1)G\_S + I, D\_S = det H\_S / det(G\_S + I), Q\_S = yᵀy − n z\_SᵀH\_S⁻¹z\_S, ν = 0.01 + n/2, a₀ = b₀ = 0.01; m(S) = (2b₀)\textasciicircum{}\{a₀\} Γ(ν) / (π\textasciicircum{}\{n/2\} Γ(a₀)) · D\_S\textasciicircum{}\{−1/2\} (2b₀ + Q\_S)\textasciicircum{}\{−ν\} (for S = ∅: D\_∅ = 1, Q\_∅ = yᵀy).
\item w(k) = ½ / C(2p−1, k); W\_K = Σ\_\{k≤K\} C(p,k) w(k); ϖ(S) = w(|S|)/W\_K.
\item Evidence Z = Σ\_\{S∈Ω\} ϖ(S) m(S). BF(S) = m(S)/m(∅).
\item Blocks C₁, …, C\_J (a fixed partition of [p] determined by X); tables r\_b(A) = BF(A), A ⊆ C\_b. γ₀(S) = ϖ(S) ∏\_b r\_b(S ∩ C\_b); Z̄₀ = Σ\_\{S∈Ω\} γ₀(S); discrepancy d(S) = log BF(S) − Σ\_b log r\_b(S ∩ C\_b).
\item Calibration 𝓒: the random object produced by the calibration phase (grid 0 = φ₀ < φ₁ < … < φ\_T = 1, T, L, N\_e, tier, candidate set and all other frozen settings). A calibration value c is \textbf{successful} if it yields a frozen grid, so that the method runs its final phase.
\item \textbf{Final run given c.} For a successful c: 4 islands r = 1, …, 4, each with N\_e particles in Ω, run with the frozen grid of c and the final streams 10–13 (island r uses stream 9 + r). Writing S\_\{r,i\}\textasciicircum{}\{(t−1)\} ∈ Ω for particle i of island r at the moment stage t's reweighting is applied (t = 1, …, T), the island evidence is (method §6.12)

Ẑ\_r = m(∅) Z̄₀ ∏\_\{t=1\}\textasciicircum{}\{T\} (1/N\_e) Σ\_\{i=1\}\textasciicircum{}\{N\_e\} exp\{(φ\_t − φ\_\{t−1\}) d(S\_\{r,i\}\textasciicircum{}\{(t−1)\})\},

and the pooled evidence is Ẑ = ¼ Σ\_\{r=1\}\textasciicircum{}\{4\} Ẑ\_r. The reported lower bound is 0.05·Ẑ.
\item P\_c, E\_c: probability and expectation over the final streams 10–13, with the calibration fixed at c. Under CAL and EX-S (and joint measurability of the final run in (c, streams)), these give a version of the conditional law and expectation given 𝓒, on the event of a successful calibration; the source text writes them P\_𝓒, E\_𝓒. All conclusions below are stated for P\_c, for every successful c.
\item \textbf{Final E2 flag} (method §6.9, recomputed after the final run as a flag only): E2\_final is raised iff the standard deviation of the four numbers log Ẑ\_1, …, log Ẑ\_4 exceeds 1 nat. The method does not fix the divisor; both the population SD (divisor 4) and the sample SD (divisor 3) are covered below.
\end{mdbul}

\mdhead{2}{2. Assumptions}

\begin{mdbul}
\item \textbf{EX-A} (declared idealisation): exact scores, sums, products, probabilities and comparisons.
\item \textbf{EX-S} (declared idealisation): independent ideal uniform random streams with the method's stream separation; the algorithm's outputs are measurable functions of the streams.
\item \textbf{CAL} (method construction): 𝓒 uses streams 0–9 only; streams 10–13 are fresh (independent of 𝓒); escalation is decided before any final randomness.
\item \textbf{(H-T12)} [taken from the sibling conjecture T12 (iii), used here only as a hypothesis]: for every successful calibration value c, E\_c[Ẑ] = Z. (T12 (i) states E\_c[Ẑ\_r] = Z for each island, which implies this.)
\end{mdbul}

\mdhead{2}{3. Conclusions to prove}

Fix any successful calibration value c. Under EX-A, EX-S, CAL and (H-T12):

\textbf{(a) Coverage.} For every α ∈ (0, 1): P\_c(Z ≥ αẐ) ≥ 1 − α. In particular P\_c(Z ≥ 0.05·Ẑ) ≥ 0.95.

\textbf{(b) Log bias.} log Ẑ is a well-defined, integrable random variable under P\_c, and E\_c[log Ẑ] ≤ log Z.

The remaining two conclusions make precise "it is not a relative-error interval" and "conditioning afterwards on final flags does not preserve this coverage automatically". They are \textbf{non-implication statements}: they show that the properties of Ẑ that yield (a) and (b) do not yield more. Let 𝓛(Z) be the set of joint laws of a vector (Ẑ\_1, …, Ẑ\_4) such that the Ẑ\_r are independent and identically distributed, each takes finitely many values, all in (0, ∞), and E[Ẑ\_r] = Z; set Ẑ = ¼ Σ\_r Ẑ\_r and define E2\_final from (Ẑ\_r) as above. Every law in 𝓛(Z) satisfies (a) and (b) (by the same proof) and has every property of Ẑ that the proof of (a)–(b) uses (positive, finitely many values, mean Z). The extra requirements in 𝓛(Z) (i.i.d. islands with islandwise mean Z) are not hypotheses about the actual final run; they only make the non-implications stronger. These clauses do not assert that any particular data set realises the laws exhibited, nor that they are not realised. This non-implication reading is the intended one; the per-data-set reading about the actual estimator would be false: for fixed data the estimator satisfies e\textasciicircum{}\{−Δ\}Z ≤ Ẑ ≤ e\textasciicircum{}\{Δ\}Z surely, with Δ = max\_Ω d − min\_Ω d, a data-dependent constant that is not among the reported quantities listed above and that no property of 𝓛(Z) supplies (proof, Remark at the end).

\textbf{(c) One-sided, not a relative-error interval.} For every c₀ > 0 and every η ∈ (0, 1) there is a law in 𝓛(Z) with P(Ẑ < c₀Z) ≥ 1 − η. Consequently, for no constants 0 < c₁ ≤ 1 ≤ c₂ < ∞ and γ > 0 does the bound P(c₁Z ≤ Ẑ ≤ c₂Z) ≥ γ follow from the properties defining 𝓛(Z) (in particular not from those used in (a)–(b)), even though (a) holds for every such law; the event \{Z ≥ αẐ\} places no upper limit on Z/Ẑ.

\textbf{(d) Conditioning on a final flag.} For every α ∈ (0, 1):

\begin{mdbul}
\item (d1) there is a law in 𝓛(Z) with P(E2\_final raised) > 0 and P(Z ≥ αẐ | E2\_final raised) = 0;
\item (d2) for every γ > 0 there is a law in 𝓛(Z) with P(E2\_final not raised) > 0 and P(Z ≥ αẐ | E2\_final not raised) < γ.
\end{mdbul}

Both hold for either SD divisor. The source's "final flags" is read existentially: one final flag (E2\_final) whose conditioning breaks the coverage suffices to show that conditioning on final flags does not preserve it automatically (a flag independent of Ẑ would preserve it). Consequently, coverage at level 1 − α conditional on either value of the final E2 flag does not follow from the properties defining 𝓛(Z), although (a) holds unconditionally for every such law.

\mdhead{2}{Approved changes to the statement}

None to hypotheses or conclusions. Exposition only (after Stage B review round 1, Opus reviewer findings 1, 2, 5): the paragraph before (c) and the sentence after (d2) were reworded to state why the non-implication reading is the intended one, that the extra requirements of 𝓛(Z) are not hypotheses about the actual run, and that "final flags" is read existentially. Round 2 (both Stage B reviewers): "the only true one" replaced by "the intended one; the per-data-set reading would be false"; "the method does not report" replaced by "not among the reported quantities listed above"; the P\_c ↔ conditional-law sentence now carries the joint-measurability proviso. Versions reviewed: review/ STATEMENT\_as\_reviewed\_r1.md (round 1), review/STATEMENT\_as\_reviewed\_r2.md (round 2).
\end{mdblock}

\mdsec{(c) Final proof, verbatim: body of PROOF\_v3.tex}
\noindent{\small\itshape The proof below is the body of \texttt{PROOF\_v3.tex} from \texttt{\textbackslash begin\{document\}} to \texttt{\textbackslash end\{document\}} without the title block. Its preamble definitions are local to this chapter; labels carry the prefix \texttt{T31:}.}\par\medskip
\mdproofbegin
% ---- preamble of PROOF_v3.tex (theorem environments carry the prefix T31) ----
\newtheorem{T31theorem}{Theorem}
\newtheorem{T31lemma}[T31theorem]{Lemma}
\newtheorem{T31corollary}[T31theorem]{Corollary}
\theoremstyle{remark}
\newtheorem{T31remark}[T31theorem]{Remark}
\newcommand{\Zh}{\widehat Z}
\newcommand{\Pc}{\mathbb{P}_c}
\newcommand{\Ec}{\mathbb{E}_c}
\newcommand{\cL}{\mathcal{L}}
\newcommand{\Etwo}{\mathrm{E2}_{\mathrm{final}}}
% ---- body of PROOF_v3.tex ----
All notation is that of STATEMENT.md. Throughout, $c$ is a fixed successful calibration value; $\Pc,\Ec$ are the probability
and expectation over the final streams 10--13 with the calibration fixed at $c$.

\section{Positivity and finiteness}

\begin{T31lemma}\label{T31:lem:pos}
Under EX-A and EX-S:
\begin{enumerate}[label=(\roman*)]
\item $m(S)\in(0,\infty)$ for every $S\in\Omega$; hence $Z\in(0,\infty)$.
\item $\bar Z_0\in(0,\infty)$ and $d(S)\in\mathbb{R}$ for every $S\in\Omega$.
\item For each $r$, $\Zh_r$ is a random variable under $\Pc$ taking finitely many values, all in $(0,\infty)$. Hence
$\Zh=\frac14\sum_r\Zh_r$ and $\log\Zh$ are random variables taking finitely many values, $\Zh>0$, and $\log\Zh$ is
bounded, in particular integrable.
\end{enumerate}
\end{T31lemma}

\begin{proof}
(i) For $S=\emptyset$, $D_\emptyset=1$ and $Q_\emptyset=y^\top y=n-1>0$. Let $S\neq\emptyset$, $k=|S|$, and take a full
singular value decomposition $X_S=U\Sigma V^\top$ with $U\in\mathbb{R}^{n\times n}$, $V\in\mathbb{R}^{k\times k}$ orthogonal
and $\Sigma\in\mathbb{R}^{n\times k}$ with diagonal entries $\sigma_1,\dots,\sigma_{\min(n,k)}\ge0$ and zeros elsewhere.
Since $G_S=X_S^\top X_S\succeq0$, both $H_S=(n+1)G_S+I$ and $G_S+I$ are positive definite, so $D_S=\det H_S/\det(G_S+I)>0$.
Further $H_S=V\{(n+1)\Sigma^\top\Sigma+I\}V^\top$, so
\[
nX_SH_S^{-1}X_S^\top=U\,\bigl[n\Sigma\{(n+1)\Sigma^\top\Sigma+I\}^{-1}\Sigma^\top\bigr]\,U^\top,
\]
and the bracket is diagonal with entries $n\sigma_j^2/\{(n+1)\sigma_j^2+1\}\in[0,n/(n+1))$ and zeros. Hence
$I_n-nX_SH_S^{-1}X_S^\top$ has all eigenvalues at least $1/(n+1)$, and
\[
Q_S=y^\top y-nz_S^\top H_S^{-1}z_S=y^\top\bigl(I_n-nX_SH_S^{-1}X_S^\top\bigr)y\ \ge\ \frac{y^\top y}{n+1}>0 .
\]
With $a_0=b_0=0.01$ and $\nu=0.01+n/2$, the constants $(2b_0)^{a_0}$, $\Gamma(\nu)$, $\Gamma(a_0)$, $\pi^{n/2}$ are finite and
positive, $D_S^{-1/2}\in(0,\infty)$ and $(2b_0+Q_S)^{-\nu}\in(0,\infty)$; so $m(S)\in(0,\infty)$. Since $1\le p$ and
$K\le p\le 2p-1$, each $\binom{2p-1}{k}$ with $k\le K$ is a positive integer, so $w(k)>0$, $W_K\in(0,\infty)$ and
$\varpi(S)>0$. As $\Omega$ is finite, $Z=\sum_\Omega\varpi(S)m(S)\in(0,\infty)$.

(ii) By (i), $\mathrm{BF}(S)=m(S)/m(\emptyset)\in(0,\infty)$ for $S\in\Omega$. The table entries that enter $\gamma_0$
and $d$ are $r_b(S\cap C_b)$ with $S\in\Omega$; since $|S\cap C_b|\le|S|\le K$, $S\cap C_b\in\Omega$ and so
$r_b(S\cap C_b)=\mathrm{BF}(S\cap C_b)\in(0,\infty)$. Hence $\gamma_0(S)\in(0,\infty)$, $\bar Z_0=\sum_\Omega\gamma_0(S)\in(0,\infty)$ ($\Omega$ finite and
nonempty), and $d(S)=\log\mathrm{BF}(S)-\sum_b\log r_b(S\cap C_b)$ is a finite real number.

(iii) Given $c$, the number of stages $T$ and the particle count $N_e$ are fixed positive integers ($T\ge1$ because
$\phi_0=0<\phi_T=1$; $N_e\ge1$ as a particle count) and the grid $\phi_0,\dots,\phi_T$ is fixed. By STATEMENT.md \S1 every
particle state $S_{r,i}^{(t-1)}$ lies in the finite set $\Omega$, so the array
$A_r=(S_{r,i}^{(t-1)})_{1\le i\le N_e,\,1\le t\le T}$ takes values in the finite set $\Omega^{N_eT}$. By the formula of
STATEMENT.md \S1, $\Zh_r=g(A_r)$ with the deterministic map
\[
g(a)=m(\emptyset)\,\bar Z_0\prod_{t=1}^T\frac1{N_e}\sum_{i=1}^{N_e}\exp\{(\phi_t-\phi_{t-1})\,d(a_{i,t})\}.
\]
By (i)--(ii) every factor is finite and positive (an average of finitely many finite positive numbers), so $g$ maps the
finite set $\Omega^{N_eT}$ into $(0,\infty)$; under EX-A this value is computed exactly. Hence, pointwise on the sample
space, $\Zh_r\in g(\Omega^{N_eT})$, a finite subset of $(0,\infty)$. Each $\Zh_r$ is an output of the algorithm, hence a
random variable under $\Pc$ by EX-S. The same then holds for $\Zh=\frac14\sum_r\Zh_r$, and $\log\Zh$ is a random
variable taking finitely many real values, so it is bounded.
\end{proof}

\section{Coverage and log bias: conclusions (a) and (b)}

The next lemma isolates the only properties of $\Zh$ that (a) and (b) use.

\begin{T31lemma}\label{T31:lem:core}
Let $W$ be a random variable on some probability space, taking finitely many values, all in $(0,\infty)$, with
$\mathbb{E}W=Z\in(0,\infty)$. Then for every $\alpha\in(0,1)$, $\mathbb{P}(Z\ge\alpha W)\ge1-\alpha$, and
$\mathbb{E}\log W\le\log Z$.
\end{T31lemma}

\begin{proof}
Since $W\ge0$, pointwise $W\ge (Z/\alpha)\mathbf 1\{W>Z/\alpha\}$; taking expectations,
$Z=\mathbb{E}W\ge(Z/\alpha)\,\mathbb{P}(W>Z/\alpha)$, so $\mathbb{P}(\alpha W>Z)\le\alpha$ (Markov's inequality), and
$\mathbb{P}(Z\ge\alpha W)=1-\mathbb{P}(\alpha W>Z)\ge1-\alpha$.

Let $v_1,\dots,v_m\in(0,\infty)$ be the distinct values of $W$ and $p_i=\mathbb{P}(W=v_i)$, so $\sum_ip_i=1$ and
$\sum_ip_iv_i=Z$. By concavity of $\log$ on $(0,\infty)$ (finite Jensen inequality),
$\mathbb{E}\log W=\sum_ip_i\log v_i\le\log\sum_ip_iv_i=\log Z$.
\end{proof}

\begin{T31theorem}[Conclusions (a) and (b)]\label{T31:thm:ab}
Under EX-A, EX-S, CAL and (H-T12), for every successful calibration value $c$:
\begin{enumerate}[label=(\alph*)]
\item $\Pc(Z\ge\alpha\Zh)\ge1-\alpha$ for every $\alpha\in(0,1)$; in particular $\Pc(Z\ge0.05\,\Zh)\ge0.95$.
\item $\log\Zh$ is a bounded (hence integrable) random variable under $\Pc$ and $\Ec\log\Zh\le\log Z$.
\end{enumerate}
\end{T31theorem}

\begin{proof}
By Lemma~\ref{T31:lem:pos}, under $\Pc$ the variable $W=\Zh$ takes finitely many values, all in $(0,\infty)$, $\log\Zh$ is
bounded, and $Z\in(0,\infty)$. By (H-T12), $\Ec\Zh=Z$. Lemma~\ref{T31:lem:core} gives (a) and the inequality in (b); the case
$\alpha=0.05$ gives $\Pc(Z\ge0.05\,\Zh)\ge0.95$.
\end{proof}

\begin{T31remark}[$\Pc$ and conditioning on $\mathcal C$]\label{T31:rem:cond}
STATEMENT.md defines $\Pc,\Ec$ as the law of the final run with the calibration frozen at $c$, and all conclusions are
proved for $\Pc$ for every successful $c$. Under CAL and EX-S the final outputs are a function $F(c,U)$ of the
calibration and the final streams $U$, which are independent of $\mathcal C$; by the standard substitution (freezing)
property of conditional expectation, on the event of a successful calibration,
$\mathbb{P}(\,\cdot\mid\mathcal C)=\mathbb{P}_{\mathcal C}$ almost surely (for jointly measurable $F$). This is how the
source's $\mathbb{P}_{\mathcal C},\mathbb{E}_{\mathcal C}$ are read; CAL is used only here.
\end{T31remark}

\begin{T31remark}\label{T31:rem:class}
Let $\cL(Z)$ be as in STATEMENT.md \S3: joint laws of $(\Zh_1,\dots,\Zh_4)$ with i.i.d.\ coordinates, each taking finitely
many values in $(0,\infty)$, with $\mathbb{E}\Zh_r=Z$. For every law in $\cL(Z)$, $\Zh=\frac14\sum_r\Zh_r$ takes finitely
many values in $(0,\infty)$ and $\mathbb{E}\Zh=Z$; so Lemma~\ref{T31:lem:core} applies and (a), (b) hold for it. The proof of
Theorem~\ref{T31:thm:ab} uses only these properties of $\Zh$ (positivity, finitely many values, mean $Z$). The extra
properties in the definition of $\cL(Z)$ (i.i.d.\ islands, islandwise mean $Z$) are not hypotheses about the actual final
run; adding them to the class only makes the non-implications (c) and (d) stronger.
\end{T31remark}

\section{The bound is one-sided: conclusion (c)}

\begin{T31lemma}[rare-tail law]\label{T31:lem:rare}
Let $c_0>0$ and $\eta\in(0,1)$. Put $c'=\min(c_0,1)/2\in(0,\tfrac12]$, $\varepsilon=\eta/4\in(0,\tfrac14)$ and
$B=\{1-(1-\varepsilon)c'\}/\varepsilon$. Let $\Zh_1,\dots,\Zh_4$ be i.i.d.\ with $\mathbb{P}(\Zh_r=c'Z)=1-\varepsilon$ and
$\mathbb{P}(\Zh_r=BZ)=\varepsilon$. This law belongs to $\cL(Z)$ and $\mathbb{P}(\Zh<c_0Z)\ge1-\eta$.
\end{T31lemma}

\begin{proof}
$(1-\varepsilon)c'<1$, so $B>0$; and $B-c'=(1-c')/\varepsilon>0$, so the two values are distinct and positive.
$\mathbb{E}\Zh_r=Z\{(1-\varepsilon)c'+\varepsilon B\}=Z$. So the law is in $\cL(Z)$. On the event
$A=\{\Zh_1=\dots=\Zh_4=c'Z\}$ we have $\Zh=c'Z<c_0Z$ (as $c'\le c_0/2$), and by independence and Bernoulli's inequality
$\mathbb{P}(A)=(1-\varepsilon)^4\ge1-4\varepsilon=1-\eta$.
\end{proof}

\begin{T31corollary}[Conclusion (c)]\label{T31:cor:c}
For every $c_0>0$ and $\eta\in(0,1)$ some law in $\cL(Z)$ has $\mathbb{P}(\Zh<c_0Z)\ge1-\eta$. Consequently, for no
constants $0<c_1\le1\le c_2<\infty$ and $\gamma>0$ does $\mathbb{P}(c_1Z\le\Zh\le c_2Z)\ge\gamma$ hold for every law in
$\cL(Z)$, although every law in $\cL(Z)$ satisfies (a).
\end{T31corollary}

\begin{proof}
The first sentence is Lemma~\ref{T31:lem:rare}. For the second, given $c_1,c_2,\gamma$ we may take $\gamma\le1$ (otherwise the
bound fails trivially). Apply Lemma~\ref{T31:lem:rare} with $c_0=c_1$ and $\eta=\gamma/2\in(0,\tfrac12]$: then
$\mathbb{P}(c_1Z\le\Zh\le c_2Z)\le\mathbb{P}(\Zh\ge c_1Z)\le\eta<\gamma$. The last clause is Remark~\ref{T31:rem:class}.
\end{proof}

\begin{T31remark}
On the event $A$ of Lemma~\ref{T31:lem:rare}, $Z/\Zh=1/c'$, which is arbitrarily large as $c_0\downarrow0$, while the coverage
event $\{Z\ge\alpha\Zh\}$ holds on $A$ for every $\alpha\in(0,1)$. So the reported $0.05\,\Zh$ is a one-sided lower
confidence bound for $Z$: the coverage event places no upper limit on $Z/\Zh$, and $[0.05\,\Zh,\infty)$ is not a
relative-error interval.
\end{T31remark}

\section{Conditioning on the final E2 flag: conclusion (d)}

\begin{T31lemma}[SD and range]\label{T31:lem:sd}
Let $x\in\mathbb{R}^4$ with range $R=\max_ix_i-\min_ix_i$ and $SS=\sum_i(x_i-\bar x)^2$. Then $SS\ge R^2/2$, so the
population SD $\sqrt{SS/4}$ and the sample SD $\sqrt{SS/3}$ are both at least $R/(2\sqrt2)$. In particular, if $R\ge3$
both exceed $1$, and if $R=0$ both equal $0$.
\end{T31lemma}

\begin{proof}
If $R=0$ all coordinates are equal and $SS=0$. Otherwise let $u=\max_ix_i$, $v=\min_ix_i$, attained at distinct indices.
For any real $m$, $(u-m)^2+(v-m)^2=\tfrac12(u-v)^2+2\{\tfrac12(u+v)-m\}^2\ge\tfrac12R^2$; with $m=\bar x$ this gives
$SS\ge R^2/2$. Then $SS/4\ge R^2/8$ and $SS/3\ge R^2/6\ge R^2/8$. If $R\ge3$, both SDs are at least
$3/(2\sqrt2)=\sqrt{9/8}>1$.
\end{proof}

\begin{T31corollary}\label{T31:cor:flag}
Suppose each $\Zh_r/Z$ takes values in a finite set $\mathcal A\subset(0,\infty)$ such that $|\log a-\log a'|\ge3$ for all
distinct $a,a'\in\mathcal A$. Then, for either SD divisor, $\Etwo$ is not raised if and only if
$\Zh_1=\Zh_2=\Zh_3=\Zh_4$.
\end{T31corollary}

\begin{proof}
$\log\Zh_r=\log Z+\log(\Zh_r/Z)$, and adding a constant does not change the SD. If all four are equal, the SD is
$0\le1$. If not, two of them take distinct values in $\mathcal A$, so the range of $(\log\Zh_r)_r$ is at least $3$ and, by
Lemma~\ref{T31:lem:sd}, the SD exceeds $1$.
\end{proof}

\begin{T31theorem}[Conclusion (d1)]\label{T31:thm:d1}
Let $\alpha\in(0,1)$, $a=\tfrac14$, $b=8/\alpha$, $\varepsilon=3\alpha/(32-\alpha)$, and let $\Zh_1,\dots,\Zh_4$ be
i.i.d.\ with $\mathbb{P}(\Zh_r=aZ)=1-\varepsilon$, $\mathbb{P}(\Zh_r=bZ)=\varepsilon$. This law is in $\cL(Z)$,
$\mathbb{P}(\Etwo\text{ raised})>0$, and $\mathbb{P}(Z\ge\alpha\Zh\mid\Etwo\text{ raised})=0$, for either SD divisor.
\end{T31theorem}

\begin{proof}
$0<3\alpha<32-\alpha$, so $\varepsilon\in(0,1)$. $\mathbb{E}\Zh_r/Z=a+\varepsilon(b-a)=\tfrac14+\frac{3}{4b-1}\cdot
\frac{4b-1}{4}=1$. So the law is in $\cL(Z)$. Further $\log(b/a)=\log(32/\alpha)>\log32>3$ (as $e^3<20.1$), so
Corollary~\ref{T31:cor:flag} applies with $\mathcal A=\{a,b\}$: $\Etwo$ is raised iff the four values are not all equal.
Hence $\mathbb{P}(\Etwo\text{ raised})=1-(1-\varepsilon)^4-\varepsilon^4$, which is positive because it is at least
the probability $\varepsilon(1-\varepsilon)^3>0$ of the outcome $(bZ,aZ,aZ,aZ)$. On the event $\{\Etwo\text{ raised}\}$ at
least one $\Zh_r$ equals $bZ$ and the others are at least $aZ>0$, so $\Zh\ge(b+3a)Z/4>bZ/4=2Z/\alpha$, i.e.\
$\alpha\Zh>2Z>Z$. Thus the coverage event fails everywhere on $\{\Etwo\text{ raised}\}$.
\end{proof}

\begin{T31theorem}[Conclusion (d2)]\label{T31:thm:d2}
Let $\alpha\in(0,1)$ and $\gamma>0$. Put $h=2/\alpha$ and let $M$ be an integer with $M\ge h$ and
$M^3\ge81/(\alpha^4\gamma)$. Let $\ell_k=he^{-3k}$ for $k=1,\dots,M$, $\bar\ell=M^{-1}\sum_{k=1}^M\ell_k$,
$q=(1-\bar\ell)/(h-\bar\ell)$, and let $\Zh_1,\dots,\Zh_4$ be i.i.d.\ with $\mathbb{P}(\Zh_r=hZ)=q$ and
$\mathbb{P}(\Zh_r=\ell_kZ)=(1-q)/M$ for each $k$. Then this law is in $\cL(Z)$, $\mathbb{P}(\Etwo\text{ not raised})>0$,
and $\mathbb{P}(Z\ge\alpha\Zh\mid\Etwo\text{ not raised})<\gamma$, for either SD divisor. Taking $\gamma=1-\alpha$, the
conditional coverage given that the flag is not raised is below $1-\alpha$.
\end{T31theorem}

\begin{proof}
\emph{The law is well defined and in $\cL(Z)$.} Using $M\ge h$,
\[
0<\bar\ell=\frac hM\sum_{k=1}^Me^{-3k}<\frac hM\cdot\frac{e^{-3}}{1-e^{-3}}=\frac hM\cdot\frac1{e^3-1}\le\frac1{e^3-1}<\frac1{19},
\]
since $e^3-1>19$ (as $e>2.718$ gives $e^3>2.718^3>20.07$). As $h=2/\alpha>2$, we have $0<\bar\ell<1<h$; so $0<1-\bar\ell<h-\bar\ell$ and $q\in(0,1)$. Moreover
$q>(1-\tfrac1{19})/h=\tfrac{9}{19}\alpha>\alpha/3$. The values $h>\ell_1>\dots>\ell_M>0$ are distinct and the
probabilities $q,(1-q)/M,\dots$ are positive and sum to $1$. Finally
$\mathbb{E}\Zh_r/Z=qh+(1-q)\bar\ell=q(h-\bar\ell)+\bar\ell=1$. So the law is in $\cL(Z)$.

\emph{The flag.} Consecutive values in $\mathcal A=\{h,\ell_1,\dots,\ell_M\}$ have log-differences exactly $3$, so any two
distinct values differ in log by at least $3$, and Corollary~\ref{T31:cor:flag} gives
$\{\Etwo\text{ not raised}\}=\{\Zh_1=\dots=\Zh_4\}$. This event is the disjoint union of $H=\{\text{all four equal }hZ\}$,
with $\mathbb{P}(H)=q^4>0$, and $L=\bigcup_k\{\text{all four equal }\ell_kZ\}$, with
$\mathbb{P}(L)=M\{(1-q)/M\}^4=(1-q)^4M^{-3}$.

\emph{Coverage on each part.} On $H$, $\Zh=hZ$ and $\alpha\Zh=2Z>Z$: coverage fails. On $\{\text{all four equal }
\ell_kZ\}$, $\alpha\Zh=2e^{-3k}Z<Z$: coverage holds. Hence, with $x=(1-q)^4M^{-3}>0$,
\[
\mathbb{P}(Z\ge\alpha\Zh\mid\Etwo\text{ not raised})=\frac{x}{q^4+x}<\frac{x}{q^4}\le\frac{M^{-3}}{q^4}
<\frac{81}{\alpha^4M^3}\le\gamma,
\]
using $(1-q)^4\le1$, $q>\alpha/3$ and $M^3\ge81/(\alpha^4\gamma)$.
\end{proof}

\begin{T31remark}[What (c) and (d) do and do not say]
By Remark~\ref{T31:rem:class}, every law used in Lemma~\ref{T31:lem:rare} and Theorems~\ref{T31:thm:d1}--\ref{T31:thm:d2} satisfies (a) and
(b), has i.i.d.\ islands with islandwise mean $Z$, and has $\Zh$ positive with finitely many values. So none of these
properties --- in particular, none of the properties of the final-run evidence used to prove (a) and (b) --- implies a
two-sided relative-error bound, or coverage conditional on either value of the final E2 flag. These are non-implication
statements about the properties used; they do not assert that a particular data set and calibration realise the laws
exhibited.

This reading is the one intended, and the per-data-set reading about the actual estimator is false, as follows. For
fixed data, put $d_{\min}=\min_\Omega d$, $d_{\max}=\max_\Omega d$ and $\Delta=d_{\max}-d_{\min}$, and let
$A=m(\emptyset)\bar Z_0>0$. The grid is strictly increasing, so $\phi_t-\phi_{t-1}>0$ and each term
$\exp\{(\phi_t-\phi_{t-1})d(S)\}$, hence each stage average, lies in
$[e^{(\phi_t-\phi_{t-1})d_{\min}},e^{(\phi_t-\phi_{t-1})d_{\max}}]$. Since $\sum_t(\phi_t-\phi_{t-1})=\phi_T-\phi_0=1$, every
$\Zh_r$, and hence their average $\Zh$, lies in $[Ae^{d_{\min}},Ae^{d_{\max}}]$. By the definitions of $\gamma_0$ and $d$,
$\varpi(S)\mathrm{BF}(S)=\gamma_0(S)e^{d(S)}$, so $Z/A=\sum_\Omega\gamma_0(S)e^{d(S)}/\bar Z_0$ is a $\gamma_0$-weighted average
of the numbers $e^{d(S)}$ and $Z$ lies in the same interval. Hence, surely, $e^{-\Delta}Z\le\Zh\le e^{\Delta}Z$. So, for
fixed data, the actual estimator does satisfy a two-sided relative-error bound, but with a data-dependent constant
$e^{\Delta}$ that is not among the reported quantities listed in STATEMENT.md (the $\Zh_r$, $\Zh$ and $0.05\,\Zh$) and that
no property of $\cL(Z)$ supplies. Similarly, the four numbers $\log\Zh_r$ lie in an interval of length at most $\Delta$, so
their sum of squared deviations is at most $\Delta^2$ (as $\sum_i(x_i-\bar x)^2\le\sum_i(x_i-m)^2$ for the midpoint $m$ of
the interval, and each of the four terms is at most $(\Delta/2)^2$), so
the population SD is at most $\Delta/2$ and the sample SD at most $\Delta/\sqrt3$; thus, if $\Delta\le\sqrt3$, $\Etwo$ is
never raised, and if $e^{\Delta}\le1/\alpha$, coverage holds surely. For such data, conditioning on the flag would
preserve coverage; (d) says only that coverage is not preserved \emph{automatically}, that is, as a consequence of the
properties used in (a)--(b). One witness flag (E2) suffices for this; a flag of positive probability independent of
$\Zh$ would preserve coverage. Whether laws like those of (c)--(d) are realised by the actual final run for some data
set is not addressed here.
\end{T31remark}
\mdproofend
\mdsec{(d) Review record}
\mdsubsec{FINAL\_REPORT.md: statement changes}
\begin{mdblock}
\mdhead{2}{Statement changes}

None to hypotheses or conclusions. The step-1 decision was (a), correct as stated. The source text is quoted verbatim in STATEMENT.md. "Under T12" became the explicit hypothesis (H-T12). The two informal cautions ("not a relative-error interval", "conditioning on final flags does not preserve coverage automatically") were written as precise non-implication statements (c), (d). After review, explanatory wording only was revised twice; the reviewed versions are kept in {\mdtt review/\allowbreak{}STATEMENT\_\allowbreak{}as\_\allowbreak{}reviewed\_\allowbreak{}r1.\allowbreak{}md} and {\mdtt \_\allowbreak{}r2.\allowbreak{}md}. Every reviewer confirmed that no hypothesis or conclusion changed and that the statement is a faithful formalisation of clause (i).
\end{mdblock}

\mdsubsec{FINAL\_REPORT.md: review history}
\begin{mdblock}
\mdhead{2}{Review history}

{\small
\begin{longtable}{|>{\raggedright\arraybackslash}p{\dimexpr(\linewidth-10\tabcolsep-6\arrayrulewidth)*65/1000\relax}|>{\raggedright\arraybackslash}p{\dimexpr(\linewidth-10\tabcolsep-6\arrayrulewidth)*362/1000\relax}|>{\raggedright\arraybackslash}p{\dimexpr(\linewidth-10\tabcolsep-6\arrayrulewidth)*90/1000\relax}|>{\raggedright\arraybackslash}p{\dimexpr(\linewidth-10\tabcolsep-6\arrayrulewidth)*103/1000\relax}|>{\raggedright\arraybackslash}p{\dimexpr(\linewidth-10\tabcolsep-6\arrayrulewidth)*379/1000\relax}|}
\hline
\textbf{Stage} & \textbf{Reviewer} & \textbf{Version} & \textbf{Verdict} & \textbf{Findings} \\ \hline\endhead
A & Astra \#1 (GPT-6 Astra, high; session 01a1641c-085b-7532-afa0-084ca7a98792) & v1 & verified & none \\ \hline
B r1 & Astra \#2 (GPT-6 Astra, high; new session 01a1641e-28c0-7c51-a27c-31ffb6ff38dc) & v1 & verified & none \\ \hline
B r1 & Opus 5.5 xhigh (read-only agent) & v1 & verified & 11 exposition (incl. the sure bound) \\ \hline
B r2 & Astra \#2 (resumed) & v2 + log & verified & 2 exposition ("only true one" overclaim; unverifiable "method does not report") \\ \hline
B r2 & Opus (continued) & v2 + log & verified & 5 exposition (same overclaim; small-Δ derivation; implicit steps; measurability proviso; file name) \\ \hline
B r3 & Astra \#2 (resumed) & v3 + log & verified & none \\ \hline
B r3 & Opus (continued) & v3 + log & verified & 3 exposition (cosmetic; recorded, not applied so the final text is the reviewed v3) \\ \hline
\end{longtable}}

No reviewer, in any round, raised a fatal or gap finding. All exposition findings and the lead's answers are in {\mdtt CLAIMS\_\allowbreak{}OF\_\allowbreak{}FATAL\_\allowbreak{}FLAWS.\allowbreak{}md}; the reviewers agreed with every answer. The Astra judge was not needed. Replies: {\mdtt review/\allowbreak{}astra\_\allowbreak{}01\_\allowbreak{}review\_\allowbreak{}r1.\allowbreak{}log}, {\mdtt review/\allowbreak{}stageB\_\allowbreak{}astra2\_\allowbreak{}r\{1,\allowbreak{}2,\allowbreak{}3\}.\allowbreak{}log}, {\mdtt review/\allowbreak{}stageB\_\allowbreak{}opus\_\allowbreak{}r\{1,\allowbreak{}2,\allowbreak{}3\}.\allowbreak{}md}.

Sanity check (not evidence of proof): a Haiku assistant computed the example laws of (c), (d1) and (d2) exactly for several α, γ, c₀, η. All 104 checks passed and the canary was exact ({\mdtt helpers/\allowbreak{}numcheck/\allowbreak{}check.\allowbreak{}py}, {\mdtt results.\allowbreak{}txt}).
\end{mdblock}

\mdsubsec{CLAIMS\_OF\_FATAL\_FLAWS.md (every entry, verbatim)}
\begin{mdblock}
\mdhead{1}{CLAIMS\_OF\_FATAL\_FLAWS — T31 (T22) part (i)}

\mdhead{2}{Fatal or gap findings}

None. Three independent reviews of PROOF\_v1 (Astra \#1, Astra \#2, Opus 5.5 xhigh) reported no fatal and no gap finding; all three returned \textbf{verified}.

\mdhead{2}{Exposition findings and the lead's answers (for the record; Opus reviewer, Stage B round 1)}

{\small
\begin{longtable}{|>{\raggedright\arraybackslash}p{\dimexpr(\linewidth-8\tabcolsep-5\arrayrulewidth)*46/1000\relax}|>{\raggedright\arraybackslash}p{\dimexpr(\linewidth-8\tabcolsep-5\arrayrulewidth)*385/1000\relax}|>{\raggedright\arraybackslash}p{\dimexpr(\linewidth-8\tabcolsep-5\arrayrulewidth)*256/1000\relax}|>{\raggedright\arraybackslash}p{\dimexpr(\linewidth-8\tabcolsep-5\arrayrulewidth)*312/1000\relax}|}
\hline
\textbf{\#} & \textbf{Finding (short)} & \textbf{Answer} & \textbf{Change in v2} \\ \hline\endhead
1 & Non-implication reading of (c)/(d) is the only true one: surely e\textasciicircum{}\{−Δ\}Z ≤ Ẑ ≤ e\textasciicircum{}\{Δ\}Z with Δ = max d − min d; statement should say so & Accepted (checked: ϖ·BF = γ₀e\textasciicircum{}\{d\}, stage factors in [e\textasciicircum{}\{Δφ d\_min\}, e\textasciicircum{}\{Δφ d\_max\}]) & STATEMENT §3 paragraph before (c) reworded; new final paragraph of the last Remark in the proof derives the bound \\ \hline
2 & "Final flags" plural vs one witness flag & Accepted: existential reading & STATEMENT sentence after (d2); last Remark \\ \hline
3 & Link P\_c ↔ conditional law given 𝓒 not written out; CAL unused & Accepted & New Remark "P\_c and conditioning on 𝓒" (substitution lemma; CAL used only there) \\ \hline
4 & Fidelity to METHODOLOGY\_v3 and T12 not checkable from the packet & Noted (scope of the packet). Lead checked: Ẑ\_r formula = METHODOLOGY\_v3 §6.12 verbatim; E2 = §6.9, "recomputed as flags only" after the final run; (H-T12) = T12(iii) first identity & None \\ \hline
5 & Preamble could be read as asserting iid islands / islandwise unbiasedness of the actual run & Accepted & STATEMENT §3 preamble reworded; Remark on 𝓛(Z) extended \\ \hline
6 & π₀ and "every move stays in Ω" not in the packet & Accepted & Lemma 1(iii) now cites STATEMENT §1 (states in Ω) only \\ \hline
7 & Measurability via "outputs", N\_e ≥ 1, T ≥ 1 & Accepted & Lemma 1(iii) reworded \\ \hline
8 & Table entries with |A| > K & Accepted & Lemma 1(ii): only A = S ∩ C\_b ∈ Ω enter \\ \hline
9 & "On this event" ambiguous & Accepted & Theorem (d1) proof \\ \hline
10 & e³ > 20 unjustified & Accepted & Theorem (d2) proof: e > 2.718 ⇒ e³ > 20.07 \\ \hline
11 & "No upper limit on Z/Ẑ" informal & Accepted as informal remark; no change needed & None \\ \hline
\end{longtable}}

\mdhead{2}{Round 2 (PROOF\_v2): rulings and new findings}

Both Stage B reviewers returned \textbf{verified} on PROOF\_v2, with no fatal or gap finding, and agreed with the answers to rows 2–11. Row 1: both agree with the sure bound but disagree with the wording "the only true one" (accepted, see R2-1). Full rulings: review/stageB\_opus\_r2.md, review/stageB\_astra2\_r2.log.

{\small
\begin{longtable}{|>{\raggedright\arraybackslash}p{\dimexpr(\linewidth-8\tabcolsep-5\arrayrulewidth)*32/1000\relax}|>{\raggedright\arraybackslash}p{\dimexpr(\linewidth-8\tabcolsep-5\arrayrulewidth)*374/1000\relax}|>{\raggedright\arraybackslash}p{\dimexpr(\linewidth-8\tabcolsep-5\arrayrulewidth)*120/1000\relax}|>{\raggedright\arraybackslash}p{\dimexpr(\linewidth-8\tabcolsep-5\arrayrulewidth)*475/1000\relax}|}
\hline
\textbf{\#} & \textbf{Finding (short; source)} & \textbf{Answer} & \textbf{Change in v3} \\ \hline\endhead
R2-1 & "Only true reading" overclaims; only the per-data-set reading is refuted (Opus R2-1; Astra \#2 finding 1) & Accepted & STATEMENT §3 and last Remark: "the intended one; the per-data-set reading would be false"; added "whether such laws are realised by the actual run for some data set is not addressed" \\ \hline
R2-2 & Small-Δ claims lack derivation and assume such data exist (Opus) & Accepted & Last Remark: SD ≤ Δ/2 (pop), ≤ Δ/√3 (sample) derived; stated conditionally ("if Δ ≤ √3", "if e\textasciicircum{}Δ ≤ 1/α"), no existence claim \\ \hline
R2-3 & Implicit steps in the sure bound; "method does not compute/report" unverifiable (Opus; Astra \#2 finding 2) & Accepted & Steps written out (positive increments, averaging, γ₀-weighted average); "not among the reported quantities listed in STATEMENT.md" \\ \hline
R2-4 & Joint-measurability proviso missing in STATEMENT §1 (Opus) & Accepted & STATEMENT §1 P\_c sentence now carries the proviso and says conclusions are stated for P\_c \\ \hline
R2-5 & Proof cites "STATEMENT.md" while packet file is STATEMENT\_v2.md (Opus) & No change: the project file is STATEMENT.md; packet copies are renamed only to keep versions apart & None \\ \hline
\end{longtable}}

\mdhead{2}{Round 3 (PROOF\_v3): confirmation}

Astra \#2 and the Opus reviewer both returned \textbf{verified} on PROOF\_v3, agreed with all Round 2 answers (R2-1..R2-5: change works), and reported no fatal or gap finding. Opus added three exposition notes (R3-1: the (d) half of the per-data remark is true for every data set by a one-line argument not on the page; R3-2: symbol A reused; R3-3: wording of "reported quantities"). Lead's answer: all three accepted as correct and cosmetic; \textbf{not applied}, so that the final text is exactly the reviewed PROOF\_v3. They concern commentary only, not any conclusion.
\end{mdblock}

\end{document}
